\documentclass[11pt,a4paper]{amsart}

% Begin input: preamble.tex
\usepackage[T1]{fontenc}
\usepackage{lmodern,microtype}
\usepackage[margin=26mm]{geometry}
\usepackage{amsmath,amssymb,amscd,mathtools,mathrsfs}
\usepackage{xcolor,enumitem,array,booktabs,longtable}
\usepackage[colorlinks=true,linkcolor=blue!35!black,citecolor=blue!35!black,urlcolor=blue!35!black]{hyperref}
\hypersetup{pdftitle={Rational Gersten kernels in ramified regular local rings},pdfauthor={OpenAI},pdfsubject={Algebraic K-theory and rational Gersten kernels}}
\definecolor{EditorialBlue}{rgb}{0.05,0.20,0.65}
\definecolor{EditorialRed}{rgb}{0.65,0.08,0.08}
\definecolor{EditorialPurple}{rgb}{0.45,0.10,0.55}
\numberwithin{equation}{section}

\newtheorem{theorem}{Theorem}[section]
\newtheorem{lemma}[theorem]{Lemma}
\newtheorem{proposition}[theorem]{Proposition}
\newtheorem{corollary}[theorem]{Corollary}
\theoremstyle{definition}
\newtheorem{definition}[theorem]{Definition}
\newtheorem{setup}[theorem]{Setup}
\newtheorem{notation}[theorem]{Notation}
\theoremstyle{remark}
\newtheorem{remark}[theorem]{Remark}
\newtheorem{example}[theorem]{Example}


% Begin input: common-macros.tex
\newcommand{\Q}{\mathbf Q}
\newcommand{\Z}{\mathbf Z}
\newcommand{\R}{\mathbf R}
\newcommand{\C}{\mathbf C}
\newcommand{\F}{\mathbf F}
\newcommand{\A}{\mathbf A}
\newcommand{\PP}{\mathbf P}
\newcommand{\OO}{\mathcal O}
\newcommand{\m}{\mathfrak m}
\newcommand{\n}{\mathfrak n}
\newcommand{\Spec}{\operatorname{Spec}}
\newcommand{\Pic}{\operatorname{Pic}}
\newcommand{\Cl}{\operatorname{Cl}}
\newcommand{\Frac}{\operatorname{Frac}}
\newcommand{\divisor}{\operatorname{div}}
\newcommand{\colim}{\operatorname*{colim}}
\newcommand{\HD}{H_{\mathcal D,\mathrm{an}}}
\newcommand{\reg}{\operatorname{reg}}
\newcommand{\Qrat}{_{\Q}}
\newcommand{\Gal}{\operatorname{Gal}}

% End input: common-macros.tex


% Begin input: core-macros.tex
\newcommand{\CthreeZ}{\mathbf{Z}}
\newcommand{\CthreeQ}{\mathbf{Q}}
\newcommand{\CthreeF}{\mathbf{F}}
\newcommand{\CthreeA}{\mathbf{A}}
\newcommand{\CthreePP}{\mathbf{P}}
\newcommand{\CthreebS}{\mathbb{S}}
\newcommand{\CthreecO}{\mathcal{O}}
\newcommand{\Cthreefm}{\mathfrak{m}}
\newcommand{\Cthreedd}{\mathrm{d}}
\newcommand{\Cthreects}{\mathrm{cts}}
\newcommand{\Cthreerel}{\mathrm{rel}}
\newcommand{\Cthreetriv}{\mathrm{triv}}
\newcommand{\Cthreetopcomp}{\mathrm{top}}
\DeclareMathOperator{\CthreeK}{K}
\DeclareMathOperator{\CthreeHH}{HH}
\DeclareMathOperator{\CthreeHC}{HC}
\DeclareMathOperator{\CthreeHP}{HP}
\DeclareMathOperator{\CthreeTHH}{THH}
\DeclareMathOperator{\CthreeTC}{TC}
\DeclareMathOperator{\CthreeSpec}{Spec}
\DeclareMathOperator{\CthreeSpf}{Spf}
\DeclareMathOperator{\CthreeFrac}{Frac}
\DeclareMathOperator{\Cthreefib}{fib}
\DeclareMathOperator{\Cthreeid}{id}
\DeclareMathOperator{\Cthreedlog}{dlog}
\newcommand{\DtwoZZ}{\mathbf{Z}}
\newcommand{\DtwoQQ}{\mathbf{Q}}
\newcommand{\DtwoRR}{\mathbf{R}}
\newcommand{\DtwoCC}{\mathbf{C}}
\newcommand{\DtwoFF}{\mathbf{F}}
\newcommand{\DtwoPP}{\mathbf{P}}
\newcommand{\DtwoAA}{\mathbf{A}}
\newcommand{\DtwocO}{\mathcal O}
\newcommand{\DtwocK}{\mathcal K}
\newcommand{\DtwocE}{\mathcal E}
\newcommand{\DtwocQ}{\mathcal Q}
\newcommand{\DtwocT}{\mathcal T}
\newcommand{\DtwocY}{\mathcal Y}
\newcommand{\DtwocP}{\mathcal P}
\newcommand{\DtwocL}{\mathcal L}
\newcommand{\Dtwofm}{\mathfrak m}
\newcommand{\Dtwofn}{\mathfrak n}
\newcommand{\Dtwowt}{\widetilde}
\newcommand{\Dtwowh}{\widehat}
\newcommand{\Dtwoct}{\widehat\otimes}
\newcommand{\DtwodR}{\mathrm{dR}}
\newcommand{\Dtworig}{\mathrm{rig}}
\newcommand{\Dtwoconv}{\mathrm{conv}}
\newcommand{\Dtwocris}{\mathrm{cris}}
\newcommand{\Dtwocts}{\mathrm{cts}}
\newcommand{\Dtwotopol}{\mathrm{top}}
\newcommand{\Dtwoar}{\mathrm{ar}}
\newcommand{\Dtwocyc}{\mathrm{cyc}}
\newcommand{\Dtwoet}{\mathrm{\acute et}}
\newcommand{\Dtwoan}{\mathrm{an}}
\newcommand{\Dtwoid}{\mathrm{id}}
\newcommand{\Dtwoim}{\mathrm{im}}
\newcommand{\Dtwofib}{\mathrm{fib}}
\newcommand{\Dtwogr}{\mathrm{gr}}
\newcommand{\Dtwored}{\mathrm{red}}
\newcommand{\DtwoSpec}{\operatorname{Spec}}
\newcommand{\DtwoProj}{\operatorname{Proj}}
\DeclareMathOperator{\DtwoFrac}{Frac}
\DeclareMathOperator{\DtwoPic}{Pic}
\DeclareMathOperator{\DtwoCl}{Cl}
\DeclareMathOperator{\DtwoTr}{Tr}
\DeclareMathOperator{\DtwoNm}{Nm}
\DeclareMathOperator{\DtwoGal}{Gal}
\DeclareMathOperator{\Dtwodivisor}{div}
\DeclareMathOperator{\Dtworank}{rank}
\DeclareMathOperator{\DtwoLi}{Li}
\DeclareMathOperator{\DtwoHH}{HH}
\DeclareMathOperator{\DtwoHC}{HC}
\DeclareMathOperator{\DtwoHP}{HP}
\DeclareMathOperator{\DtwoTHH}{THH}
\DeclareMathOperator{\DtwoTC}{TC}
\DeclareMathOperator{\Dtwoch}{ch}

% End input: core-macros.tex


% End input: preamble.tex


\title[Rational Gersten kernels]{Rational Gersten kernels in ramified regular local rings}
\author{OpenAI}
\thanks{The OpenAI byline is a user-requested attribution to preparation with ChatGPT/Codex; this is not an official OpenAI publication or institutional endorsement.}
\date{8 October 2026}
\keywords{Algebraic K-theory, Gersten conjecture, ramified regular local rings, regulators, continuous cyclic homology}

\begin{document}
\begin{abstract}

% Begin input: abstract.tex
We construct infinite-order elements in the kernels of rational generic
restriction maps for ramified regular local rings of mixed
characteristic. The basic examples have dimensions two and three,
with residue fields $\mathbf F_5(T)$ and $\mathbf F_3$, respectively.
In both cases an actual $K_3$-class remains of infinite order after
completion and vanishes integrally in the fraction field. Explicit
products of three units are given in each dimension. The proofs combine
supported coefficient classes, arithmetic divisor calculations, and
two cohomological detectors. Analytic Deligne regulators survive every
finite localization needed for the arithmetic local rings; relative
continuous cyclic characters detect the completed classes in ordinary
continuous de Rham quotients. We extend the constructions to every
residue prime and obtain arbitrarily large kernel ranks in fixed local
dimension. Cyclotomic coefficients and continuous Laurent residues give
classes in higher degrees. At every prime, fixed three-dimensional
rings, including complete rings, have infinite rational generic-kernel
rank in every degree at least three. A fixed complete two-dimensional
ring in residue characteristic five has the same property. We also
construct explicit Milnor, Beilinson-motivic, and higher-Chow classes
with zero generic restriction, keeping actual classes distinct from
their images in completed spectra.

% End input: abstract.tex

\end{abstract}
\maketitle

\setcounter{tocdepth}{1}
\tableofcontents


% Begin input: sections/01-introduction.tex
\section{Introduction}\label{sec:introduction}

For a regular local domain $R$ with fraction field $F$, consider the
rational restriction map
\[
 K_j(R)_{\Q}\longrightarrow K_j(F)_{\Q},
 \qquad K_j(R)_{\Q}=K_j(R)\otimes_{\Z}\Q.
\]
Its injectivity is the first assertion of the augmented rational
Gersten conjecture. We construct classes in its kernel for ramified
regular local rings of mixed characteristic. The classes are represented
in actual algebraic $K$-groups: they have infinite order, and their
restrictions to the fraction field vanish before tensoring with $\Q$.

The two basic rings are
\begin{equation}\label{eq:intro-basic-rings}
\begin{aligned}
 A_2&=\left(\Z_{(5)}[T,x,y]/
       (\Phi_5(1+y)+x^4+Txy^3)\right)_{(5,x,y)},\\
 A_3&=\left(\Z_{(3)}[x,y,z]/(3+xy+z^2)\right)_{(x,y,z)},
\end{aligned}
\end{equation}
where $\Phi_5(t)=1+t+t^2+t^3+t^4$.
The ring $A_2$ has dimension two, maximal ideal $(x,y)$, and residue
field $\F_5(T)$; the ring $A_3$ has dimension three, maximal ideal
$(x,y,z)$, and residue field $\F_3$.
Write $\widehat A_i$ for completion at the maximal ideal.
Theorem~\ref{d2:thm:completed} and
Theorem~\ref{thm:explicit-three-dimensional} construct classes
$c_i\in K_3(A_i)$ whose images in $K_3(\widehat A_i)$ also have
infinite order. In each of the four rings the corresponding class
vanishes after inverting $x$, and hence in the fraction field,
integrally.

In dimension three the class is the product of three explicit units:
\begin{equation}\label{eq:intro-explicit-three}
 c_3=2\left\{\frac{1+z}{2},\,1+x,\,
                   1+\frac{(1+x)y}{4}\right\}.
\end{equation}
The braces denote the product of unit classes in Quillen $K_1$.
Dennis--Stein identities prove its integral vanishing after inverting
$x$. Its infinite order is detected on a smooth formal quadric by a
convergent three-adic endpoint series. The actual vanishing of the
special-fiber group $K_4$ makes the relative-to-absolute map injective,
so this argument detects an actual $K_3$-class.
The two-dimensional construction starts with a coefficient in the
$K_3$ of a cyclotomic integer ring and pushes it forward from the
regular divisor $x=0$. A second construction identifies the explicit
unit product
\[
 \left\{2+y,\ 1-x,\
 1+(1-x)(x^3+Ty^3)y\bigl((1+y)^5+1\bigr)\right\}
 \in K_3(A_2)
\]
as an infinite-order generic-kernel class in both $A_2$ and
$\widehat A_2$; see Proposition~\ref{d2ext:dimtwoExplicitTriple}.

The same methods give families with larger rank, higher $K$-degree,
and one fixed ring supporting infinitely many independent classes.
Table~\ref{tab:intro-results} locates these extensions. Every rank
assertion concerns the rational generic kernel; the witnesses are
actual integral classes vanishing on the indicated principal open.
The parameters introduced in residue fields preserve the local
Krull dimension.

\begin{table}[htbp]
\centering
\small
\renewcommand{\arraystretch}{1.18}
\begin{tabular}{@{}p{0.10\textwidth}p{0.64\textwidth}p{0.20\textwidth}@{}}
\hline
Dim. & Result & Reference\\
\hline
$2$ & At every residue prime, tame families have unbounded kernel rank
in degree three and nonzero kernel in every odd degree at least three.
& \ref{d2ext:twodimtamefamily}, \ref{d2ext:twodimtameallodd}\\
$3$ & Eisenstein families have explicit rank formulas before completion
and arbitrarily large lower bounds after completion.
& \ref{d3:thm:eisenstein}, \ref{d3:thm:completedeisensteinrank}\\
$3$ & Explicit cyclotomic triples at every prime retain their independence
after henselization and completion.
& \ref{d3:thm:explicitwildtriples}\\
$3$ & Products of $j$ specified units give infinite-order kernel classes
for every $j\geq3$.
& \ref{d3:thm:everydegreelaurent}\\
$2,3$ & At each prime, a fixed ring has infinite kernel rank in every odd
degree at least three; in dimension three this also holds for its completion.
& \ref{d2ext:twodimstrictcoefficientinfinite}, \ref{d3:thm:strictallprimeallodd}\\
$3$ & At every prime, one fixed ring and one fixed complete ring, with
residue field $\overline{\F}_p(T)$, have infinite kernel rank in every $j\geq3$.
& \ref{d3:thm:singlefixedalldegree}, \ref{d3:thm:completedfixedalldegree}\\
$2$ & One complete ring with residue field $\overline{\F}_5(T,S)$ has
infinite kernel rank in every $j\geq3$.
& \ref{d2ext:dimtwoCompleteAllDegreesInfiniteRank}\\
$2,3$ & Explicit degree-three classes yield Milnor kernels; the
three-dimensional class also gives Beilinson-motivic and higher-Chow kernels.
& \ref{d2ext:twodimexplicitMilnor}, \ref{d3:cor:explicitmilnormotivic}\\
\hline
\end{tabular}
\caption{Families and consequences of the basic constructions.
The cited statements give the coefficient rings and all rank bounds.}
\label{tab:intro-results}
\end{table}

The geometric and completed proofs detect the supported classes by
different cohomology theories. On an arithmetic quadric, a small
resolution shows that the ruling has infinite order in the local
class group. This calculation forces each geometric boundary component
removed by an allowed denominator to have zero divisor class.
Consequently the ruling's Betti class survives every finite
localization. Multiplying by a number-field coefficient with nonzero
regulator gives a nonzero analytic Deligne class on every such open.
Filtered continuity then proves nonvanishing on the local ring.
For the quartic surface, the corresponding divisor calculation uses
the exceptional quartic curve and its Jacobian.

Completion requires a separate detector. The Beilinson fiber square
of Antieau--Mathew--Morrow--Nikolaus identifies the relevant relative
$p$-completed $K$-theory with continuous cyclic homology
\cite{AMMN}. A finite mixed Hochschild--Kostant--Rosenberg comparison
retains the ordinary continuous de Rham quotients. On a quadric, an
explicit residue operator detects the ruling class. On the quartic
family, the elliptic quotients of its Jacobian and a bounded
slope-one functional detect the weighted root class after passage to
the generic Cohen completion. In each case the proof separately
controls the relative-to-absolute map and localization of the
integral test ring. This is the step that returns a cohomological
nonvanishing statement to the actual algebraic $K$-group.

For the rank constructions, cyclotomic classes supply finite families
of independent coefficient values. Their Chern classes factor through
spectral completion, so independence passes to the continuous cyclic
point coordinates without requiring surjectivity from actual
$K$-theory. A continuous Laurent constant-term operator preserves the
ruling detector after adjoining units. Products with these units raise
the $K$-degree, while unramified coefficient extensions increase the
rank. Finite descent handles each possible relation before completion;
a separate comparison for completed essentially smooth models handles
the fixed complete rings.

We use $K(R)$ for the algebraic $K$-theory spectrum and
$K(R;\Q_p)=K(R)^{\wedge}_p[1/p]$ for its rationalized spectral
$p$-completion. This notation is distinct from tensoring each actual
$K$-group with $\Q_p$. Throughout the continuous de Rham arguments,
cohomology means the quotient by the image of the differential, with
no closure of that image.
The final mathematical section gives an explicit obstruction to the
strict simplicial-functor construction in
\cite[Lemma~13(III)]{MochizukiParameters}.
The scope of the accompanying computational and formal checks is
recorded separately at the end of the article.

% End input: sections/01-introduction.tex


% Begin input: sections/02-explicit-three-adic.tex
\section{An explicit three-adic kernel class}
\label{sec:explicit-core}
\subsection{The ring and the class}\label{c3:sec:introduction}

\medskip
\noindent\textbf{Generic restriction.}
Products of units give explicit classes with which to study restriction
from a regular local ring to its fraction field. Dennis--Stein identities
can establish the vanishing of such a class after localization
\cite[III, Definition~5.11 and Theorem~5.11.1]{Weibel}.
For nonvanishing, we use the continuous fiber comparison
\cite[Proposition~4.10]{AMMN}, together with multiplicative traces and
the module structure on cyclic homology
\cite{BGTTrace,NikolausScholze,HoyoisCircle,BokstedtOttosen}.
The example below requires both an integral calculation in algebraic
$\CthreeK$-theory and an ordinary, possibly non-Hausdorff, differential quotient.

\medskip
\noindent\textbf{An explicit class.}
For units $b_1,\ldots,b_r$ in a commutative ring $R$, write
$\{b_1,\ldots,b_r\}=[b_1]\cdots[b_r]\in\CthreeK_r(R)$ for the product of their
classes in actual Quillen algebraic $\CthreeK$-theory.

\begin{theorem}
\label{thm:explicit-three-dimensional}
Let
\[
 A_3=
 \left(\CthreeZ_{(3)}[x,y,z]/(3+xy+z^2)\right)_{(x,y,z)},
 \qquad
 \widehat A_3=\CthreeZ_3[[x,y,z]]/(3+xy+z^2).
\]
Then $\widehat A_3$ is the maximal-ideal completion of $A_3$, and both
rings are ramified regular local domains of dimension three.
For either $R=A_3$ or $R=\widehat A_3$, the class
\[
 c_R=
 2\left\{\frac{1+z}{2},\,1+x,\,
              1+\frac{(1+x)y}{4}\right\}\in\CthreeK_3(R)
\]
has infinite order and maps to zero in $\CthreeK_3(R[1/x])$ integrally.
Consequently,
\[
 \CthreeK_3(R)\otimes_{\CthreeZ}\CthreeQ
 \longrightarrow
 \CthreeK_3(\CthreeFrac(R))\otimes_{\CthreeZ}\CthreeQ
\]
is not injective.
\end{theorem}

The same unit formula therefore gives a nonzero rational class in both
the algebraic local ring and its completion. Its generic vanishing
already holds before rationalization.

\medskip
\noindent\textbf{Method.}
The integral calculation uses a Dennis--Stein symbol whose product with
twice a unit class is $c_R$.
For nonvanishing, we map to a smooth formal quadric over an Eisenstein
quadratic extension of $\CthreeZ_3$.
Proposition~\ref{c3:prop:relative-character} constructs the relative
character with its $\CthreeK$-module action and fixes its value on relative
units. Lemma~\ref{c3:lem:ordinary-residue} then detects its top component by
two endpoint values, without constructing a convergent primitive for a
closed one-form. The endpoint difference is nonzero by an estimate
valid for every term of its defining series.
Vanishing of the actual fourth $\CthreeK$-group of the special fiber finally
passes this detection from relative to absolute $\CthreeK$-theory.

Section~\ref{c3:sec:preliminaries} fixes the conventions and the continuous
comparison. Section~\ref{c3:sec:integral} proves integral vanishing.
Section~\ref{c3:sec:character} constructs and normalizes the character.
Section~\ref{c3:sec:residue} proves the residue criterion.
Section~\ref{c3:sec:nonvanishing} completes the proof.







\subsection{Preliminaries}\label{c3:sec:preliminaries}

We use Quillen algebraic $\CthreeK$-theory and its unit products with the
conventions of \cite{Quillen,Weibel}.
For continuous cyclic constructions we follow
\cite[Definition~4.1 and Proposition~4.2]{AMMN}.
We write $\mathrm{HH}$, $\mathrm{HC}$, $\mathrm{HC}^{-}$, and $\mathrm{HP}$ for Hochschild, cyclic, negative cyclic, and periodic cyclic homology, respectively. The symbols $\mathrm{THH}$ and $\mathrm{TC}$ denote topological Hochschild and topological cyclic homology.

\begin{notation}\label{c3:notation:k-theory}
An uncompleted group $\CthreeK_i(D)$ always means an actual Quillen
$\CthreeK$-group.
For an ideal $J\subseteq D$, the relative spectrum is
\[
 \CthreeK(D,J)=\Cthreefib\bigl(\CthreeK(D)\longrightarrow\CthreeK(D/J)\bigr),
 \qquad
 \CthreeK_i(D,J)=\pi_i\CthreeK(D,J).
\]
Thus relative $\CthreeK$-theory is defined by a homotopy fiber, rather than
by the kernel of a map on homotopy groups.

For a spectrum-valued invariant, the notation $;\CthreeQ_p$ means
$p$-completion of the spectrum followed by inversion of $p$.
It does not mean tensoring its actual homotopy groups with $\CthreeQ_p$.
The superscript $\Cthreects$ denotes the homotopy inverse limit over
the $p$-power reductions.
\end{notation}

Unit products are natural under ring maps, multilinear, and graded
commutative. They satisfy $\{b,1-b\}=0$ whenever both entries are units;
see \cite[III, Lemma~5.10.2, Definition~5.11, and the unit-product
discussion following Theorem~5.11.1]{Weibel}.
The distinction in Notation~\ref{c3:notation:k-theory} will be retained
when these products are mapped to continuous invariants.

\begin{setup}\label{c3:setup:coefficient-ring}
Put
\[
 W=\CthreeZ_3[u]/(u^2+3),\qquad E=W[1/3].
\]
The Eisenstein equation makes $W$ a complete discrete valuation ring
with uniformizer $u$ and residue field $\CthreeF_3$.
Its $u$-adic and $3$-adic topologies agree.
For a smooth $W$-algebra $D_0$ of relative dimension at most two, let
$D$ denote its $3$-adic completion, and put
\[
 \Omega_D^j=
 \left(\varprojlim_n\Omega^j_{D_0/W}/3^n\Omega^j_{D_0/W}\right)[1/3].
\]
All differentials on these modules are relative to $W$.
The notation $W\langle T_1,\ldots,T_r\rangle$ denotes the restricted
power series algebra, namely the $3$-adic completion of
$W[T_1,\ldots,T_r]$.

The two cases used below are $D=W\langle T\rangle$ and
\[
 B=W\langle X,Y,Z\rangle/(XY+Z^2-1).
\]
The algebra $W[X,Y,Z]/(XY+Z^2-1)$ is smooth of relative dimension two:
its partial derivatives $Y,X,2Z$ have no simultaneous zero on the
quadric.
\end{setup}

The continuous comparison will be used in its formal-scheme form.

\begin{proposition}[Continuous fiber comparison
{\cite[Proposition~4.10]{AMMN}}]\label{c3:prop:continuous-comparison}
Let $\CthreecO$ be the integer ring of a complete discretely valued
characteristic-zero field with perfect residue field $k$ of
characteristic $p$. Let $\mathfrak X$ be a smooth $p$-adic formal
$\CthreecO$-scheme. There is a natural diagram
\begin{equation}\label{c3:eq:continuous-comparison}
\begin{CD}
 \CthreeK^\Cthreects(\mathfrak X;\CthreeQ_p)
     @>>> \CthreeK(\mathfrak X_k;\CthreeQ_p)\\
 @VVV @VVV\\
 \CthreeTC^\Cthreects(\mathfrak X;\CthreeQ_p)
     @>>> \CthreeTC(\mathfrak X_k;\CthreeQ_p)\\
 @VVV @VVV\\
 \CthreeHC^{-,\Cthreects}(\mathfrak X/\CthreecO;\CthreeQ_p)
     @>>> \CthreeHP^\Cthreects(\mathfrak X/\CthreecO;\CthreeQ_p)
\end{CD}
\end{equation}
in which both squares are cartesian.
\end{proposition}

The proof of \cite[Proposition~4.10]{AMMN} identifies the horizontal
fibers after projection from absolute Hochschild homology to
Hochschild homology relative to $\CthreecO$.
For an algebraization with bounded $p$-power torsion,
\cite[Proposition~4.2]{AMMN} identifies continuous Hochschild and cyclic
constructions with their $p$-completed counterparts.
The relative version is stated immediately before
\cite[Proposition~4.10]{AMMN}.
These assertions apply to the smooth, $W$-flat algebraizations in
Setup~\ref{c3:setup:coefficient-ring}.

\subsection{Integral vanishing}\label{c3:sec:integral}

In this section, we prove the regularity and integral vanishing
assertions of Theorem~\ref{thm:explicit-three-dimensional}.

\begin{proposition}\label{c3:prop:integral}
The rings $A_3$ and $\widehat A_3$ of
Theorem~\ref{thm:explicit-three-dimensional} are ramified regular local
domains of dimension three, with $\widehat A_3$ the maximal-ideal
completion of $A_3$. For either ring $R$, all entries defining $c_R$
are units, and
\[
 c_R|_{R[1/x]}=0\quad\text{in }\CthreeK_3(R[1/x]).
\]
\end{proposition}

\begin{proof}
The ambient local ring
$\CthreeZ_{(3)}[x,y,z]_{(3,x,y,z)}$ is regular of dimension four.
The image of $3+xy+z^2$ in its maximal ideal modulo its square is the
nonzero image of $3$.
Its quotient is therefore regular of dimension three, with maximal
ideal $(x,y,z)$, and is a domain.
Completion gives the displayed presentation of $\widehat A_3$.
The same parameter calculation in $\CthreeZ_3[[x,y,z]]$ proves that this
completion is regular of dimension three.
In both rings, $3=-xy-z^2$ belongs to the square of the maximal ideal,
so the rings are ramified.
Each entry in the definition of $c_R$ has nonzero image in $\CthreeF_3$.

We use Dennis--Stein symbols with the convention that
$\langle r,s\rangle$ is defined when $1-rs$ is invertible.
In additive notation their relations are
\begin{align*}
 \langle r,s\rangle&=-\langle s,r\rangle,\\
 \langle r,s\rangle+\langle r,t\rangle
   &=\langle r,s+t-rst\rangle,\\
 \langle r,st\rangle
   &=\langle rs,t\rangle+\langle tr,s\rangle
\end{align*}
whenever the symbols involved are defined
\cite[III, Definition~5.11, relations (D1)--(D3)]{Weibel}.
If $r$ is a unit, then
$\langle r,s\rangle=\{r,1-rs\}$.
Skew-symmetry gives
$\langle r,s\rangle=\{1-rs,s\}$ when $s$ is a unit.
In particular, $\langle r,1\rangle=0$ whenever it is defined.
These identities hold in actual $\CthreeK$-groups.

In $R$, put
\[
 a=\frac{1+z}{2},\qquad t=\frac y4,\qquad
 \lambda=1+x,\qquad \mu=1+\lambda t,\qquad q=a(1-a).
\]
The defining equation of $R$ gives
\[
 q=\frac{1-z^2}{4}=1+xt.
\]
The elements $a$, $1-a$, $\lambda$, $\mu$, $q$, and $\lambda q$
are units. The second Dennis--Stein relation gives
\[
 \delta\coloneqq \langle -x,t\rangle
   =\langle -x,t\rangle+\langle -x,1\rangle
   =\langle -x,t+1+xt\rangle
   =\langle -x,\mu\rangle
   =\{\lambda q,\mu\},
\]
where the last equality uses $1+x\mu=\lambda q$.

Graded commutativity gives $2\{a,a\}=0$, and the Steinberg relation
gives $\{a,1-a\}=0$. Hence
\[
 2\{a,q\}=2\{a,a\}+2\{a,1-a\}=0.
\]
Multiplying $\delta$ by $2[a]$ now yields
\[
 2[a]\delta
   =2\{a,\lambda,\mu\}+2\{a,q,\mu\}
   =c_R.
\]
After inverting $x$, the first-entry formula for the Dennis--Stein
symbol identifies $\delta$ with $\{-x,q\}$. Consequently,
\[
 c_R|_{R[1/x]}
   =2\{a,-x,q\}
   =-2\{a,q,-x\}
   =0.
\]
No completion or rationalization of a $\CthreeK$-group enters this calculation.
\end{proof}

\subsection{The relative character and its normalization}
\label{c3:sec:character}

The goal of this section is to construct the character used to detect
$c_R$ and to determine its value on a relative unit multiplied by two
absolute unit classes.

\begin{proposition}
\label{c3:prop:relative-character}
In Setup~\ref{c3:setup:coefficient-ring}, the continuous comparison defines
a natural map from the actual relative spectrum
\[
 \chi_D\colon
 \CthreeK(D,uD)\longrightarrow
 \Sigma\CthreeHC^\Cthreects(D/W;\CthreeQ_3).
\]
This map is $\CthreeK(D)$-linear, where the target action is obtained by
restriction along
\[
 \CthreeK(D)\longrightarrow\CthreeHC^{-,\Cthreects}(D/W;\CthreeQ_3).
\]
For $v\in1+uD$, let $[v]_{\Cthreerel}\in\CthreeK_1(D,uD)$ be the relative unit
class specified by the scalar automorphism $v$ and its identity
reduction modulo $u$. Then
\[
 \chi_D([v]_{\Cthreerel})=\log v
 \quad\text{in }\CthreeHC^\Cthreects_0(D/W;\CthreeQ_3)=\Omega_D^0.
\]
For $b_1,b_2\in D^\times$, its product character satisfies
\begin{equation}\label{c3:eq:unit-character}
 \chi_D\bigl([v]_{\Cthreerel}[b_1][b_2]\bigr)_{\Cthreetopcomp}
   =\pm\log(v)\,\Cthreedlog b_1\wedge\Cthreedlog b_2
   \quad\text{in }\Omega_D^2/\Cthreedd\Omega_D^1.
\end{equation}
Here the top component is the first summand of
\begin{equation}\label{c3:eq:cyclic-degree-two}
 \CthreeHC^\Cthreects_2(D/W;\CthreeQ_3)
   =
   \bigl(\Omega_D^2/\Cthreedd\Omega_D^1\bigr)
   \oplus
   \ker\bigl(\Cthreedd\colon\Omega_D^0\longrightarrow\Omega_D^1\bigr).
\end{equation}
The sign in \eqref{c3:eq:unit-character} is the universal suspension sign.
Equation~\eqref{c3:eq:cyclic-degree-two} uses the ordinary differential image.
\end{proposition}

\begin{proof}
The character is obtained from the horizontal fibers of
\eqref{c3:eq:continuous-comparison} for $\CthreeSpf D$.
Compatible maps from actual to continuous $\CthreeK$-theory, followed by
completion of spectra and rationalization, give a map from
$\CthreeK(D,uD)$ to the fiber of the top row.
The two squares carry this fiber to the fiber of
$\CthreeHC^{-,\Cthreects}\to\CthreeHP^\Cthreects$.
For a circle object $M$, the norm sequence is
\[
 \Sigma M_{hS^1}\longrightarrow M^{hS^1}
       \longrightarrow M^{tS^1};
\]
see \cite[\S I.3 and Corollary~I.4.3]{NikolausScholze}.
By \cite[Theorem~2.1, Lemma~2.2, and Theorem~2.3]{HoyoisCircle},
its cyclic interpretation identifies the last fiber with
\[\Sigma\CthreeHC^\Cthreects(D/W;\CthreeQ_3).\]
This defines a map out of actual relative $\CthreeK$-theory.

The required module action can be followed through the construction of
the comparison. Put $H=\CthreeTHH(D)$, regarded as a commutative algebra in
cyclotomic spectra.
For a bounded-below cyclotomic spectrum $H$, the construction in
\cite[Proposition~2.8]{AMMN} tensors the cyclotomic maps
$\CthreeZ^{\Cthreetriv}\to\CthreeTHH(\CthreeF_p)$ with $H$.
It gives a cartesian square with $\CthreeTC$ in the top row and $\CthreeTC^-$ in
the bottom row.
Its horizontal fiber is the $p$-completion of
$\Sigma(H\otimes\CthreeZ_{hC_p})_{hS^1}$.
The comparison in \cite[Corollary~2.10]{AMMN} replaces the lower-right
term by a Tate construction, using the equivalence induced by
$\CthreeZ^{\Cthreetriv}\to\CthreeTHH(\CthreeF_p)$, and then inverts $p$.
Specializing to $H=\CthreeTHH(D)$ and using
$\CthreeTHH(D)\otimes\CthreeZ\to\CthreeHH(D)$ gives the square with rows
$\CthreeTC(D)\to\CthreeTC(D\otimes_{\CthreebS}\CthreeF_p)$ and
$\CthreeHC^-(D)\to\CthreeHP(D)$ in \cite[Theorem~2.12]{AMMN}.
The reduction construction of \cite[Definition~2.14]{AMMN} replaces
derived reduction by ordinary reduction through their rational
$\CthreeTC$-equivalence.
Finally, \cite[Proposition~4.10]{AMMN} replaces reduction modulo $p$
by reduction modulo a uniformizer and projects to relative
Hochschild homology.

Every object $H\otimes T$ in the first two constructions is an
$H$-module through its first factor, and every map
$\Cthreeid_H\otimes f$ is $H$-linear.
The lax monoidal structure of $\CthreeTC$ therefore makes their $\CthreeTC$ images
modules over $\CthreeTC(H)=\CthreeTC(D)$.
The homotopy-fixed-point and Tate objects carry actions of $\CthreeTC^-(H)$,
and hence of $\CthreeTC(H)$ through its canonical map to $\CthreeTC^-(H)$.
Their norm and comparison maps preserve these actions.
This compatibility follows either from the lax monoidal constructions
in \cite[\S I.3 and Corollary~I.4.3]{NikolausScholze}, or by forming the norm
inside the category of modules over the circle-equivariant algebra.

In particular, the Tate equivalence inverted in
\cite[Corollary~2.10]{AMMN} is a module map.
Its inverse is taken in the module category.
The map
$\CthreeTC(H)\to\CthreeTC(H\otimes\CthreeZ^{\Cthreetriv})$, which becomes an equivalence after
rationalization in \cite[Theorem~2.12]{AMMN}, is a module map induced
by the unit of $\CthreeZ^{\Cthreetriv}$.
After forgetting to circle objects, the maps
\[
 H\otimes\CthreeZ\longrightarrow\CthreeHH(D)
       \longrightarrow\CthreeHH(D/W)
\]
are maps of $H$-algebras.
The reduction maps
$D\otimes_{\CthreebS}\CthreeF_3\to D/3\to D/u$ are maps of $D$-algebras.
Thus the rational $\CthreeTC$-equivalences used in
\cite[Definition~2.14 and Proposition~4.10]{AMMN} are
$\CthreeTC(D)$-linear as well.

Completion, rationalization, and passage to fibers preserve these
module constructions. The resulting action on the cyclic fiber is
restriction along
\[
 \CthreeTC(D)\longrightarrow\CthreeTC^-(D)
       \longrightarrow\CthreeHC^{-,\Cthreects}(D/W;\CthreeQ_3).
\]
The multiplicative cyclotomic trace, whose composite with $\CthreeTHH$ is
the topological Dennis trace
\cite[Theorems~7.3 and~7.4]{BGTTrace}, and the compatible
actual-to-continuous maps make $\chi_D$ a $\CthreeK(D)$-module map.
The suspension contributes the usual sign for a suspended module.
Only the negative cyclic action on cyclic homology is used here;
no ring structure on $\CthreeHC$ is asserted.

For the calculation of the target, the relative dimension bound gives
a finite mixed-complex model. On the algebraization $D_0$, the
normalized Hochschild map is
\begin{equation}\label{c3:eq:finite-hkr}
 b_0\otimes\cdots\otimes b_n
   \longmapsto
   \frac{1}{n!}b_0\,\Cthreedd b_1\wedge\cdots\wedge\Cthreedd b_n
   \qquad(0\leq n\leq2),
\end{equation}
and is zero in higher degrees.
The factorials occurring here are units in $W$.
The Hochschild differential maps to zero, while the Connes operator
maps to $\Cthreedd$; beyond degree two its target is zero.
Smooth Hochschild--Kostant--Rosenberg identifies the underlying
Hochschild homology with $\Omega^n_{D_0/W}$, vanishing for $n>2$.
This identification follows from the Koszul resolution of the smooth
diagonal and is realized by \eqref{c3:eq:finite-hkr}.
Hence \eqref{c3:eq:finite-hkr} is a quasi-isomorphism of mixed complexes
to $(\Omega^\ast_{D_0/W},0,\Cthreedd)$.
Only this finite mixed complex is needed, rather than an infinite
rational splitting.

The cyclic, negative cyclic, and periodic total complexes of this
model contain only finitely many form modules in each total degree.
Those modules are $W$-flat, and their reduction towers modulo $3^n$
are termwise surjective.
Derived $3$-completion is therefore given by termwise completion
of these total complexes.
The relative continuity assertion of
\cite[Proposition~4.2 and the paragraph preceding
Proposition~4.10]{AMMN} identifies them with the continuous
constructions. Inverting $3$ gives
$\CthreeHC^\Cthreects_0(D/W;\CthreeQ_3)=\Omega_D^0$ and
\eqref{c3:eq:cyclic-degree-two}.
Indeed, cyclic degree two has $\Omega_D^2$ in column zero and
$\Omega_D^0$ in column minus one.
The differential out of the second term is $\Cthreedd$, and boundaries
into the first term form $\Cthreedd\Omega_D^1$.
These are homology groups in the ordinary derived category, with
no closure taken on a differential image.

A unit $v\in1+uD$ determines a particular relative class: its scalar
automorphism is a loop whose reduction is the identity automorphism.
This construction is natural in $D$ and maps to $[v]\in\CthreeK_1(D)$.
To determine its character, first take $D=W\langle T\rangle$ and
$v=1+uT$, and write
\[
 f=\chi_D([1+uT]_{\Cthreerel})\in E\langle T\rangle.
\]
Choose the circle orientation for which the norm is Connes' operator,
as in \cite[Lemma~2.2 and Theorem~2.3]{HoyoisCircle}.
The norm followed by projection to $\CthreeHH_1$ then sends $f$ to $\Cthreedd f$.
On the $\CthreeK$-theory side, this composite is the Dennis trace.
The scalar-unit calculation for that trace represents the class of
$b$ by $b^{-1}\otimes b$ in $\CthreeHH_1$, which
\eqref{c3:eq:finite-hkr} sends to $\Cthreedlog b$.
It follows that
\[
 \Cthreedd f=\Cthreedlog(1+uT).
\]
Evaluation at $T=0$ sends the specified relative unit to zero, so
$f(0)=0$.
The kernel of differentiation on $E\langle T\rangle$ consists of
constants: for each coefficient $c_n$ with $n\geq1$, the equation
$n c_n=0$ forces $c_n=0$.
The convergent restricted series
\[
 \log(1+uT)
   =\sum_{n\geq1}\frac{(-1)^{n+1}u^nT^n}{n}
\]
has the required derivative and constant term.
Thus $f=\log(1+uT)$.
For general $v\in1+uD$, substitution
$T=(v-1)/u\in D$ defines a continuous $W$-algebra map from
$W\langle T\rangle$ to $D$.
Naturality gives $\chi_D([v]_{\Cthreerel})=\log v$.

It remains to specify the coefficient in the product formula.
Over any commutative base ring, negative cyclic homology acts on
cyclic homology by ordinary shuffle plus the cyclic variable times
cyclic shuffle
\cite[\S2, Definition~2.1 and Remark~2.2]{BokstedtOttosen}.
The cyclic, negative cyclic, and periodic long exact sequences are
module sequences, with the usual sign on connecting homomorphisms
\cite[Proposition~2.3]{BokstedtOttosen}.
These statements precede, and do not require, the characteristic-two
hypotheses imposed in later sections of that reference.

Write the cyclic complex, with $h$ of degree $-2$, as
\[
 C^+=\bigl(E[h,h^{-1}]/hE[h]\bigr)\otimes_E C_\ast,
\]
where $C_\ast$ is the rationalized Hochschild chain complex.
The function $\log v$ is represented in column zero.
A negative cyclic representative of a unit has leading Hochschild
class $b^{-1}\otimes b$; its remaining terms contain positive powers
of $h$.
Multiplication on column zero kills all those terms, including the
cyclic-shuffle correction $h\,\mathrm{sh}'$.
The first multiplication remains in column zero, so the second
multiplication has the same property.
The degree-two shuffle has two terms.
After antisymmetrization both give the same wedge, canceling the
factor $2!$ in \eqref{c3:eq:finite-hkr}.
Changing a leading Hochschild representative by a boundary does not
change its resulting cyclic class.
The coefficient in \eqref{c3:eq:unit-character} is therefore one,
apart from the suspension sign already specified.
\end{proof}

\subsection{An ordinary residue criterion}\label{c3:sec:residue}

In this section, we detect differential classes on the formal quadric
of Setup~\ref{c3:setup:coefficient-ring} by a residue on the overlap of two
affine-plane charts.

\begin{lemma}
\label{c3:lem:ordinary-residue}
For $F\in E\langle Z\rangle$, the form initially defined where $X$ is
invertible by
\[
 \omega_F=\frac{\Cthreedd X}{X}\wedge\Cthreedd F
\]
extends to a global element of $\Omega_B^2$.
If $\omega_F$ belongs to the ordinary image $\Cthreedd\Omega_B^1$, then
$F(1)=F(-1)$.
\end{lemma}

\begin{proof}
Differentiating the equation of the quadric gives
$X\Cthreedd Y+Y\Cthreedd X+2Z\Cthreedd Z=0$.
On the locus where $Z$ is invertible,
\[
 \frac{\Cthreedd X}{X}\wedge\Cthreedd Z
   =-\frac{\Cthreedd X\wedge\Cthreedd Y}{2Z}.
\]
The loci where $X$ or $Z$ is invertible cover the quadric, since
$X=Z=0$ is incompatible with $XY+Z^2=1$.
This identity extends $\omega_F$ across its potential pole.

Before completion, let $U_+$ be the complement of the line
$X=0$, $Z=-1$, and let $U_-$ be the complement of the line
$X=0$, $Z=1$.
Put $a_+=1$ and $a_-=-1$.
Each $U_i$ is an affine plane with coordinates $X,t_i$, under
\begin{equation}\label{c3:eq:residue-charts}
 Z=a_i+Xt_i,\qquad
 Y=-t_i(2a_i+Xt_i).
\end{equation}
To see this, write $U_i=D(X)\cup D(Z+a_i)$.
The inverse coordinate is
$(Z-a_i)/X$ on $D(X)$ and $-Y/(Z+a_i)$ on $D(Z+a_i)$.
The equation $(Z-a_i)(Z+a_i)=-XY$ identifies these functions on the
overlap.
Their preimages cover the $(X,t_i)$-plane, since $2a_i$ is a unit.
Thus \eqref{c3:eq:residue-charts} is an isomorphism.
The intersection $U_+\cap U_-$ is $D(X)$.
After completion these charts are formal bidiscs, with overlap
$W\langle X,X^{-1},Z\rangle$.

On the $i$-th bidisc, define
\[
 q_i=-\bigl(F(Z)-F(a_i)\bigr)\frac{\Cthreedd X}{X}.
\]
Substitution from \eqref{c3:eq:residue-charts} makes
$F(Z)-F(a_i)$ divisible by $X$ in the restricted series algebra, so
$q_i$ is regular.
Differentiation gives $\Cthreedd q_i=\omega_F$.
On the overlap,
\begin{equation}\label{c3:eq:primitive-difference}
 q_+-q_-=\bigl(F(1)-F(-1)\bigr)\frac{\Cthreedd X}{X}.
\end{equation}

We use a coefficient calculation for closed one-forms on a bidisc.
Let
\[
 A(X,t)\Cthreedd X+D(X,t)\Cthreedd t
\]
be closed over $E\langle X,t\rangle$, and write
$A=\sum_{m,n\geq0}a_{m,n}X^mt^n$ and
$D=\sum_{m,n\geq0}b_{m,n}X^mt^n$.
The coefficients tend to zero.
For $c\in W$, substitution $t=c/X$ is defined on
$E\langle X,X^{-1}\rangle$.
The coefficient of $X^{-1}\Cthreedd X$ in the resulting form is the
convergent sum
\[
 \sum_{m\geq0}
       \bigl(a_{m,m+1}-b_{m+1,m}\bigr)c^{m+1}.
\]
Closedness means $\partial_t A=\partial_X D$, hence
\[
 (m+1)a_{m,m+1}=(m+1)b_{m+1,m}
 \qquad(m\geq0).
\]
Since $E$ has characteristic zero, each summand in the residue is
zero. This calculation does not divide an infinite series to
construct a primitive.

Assume now that $\omega_F=\Cthreedd q$ for a global
$q\in\Omega_B^1$.
The forms $q_i-q$ are closed on the two bidiscs.
Restrict them to the common slice $Z=0$.
In chart $i$, this restriction is the substitution
$t_i=-a_i/X$, with $-a_i\in W$.
Both restricted forms have residue zero by the coefficient
calculation.
Their difference, by \eqref{c3:eq:primitive-difference}, has residue
$F(1)-F(-1)$.
Therefore $F(1)=F(-1)$.
The argument concerns exactness in the ordinary quotient
$\Omega_B^2/\Cthreedd\Omega_B^1$.
\end{proof}

\subsection{Nonvanishing of the explicit class}\label{c3:sec:nonvanishing}

The goal of this section is to complete
Theorem~\ref{thm:explicit-three-dimensional} by detecting an actual
relative class on $B$ and passing to its absolute image.

\begin{proof}[Proof of Theorem~\ref{thm:explicit-three-dimensional}]
Proposition~\ref{c3:prop:integral} proves regularity and integral
vanishing.
It remains to prove that $c_{A_3}$ and $c_{\widehat A_3}$ have infinite
order.

The substitutions
\[
 x=uX,\qquad y=uY,\qquad z=uZ
\]
define a continuous homomorphism $\widehat A_3\to B$.
Indeed, every formal series in $x,y,z$ converges $u$-adically after
substitution, and the relation maps to
$3+u^2(XY+Z^2)=0$.
In $B$, put
\[
 a=\frac{1+uZ}{2},\qquad
 \lambda=1+uX,\qquad
 \mu=1+\frac{(1+uX)uY}{4},\qquad
 q=a(1-a)=\frac{1+3Z^2}{4}.
\]
These elements are units, and $a^2$ reduces to $1$ modulo $u$.
Consequently, the specified relative unit of
Proposition~\ref{c3:prop:relative-character} gives an actual class
\begin{equation}\label{c3:eq:relative-class}
 \gamma=[a^2]_{\Cthreerel}[\lambda][\mu]\in\CthreeK_3(B,uB).
\end{equation}
Its image in $\CthreeK_3(B)$ is
$[a^2][\lambda][\mu]=2[a][\lambda][\mu]=c_B$,
where $c_B$ is the image of either source class.

The map from relative to absolute $\CthreeK$-theory is injective in this
degree. To prove this using actual groups, put
\[
 C=B/uB=\CthreeF_3[X,Y,Z]/(XY+Z^2-1).
\]
A second pair of affine-plane charts describes $\CthreeSpec C$ as an
affine-line torsor over $\CthreePP^1_{\CthreeF_3}$.
The first chart has coordinates $X,t$ and equations
\[
 Z=1+Xt,\qquad Y=-t(2+Xt).
\]
It is the complement of the line $X=0$, $Z=-1$.
The second has coordinates $Y,s$ and equations
\[
 Z=-1-sY,\qquad X=-s(2+sY).
\]
It is the complement of the line $Y=0$, $Z=1$.
These charts are affine planes by the inverse-coordinate calculation
used in \eqref{c3:eq:residue-charts}.
Their omitted lines are disjoint, so the charts cover $\CthreeSpec C$.
On their intersection,
\begin{equation}\label{c3:eq:affine-line-transition}
 s=t^{-1},\qquad Y=-t^2X-2t.
\end{equation}
Thus projection to $t$ on the first chart and to $s$ on the second
glues to a morphism $\CthreeSpec C\to\CthreePP^1_{\CthreeF_3}$.
It is locally the affine-line projection, with affine-linear
transition of invertible slope $-t^2$.
In particular, it is a torsor under a line bundle.

Homotopy invariance on the two charts and their overlap, followed by
the two-open Mayer--Vietoris sequence, shows that pullback gives an
equivalence on $\CthreeK$-theory
\cite[\S7, Remark~3.5 and Proposition~4.1]{Quillen}.
The projective bundle formula
\cite[\S8, Theorem~2.1]{Quillen} identifies
$\CthreeK(\CthreePP^1_{\CthreeF_3})$ with $\CthreeK(\CthreeF_3)\oplus\CthreeK(\CthreeF_3)$.
For a finite field, the actual-group computation is
\[
 \CthreeK_{2i}(\CthreeF_q)=0,\qquad
 \CthreeK_{2i-1}(\CthreeF_q)=\CthreeZ/(q^i-1)
 \qquad(i\geq1);
\]
see \cite{QuillenFinite} and
\cite[IV, Theorem~1.12 and Corollary~1.13]{Weibel}.
Only its consequence $\CthreeK_4(\CthreeF_3)=0$ is needed here.
It follows that
\begin{equation}\label{c3:eq:actual-k-four}
 \CthreeK_4(C)=\CthreeK_4(\CthreeF_3)\oplus\CthreeK_4(\CthreeF_3)=0.
\end{equation}
The long exact sequence of the actual relative spectrum contains
\[
 \CthreeK_4(C)\longrightarrow\CthreeK_3(B,uB)\longrightarrow\CthreeK_3(B).
\]
Hence \eqref{c3:eq:actual-k-four} makes
$\CthreeK_3(B,uB)\to\CthreeK_3(B)$ injective.
This argument uses the actual fourth $\CthreeK$-group.

We detect $\gamma$ by its character.
Define the $3$-adic logarithm of $2$ by
\[
 \log2=\frac12\log4=\frac12\log(1+3),
 \qquad
 L(Z)=-\log2+\log(1+uZ).
\]
The logarithm series are additive on $1+uB$ by the formal power
series identity and their $u$-adic convergence.
Since $4\in1+uB$ and
$a^2=(1+uZ)^2/4$, this gives
$\log(a^2)=2L$.
This use of the logarithm remains valid at ramification index two.
Proposition~\ref{c3:prop:relative-character} identifies the top
component of $\chi_B(\gamma)$ with $\pm\tau$, where
\begin{equation}\label{c3:eq:character-form}
 \tau=2L\,\Cthreedlog\lambda\wedge\Cthreedlog\mu.
\end{equation}

For the differential calculation, put $t=uY/4$.
Then $q=1+uXt$, $\mu=1+\lambda t$, and
$\lambda q=1+uX\mu$.
Differentiating the last identity and wedging with $\Cthreedlog\mu$ gives,
where $X$ is invertible,
\begin{equation}\label{c3:eq:differential-identity}
 \Cthreedlog\lambda\wedge\Cthreedlog\mu
   =\frac{\Cthreedd X}{X}\wedge\Cthreedlog q
      -\Cthreedlog q\wedge\Cthreedlog\mu.
\end{equation}
More explicitly, the left-hand side plus
$\Cthreedlog q\wedge\Cthreedlog\mu$ is
\[
 \frac{u\,\Cthreedd X\wedge\Cthreedd\mu}{\lambda q}
   =\frac{u\,\Cthreedd X\wedge\Cthreedd t}{q}
   =\frac{\Cthreedd X}{X}\wedge\frac{\Cthreedd q}{q}.
\]
The form involving $\Cthreedd X/X$ is global by the extension calculation
in Lemma~\ref{c3:lem:ordinary-residue}, since $q$ depends only on $Z$.
Thus \eqref{c3:eq:differential-identity} is an identity of global regular
two-forms.

Choose $F\in E\langle Z\rangle$ satisfying
\begin{equation}\label{c3:eq:primitive-derivative}
 F'(Z)=\frac{6ZL(Z)}{1+3Z^2},
 \qquad\text{so that}\qquad
 \Cthreedd F=L\,\Cthreedlog q.
\end{equation}
Such an $F$ exists in the restricted series algebra.
With the absolute value normalized by $|3|=1/3$, both $L$ and
$(1+3Z^2)^{-1}$ converge for $|Z|<\sqrt3$.
Termwise integration divides each coefficient by a positive integer,
whose $3$-adic valuation grows at most logarithmically.
The integrated series therefore still has radius of convergence at
least $\sqrt3$, which is greater than one.
Equations \eqref{c3:eq:character-form} and
\eqref{c3:eq:differential-identity} now give the global identity
\begin{equation}\label{c3:eq:character-residue}
 \tau=
 2\frac{\Cthreedd X}{X}\wedge\Cthreedd F
       -2\Cthreedd\bigl(F\,\Cthreedlog\mu\bigr).
\end{equation}
By Lemma~\ref{c3:lem:ordinary-residue}, a nonzero difference
$F(1)-F(-1)$ makes the class of $\tau$ nonzero in the ordinary
quotient $\Omega_B^2/\Cthreedd\Omega_B^1$.

Only the odd part of $L$ contributes to this endpoint difference.
Since $u^2=-3$, that odd part is
\[
 u\sum_{k\geq0}\frac{(-3)^k Z^{2k+1}}{2k+1}.
\]
Multiply this series by $6Z/(1+3Z^2)$, integrate, and subtract the
values at $1$ and $-1$.
All series converge on a disc of radius greater than one, which
justifies termwise operations and regrouping by $j=k+l$.
The resulting formula is
\begin{equation}\label{c3:eq:endpoint-series}
 \frac{F(1)-F(-1)}{u}
   =12\sum_{j\geq0}\frac{(-3)^jS_j}{2j+3},
 \qquad
 S_j=\sum_{k=0}^j\frac1{2k+1}.
\end{equation}
Its $j=0$ term is $4$.

For completeness, an estimate for all indices controls the tail.
Normalize $v_3(3)=1$.
Every denominator in $S_j$ is at most $2j+1$, so
\[
 v_3(S_j)\geq-\lfloor\log_3(2j+1)\rfloor.
\]
For $j\geq1$, the valuation of the $j$-th summand in
\eqref{c3:eq:endpoint-series} is therefore at least
\[
 A_j=1+j-\lfloor\log_3(2j+1)\rfloor-v_3(2j+3).
\]
Direct substitution gives $A_1=1$, $A_2=2$, and $A_3=1$.
For $j\geq4$, we have
\begin{equation}\label{c3:eq:logarithmic-bound}
 \lfloor\log_3(2j+3)\rfloor\leq\lfloor j/2\rfloor.
\end{equation}
Equivalently, $2j+3<3^{\lfloor j/2\rfloor+1}$.
This holds for $j=4$ and $j=5$.
Advancing $j$ by two preserves the inequality because
$2j+7<3(2j+3)$.
Both $\lfloor\log_3(2j+1)\rfloor$ and $v_3(2j+3)$ are bounded by
$\lfloor\log_3(2j+3)\rfloor$.
Thus \eqref{c3:eq:logarithmic-bound} gives
\[
 A_j\geq1+j-2\lfloor j/2\rfloor\geq1.
\]
Moreover,
\[
 A_j\geq1+j-2\log_3(2j+3)\longrightarrow+\infty.
\]
The entire tail of \eqref{c3:eq:endpoint-series} consequently converges
in $3\CthreeZ_3$, and
\begin{equation}\label{c3:eq:endpoint-nonzero}
 \frac{F(1)-F(-1)}{u}\in4+3\CthreeZ_3.
\end{equation}
In particular, the endpoint difference is nonzero.

Equations \eqref{c3:eq:character-residue} and
\eqref{c3:eq:endpoint-nonzero}, together with
Lemma~\ref{c3:lem:ordinary-residue}, show that
$\chi_B(\gamma)$ has nonzero top component.
Its target is a $\CthreeQ_3$-vector space, so a torsion element cannot have
this image.
Hence $\gamma$ has infinite order in the actual group
$\CthreeK_3(B,uB)$.
The injection proved using \eqref{c3:eq:actual-k-four} shows that $c_B$
has infinite order in $\CthreeK_3(B)$.
As $c_B$ is the image of both $c_{\widehat A_3}$ and $c_{A_3}$, both
source classes have infinite order.

For either source ring $R$, Proposition~\ref{c3:prop:integral} gives
integral vanishing in $\CthreeK_3(R[1/x])$.
Since $R$ is a domain, $R[1/x]$ embeds in $\CthreeFrac(R)$, so the class
also vanishes in $\CthreeK_3(\CthreeFrac(R))$.
An element of an abelian group becomes zero after tensoring with
$\CthreeQ$ exactly when it is torsion.
Thus $c_R$ remains nonzero in $\CthreeK_3(R)\otimes_{\CthreeZ}\CthreeQ$, proving the
asserted failure of injectivity.
\end{proof}

\begin{remark}[Ordinary differential homology]
\label{c3:rem:ordinary-homology}
Every differential quotient used above is ordinary algebraic homology
after derived completion and rationalization.
No differential image is replaced by its closure.
The argument makes no assertion about the order of the class in an
individual finite Artin quotient, and uses no identification
\[
 \CthreeK_i(D)\otimes_{\CthreeZ}\CthreeQ_3
   =\pi_i\bigl(\CthreeK(D)^\wedge_3[1/3]\bigr).
\]
It also requires no injection from actual $\CthreeK$-theory into completed
$\CthreeK$-theory.

The distinction between an ordinary image and its closure is visible
on this quadric.
The global form
\[
 \omega_0=\frac{\Cthreedd X}{X}\wedge\Cthreedd Z
\]
is nonexact by Lemma~\ref{c3:lem:ordinary-residue}, applied to $F(Z)=Z$,
whose endpoint difference is $2$.
For every integer $j\geq0$, however, the one-form
\[
 \beta_j=-\bigl(Z-Z^{3^j}\bigr)\frac{\Cthreedd X}{X}
\]
is global and regular.
Indeed, $Z-Z^{3^j}$ vanishes at both $1$ and $-1$, and is divisible
by $Z^2-1=-XY$.
Its derivative is
\[
 \Cthreedd\beta_j
    =\omega_0-3^jZ^{3^j-1}\omega_0
    \longrightarrow\omega_0.
\]
Thus $\omega_0$ lies in the closure of the differential image while
remaining outside that image.
Passing to the Hausdorff quotient would discard the obstruction used
in the proof.
\end{remark}

% End input: sections/02-explicit-three-adic.tex


% Begin input: sections/03-quartic-core.tex
\section{The quartic construction}\label{d2:sec:introduction}

\subsection*{Supported classes and generic restriction}

A class pushed forward from a Cartier divisor vanishes after restriction to the complement of that divisor. For a regular local ring, the question is whether such a class can retain infinite order in algebraic K-theory. The localization sequence and the projection formula express this question in terms of supported classes and their products; we use these results in the forms of \cite[V.3--V.6]{Weibel}. Our examples allow two different regulators to detect one such product. The first uses the analytic regulator and its comparison with the regulator of a number ring \cite{Borel,BunkeTamme}. The second uses the relative fiber comparison with cyclic homology \cite{AMMN,BGTTrace}.

The ring is defined by a quartic equation whose constant term is the residue characteristic. Its cyclotomic divisor splits into four components after a flat coefficient extension. The coefficients of these components are conjugates of one actual class in degree three. The resulting divisor calculation is the common starting point of both proofs.

\subsection{A rational generic kernel}

The geometric argument detects the class on every finite open that occurs in the localization defining the local ring.

\begin{theorem}\label{d2:thm:geometric}
Let
\[
 A=
 \left(
 \DtwoZZ_{(5)}[T,x,y]/
 \bigl(\Phi _5(1+y)+x^4+Txy^3\bigr)
 \right)_{(5,x,y)},
 \qquad
 \Phi _5(z)=z^4+z^3+z^2+z+1.
\]
Then \(A\) is a regular local ring of dimension two with residue field
\(\DtwoFF _5(T)\). There is an actual infinite-order class
\(c\in K_3(A)\) whose image in \(K_3(\DtwoFrac A)\) is zero.
Here \(K_3\) denotes Quillen algebraic K-theory.
\end{theorem}

The class is the pushforward of a cyclotomic coefficient from \(x=0\).
After the divisor splits, its analytic regulator is a linear combination
of four first Chern classes. A blowup calculation shows that the only
relation among these classes is the diagonal relation. It also controls
every geometric component of every divisor removed by an allowed
localization.

\subsection{Persistence after completion}

The completed argument requires a coefficient with prescribed nonzero
components at both odd characters.

\begin{theorem}\label{d2:thm:completed}
Let \(A\) be the ring in Theorem~\ref{d2:thm:geometric}, and let
\(\Dtwowh A\) be its \((x,y)\)-adic completion. Both rings are regular local
domains of dimension two, with maximal ideal \((x,y)\) and residue field
\(\DtwoFF _5(T)\). There is an actual class \(c\in K_3(A)\) such that its image
\(\Dtwowh c\in K_3(\Dtwowh A)\) has infinite order and
\[
 c|_{\DtwoFrac A}=0,
 \qquad
 \Dtwowh c|_{\DtwoFrac\Dtwowh A}=0
\]
integrally. Consequently both maps
\[
 K_3(A)\otimes\DtwoQQ\longrightarrow K_3(\DtwoFrac A)\otimes\DtwoQQ,
 \qquad
 K_3(\Dtwowh A)\otimes\DtwoQQ\longrightarrow
 K_3(\DtwoFrac\Dtwowh A)\otimes\DtwoQQ
\]
have nonzero kernel.
\end{theorem}

The class in Theorem~\ref{d2:thm:completed} can also be used in the geometric
proof of Theorem~\ref{d2:thm:geometric}. We establish this identification
after constructing its arithmetic coefficient. In particular, the
passage to completion does not require choosing a second supported
class.

\subsection*{Proof architecture}

The completed detector is a continuous two-form relative to a ramified
coefficient ring. Its parameter differential is retained. A tame trace
to a constant ordinary elliptic curve, followed by a bounded contraction,
takes the weighted root-divisor class to
\(4e_3(b)\) times a nonzero one-form class on a generic Cohen completion.
The exact linear identity is proved in
Proposition~\ref{d2:prop:weighted}; it holds for every anti-invariant
coefficient \(b\). Nonvanishing of the one-form rests on scalar
Frobenius pole removal in Lemma~\ref{d2:lem:scalar} and on the
rigid-to-convergent results of \cite{CaroDAddezio2025}.

Proposition~\ref{d2:prop:cyclic} identifies the relevant relative K-theory
fiber with ordinary cyclic homology after spectral completion and
inversion of \(5\). Its product calculation detects an actual relative
class. The smooth-surface finite-coefficient bound of
\cite[Theorem~8.4]{GeisserLevineCharP} then applies twice: first to the
fiber of reduction, and then to the support term in localization.
These are distinct fiber sequences.

Section~\ref{d2:sec:preliminaries} fixes the K-theory conventions and the
supported construction. Sections~\ref{d2:sec:geometry}
and~\ref{d2:sec:analytic} prove the geometric route.
Section~\ref{d2:sec:cyclic} establishes the relative cyclic comparison.
Section~\ref{d2:sec:cohen} develops the Cohen and elliptic arguments.
Section~\ref{d2:sec:quartic} computes the weighted quartic detector.
Section~\ref{d2:sec:arithmetic} constructs its actual arithmetic coefficient.
Section~\ref{d2:sec:assembly} proves Theorem~\ref{d2:thm:completed} and identifies
the coefficient used in both routes.

\section{Preliminaries}\label{d2:sec:preliminaries}

We use algebraic K-theory, transfers, localization, and finite
coefficients with the conventions of
\cite[Chapters~IV and~V]{Weibel}.

\begin{notation}\label{d2:notation:completion}
For a ring \(H\) and an ideal \(I\), put
\[
 K(H,I)=\Dtwofib\bigl(K(H)\longrightarrow K(H/I)\bigr),
 \qquad
 \DtwocK(H)=K(H)^\wedge_5[1/5],
 \qquad
 \DtwocK(H,I)=K(H,I)^\wedge_5[1/5].
\]
The completion is the derived inverse limit of the spectra
\(K(H)/5^n\), and likewise for a relative spectrum.
It is not defined by tensoring the actual K-groups with \(\DtwoZZ _5\).
Completion and inversion of \(5\) preserve fiber sequences.

On adically completed rings, differential forms are continuous forms.
Every de Rham quotient in this article uses the ordinary image of the
differential, unless a closure is explicitly introduced in the
construction of a bounded functional.
\end{notation}

For regular noetherian rings of finite dimension, projective-module
K-theory agrees with coherent-module G-theory. Finite closed immersions
have G-theory pushforwards by restriction of scalars. These pushforwards commute with flat base change and satisfy the projection formula
\[
 i_*(i^*\alpha\cdot z)=\alpha\cdot i_*z.
\]
For an effective Cartier divisor with invertible ideal \(I\), its
two-term resolution gives \(i_*[\DtwocO]=1-[I]\).
We use \cite[V.3.3, V.3.3.2, V.3.7.2, V.3.12]{Weibel}
for these statements. In each application, the immersions are finite
between affine schemes and the base change in question is flat.
Thus the required properness and Tor-independence hold.

For a regular noetherian ring \(H\), a nonzero divisor \(s\), and a
regular quotient \(H/sH\), localization and d\'evissage give
\[
 K(H/sH)\longrightarrow K(H)\longrightarrow K(H[1/s]),
\]
where the first map is support pushforward
\cite[V.4.1, V.6.1]{Weibel}.
K-groups commute with filtered colimits of rings
\cite[IV, Section~6.4]{Weibel}. The same holds with finite coefficients,
by the cofiber construction. In particular, an equality in a filtered
localization holds after some finite further localization.

\begin{setup}\label{d2:setup:ring}
Write
\[
 \begin{split}
 P(y)&=\Phi _5(1+y)=5+10y+10y^2+5y^3+y^4,\\
 F&=P(y)+x^4+Txy^3,\\
 R_0&=\DtwoZZ _{(5)}[T,x,y]/(F),\qquad
 \Dtwofm=(5,x,y),\qquad
 S=R_0\setminus\Dtwofm,\qquad A=S^{-1}R_0.
 \end{split}
\]
Fix a primitive fifth root \(\rho\), and put
\[
 E_{\Dtwoar}=\DtwoQQ(\rho),\qquad
 W_{\Dtwoar}=\DtwoZZ _{(5)}[\rho],\qquad
 \pi=\rho-1,\qquad
 G=(\DtwoZZ/5\DtwoZZ)^\times.
\]
For \(g\in G\), let \(\sigma_g(\rho)=\rho^g\).
\end{setup}

The two coefficients \(xy^3\) and \(P(y)+x^4\) of the linear polynomial
\(F\) in \(T\) have no common irreducible factor in
\(\DtwoZZ _{(5)}[x,y]\). Neither \(x\) nor \(y\) divides \(P(y)+x^4\).
Gauss's lemma therefore shows that \(R_0\) is a domain.
Solving for \(T\) gives \(\DtwoFrac A=\DtwoQQ(x,y)\).

The ambient local ring
\(\DtwoZZ _{(5)}[T,x,y]_{(5,x,y)}\) is regular of dimension three,
with residue field \(k_0=\DtwoFF _5(T)\).
The class of \(F\) in its cotangent space is the nonzero class of \(5\).
Thus \(A\) is regular local of dimension two. The identity
\begin{equation}\label{d2:eq:ramification}
 5(1+2y+2y^2+y^3)=-(x^4+y^4+Txy^3)
\end{equation}
shows that its maximal ideal is \((x,y)\).
Completion preserves the dimension and cotangent space of a noetherian
local ring, so \(\Dtwowh A\) is also regular local of dimension two, with the
same residue field. Both rings are domains.

The polynomial \(P\) is Eisenstein at \(5\).
Consequently \(W_{\Dtwoar}\) is finite free of rank four over
\(\DtwoZZ _{(5)}\). Its reduction modulo \(5\) is
\(\DtwoFF _5[\pi]/(\pi^4)\), so it is local.
Moreover,
\begin{equation}\label{d2:eq:epsilon}
 5=\pi^4\epsilon,\qquad
 \epsilon=-\frac{1}{1+2\pi+2\pi^2+\pi^3},
 \qquad
 \epsilon\bmod\pi=-1.
\end{equation}
Thus \(W_{\Dtwoar}\) is a discrete valuation ring with uniformizer \(\pi\)
and residue field \(\DtwoFF _5\). Taking the quotient by \(x\) gives
\begin{equation}\label{d2:eq:cartier-quotient}
 A/xA=W_{\Dtwoar}[T]_{(\pi)},
 \qquad y\longmapsto\pi.
\end{equation}
This is a regular discrete valuation ring. Its completion
\(\Dtwowh A/x\Dtwowh A\) is regular as well.

\begin{definition}\label{d2:defn:supported}
Let \(i\colon\DtwoSpec(A/xA)\to\DtwoSpec A\) be the closed immersion.
For an actual class \(\beta\in K_3(W_{\Dtwoar})\), define
\[
 c(\beta)=i_*\bigl(\beta|_{A/xA}\bigr)\in K_3(A).
\]
\end{definition}

Restriction to \(A[1/x]\) kills \(c(\beta)\), because the modules in its
pushforward are supported on \(x=0\). Hence its generic image is zero
integrally. The map \(A\to\Dtwowh A\) is flat.
Flat base change identifies its completed image with the same
pushforward on \(\DtwoSpec\Dtwowh A\), so that image also vanishes over
\(\DtwoFrac\Dtwowh A\). The remaining issue in both theorems is infinite order.

\section{The divisor lattice and its finite opens}\label{d2:sec:geometry}

In this section, we prove the Picard and Betti persistence used in
Theorem~\ref{d2:thm:geometric}.

\begin{proposition}\label{d2:prop:geometry}
In Setup~\ref{d2:setup:ring}, put
\[
 B_{\Dtwoar}=A\otimes_{\DtwoZZ _{(5)}}W_{\Dtwoar},
 \qquad
 \DtwocT=\DtwoSpec B_{\Dtwoar}[1/\pi,1/y],
\]
and
\[
 \DtwocQ=\DtwoSpec E_{\Dtwoar}[T,x,y,y^{-1}]/(F),
 \qquad
 D_g=V(x,y-(\rho^g-1))\subset\DtwocQ.
\]
For \(s\in S\), write \(\DtwocQ_s=D(s)\subset\DtwocQ\).
Then \(\DtwocQ\) is a smooth geometrically integral surface,
\(\DtwocT=\DtwoSpec S^{-1}\DtwocO(\DtwocQ)\), and
\[
 \DtwoPic(\DtwocQ)=\DtwoZZ ^4/\DtwoZZ(1,1,1,1),
\]
where the standard basis maps to \(\DtwocO(D_g)\).
This calculation persists under every extension of \(E_{\Dtwoar}\), and
\(\DtwoPic(\DtwocQ)\to\DtwoPic(\DtwocT)\) is injective.

For every embedding \(E_{\Dtwoar}\to\DtwoCC\), the classes
\[
 e_g=c_1^{\Dtwotopol}(\DtwocO(D_g))
       \in H^2(\DtwocQ(\DtwoCC),\DtwoRR)
\]
have exactly the real diagonal relation. For every \(s\in S\),
restriction induces an injection
\[
 H^2(\DtwocQ(\DtwoCC),\DtwoRR)\longrightarrow H^2(\DtwocQ_s(\DtwoCC),\DtwoRR).
\]
\end{proposition}

\begin{proof}
The ring \(B_{\Dtwoar}\) is finite free over \(A\).
Its closed fiber is \(k_0[\pi]/(\pi^4)\), so it has the single maximal
ideal \(\Dtwofn=(\pi,x,y)\).
Tensoring \(A\hookrightarrow\DtwoQQ(x,y)\) with the free
\(\DtwoZZ _{(5)}\)-module \(W_{\Dtwoar}\) embeds \(B_{\Dtwoar}\) in
\(E_{\Dtwoar}(x,y)\). Thus \(B_{\Dtwoar}\) is a domain.
With \(V_{\Dtwoar}=W_{\Dtwoar}[T]_{(\pi)}\), comparison of localizations gives
\[
 B_{\Dtwoar}=
 \bigl(V_{\Dtwoar}[x,y]/(F)\bigr)_{(\pi,x,y)}.
\]
Here every polynomial whose reduction is nonzero modulo \(\pi\)
becomes a unit. Inverting \(\pi\) and \(y\) gives
\(S^{-1}E_{\Dtwoar}[T,x,y,y^{-1}]/(F)\), as asserted.

Consider the blowup
\(\DtwocY=\operatorname{Bl}_{\Dtwofn}(\DtwoSpec B_{\Dtwoar})\).
Its tangent coordinates \((U:X:Y)\) correspond to \((\pi,x,y)\).
The associated graded ring is
\[
 k_0[U,X,Y]/(X^4+Y^4+TXY^3-U^4).
\]
Indeed \(F\) has order four in the ambient regular local ring, and its
nonzero initial form generates the initial ideal of the principal
ideal \((F)\). The exceptional scheme is the plane curve
\[
 C_{\mathrm{exc}}\colon X^4+Y^4+TXY^3=U^4.
\]

A geometric singular point of this curve would satisfy \(U=0\).
It cannot have \(Y=0\), since then \(X=0\).
Putting \(s=X/Y\), the remaining equations are
\[
 s^4+Ts+1=0,\qquad 4s^3+T=0.
\]
In characteristic five these imply \(T=s^3\), \(s^4=2\), and
\(T^4=3\), contrary to the transcendence of \(T\).
Thus \(C_{\mathrm{exc}}\) is geometrically smooth.
Distinct geometric plane components would intersect by B\'ezout and
give a singular point; repeated components would also be singular.
Therefore the curve is geometrically integral.

The strict-transform equations retain the coefficient
\(\epsilon=5/\pi^4\) from \eqref{d2:eq:epsilon}:
\begin{align*}
 h_\pi&=X^4+Y^4+TXY^3+
       \epsilon(1+2\pi Y+2\pi^2Y^2+\pi^3Y^3),\\
 h_x&=1+Y^4+TY^3+
       \epsilon U^4(1+2xY+2x^2Y^2+x^3Y^3),\\
 h_y&=X^4+1+TX+
       \epsilon U^4(1+2y+2y^2+y^3).
\end{align*}
The substitutions are respectively
\((x,y)=(\pi X,\pi Y)\),
\((\pi,y)=(xU,xY)\), and
\((\pi,x)=(yU,yX)\).
The latter two ambient charts retain the relations
\(\pi=xU\) and \(\pi=yU\).

These ambient charts are regular.
For example, on the \(x\)-chart one can eliminate \(x\) or \(U\) when
the other is a unit. At a point where both vanish,
\(\pi-xU\) has the nonzero linear term \(\pi\) in a regular local
polynomial ring over \(V_{\Dtwoar}\).
The argument on the \(y\)-chart is identical with \(x\) replaced by
\(y\); the \(\pi\)-chart is a localization of \(V_{\Dtwoar}[X,Y]\).

Let \(t\) denote the exceptional parameter on any one of these charts.
It is a nonzero divisor, the ambient quotient modulo \(t\) is a
polynomial domain over \(k_0\), and the corresponding \(h\) has
nonzero reduction. If \(ta=hb\), reduction modulo \(t\) forces
\(b=tb'\), and cancellation gives \(a=hb'\).
Thus the quotient by \(h\) has no \(t\)-torsion.
The displayed equations describe the strict transform without an
additional exceptional component.

At a point of \(C_{\mathrm{exc}}\), its smoothness gives a nonzero
differential of the reduced equation in an exceptional-plane
direction. Therefore the strict transform is regular there.
Away from the exceptional divisor, where \(\pi\) is invertible, the
characteristic-zero hypersurface is smooth: if
\(F_T=xy^3=0\), then either \(x=0\), where \(P'(y)\ne0\) at each
root of \(P\), or \(y=0\), where \(x^4=-5\) and
\(F_x=4x^3\ne0\).
At a nonclosed point over \(\pi=0\), the binary form
\(x^4+y^4+Txy^3\) is geometrically square-free by the repeated-root
calculation above. Its affine zero locus is smooth away from
\((x,y)=(0,0)\), so its \(x,y\) Jacobian proves regularity there.
Consequently \(\DtwocY\) is regular everywhere and is an isomorphism off
\(\Dtwofn\). The hypersurface \(B_{\Dtwoar}\) is Cohen--Macaulay and regular
in codimension one, hence normal.

The exceptional ideal satisfies
\[
 \Dtwofn\DtwocO_{\DtwocY}=\DtwocO_{\DtwocY}(-C_{\mathrm{exc}}),
 \qquad
 \DtwocO_{\DtwocY}(C_{\mathrm{exc}})|_{C_{\mathrm{exc}}}
       =\DtwocO_{C_{\mathrm{exc}}}(-1).
\]
These identities follow from the exceptional parameters on the charts,
or from \cite[Lemmas~31.33.2 and~31.33.4, Tags~0804 and~02OS]{StacksBlowup}:
the charts are \(R[I/a]\), the extended ideal is generated by \(a\),
and the degree-minus-one tautological line is the exceptional-divisor
line.

Write \(a_g=\rho^g-1\) and \(b_g=a_g/\pi\).
Then \(b_g\) is a unit with reduction \(g\).
The strict transform \(\Dtwowt D_g\) of the prime
\((x,y-a_g)\) meets the exceptional curve only at
\[
 P_g=[1:0:g].
\]
Its parameter is \(\pi\), giving this tangent direction.
On the \(\pi\)-chart, setting \(X=0\) gives
\[
 h_\pi(0,Y)=\prod_{g=1}^4(Y-b_g),
\]
because \(P(\pi Y)=\pi^4\prod_g(Y-b_g)\).
The \(Y\)-derivative is a unit at
\((\pi,X,Y-b_g)\).
The maximal ideal there is generated by \(\pi,X\), while
\(\Dtwowt D_g\) is given by \(X=0\).
Its intersection with \(C_{\mathrm{exc}}\), given by \(\pi=0\),
has multiplicity one.

We next compute the lattice on
\(C_{\mathrm{exc}}^*=C_{\mathrm{exc}}\cap\{UY\ne0\}\).
Extend the constants to \(k=\overline{\DtwoFF}_5\), and consider
\[
 M=\DtwoSpec k[X,Z,Y,Y^{-1}]/(XZ-(1-Y^4)),
 \qquad
 T=\frac{Z-X^3}{Y^3}.
\]
This surface is smooth, since the derivative \(4Y^3\) is a unit.
The open \(X\ne0\) has coordinate ring
\(k[X^{\pm1},Y^{\pm1}]\), which is factorial.
The omitted prime divisors are
\(L_g=V(X,Y-g)\), each with multiplicity one in \(\Dtwodivisor(X)\).
They generate \(\DtwoCl(M)=\DtwoPic(M)\) by divisor localization.
A relation among them is the divisor of a function that becomes a unit
on the Laurent open. Such a unit is \(cX^rY^s\).
Since \(Y\) is already invertible on \(M\), the only relation is
\[
 \DtwoPic(M)=\DtwoZZ^4/\DtwoZZ(1,1,1,1).
\]

For every \(t\in k\), the fiber \(T=t\) is
\[
 X^4+Y^4+tXY^3=1,\qquad Y\ne0.
\]
With \(s=X/Y\) and \(w=1/Y\), its coordinate ring is
\[
 k[s,w,w^{-1}]/\bigl(w^4-(s^4+ts+1)\bigr).
\]
The polynomial \(s^4+ts+1\) is not a square in \(k(s)\).
A rational square that is polynomial is a polynomial square by unique
factorization. Comparing it with \((s^2+bs+c)^2\), after normalizing
the leading coefficient, gives \(b=c=0\) from the cubic and quadratic
coefficients, contradicting the constant term.
Since \(k\) contains the fourth roots of unity and has characteristic
different from two, adjoining a fourth root of this nonsquare has
degree four. The extension is a splitting field with Galois group
inside \(\mu_4\); a subgroup of order at most two would fix the square
of the root and make the polynomial a square.
Thus every fiber is integral and reduced.
Its single vertical prime divisor is \(\Dtwodivisor(T-t)\), with
multiplicity one.

Localization to \(k(T)\) removes exactly these vertical prime
divisors. Their classes are zero, so
\[
 \DtwoPic(M)\xrightarrow{\ \sim\ }\DtwoPic(M_{k(T)}).
\]
The generic fiber is \(C_{\mathrm{exc},k(T)}^*\), and \(L_g\)
restricts to \(P_g\) with multiplicity one.
Hence a principal combination \(\sum_g n_gP_g\) on
\(C_{\mathrm{exc}}^*\) must be diagonal: extend to \(k(T)\) and use
the lattice just computed.

The substitution \(v=x^3+Ty^3\) identifies
\[
 \DtwocQ=\DtwoSpec E_{\Dtwoar}[x,v,y,y^{-1}]/(xv+P(y)).
\]
This surface is smooth because \(P\) has four distinct roots, all in
\(E_{\Dtwoar}\). Its open \(x\ne0\) is the Laurent factorial surface
with coordinates \(x,y\).
The removed prime divisors are exactly \(D_g\), each occurring once
in \(\Dtwodivisor(x)\).
Every relation is obtained from a Laurent unit \(cx^ry^s\).
It follows that
\(\DtwoPic(\DtwocQ)=\DtwoZZ^4/\DtwoZZ(1,1,1,1)\).
This also proves geometric integrality.
The same argument works over every extension of \(E_{\Dtwoar}\).

Suppose that \(\sum_g n_gD_g\) is principal on \(\DtwocT\), with
rational function \(f\).
On the regular surface \(\DtwocY\), its divisor differs from
\(\sum_g n_g\Dtwowt D_g\) by
\[
 mC_{\mathrm{exc}}+\sum_jm_jV_j,
\]
where \(V_j\) are strict transforms of primes in \(\pi=0\) or \(y=0\).
Indeed \(\DtwocY[1/\pi,1/y]=\DtwocT\).
The strict transforms over \(\pi=0\) meet the exceptional curve only
in \(U=0\), and those over \(y=0\) meet it only in \(Y=0\), as the
chart equations show.
They therefore miss \(C_{\mathrm{exc}}^*\).
The line \(\DtwocO_{C_{\mathrm{exc}}}(-1)\) is trivial where \(U\ne0\).
Restriction of the principal-divisor identity gives
\[
 \sum_g n_g[P_g]=0\quad\text{in }\DtwoPic(C_{\mathrm{exc}}^*).
\]
The coefficients \(n_g\) are equal.
The original combination is already a multiple of \(\Dtwodivisor(x)\)
on \(\DtwocQ\), which proves Picard injectivity.

Fix an embedding into \(\DtwoCC\).
The divisors \(D_g(\DtwoCC)\) are disjoint affine lines, and their
complement is \((\DtwoCC^\times)^2\).
The Thom isomorphism gives the exact segment
\[
 H^1((\DtwoCC^\times)^2,\DtwoRR)
 \longrightarrow \bigoplus_g\DtwoRR
 \longrightarrow H^2(\DtwocQ(\DtwoCC),\DtwoRR).
\]
The second map sends the standard vectors to \(e_g\).
The angular classes of \(x,y\) form a basis of the first group.
Around each \(D_g\), \(x\) winds once because its zero is simple,
whereas \(y\) winds zero times because it is invertible.
The image of the first map is precisely the diagonal line.
This proves the asserted relation among the \(e_g\).

Let \(Z\) be an irreducible \(E_{\Dtwoar}\)-divisor in the zero locus of
some \(s\in S\).
Its restriction to \(\DtwocT\) is empty.
Picard injectivity therefore gives \([Z]=0\) on \(\DtwocQ\).
Over an algebraic closure, its geometric components
\(Z_1,\ldots,Z_d\) form one Galois orbit and each has multiplicity
one, since the characteristic is zero.
The geometric Picard group is the torsion-free lattice computed
above. Galois acts trivially on it because its generators \(D_g\)
are already defined over \(E_{\Dtwoar}\).
Thus all \([Z_j]\) equal a class \(h\), and \(dh=0\).
It follows that every \([Z_j]\) vanishes.
Geometric irreducibility persists under extension of algebraically
closed fields, so every irreducible complex component of the boundary
has zero Picard class and zero topological first Chern class.

Let \(H\) be the reduced complex divisor cut out by \(s\).
If it is empty, restriction is the identity.
Otherwise remove the finite set \(\Sigma\) of singular points and
component intersections.
Local cohomology at a point of a smooth real four-manifold is
concentrated in degree four, by the cohomology of a ball relative to
its punctured ball. Therefore
\[
 H^2_\Sigma(\DtwocQ(\DtwoCC),\DtwoRR)
 =H^3_\Sigma(\DtwocQ(\DtwoCC),\DtwoRR)=0.
\]
Put \(X=\DtwocQ(\DtwoCC)\), \(X^\circ=X\setminus\Sigma\), and
\(D=H\setminus\Sigma\). The support exact sequences give
\[
 H^2_H(X,\DtwoRR)\xrightarrow{\sim}H^2_D(X^\circ,\DtwoRR),
 \qquad
 H^2(X,\DtwoRR)\xrightarrow{\sim}H^2(X^\circ,\DtwoRR).
\]
Each irreducible component \(H_i\) of \(H\) has connected smooth
part \(D_i=H_i\setminus\Sigma\): its normalization is a connected
Riemann surface, and removing finitely many points preserves
connectedness. The \(D_i\) are disjoint closed smooth curves in
\(X^\circ\).

A tubular neighborhood is a disk neighborhood in their normal
complex line bundles; a varying radius accommodates noncompact
curves. Excision reduces cohomology with supports to the relative
cohomology of this neighborhood and its complement of the zero
section. Over a contractible trivializing open, the normal fiber
pair \((B^2,B^2\setminus\{0\})\) has relative cohomology
\(\DtwoRR\) in degree two and zero otherwise.
Changes of trivialization preserve its complex orientation and
therefore its generator. Relative Mayer--Vietoris on a good
trivializing cover gives
\[
 H^2_D(X^\circ,\DtwoRR)=H^0(D,\DtwoRR)
       =\bigoplus_i\DtwoRR.
\]
This uses ordinary cohomology of \(D\), without a compactness
assumption.

The canonical section of \(\DtwocO_X(H_i)\) is transverse to zero
along \(D_i\) on \(X^\circ\), with complex-linear normal derivative.
Pulling back the oriented fiber Thom class gives the corresponding
support generator. After forgetting supports, homotopy of this
section to the zero section identifies its class with the Euler
class of the underlying oriented real plane bundle. This equals
\(c_1(\DtwocO_X(H_i))|_{X^\circ}\): both classes are given by the
winding-number cocycle of the circle-valued transition functions,
with the complex orientation fixing the positive generator.
The isomorphism on ambient \(H^2\) above identifies its preimage
with the global class \(c_1(\DtwocO_X(H_i))\) on \(X\).
Every such class is zero by the Picard-group calculation.
The restriction sequence
\[
 H^2_H(\DtwocQ(\DtwoCC),\DtwoRR)\longrightarrow H^2(\DtwocQ(\DtwoCC),\DtwoRR)
 \longrightarrow H^2(\DtwocQ_s(\DtwoCC),\DtwoRR)
\]
now proves injectivity for every \(s\).
\end{proof}

\section{The multiplicative analytic regulator}\label{d2:sec:analytic}

The goal of this section is to prove Theorem~\ref{d2:thm:geometric} from
Proposition~\ref{d2:prop:geometry}.

\begin{proof}[Proof of Theorem~\ref{d2:thm:geometric}]
Regularity and the supported construction were established in
Section~\ref{d2:sec:preliminaries}.
Let \(\DtwocO_{E_{\Dtwoar}}\) be the integer ring of \(E_{\Dtwoar}\).
It is contained in \(W_{\Dtwoar}\): an element integral over \(\DtwoZZ\)
is integral over \(W_{\Dtwoar}\), which is integrally closed in
\(E_{\Dtwoar}\).

The field \(E_{\Dtwoar}\) has two complex-conjugate pairs of embeddings
and no real embedding. By \cite[Proposition~12.2]{Borel},
\[
 \dim_\DtwoRR\bigl(K_3(\DtwocO_{E_{\Dtwoar}})\otimes\DtwoRR\bigr)=2.
\]
The analytic regulator is injective on this real vector space by
\cite[Section~6, especially~(6.4)]{BunkeTamme}.
The comparison there with the Borel regulator is up to a nonzero
factor \(2\), which does not affect injectivity.
At a point, its weight-two target is \(\DtwoCC/\DtwoRR(2)\), where
\(\DtwoRR(p)=(2\pi i)^p\DtwoRR\).
Conjugate embeddings give conjugate values.
The interval description has point representatives
\(i^{p+1}dt\); in weight two they are purely imaginary.

Choose \(\gamma_0\in K_3(\DtwocO_{E_{\Dtwoar}})\) with nonzero regulator,
and set
\[
 \gamma=\gamma_0-\sigma_{-1}\gamma_0,\qquad
 \beta=\gamma|_{W_{\Dtwoar}}.
\]
Since \(\DtwoRR(2)\) is the real subspace of \(\DtwoCC\), conjugation negates
the unique imaginary representative of a point regulator.
Thus the regulator of \(\gamma\) is twice that of \(\gamma_0\)
and remains nonzero after restriction to \(E_{\Dtwoar}\).
The identity \(\sigma_{-1}\beta=-\beta\) holds in actual K-theory.
Put \(c=c(\beta)\).
Its generic restriction is zero by Definition~\ref{d2:defn:supported}.

On \(\DtwocQ\), define
\[
 J_g=\DtwocO_{\DtwocQ}(-D_g),\qquad
 \eta_g=1-[J_g],\qquad
 \kappa=\sum_{g\in G}(\sigma_g\beta)|_{E_{\Dtwoar}}\cdot\eta_g.
\]
The resolution of the Cartier divisor gives
\(\eta_g=[\DtwocO_{D_g}]\) in \(K_0(\DtwocQ)\).
We will detect \(\kappa\) on every \(\DtwocQ_s\).

The regulator is first defined at a finite arithmetic stage.
Choose \(N\) divisible by \(5\) and by the denominators needed for
\(s\). The roots \(a_g=\rho^g-1\) and their pairwise differences are
units over \(\DtwocO_{E_{\Dtwoar}}[1/N]\).
Indeed their norms have absolute value \(5\):
for \(a_g\) this is \(\Phi_5(1)=5\), and a difference is
\(\rho^h(\rho^{g-h}-1)\).
The scheme
\[
 X_N=\DtwoSpec
 \DtwocO_{E_{\Dtwoar}}[1/N][x,v,y,y^{-1}]/(xv+P(y))
\]
is smooth over \(\DtwocO_{E_{\Dtwoar}}[1/N]\).
At a residue point where either \(x\) or \(v\) is nonzero, one of the
corresponding partial derivatives is nonzero.
Where both vanish, \(P'(y)\ne0\), since the four roots remain distinct.
Thus \(X_N\) and \((X_N)_s\) are regular separated schemes of finite
type over \(\DtwoZZ\).
The root divisors extend as Cartier divisors: near \(y=a_g\), the
other factors of \(P\) are units and the equation expresses
\(y-a_g\) as a multiple of \(x\).
The coefficient classes \(\sigma_g\gamma\) extend there as well.

For such schemes, \cite[Theorem~2.31 and Remark~2.27]{BunkeTamme}
gives a natural regulator of ring spectra with homotopy maps
\[
 K_n(X)\longrightarrow
 \prod_{p\ge0}H_{\mathcal D,\Dtwoan}^{\,2p-n}
                  (X(\DtwoCC),\DtwoRR(p)).
\]
The target satisfies the arithmetic conjugation invariance.
After constructing this map, project to the component selected by
a fixed embedding \(\tau\colon E_{\Dtwoar}\to\DtwoCC\).
Increasing \(N\) is compatible with this construction.
Filtered-colimit continuity
\cite[IV, Section~6.4]{Weibel} therefore defines the regulator on
\(K_3(\DtwocQ_s)\); the complex component is always
\(U_s=\DtwocQ_s(\DtwoCC)\).

We use its multiplicative structure explicitly.
By \cite[Remark~2.6 and Definition~2.7]{BunkeTamme}, analytic
Deligne cohomology is the shifted cone of
\[
 (2\pi i)^pA_\DtwoRR^\bullet(U)
       \oplus F^pA_\DtwoCC^\bullet(U)
 \longrightarrow A_\DtwoCC^\bullet(U),
 \qquad (a,b)\longmapsto a-b.
\]
Here \(A^\bullet\) denotes smooth forms, and \(F^p\) means at least
\(p\) holomorphic differentials. There is no logarithmic boundary
condition. Its interval model consists of forms on
\([0,1]\times U\) whose endpoint values are respectively real-twisted
and in \(F^p\).
Multiplication is wedge product.
The map to the cone is
\[
 \theta\longmapsto
 \left(\theta|_0,\theta|_1,-\int_{[0,1]}\theta\right);
\]
see \cite[Proposition~2.8 and Corollary~2.14]{BunkeTamme}.

For a complex surface, \(F^3A_\DtwoCC^\bullet(U)=0\), because at most two
holomorphic coordinate differentials can occur.
The cone sequence consequently identifies
\begin{equation}\label{d2:eq:analytic-quotient}
 H_{\mathcal D,\Dtwoan}^{3}(U,\DtwoRR(3))
   =H^2(U,\DtwoCC)/H^2(U,\DtwoRR(3)).
\end{equation}
The next possible term is the kernel of
\(H^3(U,\DtwoRR(3))\to H^3(U,\DtwoCC)\), which is zero by extension of
scalars. This argument does not require properness.

Let \(r_g\in i\DtwoRR\) be the imaginary representative of the point
regulator of \(\sigma_g\beta\) at \(\tau\).
At least one \(r_g\) is nonzero, since the embeddings
\(\tau\sigma_g\) exhaust all embeddings of \(E_{\Dtwoar}\).
Antisymmetry and commutativity of \(G\) give \(r_{-g}=-r_g\).
The vector \((r_g)\) is therefore nonzero and not diagonal.

The interval representative of this point class is \(-r_gdt\).
Choose a Hermitian metric and Chern connection on \(\DtwocO(D_g)\),
and let \(\omega_g\) be the resulting real closed representative of
\(e_g\). Its weight-one regulator is the interval-constant form
\(2\pi i\omega_g\).
This follows from the unitary, type-\((1,1)\) curvature and the
degree-two term of the unnormalized character
\(\DtwoTr\exp(-\mathrm{curvature})\) in
\cite[Definitions~2.16--2.20]{BunkeTamme}.
The weight-one class of \(\eta_g\) is the same, by additivity of
the first Chern class.
Wedge multiplication followed by minus integration gives
\[
 (-r_gdt)\wedge(2\pi i\omega_g)
 \longmapsto 2\pi i r_g\omega_g.
\]
Thus the weight-three regulator of \(\kappa|_{\DtwocQ_s}\) is
\begin{equation}\label{d2:eq:analytic-detector}
 \left[2\pi i\sum_g r_g e_g|_{U_s}\right]
 \in H^2(U_s,\DtwoCC)/H^2(U_s,\DtwoRR(3)).
\end{equation}
No other point weight contributes: point Deligne cohomology is
nonzero only in degrees zero and one, and for \(K_3\) the equation
\(2p-3=1\) forces \(p=2\).

Write \(r_g=it_g\) with \(t_g\in\DtwoRR\).
Proposition~\ref{d2:prop:geometry} shows that
\(\sum_g t_ge_g\ne0\), since \((t_g)\) is not diagonal.
It stays nonzero on every \(U_s\).
The representative in \eqref{d2:eq:analytic-detector} is the nonzero
real class
\[
 -2\pi\sum_g t_ge_g|_{U_s}.
\]
By contrast, \(H^2(U_s,\DtwoRR(3))=iH^2(U_s,\DtwoRR)\) is purely imaginary.
The real and purely imaginary subspaces intersect in zero.
Hence \eqref{d2:eq:analytic-detector} is nonzero.
Its target is a real vector space, so
\(\kappa|_{\DtwocQ_s}\) has infinite order for every \(s\in S\).

The map \(A\to B_{\Dtwoar}[1/\pi,1/y]\) is flat.
Over it, the divisor \(x=0\) splits into the disjoint components
\(D_g\cap\DtwocT\).
On the \(g\)-th component, the coefficient map sends
\(\rho=1+y\) to \(\rho^g\).
Flat base change and the projection formula give the equality of
actual classes
\begin{equation}\label{d2:eq:analytic-splitting}
 c|_{\DtwocT}
 =\sum_g(\sigma_g\beta)|_{E_{\Dtwoar}}\,[\DtwocO_{D_g\cap\DtwocT}]
 =\kappa|_{\DtwocT}.
\end{equation}
The ring of \(\DtwocT\) is regular, being a localization of the smooth
surface \(\DtwocQ\). No identification of K-theory and G-theory on the
possibly singular ring \(B_{\Dtwoar}\) is involved.

Finally,
\[
 K_3(\DtwocT)=\varinjlim_{s\in S}K_3(\DtwocQ_s).
\]
If a nonzero integer killed \(c\), it would kill \(\kappa\) in this
colimit by \eqref{d2:eq:analytic-splitting}.
The resulting equality would hold at a finite stage, contrary to
\eqref{d2:eq:analytic-detector}. Thus \(c\) has infinite order.
\end{proof}

\section{Relative cyclic detection}\label{d2:sec:cyclic}

In this section, we prove the comparison and product formula that will
detect the completed supported class.

\begin{proposition}\label{d2:prop:cyclic}
Let \(k=\overline{\DtwoFF}_5\), \(W_0=W(k)\), and let \(R\) be the
integer ring of a finite totally ramified extension of \(W_0[1/5]\).
Write \(\pi_R\) for a uniformizer and \(E_R=\DtwoFrac R\).
Let \(B_0\) be a localization of a smooth \(R\)-algebra whose relative
cotangent module is finite projective of constant rank two.
Assume that \(B=\varprojlim_nB_0/5^nB_0\ne0\).
There is a natural equivalence
\begin{equation}\label{d2:eq:relative-cyclic}
 \DtwocK(B,\pi_RB)\simeq\Sigma\DtwoHC^\Dtwocts(B/R;\DtwoQQ_5)
\end{equation}
compatible with the \(K(B)\)-module action through the negative cyclic
Chern character.

The degree-two cyclic group has the canonical summand
\[
 \Omega^2_{B/R}[1/5]\big/d\Omega^1_{B/R}[1/5].
\]
For an actual \(\beta\in K_3(R,\pi_RR)\), let \(b\in E_R\) be its
relative point coordinate, so its point character is \(t^{-1}b\)
in the mixed-HKR convention with \(|t|=-2\).
There is a scalar \(\lambda_R\in E_R^\times\), independent of
\(B\), \(J\), and \(\beta\), such that for every invertible
\(B\)-module \(J\), the top-form component of
\(\beta_B(1-[J])\) under \eqref{d2:eq:relative-cyclic} is
\[
 -\lambda_R b\,c_1^\DtwodR(J).
\]
Nonvanishing of \(-b\,c_1^\DtwodR(J)\), or of a finite sum of
these forms, proves nonvanishing of the corresponding actual
relative class after rationalized spectral completion: the whole
sum is multiplied by the same nonzero \(\lambda_R\).
No injectivity from actual K-theory to its completion is asserted.
\end{proposition}

\begin{proof}
The fiber theorem \cite[Theorem~A]{AMMN} applies to a commutative
ring henselian along \((5)\). Its coefficients \(\DtwoQQ_5\) mean
spectral completion followed by inversion of \(5\).
It gives a cartesian square with rows
\[
 K(T;\DtwoQQ_5)\longrightarrow K(T/5;\DtwoQQ_5),
 \qquad
 \DtwoHC^-(T;\DtwoQQ_5)\longrightarrow\DtwoHP(T;\DtwoQQ_5),
\]
and hence identifies the relative K-theory fiber with
\(\Sigma\DtwoHC(T;\DtwoQQ_5)\).
The continuity result
\cite[Proposition~4.2]{AMMN} applies to bounded \(5\)-power
torsion and compares completed Hochschild and cyclic objects with
their infinitesimal inverse limits.
The paragraph preceding \cite[Proposition~4.10]{AMMN}
also treats relative Hochschild and cyclic theories.
We extend that relative argument to \(B_0\) by bounds in each finite
degree range; no topological finite-presentation assertion about
\(B\) is needed.

Both \(B_0\) and \(B\) are \(5\)-torsion-free, by flatness over the
discrete valuation ring and flatness of noetherian adic completion.
Their continuous relative forms are the completions of
\(\Omega^j_{B_0/R}\). Completion is exact on these finite projective
modules, and they vanish for \(j>2\).

For every \(n\), the ring \(W_0/5^n\) is the filtered union of the
finite \'etale Witt extensions of \(\DtwoZZ/5^n\).
Its relative cotangent complex over \(\DtwoZZ/5^n\) is zero.
Flat base change therefore gives
\({L_{W_0/\DtwoZZ}}^\wedge_5=0\).
Choose an Eisenstein presentation \(R=W_0[\pi_R]\), with polynomial
\(f\). Then
\[
 {L_{R/\DtwoZZ}}^\wedge_5
   \simeq [R\xrightarrow{f'(\pi_R)}R]
\]
in cohomological degrees \(-1,0\).
Separability in characteristic zero gives \(f'(\pi_R)\ne0\).
Choose \(a\ge0\) with \(5^a\in f'(\pi_R)R\).
Multiplication by \(5^a\) on this complex is null-homotopic.

Completed cotangent transitivity gives a triangle
\[
 M\longrightarrow \Dtwowh L\longrightarrow P,
 \qquad
 M=[B\xrightarrow{f'(\pi_R)}B],\quad
 \Dtwowh L={L_{B_0/\DtwoZZ}}^\wedge_5,\quad
 P=(\Omega^1_{B_0/R})^\wedge_5.
\]
The module \(P\) is projective in degree zero and \(M\) lies in
degrees \(-1,0\).
The connecting map \(P\to M[1]\) is zero, so the triangle splits
as a triangle of \(B\)-complexes.
The kernel terms in the \(i\)-th derived exterior power of the
projection are sums of
\[
 \bigl(L\!\Lambda^rM\bigr)\otimes\Lambda^{i-r}P,
 \qquad 1\le r\le i.
\]
The \(r\)-th exterior-power functor takes the null-homotopy of
\(5^a\Dtwoid_M\) to one of \(5^{ar}\).
Thus each displayed term is killed as an object by \(5^{ar}\).

The Hochschild--Kostant--Rosenberg filtration has graded pieces
\((L\!\Lambda^iL)[i]\), and its tails become increasingly connective.
Only finitely many pieces affect any fixed homotopy range.
This remains true after derived completion.
For a \(c\)-connective spectrum, the only possible additional bottom
group is the derived inverse limit of the quotients of its
\(c\)-th homotopy group; that quotient tower is surjective.
Completion therefore preserves connectivity.
Combining the exterior-power bounds along the finitely many
extensions bounds the \(5\)-power torsion in the comparison cone in
each fixed range. Hence
\begin{equation}\label{d2:eq:hh-comparison}
 \DtwoHH(B_0;\DtwoZZ_5)[1/5]
 \xrightarrow{\ \sim\ }
 \DtwoHH(B_0/R;\DtwoZZ_5)[1/5]
\end{equation}
is an equivalence of spectra with circle action.

Ordinary cyclic homology is the circle homotopy-orbit construction.
For a connective spectrum, the omitted part of the skeletal
filtration of its bar construction becomes arbitrarily connective.
Each finite stage uses finite colimits and tensors with finite cell
spaces, which commute with spectral completion.
The finite-range argument proving \eqref{d2:eq:hh-comparison} thus
also proves the ordinary cyclic comparison.
Equivalently, the circle norm sequence makes the square comparing
the absolute and relative maps \(\DtwoHC^-\to\DtwoHP\) cartesian.
We do not need the separate vertical maps on \(\DtwoHC^-\) and \(\DtwoHP\)
to be equivalences; those constructions involve infinite limits.

The complete ring \(B\) is henselian along \((5)\).
Applying the fiber theorem to the ordinary ring \(B\), and replacing
the cyclic fiber by the relative one just computed, gives the desired
comparison with \(\DtwocK(B,5B)\).
The map \(B/5\to B/\pi_R\) has nilpotent kernel.
By \cite[Corollary~3.8]{AMMN}, a nilpotent relative extension in
a ring with nilpotent \(5\) has rationalized completed relative
K-theory zero. Exactness therefore identifies
\(\DtwocK(B,5B)\) with \(\DtwocK(B,\pi_RB)\).
This proves \eqref{d2:eq:relative-cyclic}.
Precomposing with the map from actual relative K-theory defines the
actual character.

The module structure requires more than naturality.
Put \(H=\DtwoTHH(B)\), a commutative algebra object in cyclotomic spectra.
The coefficient comparisons in
\cite[Proposition~2.8 and Corollary~2.10]{AMMN}
are tensored with \(H\); their maps are
\(\Dtwoid_H\) tensored with the coefficient maps.
They are consequently maps of \(H\)-modules.
The maps from \(H\), the linearization
\(H\otimes H\DtwoZZ\to\DtwoHH(B)\) used in
\cite[Theorem~2.12]{AMMN}, and projection to
\(\DtwoHH(B/R)\) are \(H\)-linear.

The functor \(\DtwoTC\) is lax symmetric monoidal, so these diagrams
become diagrams of \(\DtwoTC(B)\)-modules.
Circle invariants, orbits, and the norm fiber carry the corresponding
actions. The norm projection formula identifies the fiber action
with the usual \(\DtwoHC^-\)-action on \(\Sigma\DtwoHC\).
The equivalences used above may therefore be inverted in the module
category.
The relative K-to-\(\DtwoTC\) map is induced by the multiplicative
cyclotomic trace, whose \(E_\infty\) refinement is
\cite[Theorem~7.4]{BGTTrace}.
Restricting the module action along this trace proves the asserted
\(K(B)\)-linearity through the negative cyclic character.

It remains to compute the ordinary cyclic component.
In relative dimension two the normalized HKR map is integral,
because \(2\) is invertible:
\[
 a_0[a_1|\cdots|a_j]\longmapsto
 \frac1{j!}a_0\,da_1\wedge\cdots\wedge da_j
 \quad (0\le j\le2),
\]
and higher chains map to zero.
The Hochschild boundary maps to zero.
The alternating sum defining the Connes operator contributes
\(j+1\) before normalization, so the factors \(1/j!\) make it
correspond to \(d\).
For \(j=2\), the target of \(d\) is zero.
The antisymmetrized smooth HKR representative maps to the original
form, and higher relative Hochschild homology vanishes.
Thus this is a quasi-isomorphism of mixed complexes.
The completed comparison gives the finite continuous form modules
already identified.

Give a bookkeeping variable \(t\) homological degree \(-2\), and give
a \(j\)-form degree \(j\).
The mixed differential is \(td\).
The ordinary cyclic model retains the powers
\(1,t^{-1},t^{-2},\ldots\) and discards positive powers.
In degree two it has \(\Omega^2\oplus t^{-1}B\).
The differential sends \(t^{-1}b\) to \(db\).
The degree-three boundaries in the top-form term are exactly
\(d\Omega^1\).
Consequently
\begin{equation}\label{d2:eq:hc-two}
 \DtwoHC_2^\Dtwocts(B/R;\DtwoQQ_5)=
 \frac{\Omega^2_{B/R}[1/5]}{d\Omega^1_{B/R}[1/5]}
 \oplus
 t^{-1}\ker\bigl(d\colon B[1/5]\to\Omega^1_{B/R}[1/5]\bigr).
\end{equation}
In particular the image in the first summand is the ordinary
differential image.

At the coefficient point, the degree-two group is \(t^{-1}E_R\).
The class \(\beta\) has representative \(t^{-1}b\), and \(b\) is
constant for differentiation relative to \(R\).
We retain this point coordinate without further normalization.
The character comparison in
\cite[v2, Constructions~4.7--4.8 and~4.11,
Propositions~4.10 and~4.12, pp.~21--23]{AMMN}
gives a common nonzero scalar, whose image in \(E_R\) we denote
by \(\lambda_R\), for which the first Chern component of the
periodic character is \(\lambda_R c_1^\DtwodR\).
Here equation~(28) of that reference gives the special-fiber
character, equation~(24) gives the crystalline/de Rham comparison,
and equation~(25) is the fiber square.
The comparison preserves first Chern classes because their
transition functions give the same \(d\log\) cocycle.
The scalar is common to all smooth formal finite-type models
over the fixed coefficient ring; neither scalar one nor
\(\DtwoQQ_5\)-rationality is required.

This formula applies to every invertible \(B\)-module \(J\) here.
The finite presentation of \(J/\pi_RJ\), together with the finite
equations witnessing rank-one projectivity, descends to the
special fiber of a finite localization of a smooth relative
rank-two \(R\)-algebra. Complete this finite model.
An idempotent presenting the descended module lifts: starting
with a matrix lift \(e\), iterate \(e\mapsto3e^2-2e^3\).
If \(\delta=e^2-e\), the new error is
\(\delta^2(4\delta-3)\), so the errors tend to zero adically
and the matrices converge to an idempotent.
Its projective module has rank one since its reduction does
and \(\pi_R\) lies in the Jacobson radical.
The completed model maps naturally to \(B\).
After pullback, an isomorphism modulo \(\pi_R\) with \(J/\pi_RJ\)
and its inverse lift between the projective modules.
Their composites are invertible because they reduce to the
identity modulo the Jacobson radical, so the pulled-back line
is isomorphic to \(J\).
Naturality therefore carries the finite-model character formula
to \(B\).

In the rank-two mixed complex, both degree-zero negative and
periodic cyclic homology are
\[
 \ker\bigl(d\colon B[1/5]\longrightarrow
                    \Omega^1_{B/R}[1/5]\bigr)
 \ \oplus\
 t\bigl(\Omega^2_{B/R}[1/5]/d\Omega^1_{B/R}[1/5]\bigr).
\]
The comparison is the identity on these summands: total degree
zero contains functions and \(t\)-times two-forms, total degree
one supplies the boundaries \(t\,d\Omega^1\), and higher forms
vanish. Thus the negative-to-periodic map is injective in this
degree. The commutative fiber square identifies the negative
cyclic character of \(J\) with
\(1+\lambda_Rt\,c_1^\DtwodR(J)\).
The character of \(1-[J]\) is consequently
\(-\lambda_Rt\,c_1^\DtwodR(J)\), whose action on \(t^{-1}b\)
has top component \(-\lambda_Rb\,c_1^\DtwodR(J)\).
At the normalized bar level, the shuffle term gives this wedge
product, while cyclic-shuffle corrections insert the relative
scalar \(b\) in a bar slot and vanish.
This is the module action established above.
The common nonzero scalar preserves nonvanishing of every
finite sum.

A nonzero top-form component cannot be the image of a zero
rationalized completed relative class.
The same calculation is additive for finite sums.
It makes no assertion that actual K-theory injects into completion.
\end{proof}

\section{Cohen completion and elliptic contraction}\label{d2:sec:cohen}

We prove that the graph-divisor class survives Cohen completion.

\subsection{Scalar Frobenius pole removal}

\begin{lemma}\label{d2:lem:scalar}
Let \(p\) be a prime, \(k=\overline{\DtwoFF}_p\), \(W=W(k)\), and
\(K=W[1/p]\).
Let \(D\) be a smooth affine \(W\)-algebra of relative dimension one
with connected nonempty special fiber \(U=\DtwoSpec(D/pD)\).
Put
\[
 S_D=\{s\in D\mid s\bmod p\ne0\},\qquad
 \Dtwowh D=\varprojlim_nD/p^nD,\qquad
 C=\varprojlim_nS_D^{-1}D/p^nD,\qquad F_C=C[1/p].
\]
Let \(D^\dagger\) be the weak completion.
Assume that an integral Frobenius lift
\(\phi\colon D^\dagger\to D^\dagger\) restricts to the \(d\)-th Witt
Frobenius on \(W\), where \(d\ge1\), and reduces to the
\(q=p^d\) power map.
If \(\omega\in\Omega^1_{D^\dagger/W}[1/p]\) has nonzero rigid
cohomology class and
\[
 \phi^*\omega-a\omega=dh
 \qquad(a\in W,\ h\in D^\dagger[1/p]),
\]
then its image in
\[
 \Omega^1_{C/W}[1/p]\big/dF_C
\]
is nonzero.
\end{lemma}

\begin{proof}
A smooth connected curve has disjoint irreducible components, since
its local rings are domains. Therefore \(D_0=D/pD\) is an integral
normal domain. Write \(L=\DtwoFrac D_0\).
The localization \(S_D^{-1}D=D_{pD}\) is a regular local ring of
dimension one with parameter \(p\) and residue field \(L\).
Thus \(C\) is a complete discrete valuation ring with these same
data.

For every \(n\), the map
\(D/p^nD\to S_D^{-1}D/p^nD\) is injective.
If \(sa\in p^nD\) with \(s\in S_D\), reduction modulo \(p\) gives
\(a\in pD\). Write \(a=pa_1\), cancel \(p\) by flatness over \(W\),
and repeat.
Taking limits embeds \(\Dtwowh D\) in \(C\) and gives
\begin{equation}\label{d2:eq:integral-intersection}
 \Dtwowh D\cap p^nC=p^n\Dtwowh D.
\end{equation}
The continuous forms on \(C\) are the completions of the forms on
\(S_D^{-1}D\).

We first determine its continuous constants:
\begin{equation}\label{d2:eq:cohen-constants}
 \ker\bigl(d\colon F_C\to\Omega^1_{C/W}[1/p]\bigr)=K.
\end{equation}
After inverting an element of \(S_D\), smoothness supplies an
\'etale coordinate \(t\) and a ring \(D'\) \'etale over \(W[t]\).
The residue of \(t\) is a separating parameter for \(L\), and the
generic-special completion of \(D'\) is still \(C\).
Formal \'etaleness extends \(t\mapsto t+z\) to a Taylor map
\(D'\to C[[z]]\), successively modulo powers of \(z\).
Further denominators have unit constant coefficient, so this map
extends to the localization and then to \(C\), coefficient by
coefficient.
Write its coefficients as integral continuous \(W\)-linear operators
\(H_j\colon C\to C\).

Uniqueness of the \'etale lift identifies successive translations by
\(z,w\) with translation by \(z+w\).
With \(\partial=H_1\), this gives
\(j!H_j=\partial^j\) and \(\partial(t)=1\).
The continuous cotangent module is free on \(dt\).
If \(f\in C\) satisfies \(df=0\), then
\(j!H_j(f)=0\) for every \(j>0\).
The ring \(C\) is torsion-free over \(\DtwoZZ\), so \(H_j(f)=0\).

On reduction modulo \(p\), all positive Hasse operators annihilate
\(\bar f\in L\).
For each \(r\ge1\),
\[
 1,t,\ldots,t^{p^r-1}
\]
is a basis of \(L\) over \(L^{p^r}\).
Indeed finite separability of \(L/k(t)\) gives
\([L:L^p]=p\), and \(dt\ne0\) implies \(t\notin L^p\).
Iteration gives \([L:L^{p^r}]=p^r\).
A smaller purely inseparable degree for \(t\) over \(L^{p^r}\)
would force \(t\) to be a \(p\)-th power in \(L\).

Expand \(\bar f=\sum_{i=0}^{p^r-1}b_i^{p^r}t^i\).
The Taylor series of \(b_i^{p^r}\) has no terms of degrees
\(1,\ldots,p^r-1\).
Comparing the coefficients of \(z^{p^r-1},z^{p^r-2},\ldots,z\)
successively gives \(b_i=0\) for every \(i>0\).
Thus \(\bar f\in\bigcap_rL^{p^r}=k\).
For the last equality, the valuation of a nonzero element of the
intersection at every point of the smooth proper model of \(L\) is
divisible by all \(p^r\), hence zero.
Such an element is a global regular function on a proper connected
curve and belongs to \(k\); conversely \(k\) is perfect.

Choose \(w_0\in W\) lifting \(\bar f\).
Then \((f-w_0)/p\) is integral and again has zero differential.
Repeating gives \(f=\sum_{i\ge0}p^iw_i\), with \(w_i\in W\).
This converges in \(W\) and in \(C\), so \(f\in W\).
Scaling by a power of \(p\) proves \eqref{d2:eq:cohen-constants}.

The integral map \(\phi\) fixes \(p\) and extends continuously to
\(\Dtwowh D\). Its reduction preserves nonzero residues, so inversion
of \(S_D\) and completion extend it to \(C\).
We claim that
\begin{equation}\label{d2:eq:pole-removal}
 f\in C,\quad \phi(f)-af\in\Dtwowh D
 \quad\Longrightarrow\quad f\in\Dtwowh D.
\end{equation}
Write \(\phi(f)-af=r\).
A pole of \(\bar f\) on \(U\) of order \(m>0\) would give pole
order \(qm\) in \(\bar f^q\) and at most \(m\) in
\(\bar a\bar f\). These cannot cancel.
Hence \(\bar f\) has no poles on \(U\).
Normality gives \(\bar f\in D_0\).

Lift \(\bar f\) to \(g_0\in D\), and put
\[
 f_1=(f-g_0)/p,\qquad
 r_1=(r-\phi(g_0)+ag_0)/p.
\]
The numerator defining \(r_1\) belongs to
\(\Dtwowh D\cap pC=p\Dtwowh D\) by \eqref{d2:eq:integral-intersection}.
Thus \(r_1\in\Dtwowh D\) and \(\phi(f_1)-af_1=r_1\).
Iteration yields
\[
 f=\sum_{i=0}^{n-1}p^ig_i+p^nf_n,
 \qquad g_i\in D,\quad f_n\in C.
\]
The partial sums converge in \(\Dtwowh D\), and the remainder tends to
zero in \(C\). This proves \eqref{d2:eq:pole-removal}.
No \(W\)-linearity of \(\phi\) was used.

Suppose now that \(\omega=df\) in \(F_C\).
The Frobenius equation gives
\(d(\phi(f)-af-h)=0\).
By \eqref{d2:eq:cohen-constants},
\(\phi(f)-af=h+c\) for \(c\in K\).
Multiply by one sufficiently large power of \(p\) to make
\(f,h,c\) integral in \(C,\Dtwowh D,W\), respectively.
The equation is preserved because \(\phi(p)=p\).
Equation~\eqref{d2:eq:pole-removal} then puts the scaled \(f\) in
\(\Dtwowh D\). Undoing the scaling makes \(\omega\) exact in the
rational completed de Rham complex.

For the smooth finite-type curve \(U\) and its constant overconvergent
Frobenius coefficient, the map
\[
 H^1_\Dtworig(U/K)\longrightarrow H^1_\Dtwoconv(U/K)
\]
is injective by \cite[Theorem~3.2.1]{CaroDAddezio2025}.
The weakly completed de Rham complex computes the first group, and
the completed de Rham complex computes the second, with comparison
induced by \(D^\dagger\to\Dtwowh D\);
see also \cite[Theorem~2.1.2]{CaroDAddezio2025}.
Both use ordinary differential images in degree one.
The primitive just constructed contradicts this injectivity.
Only the original finite-type curve \(U\) enters this external
comparison.
\end{proof}

\subsection{A bounded elliptic contraction}

\begin{setup}\label{d2:setup:elliptic}
For the remainder of this section, let
\(k=\overline{\DtwoFF}_5\), \(W=W(k)\), and \(K=W[1/5]\).
Let
\[
 \DtwocE/W\colon v^2=z^3+z
\]
be the smooth projective elliptic curve with origin \(O\), and put
\[
 \DtwocE^\circ=\DtwocE\setminus\{O,(0,0)\},
 \qquad
 E_0=\DtwocE\otimes_Wk,
 \qquad
 A_{\mathrm{ell}}=\Gamma(\DtwocE^\circ,\DtwocO).
\]
Let \(D\) be a smooth affine \(W\)-algebra of relative dimension one
with connected nonempty special fiber. Let \(C\) be its
generic-special completion, and put
\[
 B_{\mathrm{ell}}=
 \bigl(\Dtwowh A_{\mathrm{ell}}\otimes_WC\bigr)^\wedge_5.
\]
\end{setup}

\begin{proposition}\label{d2:prop:contraction}
Choose rational dagger one-forms \(\omega,\eta\) on
\(\DtwocE^\circ\) representing nonzero generators of the proper
slope-one and slope-zero crystalline lines, respectively.
There is a bounded \(K\)-linear functional
\[
 \ell\colon\Omega^1_{\Dtwowh A_{\mathrm{ell}}/W}[1/5]\longrightarrow K
\]
such that
\[
 \ell(\omega)=1,\qquad
 \ell(\eta)=0,\qquad
 \ell(d\Dtwowh A_{\mathrm{ell}}[1/5])=0.
\]
It extends to a continuous degree-minus-one chain operator
\[
 L\colon\Omega^\bullet_{B_{\mathrm{ell}}/W}[1/5]
       \longrightarrow\Omega^\bullet_{C/W}[1/5][-1]
\]
with \(L_2(\alpha\wedge\gamma)=\ell(\alpha)\gamma\), where the
elliptic form is placed first.
The map
\[
 J\colon H^1_\DtwodR(C/W)[1/5]\longrightarrow
          H^2_\DtwodR(B_{\mathrm{ell}}/W)[1/5],
 \qquad [\gamma]\longmapsto[\omega\wedge\gamma],
\]
is split injective, with left inverse induced by \(L\).
In particular, a class represented by
\(\omega\wedge\gamma+\eta\wedge\delta\) is nonzero whenever
\([\gamma]\ne0\) in the ordinary one-form quotient of \(C\).
\end{proposition}

\begin{proof}
The cubic has distinct roots in characteristic five, so \(E_0\) is
smooth. The coefficient of \(z^4\) in \((z^3+z)^2\) is \(2\),
which is nonzero modulo five.
The Hasse-invariant criterion
\cite[Theorem~V.4.1(a), p.~148]{SilvermanAEC2009} gives ordinarity.
Proper rational crystalline \(H^1\) therefore has two lines of
slopes zero and one.
Restriction to \(E_0^\circ\) is injective in rigid \(H^1\).
The removed reduced finite set \(Z=\{O,(0,0)\}\) has codimension
one in the smooth curve \(E_0\).
The support-vanishing assertion of
\cite[Theorem~1.1, third assertion, p.~63]{Petrequin2003}
gives \(H^1_{Z,\Dtworig}(E_0/K)=0\), since \(1<2\).
The localization sequence
\cite[Proposition~1.3, p.~64]{Petrequin2003} therefore gives
\[
 0\longrightarrow H^1_\Dtworig(E_0/K)
       \longrightarrow H^1_\Dtworig(E_0^\circ/K).
\]
This natural restriction respects Frobenius.

For unramified Witt coefficients,
\cite[Theorem~7.3.1]{CaroDAddezio2025} identifies slope-one rigid and
convergent cohomology in degree one.
For smooth affine schemes with unit-root coefficients, the
slope-one convergent part meets the closure of zero trivially
\cite[Theorem~7.3.6]{CaroDAddezio2025}.
The slopes in \([0,1)\) lie in the closure of zero by
\cite[Theorem~7.2.1]{CaroDAddezio2025}.
All hypotheses apply to \(E_0^\circ\), constant coefficients, and
the unramified ring \(W\).

Put \(M_0=\Omega^1_{\Dtwowh A_{\mathrm{ell}}/W}\), \(M=M_0[1/5]\),
and let \(N\) be the closure of
\(d\Dtwowh A_{\mathrm{ell}}[1/5]\) in \(M\).
All one-forms are closed, since the relative dimension is one.
The topology and convergent-cohomology description in
\cite[Definition~7.1.1]{CaroDAddezio2025} identify the separated
quotient with \(M/N\).
The class of \(\omega\) in this quotient is nonzero, whereas
\(\eta\in N\).
We use this separated quotient only to construct \(\ell\).

The finite projective, complete, separated, torsion-free module
\(M_0\) defines a Banach norm on \(M\) with discrete values.
The quotient \(V=M/N\) is Banach for the quotient norm.
For a nonzero quotient vector, the infimum of the norms of its
representatives is positive and attained, because the norm values
are discrete and \(N\) is closed.
Completeness follows by taking a Cauchy subsequence, lifting its
successive differences with norms tending to zero, and summing
those lifts in \(M\).

Let \(V_0\) be the unit ball, and lift a basis of the \(k\)-vector
space \(V_0/5V_0\) to elements \(e_i\in V_0\).
Every \(v\in V_0\) has a unique expansion
\[
 v=\sum_i a_ie_i,\qquad a_i\in W,
\]
where the coefficients tend to zero off finite subsets.
To obtain the expansion, reduce modulo five, subtract a finite
lifted basis expansion, divide by five, and repeat.
The resulting series converges by completeness.
For uniqueness, a nonzero coefficient family has a largest absolute
value. Scaling to make that value one and reducing modulo five
contradicts independence of the finitely many largest coefficients.
Scaling also gives such expansions over \(K\) for all \(v\in V\).
The coordinate maps are consequently bounded and \(K\)-linear.

Some coordinate is nonzero on the image of \(\omega\).
Divide that coordinate by its value on \(\omega\) and compose with
\(M\to V\). This gives \(\ell\).
It annihilates \(N\), hence both \(\eta\) and exact forms.
For some \(e\ge0\), boundedness gives \(5^e\ell(M_0)\subset W\).

Set \(\ell_0=5^e\ell\).
Reduce this integral map modulo \(5^n\), tensor with \(C/5^nC\),
and pass to the inverse limit.
This defines \(M_0\Dtwoct_WC\to C\).
Tensoring instead with \(\Omega^1_{C/W}/5^n\) defines
\(M_0\Dtwoct_W\Omega^1_{C/W}\to\Omega^1_{C/W}\).
After rationalization divide by \(5^e\).
These are limits of compatible operators, with no interchange
between cohomology and inverse limits.

At each finite precision, the tensor-product formula for
differentials gives
\begin{align}
 \Omega^1_{B_{\mathrm{ell}}/W}
 &=(M_0\Dtwoct_WC)\oplus
   (\Dtwowh A_{\mathrm{ell}}\Dtwoct_W\Omega^1_{C/W}),\label{d2:eq:product-one}\\
 \Omega^2_{B_{\mathrm{ell}}/W}
 &=M_0\Dtwoct_W\Omega^1_{C/W}.\label{d2:eq:product-two}
\end{align}
Higher forms vanish.
The two factors have cotangent rank one; for \(C\), this follows
from completion of an \'etale-coordinate localization of \(D\).
The form modules are finite projective.
Flatness over \(W\) ensures that completed tensor products reduce to
the ordinary tensor products modulo \(5^n\), so the displayed limits
are valid.

Define \(L_0=0\).
Define \(L_1\) by the extended \(\ell\) on the first summand of
\eqref{d2:eq:product-one} and zero on the second.
Define \(L_2\) by the second tensor extension.
On functions, \(L_1d=0\).
For an elliptic one-form \(\alpha\) and a base function \(f\),
\[
 d(\alpha f)=-\alpha\wedge d_Cf,\qquad
 L_2d(\alpha f)=-\ell(\alpha)d_Cf=-d_CL_1(\alpha f).
\]
For an elliptic function \(a\) and a base one-form \(\gamma\),
\(d(a\gamma)=d_Aa\wedge\gamma\), so its image under \(L_2\) is
zero, as is \(d_CL_1(a\gamma)\).
Finite sums of pure tensors are dense, and all operators are
continuous. Thus \(L_2d=-d_CL_1\) on the completed complex, which is
the chain identity for the shift \([-1]\).

The form \(\omega\wedge\gamma\) is closed.
If \(\gamma=d_Cf\), then
\(\omega\wedge\gamma=-d(f\omega)\).
Hence \(J\) is defined on ordinary cohomology.
Finally \(L_2(\omega\wedge\gamma)=\gamma\).
This proves the splitting, and \(\ell(\eta)=0\) proves the last
assertion.
\end{proof}

\subsection{Graph divisors}

\begin{proposition}\label{d2:prop:graph}
In Setup~\ref{d2:setup:elliptic}, assume that
\(\Dtwowt Q\colon\DtwoSpec D\to\DtwocE^\circ\) has nonconstant special fiber
\(Q\colon U\to E_0^\circ\).
Let \(\DtwocL_Q\) be the pullback to \(\DtwoSpec B_{\mathrm{ell}}\) of
\[
 \DtwocO(\Gamma_{\Dtwowt Q}-\Gamma_{-\Dtwowt Q})
\]
on \(\DtwocE^\circ\times_W\DtwoSpec D\).
Then
\[
 c_1^\DtwodR(\DtwocL_Q)\ne0
 \quad\text{in}\quad
 \Omega^2_{B_{\mathrm{ell}}/W}[1/5]\big/
       d\Omega^1_{B_{\mathrm{ell}}/W}[1/5].
\]
Nonvanishing persists under every finite extension of \(K\), with
continuous relative forms over the extended integer ring.
\end{proposition}

\begin{proof}
Choose the proper slope-zero generator \(\eta\) and slope-one
generator \(\omega\) over \(\DtwoQQ_5\), and normalize
\(\DtwoTr(\eta\cup\omega)=1\).
The curve is defined over \(\DtwoFF_5\).
Proper smooth crystalline/de Rham comparison gives corresponding
algebraic de Rham classes on \(\DtwocE_K\).
Their restrictions to \(\DtwocE_K^\circ\) have algebraic one-form
representatives because the curve is affine.
These also define dagger and continuous representatives, so
Proposition~\ref{d2:prop:contraction} applies.

With the principal-polarization convention, the normalized
Poincar\'e bundle is
\[
 \DtwocP=\DtwocO(\Delta-\DtwocE\times O-O\times\DtwocE).
\]
Its restrictions to the origin slices are trivial, so its Chern
class has only a mixed term in the proper K\"unneth decomposition.
The diagonal acts as the identity on \(H^1\).
With our pairing convention, the mixed tensor is
\[
 \omega\otimes\eta-\eta\otimes\omega.
\]
Multiplying by \(\eta\) on the first factor and tracing returns
\(\eta\) on the second. Multiplying by \(\omega\) returns
\(\omega\). Perfectness of the pairing determines the tensor.
The trace, duality, K\"unneth, and proper crystalline compatibilities
used here are those of \cite[\S1.7(i), \S2.3(i), Theorems~2.4 and~3.2, and \S3.3(i)]{Berthelot1997Duality}.

The graph line for a map \(q\) is the pullback of the diagonal line
by \(\Dtwoid\times q\).
In the difference for \(\Dtwowt Q\) and \(-\Dtwowt Q\), the first-factor
origin line cancels, and the second-factor terms cancel because
\([-1]^*\DtwocO(O)\simeq\DtwocO(O)\).
Inversion acts by \(-1\) on proper \(H^1\).
Thus
\begin{equation}\label{d2:eq:graph-class}
 c_1^\DtwodR(\DtwocL_Q)
 =
 \left[2\bigl(\omega\wedge\Dtwowt Q^*\eta
             -\eta\wedge\Dtwowt Q^*\omega\bigr)\right].
\end{equation}
This equality first holds on the finite-type affine product over
\(K\). The difference from another Chern representative is the
differential of an algebraic one-form there.
It remains exact under weak completion, generic localization, and
continuous completion. No K\"unneth formula for the completed
product is being used.

Compactify \(U\) to its smooth proper connected curve \(\bar U\).
Properness extends \(Q\) to
\(\bar Q\colon\bar U\to E_0\).
The map \(\bar Q\) is finite and flat of rank
\(m=[k(\bar U):k(E_0)]>0\): its finite direct image is
torsion-free over the local discrete valuation rings of \(E_0\).
This rank includes the inseparable degree.
Write \(\xi_C\) for the rigid homological fundamental class of a
smooth proper curve \(C\), viewed as its top-degree trace functional.
Proper cycle pushforward gives
\(\bar Q_*[\bar U]=m[E_0]\); compatibility of cycle classes with
proper pushforward therefore gives
\(\bar Q_*\xi_{\bar U}=m\xi_{E_0}\)
\cite[Proposition~2.2, Section~2.1, and Proposition~6.4]{Petrequin2003}.
The homological pushforward is the transpose of cohomological
pullback. Consequently, for every
\(u\in H^2_{\mathrm{rig}}(E_0/K)\),
\[
 \DtwoTr_{\bar U}(\bar Q^*u)=m\DtwoTr_{E_0}(u).
\]
Taking \(u=a\cup b\) and using the naturality of cup products gives
\[
 \DtwoTr_{\bar U}(\bar Q^*a\cup\bar Q^*b)
   =m\,\DtwoTr_{E_0}(a\cup b).
\]
The pairing is perfect and \(m\) is invertible in \(K\).
Consequently \(\bar Q^*\) is injective on degree-one cohomology,
without a separability assumption.
Restriction from \(\bar U\) to \(U\) is also injective.
If the reduced finite complement \(Z=\bar U\setminus U\) is
nonempty, it has codimension one in the smooth curve \(\bar U\).
Support vanishing gives \(H^1_{Z,\Dtworig}(\bar U/K)=0\), and
the localization sequence then makes
\(H^1_\Dtworig(\bar U/K)\to H^1_\Dtworig(U/K)\) injective
\cite[Theorem~1.1, third assertion, and Proposition~1.3,
pp.~63--64]{Petrequin2003}.
For an empty complement restriction is the identity.
Thus \(Q^*\eta\ne0\) in \(H^1_\Dtworig(U/K)\), represented
by \(\Dtwowt Q^*\eta\).

The special curve and map descend to some finite field
\(\DtwoFF_{5^d}\).
The four \(\DtwoFF_5\)-points of \(E_0\) are \(O\) and the three
points with \(v=0\), at \(z=0,2,3\).
At \(z=1,4\) the cubic is a nonsquare.
The Frobenius polynomial is therefore \(X^2-2X+5\).
Its reduction has simple roots \(0,2\), so its two roots lie in
\(\DtwoZZ_5\), one a unit \(\alpha\).
We choose \(\eta\) with Frobenius eigenvalue \(\alpha\).

The weak completion \(D^\dagger\) admits an integral, Witt-semilinear
Frobenius lift by the smooth affine lifting statement in
\cite[Introduction, p.~198]{DLZ2011}.
Iterate it \(d\) times to obtain \(\phi\).
Its action on rigid cohomology is canonical and independent of the
lift.
Descent of \(Q\) gives
\(\phi^*[Q^*\eta]=\alpha^d[Q^*\eta]\).
On the affine dagger complex this means
\[
 \phi^*(\Dtwowt Q^*\eta)-\alpha^d\Dtwowt Q^*\eta=dh
 \qquad(h\in D^\dagger[1/5]).
\]
Lemma~\ref{d2:lem:scalar} applies, since \(\alpha^d\in W\).
Hence the image \([\Dtwowt Q^*\eta]\) in the ordinary Cohen one-form
quotient is nonzero.

Apply \(L\) from Proposition~\ref{d2:prop:contraction} to
\eqref{d2:eq:graph-class}.
The result is \(2[\Dtwowt Q^*\eta]\), which is nonzero.
Because \(L\) takes exact forms to exact forms, the graph class
cannot be exact.

For a finite extension \(K'/K\), its integer ring \(W'\) is finite
free over \(W\) and complete.
Scalar extension therefore commutes with the adic limits defining
the product and its finite projective form modules.
Relative forms over \(W'\) are the old forms tensored with \(W'\).
After inversion of five, the complexes are tensored with \(K'\).
Field extension is exact and faithfully flat, so it commutes with
kernels and ordinary images and preserves the nonzero class.
\end{proof}

\section{The quartic quotient and its weight}\label{d2:sec:quartic}

In this section, we construct the elliptic quotient and compute the
linear functional on the four root-divisor classes.

\begin{proposition}\label{d2:prop:weighted}
Let \(k=\overline{\DtwoFF}_5\), \(W=W(k)\), \(R=W[\rho]\), and
\[
 V=\varprojlim_n W[T]_{(5)}/5^n.
\]
There is a finite unramified extension \(C/V\), obtained as the
generic-special completion of a smooth affine \(W\)-curve \'etale
over an open of the \(T\)-line, with the following property.
Put \(C'=C\otimes_WR\) and
\[
 B=C'\langle X,Y\rangle/
 \left(X^4+\frac{\Phi_5(1+\pi Y)}{\pi^4}+TXY^3\right).
\]
For \(g\in G=\DtwoFF_5^\times\), set
\[
 a_g=\frac{\rho^g-1}{\pi},\qquad
 J_g=(X,Y-a_g),\qquad
 \gamma_g=-c_1^\DtwodR(J_g).
\]
The ideals \(J_g\) are invertible.
Let \(E=\DtwoQQ_5(\rho)\), let
\(\omega_T\colon G\to\DtwoZZ_5^\times\) be the Teichm\"uller character,
and put
\[
 e_r=\frac14\sum_{g\in G}\omega_T(g)^{-r}\sigma_g
 \qquad(r=1,3)
\]
on \(\DtwoQQ_5\)-vector-space targets.
If \(b\in E\) satisfies
\(\sigma_{-1}b=-b\) and \(e_1b,e_3b\ne0\), then
\[
 \sum_{g\in G}\sigma_g(b)\gamma_g\ne0
 \quad\text{in}\quad
 \Omega^2_{B/R}[1/5]\big/d\Omega^1_{B/R}[1/5].
\]
The relative forms retain \(dT\).
\end{proposition}

\begin{proof}
Over \(k(T)\), consider
\[
 H\colon X^4+Y^4+TXY^3=Z^4,\qquad
 H^\circ=H\cap\{Z\ne0\},\qquad
 P_g=[0:g:1].
\]
The singular-point calculation in
Proposition~\ref{d2:prop:geometry} applies with \(Z\) in place of \(U\).
It gives \(T^4=3\) at a hypothetical singularity.
Thus \(H\) is smooth and geometrically integral.

The same auxiliary surface
\[
 \DtwoSpec k[X,U,Y,Y^{-1}]/(XU-(1-Y^4)),
 \qquad T=(U-X^3)/Y^3,
\]
has root-divisor lattice \(\DtwoZZ^4/\DtwoZZ(1,1,1,1)\).
For clarity, the mechanisms remain unchanged here: inversion of \(X\)
gives the factorial Laurent ring; its units give only the diagonal
relation; each fiber becomes
\(w^4=s^4+ts+1\), whose right side is not a square by comparison
with \((s^2+bs+c)^2\); hence each vertical divisor is the reduced
principal divisor \(T-t\).
Divisor localization therefore gives this same lattice on
\(H^\circ\cap\{Y\ne0\}\).
A torsion relation on \(P_1-P_{-1}\) in the proper Jacobian would
restrict to a torsion relation in this lattice, which is impossible.
Finite field extension preserves nontorsion: norm after pullback multiplies divisor classes by the extension degree.

We construct the quotients over an \'etale parameter cover.
Set \(\Delta=256-27T^4\), and order the roots
\(b_1,\ldots,b_4\) of \(s^4+Ts+1\) over a finite \'etale cover of
\(W[T,1/\Delta]\).
Put
\[
 a=\frac{b_2-b_4}{b_2-b_1},\qquad
 x=a\frac{s-b_1}{s-b_4},\qquad
 \lambda=a\frac{b_3-b_1}{b_3-b_4}.
\]
The elements \(a,a-1,a-\lambda,\lambda,\lambda-1\) are units,
as follows from the pairwise root differences.
Adjoin \(c\) with \(c^4=a(a-1)(a-\lambda)\).
Replacing \(w=Z/Y\) by \(cw/(s-b_4)\) gives
\begin{equation}\label{d2:eq:cyclic-quartic}
 w^4=x(x-1)(x-\lambda).
\end{equation}
Indeed the product of the three transformed factors is
\(a(a-1)(a-\lambda)\prod_{j=1}^3(s-b_j)/(s-b_4)^3\),
which is \(c^4(s^4+Ts+1)/(s-b_4)^4\).

Adjoin \(r\) with \(r^2=\lambda\) and
\(\alpha_\pm\) with \(\alpha_\pm^2=r\pm1\).
All these extensions are \'etale, because the radicands and \(2,4\)
are units.
Set
\[
 z=\frac{w^2}{x},\qquad v_\pm=\frac{w(x\pm r)}x.
\]
Equation~\eqref{d2:eq:cyclic-quartic} gives
\begin{equation}\label{d2:eq:elliptic-quotients}
 z^2=x+\frac{r^2}{x}-1-r^2,\qquad
 v_\pm^2=z^3+(r\pm1)^2z.
\end{equation}
After \(z=\alpha_\pm^2\zeta\), \(v_\pm=\alpha_\pm^3v\),
the target is the constant elliptic curve \(v^2=\zeta^3+\zeta\).

Let \(\tau(x,w)=(x,-w)\) and
\[
 \iota(x,w)=\left(\frac{r^2}{x},\frac{rw}{x}\right).
\]
The maps \(h_+,h_-\) are the quotients by \(\iota\) and
\(\tau\iota\), respectively.
Their degree is two: \(x\) satisfies a quadratic over the target
functions by \eqref{d2:eq:elliptic-quotients}, \(w\) is recovered from
\(x,v_\pm\), and the displayed nontrivial involution fixes the
target functions.
They extend between the proper smooth curves.

The boundary \(Z=0\) consists of the four points
\(x=0,1,r^2,\infty\).
At \(0,\infty\), the orders of \(z,v_\pm\) are \((-2,-3)\), so
the image is \(O\) with local degree one.
At \(1,r^2\), these orders are \((2,1)\), so the image is
\((0,0)\), again with local degree one.
Everywhere else \(z\) is finite and nonzero.
The exact inverse image of \(\{O,(0,0)\}\) is therefore the boundary.
Both maps restrict to finite maps of all of \(H^\circ\) to
\(E_0^\circ\).

The equations are defined over \(W\).
Shrinking the parameter open by elements nonzero at its generic
special point makes these finite flat quotient maps between smooth
relative curves.
Choose a connected component of the special fiber of the finite
\'etale parameter cover.
It can be isolated in a smooth affine \(W\)-model by inverting a
lift of its characteristic idempotent.
Write \(D\) for the resulting smooth \(W\)-algebra.
Its generic-special completion \(C\) is the corresponding finite
unramified extension of \(V\); subsequent open restrictions are
units in \(C\).

The involutions \(\tau,\iota,\tau\iota\) form a Klein four group.
Its total quotient is rational with coordinate \(z=w^2/x\).
Recovering \(x\) and then \(w\) gives degree at most four, and the
four automorphisms fix \(z\), so \(k(T)(z)\) is the invariant field.
If \(q\) is the quotient by \(\tau\), the pullback and norm identities
on Jacobians give
\begin{equation}\label{d2:eq:klein-identity}
 q^*q_*+h_+^*h_{+*}+h_-^*h_{-*}
   =2\Dtwoid+\sum_{g\in\{1,\tau,\iota,\tau\iota\}}g
   =2\Dtwoid.
\end{equation}
The full group norm is zero because it factors through the Jacobian
of the rational total quotient.
The class \(P_1-P_{-1}\) is killed by \(q_*\).
If its images under both \(h_+\) and \(h_-\) were torsion,
\eqref{d2:eq:klein-identity} would make it torsion.
Choose a quotient \(h\) for which the image is nontorsion.
Since \(h\tau=[-1]h\), this image is \(2Q\), where \(Q=h(P_1)\).
Every point in \(E_0(k)\) is defined over a finite field and is
torsion. Hence \(Q\) is nonconstant.
The lifted root section gives
\(\Dtwowt Q\colon\DtwoSpec D\to\DtwocE^\circ\) with this reduction.

Let \(i\in W\) be the fourth root of unity reducing to \(2\).
The rotation
\([X:Y:Z]\mapsto[iX:iY:Z]\) fixes \(s\), sends the transformed
\(w\) to \(-iw\), sends \(z\) to \(-z\), and sends \(v_\pm\) to
\(-iv_\pm\).
Thus it induces \([-i]\) on the elliptic target.
It sends \(P_1\) to \(P_i\), so
\begin{equation}\label{d2:eq:root-equivariance}
 h(P_g)=[g^{-1}]Q
\end{equation}
for all fourth roots \(g\), including the lifted Teichm\"uller
sections.

We need an integral trace on total forms, including base
differentials.
For a tame double quotient of smooth relative curves, the trace is
the \'etale trace away from ramification.
At ramification, choose an anti-invariant fiber parameter \(t\);
\'etale-locally the quotient parameter is \(s=t^2\).
Invariant fiber one-forms are multiples of
\(t\,dt=ds/2\).
Invariant base forms and their wedges also descend.
Thus invariant total forms are exactly the forms on the quotient.
Summing \(1+\iota^*\) and descending defines a chain map
\[
 \DtwoTr_hh^*=2,\qquad h^*\DtwoTr_h=1+\iota^*.
\]
It is integral and continuous, and extends through localization and
completion, first relative to \(W\) and then relative to \(R\).
For an invertible module \(L\), the Cartier-divisor norm identity is
\(h^*\DtwoNm_hL=L\otimes\iota^*L\), including ramification
multiplicities.
Since \(h^*\) is injective on rational cohomology by
\(\DtwoTr_hh^*=2\), it follows that
\begin{equation}\label{d2:eq:trace-chern}
 \DtwoTr_h(c_1^\DtwodR(L))=c_1^\DtwodR(\DtwoNm_hL).
\end{equation}
This is an identity on ordinary de Rham quotients, because the trace
is an actual chain map.

Consider first
\[
 B_{\Dtwocyc}=
 C'\langle X,Y\rangle/(X^4+Y^4+TXY^3-1),
 \qquad
 \gamma_g^\Dtwocyc=c_1^\DtwodR(\DtwocO(P_g)).
\]
The norm of its root line under \(h\) is the graph line of
\([g^{-1}]\Dtwowt Q\).
Choose \(\eta,\omega\) with the polarization convention of
\eqref{d2:eq:graph-class}.
The unit-root character can be determined exactly.
On the elliptic curve,
\[
 [i](z,v)=(-z,iv),\qquad
 [i]^*\!\left(\frac{dz}{2v}\right)
       =i\,\frac{dz}{2v}.
\]
By the perfect pairing, the character on \(H^1(\DtwocO)\) is
\(i^{-1}\).
A primitive integral generator of the unit-root crystalline line
has nonzero image in \(H^1(\DtwocO)\) modulo five.
Otherwise its nonzero reduction would lie in the holomorphic
subspace, which crystalline Frobenius kills modulo five,
contradicting its unit Frobenius eigenvalue.
The two order-four eigenvalues have distinct reductions.
Consequently the character on the unit-root line is exactly
\(\omega_T^3\), rather than an unspecified odd character.

Let \(\lambda_g=\sigma_g(b)\) for any anti-invariant \(b\).
Then \(\lambda_{-g}=-\lambda_g\).
Pair opposite roots and apply \eqref{d2:eq:trace-chern}, followed by
the contraction \(L\) of Proposition~\ref{d2:prop:contraction}.
The graph-difference formula \eqref{d2:eq:graph-class} contributes
twice the pullback of \(\eta\).
Equation~\eqref{d2:eq:root-equivariance} gives
\([g^{-1}]^*\eta=\omega_T(g)^{-3}\eta\).
Hence the exact linear identity is
\begin{equation}\label{d2:eq:linear-weight}
 L\DtwoTr_h\left(\sum_{g\in G}\sigma_g(b)\gamma_g^\Dtwocyc\right)
   =
 \left(\sum_{g\in G}\sigma_g(b)\omega_T(g)^{-3}\right)
       [\Dtwowt Q^*\eta]
   =4e_3(b)[\Dtwowt Q^*\eta].
\end{equation}
The factor two from the graph difference is exactly the contribution
of the two opposite roots in the full sum.
The polarization convention fixes the common sign in this formula.

Proposition~\ref{d2:prop:graph}, using
Lemma~\ref{d2:lem:scalar}, shows that \([\Dtwowt Q^*\eta]\) is nonzero
in the ordinary Cohen one-form quotient.
Its finite-constant-extension assertion preserves this class over
\(R\).
Thus \eqref{d2:eq:linear-weight} is nonzero when \(e_3(b)\ne0\);
in particular it is nonzero under the two nonvanishing assumptions
of the proposition.
Both operators preserve exactness, so the weighted class on
\(B_{\Dtwocyc}\) is nonzero.
Formula~\eqref{d2:eq:linear-weight} remains valid for all anti-invariant \(b\), even if an odd component vanishes.

Finally we transport all four root lines to \(B\).
The cyclotomic identity gives \(5/\pi^4\equiv-1\bmod\pi\), and all
other coefficients below degree four in
\(\Phi_5(1+\pi Y)/\pi^4\) are divisible by \(\pi\).
Thus \(B\) and \(B_{\Dtwocyc}\) have the same special fiber
\[
 X^4+Y^4+TXY^3=1.
\]
For \(G_0=X^4+Y^4+TXY^3\), Euler's identity
\(X(G_0)_X+Y(G_0)_Y=4G_0\) proves smoothness on this fiber.
The two equations are monic in \(Y\), so the models are
\(\pi\)-torsion-free. Their completions are formally smooth over
\(C'\).

Lift the identity modulo \(\pi\) successively modulo \(\pi^n\).
This gives a continuous \(C'\)-linear map between the two completed
models.
Successive approximation proves surjectivity.
An element of its kernel lies in \(\pi\) times the source;
torsion-freeness of the target permits division of the relation by
\(\pi\). Iteration and separatedness prove injectivity.
Thus the map is an isomorphism.

On the support of \(J_g\), the \(Y\)-derivative is
\(\prod_{h\ne g}(a_g-a_h)\), a unit.
The equation expresses \(Y-a_g\) locally as a multiple of \(X\);
elsewhere \(J_g\) is the unit ideal.
Hence \(J_g\) is invertible.
Its reduction is \((X,Y-g)\), equal to the reduction of the cyclic
root ideal.
In a complete ring, \(\pi\) lies in the Jacobson radical.
An isomorphism between the reductions of two projective rank-one
modules lifts by projectivity; Nakayama at every maximal ideal
makes the lift an isomorphism.
Therefore the formal replacement identifies every root Picard class.

The replacement fixes \(C'\), so it fixes all \(\sigma_g(b)\).
It induces an isomorphism of continuous forms relative to \(R\),
with \(dT\) retained.
The weighted nonzero class, and the transported linear functional
\eqref{d2:eq:linear-weight}, therefore give the assertion for \(B\).
\end{proof}

\section{An actual arithmetic coefficient}\label{d2:sec:arithmetic}

The goal of this section is to construct one actual coefficient whose
two odd point-character components are nonzero.

\begin{proposition}\label{d2:prop:coefficient}
Put \(W_E=\DtwoZZ_5[\rho]\), \(E=\DtwoQQ_5(\rho)\), and
\(\pi=\rho-1\).
Define
\[
 r_E\colon K_3(W_E)\longrightarrow E
\]
by passage to \(\pi_3\DtwocK(W_E)\), the inverse of
\(\pi_3\DtwocK(W_E,(\pi))\to\pi_3\DtwocK(W_E)\), and the relative point
equivalence with \(\DtwoHC_2^\Dtwocts(W_E/W_E;\DtwoQQ_5)=E\).
There is an actual class
\[
 \beta=(1-\sigma_{-1})\beta_0\in K_3(W_{\Dtwoar})
\]
whose reduction in \(K_3(\DtwoFF_5)\) is zero, such that
\(\sigma_{-1}\beta=-\beta\) integrally and
\[
 b=r_E(\beta|_{W_E})
 \quad\text{satisfies}\quad e_1b\ne0,\qquad e_3b\ne0.
\]
The projectors \(e_r\) act only on the indicated
\(\DtwoQQ_5\)-vector-space targets.
\end{proposition}

\begin{proof}
For \(r=1,3\), set
\[
 \tau_r=\sum_{g=1}^4\omega_T(g)^{-r}\rho^g.
\]
Its \(\pi\)-adic valuation is \(r\).
To see this, expand \(\rho^g=(1+\pi)^g\).
For \(j<r\), the coefficient
\(\sum_g\omega_T(g)^{-r}\binom gj\) vanishes modulo five:
\(\binom gj\) has degree \(j\) in \(g\), and the power sums on
\(\DtwoFF_5^\times\) vanish unless the exponent is divisible by four.
For \(j=r\), the coefficient is \(4/r!\bmod5\), which is nonzero.
Since \(5\) is a unit times \(\pi^4\), terms with \(j<r\)
have valuation at least \(4+j>r\), while terms with \(j>r\)
have valuation at least \(j>r\).
The degree-\(r\) term is the unique one of minimal valuation.

Let \(\DtwoLi_2\) be the Coleman dilogarithm normalized by
\(\sum_{m\ge1}z^m/m^2\) on the open unit disc.
Define
\[
 S_j(z)=\left(z\frac d{dz}\right)^j\frac1{1-z},
 \qquad
 F_2(z)=\DtwoLi_2(z)-5^{-2}\DtwoLi_2(z^5).
\]
Partitioning the indices prime to five into \(a+5m\), and expanding
\((a+5m)^{-2}\), gives on that disc
\begin{equation}\label{d2:eq:depletion}
 F_2(z)=
 \sum_{j\ge0}(-1)^j(j+1)5^jS_j(z^5)
       \sum_{a=1}^4\frac{z^a}{a^{j+2}}.
\end{equation}
The right side converges uniformly on
\(|z|\le1,\ |1-z^5|=1\).
Each \(S_j\) has integral coefficients and unit denominator there,
so its \(j\)-th term has norm at most \(5^{-j}\).
This affinoid is connected: its reduction is the integral curve
\(\DtwoSpec\DtwoFF_5[z,(1-z)^{-1}]\), and its integral affinoid algebra
has no nontrivial idempotent.

By \cite[Theorem~2.7(3)]{BesserDeJeuAlgorithm}, the depleted polylogarithm
\(\DtwoLi_n(z)-p^{-n}\DtwoLi_n(z^p)\) is analytic outside the closed disc
\(|z-1|\le p^{-1/(p-1)}\), with a convergent expression in
\(1/(1-z)\).
For \(p=5\), this region contains the affinoid just used.
Both sides of \eqref{d2:eq:depletion} are analytic there.
Their equality on the open unit disc extends by the identity
principle.

At \(z=-\rho\), \(z^5=-1\).
The value \(\DtwoLi_2(-1)\) is fixed by \(G\), so its odd projections
vanish.
Permuting the summation indices gives the exact identity
\[
 e_r(\rho^a)=\frac{\omega_T(a)^r}{4}\tau_r.
\]
Applying it to \eqref{d2:eq:depletion} yields
\begin{equation}\label{d2:eq:odd-dilogarithms}
 \frac{4e_r\DtwoLi_2(-\rho)}{\tau_r}
 =
 \sum_{j\ge0}(-1)^j(j+1)5^jS_j(-1)
    \sum_{a=1}^4\frac{(-1)^a\omega_T(a)^r}{a^{j+2}}.
\end{equation}
All \(S_j(-1)\) are \(5\)-integral.
Modulo five, only the \(j=0\) term remains:
\[
 \frac12\sum_{a=1}^4(-1)^aa^{r-2}.
\]
For \(r=1\) it is
\(\frac12(-1+3-2+4)=2\) in \(\DtwoFF_5\).
For \(r=3\) it is
\(\frac12(-1+2-3+4)=1\).
Thus both odd projections of \(\DtwoLi_2(-\rho)\) are nonzero.
Division by \(\tau_r\) is legitimate by its valuation, and the
projection identity is exact, so no precision loss is involved.

The special-unit symbol construction
\cite[Theorems~1.6 and~1.10, Remark~1.7]{BesserDeJeu}
maps the degree-one, weight-two symbol complex of a localized
number ring to rational K-theory.
Its syntomic regulator is
\([x]_n\mapsto\pm(n-1)!L_{\mathrm{mod},n}(x)\);
the sign is the single convention in the cited remark.
In weight two there is no conjectural vanishing assumption.

Take \(x=-\rho\) in \(W_{\Dtwoar}\).
Both \(x\) and \(1-x=1+\rho\) are units, since the latter reduces
to \(2\) modulo \(\pi\).
The symbol \([-\rho]_2\) is a rational degree-one cocycle:
its differential is \((1+\rho)\otimes(-\rho)\), or the exterior
version, and the root-of-unity factor is zero after rationalization.
It gives \(\beta_{\mathrm{rat}}\in K_3(W_{\Dtwoar})\otimes\DtwoQQ\).
Its syntomic regulator is \(\pm\DtwoLi_2(-\rho)\), because the
modification terms contain \(\log(-\rho)=0\).
Choose a positive integer \(N\) and an actual \(\beta_0\)
representing \(N\beta_{\mathrm{rat}}\), and set
\(\beta=(1-\sigma_{-1})\beta_0\).
This is an operation on actual groups.
The automorphism \(\sigma_{-1}\) induces the identity on the residue
field, so \(\beta\) reduces to zero.
Its integral antisymmetry follows by squaring \(\sigma_{-1}\).
Equivariance of the regulator and
\eqref{d2:eq:odd-dilogarithms} show that its two odd syntomic
components are nonzero.

For the possibly ramified field \(E\), the syntomic-to-\'etale
comparison in \cite[Section~2]{HuberKings} has no ramification
restriction.
The degree-one syntomic coordinate for weight \(n>1\) is \(E\)
\cite[Example~2.2.4]{HuberKings}.
It agrees with the definition used above
\cite[Proposition~2.2.7]{HuberKings}.
Its Chern-compatible map to \(H^1(E,\DtwoQQ_5(n))\) is the
Bloch--Kato exponential
\cite[Propositions~2.2.9 and~2.3.4]{HuberKings}, which is an
isomorphism for \(n>1\), as recalled in
\cite[Section~1.3]{HuberKings}.
These maps are natural under field automorphisms.
Thus the actual \'etale Chern class
\(c_{2,1}(\beta)\in H^1(E,\DtwoQQ_5(2))\) has both odd components
nonzero.

There is a natural \(G\)-equivariant \(\DtwoQQ_5\)-linear map
\begin{equation}\label{d2:eq:completed-chern}
 \pi_3\DtwocK(W_E)\longrightarrow H^1(E,\DtwoQQ_5(2))
\end{equation}
through which the Chern class of every actual element factors.
To construct this map, first work before inverting five.
Map from \(\pi_3(K(W_E)^\wedge_5)\) to \(\pi_3(K(E)^\wedge_5)\), and then to
\(\varprojlim_nK_3(E;\DtwoZZ/5^n)\).
Apply the finite-coefficient Chern classes
\[
 K_3(E;\DtwoZZ/5^n)\longrightarrow
 H^1(E,\mu_{5^n}^{\otimes2}).
\]
Naturality, compatibility with reduction of integral classes,
additivity for odd coefficients, and compatibility with coefficient
transitions are respectively
\cite[Proposition~2.1.1, Lemma~2.3, Propositions~2.4 and~2.8]{WeibelChern}.
The resulting compatible additive maps recover the actual Chern
class. Each finite target is killed by \(5^n\), so the inverse-limit
map is \(\DtwoZZ_5\)-linear.
The groups \(H^0(E,\mu_{5^n}^{\otimes2})\) are finite, hence
Mittag--Leffler. The continuous-cohomology inverse-limit sequence
therefore identifies the target limit with \(H^1(E,\DtwoZZ_5(2))\).
Inverting five gives \eqref{d2:eq:completed-chern}.
No derived inverse-limit vanishing on the K-theory side is needed.

The finite-field calculation \cite{QuillenFinite}\cite[IV, Corollary~1.13]{Weibel} gives
\[
 K_{2j}(\DtwoFF_5)=0,\qquad
 K_{2j-1}(\DtwoFF_5)=\DtwoZZ/(5^j-1)\qquad(j\ge1).
\]
Multiplication by five is invertible on these positive groups.
The finite-coefficient exact sequences show that all positive
homotopy groups of \(K(\DtwoFF_5)^\wedge_5\) vanish.
In particular the groups in degrees three and four of
\(\DtwocK(\DtwoFF_5)\) vanish.
The relative fiber sequence identifies
\(\pi_3\DtwocK(W_E,(\pi))\) with \(\pi_3\DtwocK(W_E)\).

The point case of
\cite[Theorem~A, Proposition~4.10, Corollary~3.8]{AMMN}
identifies this group with
\(\DtwoHC_2^\Dtwocts(W_E/W_E;\DtwoQQ_5)=E\).
Here \(W_E\) is a complete mixed-characteristic discrete valuation
ring with perfect finite residue field, its formal spectrum is
smooth over itself, and changing \((5)\) to \((\pi)\) changes the
special fiber only by a nilpotent thickening.
The comparison is natural for automorphisms of the coefficient
point, hence \(G\)-equivariant.

Transporting \eqref{d2:eq:completed-chern} through this point
identification gives a \(G\)-equivariant \(\DtwoQQ_5\)-linear map from
\(E\) to \(H^1(E,\DtwoQQ_5(2))\), carrying \(b\) to
\(c_{2,1}(\beta)\).
If \(e_rb=0\), the corresponding Chern component would vanish.
This contradicts the syntomic calculation for \(r=1,3\).
Equivariance also gives \(\sigma_{-1}b=-b\).
No equality between a chosen syntomic normalization and the point
coordinate was required.
\end{proof}

\section{Assembly over the local ring and its completion}\label{d2:sec:assembly}

In this section, we prove Theorem~\ref{d2:thm:completed} and show that its
actual coefficient also gives the geometric detector.

\begin{proof}[Proof of Theorem~\ref{d2:thm:completed}]
Choose the actual \(\beta\) of Proposition~\ref{d2:prop:coefficient},
and set \(c=c(\beta)\).
Section~\ref{d2:sec:preliminaries} proves regularity of \(A,\Dtwowh A\)
and both integral generic vanishings.
We prove that the image of this same \(c\) in \(K_3(\Dtwowh A)\) has
infinite order.

Put \(k=\overline{\DtwoFF}_5\), \(W=W(k)\), and \(R=W[\rho]\).
Choose \(C\), arising from a smooth affine curve \(D\), as in
Proposition~\ref{d2:prop:weighted}.
Write \(L\) for the function field of \(D/5D\).
It is a finite separable extension of \(k(T)\).
The coefficient ring \(C'=C\otimes_WR\) is a complete discrete
valuation ring with uniformizer \(\pi\) and residue field \(L\).
Use the test algebra
\begin{equation}\label{d2:eq:test-algebra}
 B=C'\langle X,Y\rangle/
 \left(X^4+\frac{\Phi_5(1+\pi Y)}{\pi^4}+TXY^3\right).
\end{equation}
The polynomial is integral over \(R\): every coefficient below
degree four before division has \(\pi\)-valuation at least four.
It is monic in \(Y\), so the polynomial quotient and its completion
are \(\pi\)-torsion-free.

Its special fiber is
\begin{equation}\label{d2:eq:test-special}
 B/\pi B=L[X,Y]/(X^4+Y^4+TXY^3-1).
\end{equation}
For \(G_0=X^4+Y^4+TXY^3\), Euler's identity gives
\(X(G_0)_X+Y(G_0)_Y=4G_0=4\) on this fiber.
It is therefore smooth over \(L\).
Its projective closure is the smooth geometrically integral quartic
of Proposition~\ref{d2:prop:weighted}, so the fiber is a domain.

The ring \(B\) is noetherian by its monic presentation and adic
completion. Since \(\pi\) lies in its Jacobson radical, all maximal
ideals contain \(\pi\).
At each one, the quotient by \(\pi\) is regular and \(\pi\) is a
nonzero divisor. Lifting a regular system of parameters from the
quotient and adjoining \(\pi\) proves regularity of that local ring.
Hence \(B\) is regular.
Its dimension is two: the restricted two-variable ring over \(C'\)
has dimension three, and the equation is a nonzero monic equation.
Equivalently, the parameter argument gives dimension at most two
at every maximal ideal and equality at closed points of the special
curve.
Separatedness, torsion-freeness, and the domain special fiber show
that \(B\) is a domain: the first nonzero \(\pi\)-adic terms of a
product of two nonzero elements have nonzero product.

We construct the essentially smooth rank-two model required in
Proposition~\ref{d2:prop:cyclic}.
Let \(U_D=D_{5D}\), whose completion is \(C\), and put
\[
 H=
 \bigl((U_D\otimes_WR)[X,Y]\bigr)/
 \left(X^4+\frac{\Phi_5(1+\pi Y)}{\pi^4}+TXY^3\right).
\]
Its completion is \(B\).
The two relative partial derivatives generate the unit ideal
modulo \(\pi\) by the Euler calculation.
Choose an element \(j\) of their generated ideal with
\(j\equiv1\bmod\pi\).
Then \(H[1/j]\) is smooth of relative dimension one over
\(U_D\otimes_WR\).
Inverting \(j\) does not change the completion.

The ring \(U_D\otimes_WR\) is a localization of the smooth
\(R\)-algebra \(D\otimes_WR\), of relative dimension one.
Transitivity of smoothness makes \(H[1/j]\) a localization of a
smooth \(R\)-algebra with relative cotangent rank two.
The finite coefficients and the choice of \(j\) spread over a
further localization of \(D\).
This gives the required \(B_0\).
The \(5\)-adic and \(\pi\)-adic completions agree because \(5\)
is a unit times \(\pi^4\).
In particular, relative forms over \(R\) retain the differential
\(dT\).

Write \(b=r_E(\beta|_{W_E})\).
Proposition~\ref{d2:prop:weighted} gives invertible ideals
\[
 J_g=\left(X,Y-\frac{\rho^g-1}{\pi}\right)
\]
and the nonzero class
\begin{equation}\label{d2:eq:completed-form}
 \Xi=\sum_{g\in G}\sigma_g(b)\bigl(-c_1^\DtwodR(J_g)\bigr)
 \ne0
 \quad\text{in}\quad
 \Omega^2_{B/R}[1/5]\big/d\Omega^1_{B/R}[1/5].
\end{equation}

The reduction of \(\beta|_{W_E}\) in \(K_3(\DtwoFF_5)\) is zero.
It lifts to an actual
\(\beta_{\mathrm{rel}}\in K_3(W_E,(\pi))\).
This lift is unique, since \(K_4(\DtwoFF_5)=0\) by
\cite{QuillenFinite}\cite[IV, Corollary~1.13]{Weibel}.
Naturality makes the lifts of all conjugates compatible.
Pull them through \(W_E\to R\).
The relative point coordinates are \(\sigma_g(b)\) in \(R[1/5]\):
the point comparison is natural for this coefficient extension,
and the map on degree-two point cyclic homology is inclusion of
the coefficient fields.

Form the actual relative class
\[
 \delta_{\mathrm{rel}}
   =\sum_{g\in G}
      (\sigma_g\beta_{\mathrm{rel}})|_B\cdot(1-[J_g])
       \in K_3(B,\pi B),
\]
and let \(\delta\) be its image in actual \(K_3(B)\).
Proposition~\ref{d2:prop:cyclic}, applied to the rank-two model above,
identifies its top-form component with \(\lambda_R\Xi\), where
\(\lambda_R\ne0\).
Thus the image of \(\delta_{\mathrm{rel}}\) in
\(\pi_3\DtwocK(B,\pi B)\) is nonzero.

We establish the special-fiber vanishing before using either fiber
sequence.
The ring in \eqref{d2:eq:test-special} is the filtered colimit of
\[
 (D/5D)[1/s][X,Y]/(X^4+Y^4+TXY^3-1),
 \qquad 0\ne s\in D/5D.
\]
These are smooth affine surfaces over the perfect field \(k\).
The base is a smooth affine curve, and the Euler identity makes
the displayed curve over it smooth.
For a smooth \(d\)-dimensional variety over a perfect field of
characteristic \(p\), the finite-coefficient theorem
\cite[Theorem~8.4]{GeisserLevineCharP} gives
\(K_j(-;\DtwoZZ/p^n)=0\) for \(j>d\).
Its groups are the homotopy groups of the cofiber spectrum, with
their torsion contribution, rather than mere tensor products of
actual K-groups.
The theorem follows there from the motivic-to-K spectral sequence
and the vanishing of logarithmic de Rham--Witt sheaves above the
dimension.
Applying its smooth-surface form and then filtered continuity gives
\[
 K_j(B/\pi B;\DtwoZZ/5^n)=0
 \qquad(j>2,\ n\ge1).
\]
For \(j\ge3\), both the degree-\(j\) and degree-\((j+1)\) groups
of the finite-coefficient tower vanish.
The Milnor exact sequence therefore gives
\(\pi_jK(B/\pi B)^\wedge_5=0\), and inversion of five preserves
this vanishing. In particular,
\begin{equation}\label{d2:eq:special-vanishing}
 \pi_3\DtwocK(B/\pi B)=\pi_4\DtwocK(B/\pi B)=0.
\end{equation}

First use the reduction fiber sequence
\begin{equation}\label{d2:eq:reduction-fiber}
 \DtwocK(B,\pi B)\longrightarrow\DtwocK(B)
                  \longrightarrow\DtwocK(B/\pi B).
\end{equation}
Equation~\eqref{d2:eq:special-vanishing} makes the first map an
isomorphism on degree three.
The image of \(\delta\) in \(\pi_3\DtwocK(B)\) is consequently nonzero.

Next use regularity of \(B\), regularity of \(B/\pi B\), and
the nonzero divisor \(\pi\).
Localization and d\'evissage
\cite[V.4.1, V.6.1]{Weibel} give the different fiber sequence
\begin{equation}\label{d2:eq:localization-fiber}
 \DtwocK(B/\pi B)\longrightarrow\DtwocK(B)
                         \longrightarrow\DtwocK(B[1/\pi]).
\end{equation}
The left map is support pushforward, not reduction.
Its source has zero degree-three group, so the second map is
injective on degree three.
Thus \(\delta|_{B[1/\pi]}\) has nonzero image in
\(\pi_3\DtwocK(B[1/\pi])\).
That target is a \(\DtwoQQ_5\)-vector space, so the actual class
\(\delta|_{B[1/\pi]}\) has infinite order.
Only the fact that torsion maps to zero in a characteristic-zero
vector space is used here.

There is a ring map \(A\to B\) defined by
\[
 x\longmapsto\pi X,\qquad y\longmapsto\pi Y,
 \qquad T\longmapsto T.
\]
The defining equation maps to \(\pi^4\) times the equation of \(B\).
Every element inverted in forming \(A\) reduces to a nonzero
polynomial in \(T\) modulo \((5,x,y)\).
Its image modulo \(\pi\) is a nonzero element of \(L\), hence a
unit. Since \(\pi\) lies in the Jacobson radical, its image is
a unit in \(B\).
This verifies the map on the localization.
It takes \((x,y)\) into \(\pi B\), so completeness extends it
continuously to \(\Dtwowh A\to B\).

We match the actual supported class only after a flat splitting.
Set
\[
 S_{\mathrm{spl}}=
 (A\otimes_{\DtwoZZ_{(5)}}W_{\Dtwoar})[1/5].
\]
This is flat over \(A\), and regular because it is finite \'etale
over \(A[1/5]\).
The polynomial \(P(y)\) splits there with roots \(\rho^g-1\),
whose differences are units.
The divisor \(x=0\) is the disjoint union of the Cartier divisors
with ideals
\[
 I_g=(x,y-(\rho^g-1)).
\]
Near a component, the other factors of \(P\) are units, and the
equation expresses \(y-(\rho^g-1)\) as a multiple of \(x\).
Away from that component the ideal is the unit ideal.
The element \(x\) remains a nonzero divisor by flatness, and the
component rings are localizations of \(E_{\Dtwoar}[T]\), hence regular.

On the \(g\)-th component, the original coefficient map sends
\(\rho=1+y\) to \(\rho^g\).
Flat base change and the projection formula give the equality of
actual classes
\begin{equation}\label{d2:eq:flat-splitting}
 c|_{S_{\mathrm{spl}}}
   =\sum_{g\in G}(\sigma_g\beta)|_{S_{\mathrm{spl}}}
                         \cdot(1-[I_g]).
\end{equation}
No scalar projector acts on an actual K-group in this identity.

The map \(A\to B\), together with the coefficient embedding,
gives \(S_{\mathrm{spl}}\to B[1/\pi]\).
There
\[
 x=\pi X,\qquad
 y-(\rho^g-1)=\pi(Y-a_g).
\]
The pulled-back invertible module \(I_g\) identifies with
\(J_g[1/\pi]\).
Locally on its support it is generated by \(x\), whose image
\(\pi X\) is a nonzero divisor in the domain \(B[1/\pi]\);
off the support it is generated by a unit.
Thus this identification is an isomorphism of invertible modules,
not only an equality of generated ideals.

Pulling back \eqref{d2:eq:flat-splitting} now gives
\[
 c|_{B[1/\pi]}
 =\sum_g(\sigma_g\beta)|_{B[1/\pi]}
                          (1-[J_g[1/\pi]])
 =\delta|_{B[1/\pi]}.
\]
The right side has infinite order.
The map factors through \(\Dtwowh A\), so it is also the image of
\(\Dtwowh c\). Hence \(\Dtwowh c\), and therefore \(c\), has infinite order.

An element of an abelian group maps to zero after tensoring with
\(\DtwoQQ\) precisely when it is torsion.
The two integral generic vanishings established at the start thus
give the two asserted nonzero rational kernels.

We finish by identifying this coefficient with one allowed in the
geometric proof. The inclusions
\(\DtwoZZ[\rho]\subset\DtwocO_{E_{\Dtwoar}}\subset W_{\Dtwoar}\) show that
\(W_{\Dtwoar}\) is the localization of \(\DtwocO_{E_{\Dtwoar}}\) obtained by
inverting integers prime to five.
Localization, filtered continuity, and
\(K_2\) of finite fields being zero give a surjection
\[
 K_3(\DtwocO_{E_{\Dtwoar}})\longrightarrow K_3(W_{\Dtwoar});
\]
see \cite[V.6.1]{Weibel},
\cite[IV, Section~6.4]{Weibel}, and
\cite{QuillenFinite}\cite[IV, Corollary~1.13]{Weibel}.
Lift the actual \(\beta_0\) in
Proposition~\ref{d2:prop:coefficient} to
\(\gamma_0\in K_3(\DtwocO_{E_{\Dtwoar}})\).
Then \(\gamma=(1-\sigma_{-1})\gamma_0\) maps exactly to \(\beta\).

The nonzero point coordinate makes \(\beta\) nontorsion, so
\(\gamma\) is nontorsion.
The regulator injectivity of
\cite[Section~6]{BunkeTamme}, with the rank calculation of
\cite[Proposition~12.2]{Borel}, gives \(\gamma\) a nonzero
complex regulator.
All calculations in Section~\ref{d2:sec:analytic} therefore apply to
this same \(\beta\).
Thus both detectors apply to the actual class \(c\) just constructed,
and its image under completion is exactly \(\Dtwowh c\).
\end{proof}

% End input: sections/03-quartic-core.tex


% Begin input: sections/04-three-dimensional-families.tex
\section{Arithmetic quadrics and geometric detection}\label{d3:sec:geometry}
The supported classes on an arithmetic quadric are detected by a ruling
line bundle. The geometric argument has two steps: the ruling survives
every finite localization permitted by the arithmetic local ring, and
its regulator pairs nontrivially with a number-field coefficient.

\subsection{The arithmetic quadric and its coefficient class}

For an abelian group $M$, write $M\Qrat=M\otimes_{\Z}\Q$. All $K$-groups in this article are Quillen algebraic $K$-groups. A coherent sheaf on a regular noetherian scheme is viewed in $K$-theory through a finite locally free resolution, or equivalently through the comparison $K\simeq G$.

The injectivity part of the rational Gersten conjecture asserts that, for every regular local ring $R$ with fraction field $F$ and every $n\geq 0$,
\[
K_n(R)\Qrat\longrightarrow K_n(F)\Qrat
\]
is injective. Its failure is sufficient to disprove the augmented rational Gersten conjecture.

\begin{theorem}\label{d3:thm:main}
Let $p$ be an odd prime and set
\[
A=\left(\Z_{(p)}[x,y,z]/(p+xy+z^2)\right)_{(x,y,z)}.
\]
There exists an integral class $c\in K_3(A)$ of infinite order whose image in $K_3(\Frac A)$ is zero. In particular,
\[
\ker\bigl(K_3(A)\Qrat\longrightarrow K_3(\Frac A)\Qrat\bigr)\neq 0.
\]
\end{theorem}

Only $p=3$ is needed for a counterexample. The proof is given for odd $p$ because its algebraic geometry is identical for all such primes.

In the case $p=3$, Section~\ref{d3:sec:explicit} gives the following fully specified class:
\begin{equation}\label{d3:eq:explicitfront}
\boxed{\quad
c_{\mathrm{exp}}=2\left\{\frac{1+z}{2},\,1+x,\,
1+\frac{(1+x)y}{4}\right\}\in K_3(A).
\quad}
\end{equation}
Braces mean the product of the three indicated unit classes in Quillen $K_1$. All entries are units in $A$. The factor two makes generic vanishing an integral assertion.


Put
\[
V=\Z_{(p)},\qquad W=V[u]/(u^2+p),\qquad E=\Q(u),\qquad \sigma(u)=-u.
\]
Thus $W$ is a discrete valuation ring (DVR) with uniformizer $u$ and residue field $\F_p$, and $E$ is an imaginary quadratic number field. Write
\[
A_0=V[x,y,z]/(p+xy+z^2),\quad \m=(x,y,z),\quad S=A_0\setminus\m.
\]
Then $A=S^{-1}A_0$. Let
\begin{align*}
B_0&=A_0\otimes_VW=W[x,y,z]/(xy+z^2-u^2),\\
B&=A\otimes_VW=S^{-1}B_0,\\
Q&=\Spec E[x,y,z]/(xy+z^2-u^2),\\
T&=\Spec B[1/u].
\end{align*}
For $s\in S$, let $Q_s$ denote the corresponding principal open of $Q$.

We use the following standard inputs.

\begin{enumerate}
\item Quillen localization and d\'evissage, the projection formula for a regular closed immersion, and flat base change for the associated pushforward. These may be formulated in $G$-theory and transported to $K$-theory on the regular schemes involved; see \cite[Sections~5--7]{Quillen}.
\item $K_2(\F_p)=0$ \cite{QuillenFinite}\cite[IV, Corollary~1.13]{Weibel}. Consequently $K_3(W)\to K_3(E)$ is surjective, integrally.
\item Borel's theorem gives a class in $K_3(E)$ with nonzero regulator at a complex embedding. The comparison of the Borel and Beilinson regulators identifies this, up to a nonzero scalar, with a class in
\[
\HD^1(\Spec\C,\R(2))=\C/\R(2),\qquad \R(j)=(2\pi i)^j\R.
\]
For an imaginary quadratic field, conjugating the embedding negates this regulator. See \cite[Proposition~12.2]{Borel} and \cite[Section~6]{BunkeTamme}.
\item The analytic version of the Beilinson regulator is natural and multiplicative and has the usual first Chern-class normalization. A precise reference is \cite[Theorem~2.31]{BunkeTamme}. We use its weight-three component
\[
r_3:K_3(U)\longrightarrow\HD^3(U,\R(3))
\]
on smooth complex fibers, and its weight-two component on number fields.
\item Algebraic $K$-theory of rings commutes with filtered colimits \cite[IV, Section~6.4]{Weibel}. In particular,
\begin{equation}\label{d3:eq:continuity}
K_3(T)=\colim_{s\in S} K_3(Q_s).
\end{equation}
\end{enumerate}

\begin{remark}\label{d3:rem:spreading}
The site in \cite{BunkeTamme} consists of regular separated schemes of finite type over $\Z$. Although $E$ and $Q_s$ are not finite type over $\Z$, each finite collection of coefficients, $K$-classes, line bundles, and morphisms used below spreads to a smooth model over $\OO_E[1/N]$ for a suitable positive integer $N$. These models are regular and of finite type over $\Z$, and all further integer localizations have the same complex fiber. Regulator naturality and filtered-colimit compatibility therefore give precisely the maps used here. Projection from the conjugation-invariant collection of complex components to a chosen embedding preserves products. No regulator on the infinite scheme $T$ is required.
\end{remark}

\subsection{Regularity and the supported class}

\begin{lemma}\label{d3:lem:regularA}
The ring $A$ is a three-dimensional regular local domain with residue field $\F_p$. Its maximal ideal is generated by $x,y,z$, its fraction field is $\Q(x,z)$, and $p\in\m^2$.
\end{lemma}
\begin{proof}
The ambient ring $V[x,y,z]_{(p,x,y,z)}$ is regular local of dimension four. The element $p+xy+z^2$ belongs to its maximal ideal but not to its square: its initial linear term is $p$. Quotienting by this element gives a regular local ring of dimension three. The relation $p=-xy-z^2$ gives both the asserted generators of the maximal ideal and $p\in\m^2$. Solving for $y$ after inverting $x$ gives the fraction field $\Q(x,z)$. Alternatively, the defining polynomial is primitive and irreducible as a polynomial of degree one in $y$, so the quotient is a domain.
\end{proof}

The regular principal divisor $x=0$ has coordinate ring
\[
R=A/(x)=\left(V[y,z]/(z^2+p)\right)_{(y,z)}
\cong W[y]_{(u,y)},
\]
where the map $W\to R$ sends $u$ to $z$. Let
\[
i:\Spec R\hookrightarrow\Spec A
\]
be the closed immersion.

\begin{lemma}\label{d3:lem:beta}
There is an integral class $\beta\in K_3(W)$ such that
\[
\sigma\beta=-\beta\quad\text{integrally},\qquad r_2(\beta|_E)\neq 0
\]
at a fixed complex embedding $E\hookrightarrow\C$.
\end{lemma}
\begin{proof}
Choose a class of $K_3(E)$ with nonzero Borel--Beilinson regulator. Localization gives
\[
K_3(W)\longrightarrow K_3(E)\longrightarrow K_2(\F_p)=0,
\]
so choose an integral lift $b_0\in K_3(W)$. Set $\beta=b_0-\sigma b_0$. Since $\sigma^2=1$, we have $\sigma\beta=-\beta$ without rationalization. The nontrivial embedding is complex conjugation, and on $\C/\R(2)$ this negates the nonzero regulator. Thus $r_2(\beta)=2r_2(b_0)\neq 0$.
\end{proof}

Define
\begin{equation}\label{d3:eq:c}
c=i_*(\beta|_R)\in K_3(A).
\end{equation}
The pushforward exists because $A$ and $R$ are regular, and the closed immersion is a regular Cartier immersion.

\begin{lemma}\label{d3:lem:genericzero}
The image of $c$ in $K_3(\Frac A)$ is zero integrally.
\end{lemma}
\begin{proof}
By construction, $c$ comes from a class with support on $V(x)$. The localization sequence therefore sends it to zero in $K_3(A[1/x])$, and hence in the fraction field.
\end{proof}

The remaining question is whether $c$ has infinite order. Merely knowing that $V(x)$ is principal does not answer this: a trivial normal bundle implies the vanishing of the self-intersection operator $i^*i_*$, not the vanishing of $i_*$ itself.

\subsection{The arithmetic node and its divisor class}\label{d3:sec:node}

\begin{lemma}\label{d3:lem:normalB}
The ring $B$ is a normal local domain of dimension three, with maximal ideal
\[
\n=(u,x,y,z).
\]
It is regular away from its closed point, and
\[
\operatorname{gr}_{\n}B
\cong\F_p[U,X,Y,Z]/(XY+Z^2-U^2).
\]
\end{lemma}
\begin{proof}
The map $A\to B$ is finite free of rank two. Its closed fiber is $\F_p[u]/(u^2)$, with a unique maximal ideal, so $B$ is local. In particular, $B=(B_0)_{\n}$, not merely an unspecified localization of $B_0$. The polynomial $u^2+p$ is irreducible over the purely transcendental extension $\Q(x,z)$, hence $B$ is a domain.

At a point where $u$ is invertible, the ring is on a smooth affine quadric over $E$. At a point of the special fiber other than the origin, one of $x,y,z$ is invertible, and the relative gradient $(y,x,2z)$ has a unit component. Thus the punctured spectrum is regular. Being a hypersurface, $B$ satisfies $S_2$; regularity in codimension one gives normality. The displayed associated graded ring follows from the nonzero quadratic initial form of the equation in the regular ambient local ring $W[x,y,z]_{(u,x,y,z)}$.
\end{proof}

Consider the height-one primes
\[
P_+=(x,z-u),\qquad P_-=(x,z+u).
\]
Both have multiplicity one in $\divisor_B(x)$, so
\begin{equation}\label{d3:eq:divx}
\divisor_B(x)=P_++P_-.
\end{equation}

\begin{lemma}\label{d3:lem:smallresolution}
Let $W'$ be a DVR with uniformizer $\varpi$ and residue field $k$, and let $0\neq d\in\varpi W'$. In the normal local domain
\[
C=\left(W'[x,y,w]/(xy+w(w+d))\right)_{(\varpi,x,y,w)},
\]
the height-one prime $P=(x,w)$ has infinite order in $\Cl(C)$. There is no restriction on the residue characteristic or on the positive valuation of $d$.
\end{lemma}
\begin{proof}
The defining polynomial is irreducible, and $C$ is a hypersurface. Its generic fiber is smooth since $d\neq0$. On the special fiber, a prime containing $x$ and $y$ also contains $w$; hence every nonclosed point of that fiber has $x$ or $y$ invertible. Thus $C$ is regular off its closed point, and $S_2$ and $R_1$ give normality.

Blow up the ideal $I=(x,w)$:
\[
\rho:Y=\operatorname{Bl}_{I}(\Spec C)\longrightarrow\Spec C.
\]
Its two standard charts are localizations of the following polynomial rings:
\begin{align*}
W'[x,t],&\qquad w=xt,\quad y=-t(xt+d),\\
W'[q,y],&\qquad w=-qy-d,\quad x=-q(qy+d).
\end{align*}
These are the actual Rees charts $C[w/x]$ and $C[x/w]$ inside the fraction field. For example, introducing $t=w/x$ gives $x(y+t(xt+d))=0$, and $x$-saturation gives exactly the first polynomial ring. The second is obtained by $w$-saturation. There are no extra equations because each indicated pair is a transcendence basis over $\Frac(W')$.

Consequently $Y$ is regular. The closed fiber has equations $(\varpi,x)$ in the first chart and $(\varpi,y)$ in the second. Its two affine charts are $k[t]$ and $k[q]$, with $q=t^{-1}$ on their overlap. Hence it is a reduced curve
\[
F\cong\PP^1_k.
\]

The ideal $I$ is invertible off the closed point. Outside $V(I)$ this is immediate. At a nonclosed prime containing $I$, either $y$ is a unit, giving $I=(w)$, or $\varpi$ is a unit; then $d$ and $w+d$ are units and $I=(x)$. The blowup is therefore an isomorphism off the closed point. Its only exceptional locus is the codimension-two curve $F$, so there are no exceptional prime divisors.

The standard tautological identity for a blowup is $I\OO_Y=\OO_Y(1)$, whose restriction to $F$ is $\OO_{\PP^1_k}(1)$; see \cite[Lemma~31.33.4, Tag~02OS]{StacksBlowup}. Since $I$ has multiplicity one at the generic point of $P$ and there are no exceptional divisors,
\[
I\OO_Y=\OO_Y(-\widetilde P),\qquad
\OO_Y(\widetilde P)|_F=\OO_{\PP^1_k}(-1).
\]

If $\divisor_C(f)=NP$ for a nonzero integer $N$, then the same rational function on $Y$ has divisor exactly $N\widetilde P$. Restricting its trivial divisor line bundle to $F$ gives
\[
\OO_{\PP^1_k}(-N)\cong\OO_{\PP^1_k},
\]
which forces $N=0$, a contradiction.
\end{proof}

\begin{corollary}\label{d3:lem:clnonzero}
The class $[P_+]$ is not torsion in $\Cl(B)$.
\end{corollary}
\begin{proof}
Set $w=z-u$ and $d=2u$. Then $B$ is the ring in Lemma~\ref{d3:lem:smallresolution}, with $W'=W$ and $\varpi=u$, and $P_+=(x,w)$.
\end{proof}

\begin{remark}
Blowing up the maximal ideal instead gives the regular resolution with exceptional quadric $\PP^1_{\F_p}\times\PP^1_{\F_p}$. Its normal bundle is $\OO(-1,-1)$, while the strict transform of $P_+$ restricts to one ruling, $\OO(1,0)$. If $NP_+$ were principal, its divisor upstairs would have the form $N\widetilde P_++aF$, whose restriction would be $\OO(N-a,-a)$. Triviality again forces $N=0$. The small resolution above avoids this larger exceptional divisor and also works in residue characteristic two.
\end{remark}

\begin{lemma}\label{d3:lem:invertu}
Restriction induces an isomorphism
\[
\Cl(B)\xrightarrow{\sim}\Cl(B[1/u]).
\]
In particular, the image of $[P_+]$ is not torsion in $\Pic(T)$.
\end{lemma}
\begin{proof}
The quotient
\[
B/(u)=\F_p[x,y,z]_{(x,y,z)}/(xy+z^2)
\]
is a domain. The only height-one prime removed by inverting $u$ is therefore the principal prime $(u)$. The localization sequence for divisor class groups is surjective with kernel generated by the removed prime classes, so its kernel is zero. Since $B[1/u]$ is regular, its class group equals its Picard group.
\end{proof}

\begin{remark}
One can also compute $\Cl(B)=\Z$. Indeed, $B[1/x]$ is a localization of the unique factorization domain $W[x,x^{-1},z]$, so the class group is generated by $P_+$ and $P_-$. Equation~\eqref{d3:eq:divx} gives their sum as zero, and Lemma~\ref{d3:lem:clnonzero} excludes any further relation. The argument needs only non-torsion, not this complete computation.
\end{remark}

\subsection{Persistence of Picard and Betti classes}\label{d3:sec:persistence}

Let $D_+$ and $D_-$ denote the divisors of $Q$ with equations
\[
D_+=(x,z-u),\qquad D_-=(x,z+u).
\]
They are disjoint affine lines over $E$.

\begin{lemma}\label{d3:lem:quadricpic}
We have $\Pic(Q)=\Z[D_+]$, and the same statement holds after extension to $\overline E$ or $\C$. The absolute Galois group of $E$ acts trivially on the geometric Picard group. At a complex embedding,
\[
H^2(Q(\C),\Z)=\Z,
\]
with $c_1(\OO_Q(D_+))$ a primitive generator after the usual Tate normalization.
\end{lemma}
\begin{proof}
The projective closure
\[
\overline Q=\{xy+z^2=u^2t^2\}\subset\PP^3_E
\]
is split: it has the rank-one matrix presentation
\[
\det\begin{pmatrix}x&ut-z\\ut+z&y\end{pmatrix}=0.
\]
Thus $\overline Q\cong\PP^1\times\PP^1$. The divisor at infinity is a smooth divisor of type $(1,1)$; after an automorphism of one factor it is the diagonal. Hence
\[
Q\cong (\PP^1\times\PP^1)\setminus\Delta,
\qquad
\Pic(Q)=\Z^2/\Z(1,1)=\Z.
\]
The closure of $D_+$ is a ruling, and gives a generator. This calculation is unchanged over algebraically closed extensions. Because the generator is defined over $E$, the absolute Galois action is trivial.

Projection to one factor makes $Q(\C)\to\PP^1(\C)$ a locally trivial affine-line bundle, so it is a homotopy equivalence. Equivalently, the topological localization sequence for the diagonal gives the same quotient $\Z^2/\Z(1,1)$ in degree two. The first Chern class of a ruling maps to its primitive generator.
\end{proof}

\begin{lemma}\label{d3:lem:Picinject}
The restriction map $\Pic(Q)\to\Pic(T)$ is injective.
\end{lemma}
\begin{proof}
The generator $[D_+]$ maps to the image of $[P_+]$ in $\Cl(B[1/u])$. The latter has infinite order by Lemmas~\ref{d3:lem:clnonzero} and \ref{d3:lem:invertu}.
\end{proof}

\begin{lemma}\label{d3:lem:components}
For each $s\in S$, every geometrically irreducible curve in $Q_{\C}\setminus(Q_s)_{\C}$ has zero class in $\Pic(Q_{\C})$.
\end{lemma}
\begin{proof}
Let $D$ be an irreducible divisor over $E$ in $Q\setminus Q_s$. It is absent from $T$, so $\OO_Q(D)|_T$ is trivial. Lemma~\ref{d3:lem:Picinject} implies $[D]=0$ in $\Pic(Q)$.

In characteristic zero the reduced divisor $D$ becomes a reduced union of finitely many geometric irreducible components. These components are permuted transitively by the absolute Galois group of $E$. Its action on $\Pic(Q_{\overline E})=\Z$ is trivial, so all components have the same integer class, say $a$. If there are $r>0$ components, their sum has class $ra$. This sum is the pullback of $[D]=0$, so $ra=0$ in the torsion-free group $\Z$, and therefore $a=0$. Extending further to $\C$ preserves the calculation.
\end{proof}

\begin{remark}
A geometric boundary component $\Gamma$ is defined over some finite
extension $E'/E$. Put $L=\OO_Q(D_+)$ and write
$[\OO(\Gamma)]=a[L_{E'}]$ in $\Pic(Q_{E'})=\Z[L_{E'}]$.
This line bundle is trivial on $T_{E'}$. Its finite-flat norm
(equivalently, determinants after restriction of scalars) therefore
makes $L_T^{a[E':E]}$ trivial. The infinite order of $L_T$ forces
$a=0$, giving a second proof that each geometric boundary component has zero Picard class.
\end{remark}

\begin{remark}
It is not enough that $\divisor(s)$, as a whole, is principal. For instance $\divisor(x)=D_++D_-$ is principal, while its two components have opposite nonzero classes. Removing $D_+\cup D_-$ kills $H^2(Q,\R)$. The element $x$ is not in the permitted multiplicative set $S$.
\end{remark}

\begin{lemma}\label{d3:lem:Betti}
For every $s\in S$, restriction is injective:
\[
H^2(Q_{\C},\R)\longrightarrow H^2((Q_s)_{\C},\R).
\]
\end{lemma}
\begin{proof}
Let $Z=Q_{\C}\setminus(Q_s)_{\C}$ with its reduced support. The kernel is the image of the support map
\[
H_Z^2(Q_{\C},\R)\longrightarrow H^2(Q_{\C},\R).
\]
Remove from the ambient surface the finite set of singular points and intersections of the reduced curve components of $Z$. Support cohomology in finitely many points of a real four-manifold vanishes below degree four, so this changes neither the degree-two support group nor the ambient $H^2$. The remaining components are disjoint smooth connected complex curves.
The tubular-neighborhood calculation in the proof of
Proposition~\ref{d2:prop:geometry} identifies degree-two support
cohomology with one copy of $\R$ for each component.
There the canonical transverse section of the divisor line bundle
identifies the image of its support generator with the global first
Chern class of the component. The isomorphism on ambient $H^2$
transports this identification across the deleted finite set.
Every such class is zero by Lemma~\ref{d3:lem:components}.
Thus the support map is zero.
\end{proof}

\subsection{Analytic Deligne detection}\label{d3:sec:Delignedetection}

Throughout this section, $\HD$ denotes analytic Deligne cohomology.

\begin{lemma}\label{d3:lem:Deligne}
For a smooth complex surface $U$, there is a natural isomorphism
\begin{equation}\label{d3:eq:Delignequotient}
\HD^3(U,\R(3))
\cong H^2(U,\C)/H^2(U,\R(3)).
\end{equation}
\end{lemma}
\begin{proof}
The analytic Deligne complex is
\[
[\R(3)\longrightarrow\OO_U\longrightarrow\Omega_U^1\longrightarrow\Omega_U^2].
\]
The holomorphic de Rham complex resolves the constant sheaf $\C$ by the holomorphic Poincar\'e lemma, and there are no holomorphic forms in degree three. The Deligne complex is consequently quasi-isomorphic to the constant sheaf $\C/\R(3)$ placed in degree one. Equivalently, it is the shifted cone of
\[
R\Gamma(U,\R(3))\longrightarrow R\Gamma(U,\C).
\]
Since real cohomology injects into its complexification in every degree, the long exact sequence of this cone gives \eqref{d3:eq:Delignequotient}. All constructions are natural for restriction to open subsets. Properness is not used.
\end{proof}

\begin{remark}
On a smooth normal-crossings compactification, the logarithmic de Rham complex also has $F^3=0$. Thus Deligne--Beilinson cohomology agrees with the analytic group in \eqref{d3:eq:Delignequotient}; compare \cite[Remark~2.13]{EsnaultViehweg}. The direct analytic formulation above already suffices for the proof.
\end{remark}

\begin{lemma}\label{d3:lem:Deligneinject}
For every $s\in S$, restriction induces an injection
\[
\HD^3(Q_{\C},\R(3))\longrightarrow\HD^3((Q_s)_{\C},\R(3)).
\]
\end{lemma}
\begin{proof}
Combine Lemmas~\ref{d3:lem:Betti} and \ref{d3:lem:Deligne}. Explicitly, an injection of real vector spaces $V\hookrightarrow W$ induces an injection
\[
\frac{V\otimes_{\R}\C}{(2\pi i)^3V}
\longrightarrow
\frac{W\otimes_{\R}\C}{(2\pi i)^3W},
\]
because the intersection of the image of $V\otimes\C$ with $(2\pi i)^3W$ is exactly $(2\pi i)^3V$.
\end{proof}

Set
\[
\eta=[\OO_{D_+}]\in K_0(Q),\qquad
\gamma=(\beta|_E)\eta\in K_3(Q).
\]
The structure sheaf class is a perfect-complex class, with
\[
\eta=1-[\OO_Q(-D_+)],\quad \operatorname{rank}(\eta)=0,\quad
\operatorname{ch}_1(\eta)=c_1(\OO_Q(D_+)).
\]

\begin{lemma}\label{d3:lem:product}
The class $\gamma$ has nonzero weight-three analytic Deligne regulator.
\end{lemma}
\begin{proof}
Naturality and multiplicativity give
\begin{equation}\label{d3:eq:product}
r_3(\gamma)=r_2(\beta|_E)\cup c_1(\OO_Q(D_+)).
\end{equation}
There are no additional terms of total weight three: a degree-three $K$-class pulled back from a point has the relevant regulator only in $H^1_{\mathcal D}(\mathrm{pt},\R(2))$, and the rank of $\eta$ is zero.

Let $h$ be a real primitive generator of $H^2(Q_{\C},\R)$, normalized so that $c_1(\OO_Q(D_+))$ has Betti image $2\pi i\,h$. A nonzero element of $\C/\R(2)$ is represented by $it$ for a nonzero real $t$. Under \eqref{d3:eq:Delignequotient}, the product in \eqref{d3:eq:product} is represented, up to the harmless convention for the sign of $h$, by
\[
(it)(2\pi i)h=-2\pi t\,h.
\]
This is a nonzero real multiple of $h$, whereas $H^2(Q_{\C},\R(3))$ is the imaginary real line generated by $h$. The quotient class is therefore nonzero.

One may equivalently compute this cup product on $\PP^1$: the projection of the affine-line torsor $Q\to\PP^1$ induces the required degree-two Betti isomorphism, and cup product with the hyperplane class is the usual projective-bundle map.
\end{proof}

\begin{proposition}\label{d3:prop:infinite}
The restriction $\gamma|_T\in K_3(T)$ has infinite order.
\end{proposition}
\begin{proof}
By regulator naturality and Lemma~\ref{d3:lem:Deligneinject},
\[
r_3(\gamma|_{Q_s})\neq 0
\qquad\text{for every }s\in S.
\]
If a positive integer multiple $N\gamma$ vanished in $K_3(T)$, continuity \eqref{d3:eq:continuity} would make it vanish already in $K_3(Q_s)$ for some $s$. Applying the regulator there would yield $N r_3(\gamma|_{Q_s})=0$ in a real vector space, which is impossible.
\end{proof}

\subsection{Flat base change and nonvanishing}

\begin{lemma}\label{d3:lem:basechange}
The image of the supported class \eqref{d3:eq:c} in $K_3(T)$ equals $2\gamma|_T$.
\end{lemma}
\begin{proof}
The map $A\to B[1/u]$ is flat, being a finite free base change followed by localization. Its pullback of the divisor $x=0$ is the disjoint union
\[
(D_+\cap T)\amalg(D_-\cap T).
\]
The original coefficient map $W\to R$ sends $u$ to $z$. Therefore its two restrictions after the splitting are the coefficient maps $u\mapsto u$ and $u\mapsto -u$. Flat base change and the projection formula give
\begin{align*}
c|_T
&=(\beta|_E)[\OO_{D_+\cap T}]+(\sigma\beta|_E)[\OO_{D_-\cap T}]\\
&=(\beta|_E)\bigl([\OO_{D_+\cap T}]-[\OO_{D_-\cap T}]\bigr).
\end{align*}
The two divisors are disjoint since $2u$ is a unit on $T$. Moreover their union is the principal divisor $V(x)$. The resolution by multiplication by $x$ gives
\[
[\OO_{D_+\cap T}]+[\OO_{D_-\cap T}]
=[\OO_{T/(x)}]=0\quad\text{in }K_0(T).
\]
Consequently $c|_T=2(\beta|_E)[\OO_{D_+\cap T}]=2\gamma|_T$.
\end{proof}

\begin{proof}[Proof of Theorem~\ref{d3:thm:main}]
Lemma~\ref{d3:lem:regularA} gives the required regular local ring. By Proposition~\ref{d3:prop:infinite} and Lemma~\ref{d3:lem:basechange}, the pullback of $c$ to $K_3(T)$ has infinite order; hence $c$ itself has infinite order. Lemma~\ref{d3:lem:genericzero} shows that its generic image is zero integrally. Tensoring with $\Q$ therefore gives the asserted nonzero rational kernel.
\end{proof}

The nonzero-product assertion in this proof is not deduced from a nonzero Picard class alone. The essential extra step is that its regulator remains nonzero \emph{after every permitted finite localization}. This is what makes the filtered-colimit argument valid.

\subsection{The unsplit generic fiber}

The following calculation identifies the coefficient of the supported class on the unsplit generic fiber. Put
\[
X=\Spec\Q[x,y,z]/(xy+z^2+p).
\]
Its divisor $Z=V(x)$ is $\A^1_E$, while
\[
X\setminus Z\cong\mathbb G_{m,\Q}\times\A^1_{\Q}.
\]
For every $n\geq 0$, localization, homotopy invariance, and the fundamental theorem give
\[
K_n(E)\longrightarrow K_n(X)\longrightarrow
K_n(\Q)\oplus K_{n-1}(\Q)
\xrightarrow{\partial}K_{n-1}(E).
\]
The boundary is zero on $K_n(\Q)$ and is restriction $\operatorname{res}_{E/\Q}$ on the second summand: that summand is represented by $\{x\}\cdot b$, and $x$ cuts $Z$ with multiplicity one. The preceding boundary has image $\operatorname{res}K_n(\Q)$ for the same reason.

After rationalization, restriction is injective because transfer composed with restriction is multiplication by two. Constants split the resulting short exact sequence, giving
\begin{equation}\label{d3:eq:genericdecomposition}
K_n(X)\Qrat\cong K_n(\Q)\Qrat\oplus K_n(E)\Qrat^{-}.
\end{equation}
Here the second summand is identified using $\operatorname{res}\operatorname{Tr}=1+\sigma$ and enters through $Z$-Gysin. In degree three, Borel's calculation \cite[Proposition~12.2]{Borel} yields
\[
K_3(\Q)\Qrat=0,\qquad K_3(E)\Qrat^{-}\cong\Q.
\]
Thus the generic-fiber Gysin map identifies $K_3(E)\Qrat$ with $K_3(X)\Qrat$. The supported class is already a nonzero one-dimensional summand before localization at $\m$. The main proof shows that this particular summand cannot disappear at that localization.

\section{The explicit three-adic class and its complex period}\label{d3:sec:explicit}

This section gives a geometric proof of the explicit class at $p=3$. It uses the same geometric persistence lemmas, but obtains nonvanishing by an explicit period rather than by selecting a Borel class. Put
\[
A=\left(\Z_{(3)}[x,y,z]/(3+xy+z^2)\right)_{(x,y,z)},
\]
and define
\[
a=\frac{1+z}{2},\qquad s_0=a(1-a)=1+\frac{xy}{4},\qquad
\lambda=1+x,\qquad \mu=1+\frac{(1+x)y}{4}.
\]
Every one of $a,1-a,s_0,\lambda,\mu$ is a unit of $A$, and
\begin{equation}\label{d3:eq:unitsidentity}
1+x\mu=\lambda s_0.
\end{equation}

\subsection{Dennis--Stein notation and integral generic vanishing}
We use the modern Dennis--Stein convention: $\langle r,t\rangle$ is defined if $1-rt$ is a unit, and, if $r$ is a unit,
\[
\langle r,t\rangle=\{r,1-rt\}.
\]
If $t$ is a unit, antisymmetry gives $\langle r,t\rangle=\{1-rt,t\}$. The relation used below is
\begin{equation}\label{d3:eq:DS2}
\langle r,t\rangle+\langle r,v\rangle
=\langle r,t+v-rtv\rangle,
\qquad \langle r,1\rangle=0.
\end{equation}
The modern convention and relations are those of \cite[III, Definition~5.11 and (D1)--(D3)]{Weibel}; the historical source is \cite{DennisStein}. We do not need a presentation theorem for $K_2$.

Let
\[
\delta=\langle -x,y/4\rangle\in K_2(A),\qquad
\alpha=[a]\delta\in K_3(A).
\]
On $A[1/x]$, the symbol is $\{-x,s_0\}$. Since $a$ and $1-a$ are units, the Steinberg relation and graded commutativity give
\[
2\{a,s_0\}=2\{a,a(1-a)\}
=2[a]^2+2\{a,1-a\}=0.
\]
Reordering products therefore gives
\begin{equation}\label{d3:eq:explicitgeneric}
2\alpha|_{A[1/x]}=2\{a,-x,s_0\}=0
\qquad\text{integrally}.
\end{equation}

The Dennis--Stein class can also be rewritten using units alone. By \eqref{d3:eq:DS2},
\[
\delta=\langle -x,\mu\rangle=\{1+x\mu,\mu\}
=\{\lambda s_0,\mu\}.
\]
Multiplying by $2[a]$ kills the $\{a,s_0,\mu\}$ term and yields
\begin{equation}\label{d3:eq:explicitDS}
2\alpha=2\{a,\lambda,\mu\}=c_{\mathrm{exp}}.
\end{equation}
Thus the explicitly displayed class \eqref{d3:eq:explicitfront} has zero generic image integrally.

\subsection{A finite open on which to calculate the period}
Keep $E=\Q(u)$ with $u^2=-3$, and fix $u=i\sqrt3$. Let
\[
Q=\Spec E[x,y,z]/(3+xy+z^2),\qquad
U=Q[(1-z^2)^{-1}].
\]
The classes $\delta$ and $\alpha$ are defined on $U$, since $a,1-a,s_0$ are units there. Both lines $D_+$ and $D_-$ lie entirely inside $U$, because on them $1-z^2=4$. Equation~\eqref{d3:eq:explicitgeneric} holds on $U[1/x]$ as well.

Localization therefore provides a class supported on $D=V(x)$ that maps to $2\alpha\in K_3(U)$. Since the entire divisor $D$ lies in $U$, excision identifies the supported groups for $Q$ and $U$. Equivalently, apply the localization sequence with $K_3(D)$ and the naturality of its pushforward. Thus there is a class
\begin{equation}\label{d3:eq:explicitlift}
\widetilde c\in K_3(Q),\qquad \widetilde c|_U=2\alpha.
\end{equation}
No coefficient of this supported lift has to be chosen or computed explicitly.

\subsection{The curvature and the Deligne product}
The weight-two analytic regulator of $\delta$ has the holomorphic curvature
\begin{equation}\label{d3:eq:omega}
\omega=\frac{dx\wedge d(y/4)}{1+xy/4}.
\end{equation}
Indeed, where $x\neq0$, multiplicativity for the two unit classes gives
\[
d\log(-x)\wedge d\log s_0
=\frac{dx\wedge d(y/4)}{s_0}.
\]
The Hodge endpoint is a canonical holomorphic two-form: in this degree $F^2A^1=0$, so there is no preceding filtered form that could change it by a boundary. Both sides therefore extend as actual holomorphic two-forms on $U$, and agreement on the dense open proves \eqref{d3:eq:omega}. The logarithmic normalization follows from the ordinary unit regulator and \cite[Theorem~2.31 and Lemma~4.5]{BunkeTamme}. An independent algebraic differential calculation for Dennis--Stein characters over smooth commutative $\Q$-algebras, with its own symbol conventions, appears in \cite[Section~3.2, Theorem~3.4 and Corollary~3.6]{Ginot}; the preceding dense-open calculation fixes the regulator sign and gives the needed form.

Suppose that $a$ has a single-valued holomorphic logarithm $L$ on an analytic open $N\subset U(\C)$. Then the regulator of $\alpha=[a]\delta$ on $N$ is represented in
\[
\HD^3(N,\R(3))=H^2(N,\C/\R(3))
\]
by the closed holomorphic two-form
\begin{equation}\label{d3:eq:Lomega}
-L\omega.
\end{equation}
We use the cone-coordinate convention of \cite[Proposition~2.8]{BunkeTamme}: an interval form $\Theta$ maps to
\[
\bigl(\Theta|_0\oplus\Theta|_1,-\textstyle\int_{[0,1]}\Theta\bigr).
\]
Let $v$ denote the auxiliary interval coordinate and put $g=i\operatorname{Im}L+v\operatorname{Re}L$.
By \cite[Lemma~4.5]{BunkeTamme}, the unit class $[a]$ is represented by $dg$.
We may replace this by $d(vL)$ in the same absolute interval complex:
\[
dg-d(vL)=d\bigl((1-v)i\operatorname{Im}L\bigr).
\]
The primitive is allowed in $IDR(1)$: its endpoint at zero is
$\R(1)$-valued, and its endpoint at one is zero, as required by
$F^1A^0=0$. Thus this replacement does not substitute a different
relative class for the original absolute regulator.
Choose a closed interval representative $\Theta_\delta$ whose Hodge endpoint is $\omega$.
The product $d(vL)\wedge\Theta_\delta$ has zero endpoint at zero and
endpoint $dL\wedge\omega=0$ at one, since $N$ is a complex surface.
Multiplicativity and Stokes' formula therefore give
\[
-\int_{[0,1]}d(vL)\wedge\Theta_\delta
=-L\omega+d\int_{[0,1]}vL\Theta_\delta.
\]
Thus the cone coordinate is \eqref{d3:eq:Lomega}, for every chosen
closed representative of the original absolute class. The form
$L\omega$ is closed because $N$ has complex dimension two.

\subsection{A strictly positive real period}
Inside $U(\C)$ consider the smooth sphere
\[
\Sigma:\quad z=i\sqrt3\,\tau,\quad
x=\sqrt3\sqrt{1-\tau^2}\,e^{i\vartheta},\quad
y=-\overline x,\qquad -1\leq\tau\leq1.
\]
The polar coordinates degenerate only at the poles; intrinsically this is the sphere $|x|^2+(\operatorname{Im}z)^2=3$. Along it,
\[
a=\frac{1+i\sqrt3\,\tau}{2},\qquad
s_0=\frac{1+3\tau^2}{4}>0.
\]
There is an analytic neighborhood $N$ of $\Sigma$ on which $\operatorname{Re}a>0$, so the principal $L=\log a$ is single-valued. On $\Sigma$,
\[
L=\frac12\log\frac{1+3\tau^2}{4}+i\arctan(\sqrt3\,\tau),
\qquad
\omega=-\frac{6i\tau}{1+3\tau^2}\,d\tau\wedge d\vartheta.
\]
Consequently,
\begin{equation}\label{d3:eq:positiveperiod}
\int_{\Sigma}\operatorname{Re}(L\omega)
=2\pi\int_{-1}^{1}
\frac{6\tau\arctan(\sqrt3\,\tau)}{1+3\tau^2}\,d\tau>0.
\end{equation}
The integrand is positive except at $\tau=0$. With the specified Bunke--Tamme cone convention, the real period of $r_3(2\alpha)$ is the negative double of this positive integral. Since $\R(3)$ is the imaginary axis, this detects a nonzero class in $H^2(N,\C/\R(3))$. Thus $r_3(2\alpha)\neq0$ on $U$.

For comparison, integration by parts expresses \eqref{d3:eq:positiveperiod} as
\[
4\pi D(e^{i\pi/3}),
\]
with the displayed orientation and logarithmic convention. Strict positivity alone is sufficient; neither the special value nor Borel's rank theorem is needed in this argument.

\subsection{A normalized supported coefficient}\label{d3:sec:exactcoefficient}
The exact coefficient can also be specified rationally. Put $\zeta=(1+u)/2$, and define $\beta_{\mathrm{BT}}=\tfrac12\operatorname{bl}(2[\zeta])\in K_3(E)_{\Q}$ using the Bloch symbol map of \cite[Theorem~4.9 and (4.8)]{BunkeTamme}, first on the full integer ring $\OO_E=\Z[\zeta]$ and then by restriction to $E$. Both $\zeta$ and $1-\zeta=\zeta^{-1}$ are units. The doubled wedge boundary is zero even with the antisymmetric-tensor convention, so this definition avoids any integral two-torsion convention for the Bloch group. At the chosen embedding its cone-coordinate regulator is $-iD(\zeta)$. Borel injectivity \cite[Section~6]{BunkeTamme} makes this rational class independent of the integral lift.

A supported lift as in \eqref{d3:eq:explicitlift} has the form $\widetilde c=\gamma\eta_+$, with $\gamma\in K_3(E)$: use homotopy invariance on the two affine lines and $\eta_-=-\eta_+$. The two generators of $I=(x,z-u)$ give a map $f:Q\to\mathbf P^1_E$ with $f^*\OO(1)=I$. On $\Sigma$ its coordinate is
\[
\frac{z-u}{x}=-i\sqrt{\frac{1-\tau}{1+\tau}}e^{-i\vartheta}.
\]
Both the radial derivative and the angular derivative are negative, so this is a map of degree $+1$ for the displayed orientation. Hence
\[
\int_\Sigma c_1(I)=2\pi i,\qquad
\int_\Sigma\operatorname{ch}_1(\eta_+)=-2\pi i.
\]
Multiplicativity gives period $(-iD(\zeta))(-2\pi i)=-2\pi D(\zeta)$ for $\beta_{\mathrm{BT}}\eta_+$, whereas \eqref{d3:eq:positiveperiod} gives period $-8\pi D(\zeta)$ for $\widetilde c$. Borel injectivity on $K_3(E)_{\Q}$ therefore gives
\begin{equation}\label{d3:eq:exactcoefficient}
\widetilde c=4\beta_{\mathrm{BT}}\eta_+
\quad\text{in }K_3(Q)_{\Q}.
\end{equation}
This is a regulator-normalized identity for every supported lift. It does not identify $\beta_{\mathrm{BT}}$ with the Besser--de Jeu class used in Section~\ref{d3:sec:explicitcomplete}; that latter comparison will only use a nonzero rational scalar.

\subsection{Persistence and the explicit conclusion}
By \eqref{d3:eq:explicitlift}, the regulator of $\widetilde c$ is nonzero on $Q$. Lemma~\ref{d3:lem:Deligneinject} keeps it nonzero on every $Q_s$ for $s\in S$. The element $1-z^2$ belongs to $S$, so $T$ lies in $U$ and the restrictions of $\widetilde c$ and $2\alpha$ agree on $T$. By \eqref{d3:eq:explicitDS}, this is also the image of $c_{\mathrm{exp}}\in K_3(A)$.

If a positive multiple of $c_{\mathrm{exp}}$ vanished in $K_3(A)$, the same multiple of $\widetilde c$ would vanish in $K_3(T)$, and hence on some finite $Q_s$ by continuity. Its nonzero real regulator rules this out. We have therefore obtained the explicit version of the counterexample:
\begin{equation}\label{d3:eq:explicitconclusion}
\boxed{\quad
0\neq 2\left\{\frac{1+z}{2},\,1+x,\,
1+\frac{(1+x)y}{4}\right\}\in
\ker\bigl(K_3(A)\longrightarrow K_3(\Frac A)\bigr),
\quad}
\end{equation}
where nonzero is strengthened to infinite order by the period argument. This conclusion depends on the geometric persistence proof in Sections~\ref{d3:sec:node}--\ref{d3:sec:Delignedetection}, in addition to the explicit calculation in this section.

\subsection{An independent algebraic detection modulo three}\label{d3:sec:modthree}

There is a short independent test that proves $c_{\mathrm{exp}}$ is nonzero integrally, and even that it is not divisible by three. This test does not prove infinite order, so it does not replace the regulator argument for the rational conjecture.

Let
\[
R=A[T]_{\m A[T]},\qquad V'=\Z_{(3)}[T]_{(3)}.
\]
Then $R$ is a local $V'$-algebra with residue field $\F_3(T)$. Its module of relative differentials is presented by
\[
\Omega^1_{R/V'}=
\frac{R\,dx\oplus R\,dy\oplus R\,dz}
 {R(y\,dx+x\,dy+2z\,dz)},
\]
so
\[
\Omega^3_{R/V'}=
\bigl(R/(x,y,2z)\bigr)\,dx\wedge dy\wedge dz
\cong\F_3(T)\,dx\wedge dy\wedge dz.
\]

For a local ring with infinite residue field, the usual Milnor symbol product
\[
\iota:K^M_3(R)\longrightarrow K_3(R)
\]
has a comparison homomorphism $\phi:K_3(R)\to K^M_3(R)$ satisfying $\phi\iota=2$. This is the degree-three case of the Nesterenko--Suslin comparison recorded for general local rings in \cite[Proposition~4.2.8(5),(6)]{KerzThesis}, after identifying ordinary and improved Milnor $K$-theory for the infinite residue field of $R$.

Write $c_M=2\{a,\lambda,\mu\}\in K^M_3(R)$. The logarithmic differential is well defined on Milnor symbols, since it kills every Steinberg relation. A direct calculation gives
\[
d\log(c_M)
=\frac{1}{2(1+z)\mu}\,dz\wedge dx\wedge dy.
\]
Consequently the composite additive homomorphism
\[
K_3(A)\longrightarrow K_3(R)
\xrightarrow{\phi}K^M_3(R)
\xrightarrow{d\log}\Omega^3_{R/V'}
\]
sends $c_{\mathrm{exp}}$ to
\[
\frac{1}{(1+z)\mu}\,dx\wedge dy\wedge dz
=dx\wedge dy\wedge dz\neq0.
\]
The equality is in $\Omega^3_{R/V'}$, where $x,y,z$ vanish in the coefficient ring and $\mu=1$. Since this target is killed by three, $c_{\mathrm{exp}}\notin3K_3(A)$.

In particular, the coefficient exact sequence injects the nonzero image of this class in $K_3(A)/3$ into $K_3(A;\Z/3)$. Its generic image is zero by the integral calculation in \eqref{d3:eq:explicitgeneric}. Thus this explicit example also fails Gersten injectivity with $\Z/3$-coefficients, independently of the analytic regulator.

\begin{remark}
The same conclusion is already visible in
\[
C=A/\m^2\cong
\F_3[\epsilon_x,\epsilon_y,\epsilon_z]/(\epsilon_x,\epsilon_y,\epsilon_z)^2.
\]
To check this without assuming that the finite residue field is large enough for an ordinary Milnor comparison, pass to
\[
D=C[T]_{\m_C C[T]}
\cong k[\epsilon_x,\epsilon_y,\epsilon_z]/(\epsilon_x,\epsilon_y,\epsilon_z)^2,
\qquad k=\F_3(T).
\]
Let $M=k e_x\oplus k e_y\oplus k e_z$ be a $D$-module through the residue map. The $k$-derivation $d:D\to M$ given by $d\epsilon_i=e_i$ induces
\[
K_3^M(D)\longrightarrow\bigwedge_k^3 M,\qquad
\{g_1,g_2,g_3\}\longmapsto
\frac{dg_1}{g_1}\wedge\frac{dg_2}{g_2}\wedge\frac{dg_3}{g_3}.
\]
This kills the Steinberg relations because $dg/g$ and $d(1-g)/(1-g)$ are proportional. The three entries of $c_{\mathrm{exp}}$ map to $2(1+\epsilon_z)$, $1+\epsilon_x$, and $1+\epsilon_y$. After the comparison $K_3(D)\to K_3^M(D)$, the image of $c_{\mathrm{exp}}$ is therefore
\[
4e_z\wedge e_x\wedge e_y=e_x\wedge e_y\wedge e_z\ne0.
\]
In fact its image in $K_3(C)$ has exact order three: it is nonzero by this test, and $(1+\epsilon_x)^3=1$ kills its triple product by three. Both the henselization and the completion of $A$ have the same quotient $C$. Hence the explicit class also remains nonzero and not divisible by three in their $K_3$-groups, with zero generic image. This observation alone does not prove infinite order in either ring.
\end{remark}

\section{Quadratic and Eisenstein families}\label{d3:sec:families}

\subsection{Residue characteristic two}\label{d3:sec:two}

The small resolution in Lemma~\ref{d3:lem:smallresolution} allows the same argument at the prime two. This variant is not needed for the explicit counterexample in \eqref{d3:eq:explicitfront}.

\begin{proposition}\label{d3:prop:two}
For the regular local ring
\[
A_2=\left(\Z_{(2)}[x,y,z]/(1+xy+z^2)\right)_{(2,x,y,z-1)},
\]
the map $K_3(A_2)\Qrat\to K_3(\Frac A_2)\Qrat$ has nonzero kernel.
\end{proposition}
\begin{proof}
Put $v=z-1$. The defining equation is $2+2v+v^2+xy$, whose linear term in the regular ambient local ring is $2$. Thus $A_2$ is regular of dimension three, with maximal ideal generated by $x,y,v$ and residue field $\F_2$.

Take $W_2=\Z_{(2)}[i]$, $E_2=\Q(i)$, and $\varpi=1-i$. The equation $\varpi^2-2\varpi+2=0$ is Eisenstein, so $W_2$ is a DVR. The divisor $x=0$ in $\Spec A_2$ is regular, with coordinate ring $W_2[y]_{(\varpi,y)}$. Choose an integral anti-invariant Borel class $\beta\in K_3(W_2)$ exactly as in Lemma~\ref{d3:lem:beta}, and push it forward from this divisor.

After the finite free quadratic base change, put $w=z-i$. The local ring becomes
\[
\left(W_2[x,y,w]/(xy+w(w+2i))\right)_{(\varpi,x,y,w)}.
\]
Here $d=2i=-\varpi^2$ is nonzero and has positive valuation. Lemma~\ref{d3:lem:smallresolution} proves that $(x,w)$ has infinite order in the class group. The quotient by $\varpi$ is the domain
\[
\F_2[x,y,w]_{(x,y,w)}/(xy+w^2),
\]
so inverting $\varpi$ preserves this class group.

The generic fiber over $E_2$ is the split smooth quadric $xy+z^2+1=0$. Its Picard group, geometric Galois action, and Betti cohomology are exactly those used in Section~\ref{d3:sec:persistence}. The argument of Lemma~\ref{d3:lem:components} therefore proves persistence under every denominator from the displayed localization defining $A_2$. Flat base change along $A_2\to A_2\otimes_{\Z_{(2)}}W_2$, followed by inversion of $\varpi$, gives twice the product of $\beta$ with the ruling class, and the analytic regulator in twist three proves that this product has infinite order. Its generic vanishing follows from its original support on $x=0$, as before.
\end{proof}

\subsection{Every higher odd $K$-degree}

The regulator argument also has the following extension. It is stated separately to keep the degree-three counterexample independent of extra indexing conventions.

\begin{proposition}\label{d3:prop:odddegrees}
Let $p$ be an odd prime and let $n=2m-1\geq3$. Set $\varepsilon=(-1)^{m+1}$ and
\[
A_{p,m}=\left(\Z_{(p)}[x,y,z]/(xy+z^2-\varepsilon p)\right)_{(x,y,z)}.
\]
There is an integral element of infinite order in
\[
\ker\bigl(K_n(A_{p,m})\longrightarrow K_n(\Frac A_{p,m})\bigr).
\]
Thus the family with $xy+z^2+p$ works in degrees $3,7,11,\ldots$, and the family with $xy+z^2-p$ works in degrees $5,9,13,\ldots$.
\end{proposition}
\begin{proof}
Take $W=\Z_{(p)}[u]/(u^2-\varepsilon p)$, $E=\Q(u)$, and $\sigma(u)=-u$. The regularity and small-resolution arguments above are unchanged: after base change the equation is $xy+z^2-u^2$, with $d=2u$ in Lemma~\ref{d3:lem:smallresolution}.

If $m$ is even, $E$ is imaginary quadratic. Borel's regulator identifies $K_{2m-1}(E)\otimes\R$ with a one-dimensional real space, and the involution acts by $-1$: its target $\C/\R(m)$ is the imaginary quotient. If $m$ is odd, $E$ is real quadratic. The regulator target is a pair of real lines indexed by the two real embeddings, and $\sigma$ exchanges the two coordinates. Its anti-invariant part is again one-dimensional. These regulator ranks and conjugation conventions are those of \cite[Proposition~12.2]{Borel} and \cite[Section~6]{BunkeTamme}.

In either case choose $b\in K_n(E)$ with nonzero anti-invariant regulator. Since $n-1>0$ is even, $K_{n-1}(\F_p)=0$ and localization lifts $b$ to $K_n(W)$. Antisymmetrizing gives an integral $\beta$ with $\sigma\beta=-\beta$ and
\[
r_m(\beta)\neq0\quad\text{in}\quad\HD^1(\Spec\C,\R(m)).
\]
Push $\beta$ forward from the regular divisor $x=0$ of $\Spec A_{p,m}$. The resulting class is zero in the fraction field. On the split quadric localization it becomes $2\beta\eta$.

The relevant regulator component now has target
\[
\HD^{2(m+1)-n}(Q_s,\R(m+1))
=\HD^3(Q_s,\R(m+1)).
\]
Since $m+1\geq3>\dim_{\C}Q_s$, the same holomorphic Poincar\'e-lemma calculation gives
\[
\HD^3(Q_s,\R(m+1))
=H^2(Q_s,\C)/H^2(Q_s,\R(m+1)).
\]
The proof of Lemma~\ref{d3:lem:Betti}, applied to each geometric boundary curve, therefore gives the same injection. Multiplication by $c_1(D_+)=2\pi i\,h$ identifies the nonzero coefficient in $\C/\R(m)$ with a nonzero coefficient in $\C/\R(m+1)$. No other source regulator component contributes: for a complex point and a positive twist, analytic Deligne cohomology is concentrated in degree one. The target regulator is nonzero at every finite stage, and $K$-theoretic continuity proves infinite order as before.
\end{proof}

\begin{remark}
\label{d3:rem:fixedgaussianallodd}
The same argument gives a \emph{single} arithmetic local
ring with a nonzero rational generic kernel in every odd degree
\(2n-1\ge3\). Fix an odd prime \(p\), a prime \(\mathfrak v\)
of \(\Z_{(p)}[i]\) above \(p\), and put
\[
 F=\Q(i),\quad V=(\Z_{(p)}[i])_{\mathfrak v},\quad
 R=\left(V[x,y,z]/(p+xy+z^2)\right)_{(x,y,z)}.
\]
Here \(V\) is an unramified DVR with uniformizer \(p\), so \(R\)
is regular of dimension three. Set \(E=F(u)\), \(u^2=-p\),
\(W=V[u]\), and let \(\sigma\) fix \(F\) and send \(u\) to \(-u\).
Eisenstein's criterion makes \(W\) a DVR and gives \([E:F]=2\).
Both \(F\) and \(E\) are totally imaginary, with respectively
one and two complex places. Borel's theorem \cite[Proposition~12.2]{Borel} gives
\(\dim_{\Q}K_{2n-1}(F)_{\Q}=1\) and
\(\dim_{\Q}K_{2n-1}(E)_{\Q}=2\), for every \(n\ge2\).
Restriction and transfer satisfy
\(\mathrm{res}\,\mathrm{tr}=1+\sigma\) and
\(\mathrm{tr}\,\mathrm{res}=2\).
Thus the \(\sigma\)-fixed subspace is exactly the restriction
of \(K_{2n-1}(F)_{\Q}\), and the anti-invariant subspace has
dimension one.

Choose a nonzero class there, clear denominators, lift it to
\(K_{2n-1}(W)\), and antisymmetrize. The lift exists because
\(K_{2n-2}(k_V)=0\)
\cite[Chapter~IV, Corollary~1.13]{Weibel}. This produces an actual integral
\(\beta_n\) with \(\sigma\beta_n=-\beta_n\) and a nonzero
regulator at some complex embedding. Its pushforward from
\(R/(x)=W[y]_{(u,y)}\) is generically zero and becomes
\(2\beta_n(1-[\OO(-D_+)])\) on the split generic quadric.
The small resolution and the persistence proved in
Lemma~\ref{d3:lem:Betti} are unchanged. In twist \(n+1>2\), the detecting target is
\[
 \HD^3(Q_s,\R(n+1))
   =H^2(Q_s,\C)/H^2(Q_s,\R(n+1)).
\]
Cup with the ruling multiplies the point coefficient
\(\C/\R(n)\) by \(2\pi i\), giving an isomorphism onto the
ruling line modulo \(\R(n+1)\). The product is nonzero at
every finite stage, hence has infinite order after localization.
Here \(\sigma\) is \emph{not} complex conjugation: it fixes \(i\).
The invariant-rank argument, and multiplication by \(2\pi i\),
work for both parities of \(n\) \cite{Borel,BunkeTamme,Quillen}.
\end{remark}


\subsection{Eisenstein polynomials and arbitrary kernel rank}\label{d3:sec:eisenstein}

The quadratic construction is one instance of a root-lattice calculation. This section records the stronger statement and its geometric proof.

\begin{theorem}\label{d3:thm:eisenstein}
Let $p$ be any prime, let $f\in\Z_{(p)}[z]$ be a monic Eisenstein polynomial of degree $N\geq2$, and put
\[
A_f=\left(\Z_{(p)}[x,y,z]/(xy+f(z))\right)_{(x,y,z)},\qquad
W=\Z_{(p)}[\theta]/(f(\theta)),\qquad K=\Q(\theta).
\]
The ring $A_f$ is regular local of dimension three and has fraction field $\Q(x,z)$. Let $(r_1,r_2)$ be the signature of $K$. For $n\geq3$ odd, consider
\[
g_n:K_n(W)_{\Q}\longrightarrow K_n(A_f)_{\Q},\qquad
\beta\longmapsto i_*(\beta|_{A_f/(x)}),
\]
where $i$ is the regular principal divisor $x=0$. Its image is contained in the generic kernel. Moreover,
\[
\begin{array}{c|c|c}
n&\ker g_n&\dim_{\Q}\operatorname{im}g_n\\ \hline
3\pmod4&0&r_2\\
1\pmod4&\operatorname{res}_{K/\Q}K_n(\Q)_{\Q}&r_1+r_2-1.
\end{array}
\]
In the second row $n\geq5$, and $K_n(W)_{\Q}$ is identified with $K_n(K)_{\Q}$ by localization.
\end{theorem}

\subsubsection{The divisor lattice after splitting}

Let $O$ be a DVR with uniformizer $\pi$, residue field $k$, and fraction field $L$. Suppose that $a_1,\ldots,a_N$ are distinct elements of $\pi O$. Define
\[
B=\left(O[x,y,z]/\left(xy+\prod_{i=1}^N(z-a_i)\right)\right)_{(\pi,x,y,z)},
\qquad P_i=(x,z-a_i).
\]
The same normality argument as in Lemma~\ref{d3:lem:smallresolution} shows that $B$ is normal and regular off its closed point. Since $B[1/x]$ is a localization of $O[x,x^{-1},z]$, the primes $P_i$ generate $\Cl(B)$, with the relation
\[
\divisor(x)=P_1+\cdots+P_N.
\]
To prove that there are no further relations, successively blow up the strict transforms of $P_1,\ldots,P_{N-1}$. The resulting regular scheme $Y$ has $N$ charts, each a localization of $O[u_i,v_i]$, on which
\begin{align*}
z&=a_i+u_i v_i,\\
x&=u_i\prod_{k<i}(z-a_k),&
y&=-v_i\prod_{k>i}(z-a_k).
\end{align*}
For completeness, before the $i$th blowup the unresolved chart is
\[
O[s_{i-1},y,z]/\left(s_{i-1}y+\prod_{k\geq i}(z-a_k)\right),
\quad s_0=x.
\]
Blowing up $(s_{i-1},z-a_i)$ gives the displayed $i$th polynomial chart and the next unresolved chart, where $s_i=s_{i-1}/(z-a_i)$. These are the actual Rees charts inside the common function field. A future strict transform $P_j$, $j>i$, is absent from the resolved chart, since its generic point has $z=a_j$ and cannot satisfy $u_i=0$, $z=a_i$. Thus the successive centers glue globally and do not modify previous resolved charts.

The adjacent transition is
\[
u_{i+1}=v_i^{-1},\qquad
v_{i+1}=v_i(a_i-a_{i+1}+u_i v_i).
\]
The exceptional fiber is a reduced chain $C_1,\ldots,C_{N-1}$ of projective lines. Its ideal is $(\pi,u_1)$ on the first chart, $(\pi,v_N)$ on the last, and $(\pi,u_i v_i)$ on an interior chart. In particular it has dimension one, and the resolution has no exceptional divisors. The two charts on $C_j$ have coordinates $v_j,u_{j+1}$ with $u_{j+1}=v_j^{-1}$.

Write $\widetilde P_\ell$ for the strict transform of $P_\ell$. On the $i$th chart its defining equation is
\[
d_{\ell i}=\begin{cases}
z-a_\ell,&\ell<i,\\
u_i,&\ell=i,\\
1,&\ell>i.
\end{cases}
\]
These equations also show that no vertical divisor has been added: when $\ell<i$, the quotient by $u_i v_i+a_i-a_\ell$ is a domain. Taking transition functions of $\OO_Y(\widetilde P_\ell)$ gives
\[
\deg\bigl(\OO_Y(\widetilde P_\ell)|_{C_j}\bigr)
=\delta_{\ell,j+1}-\delta_{\ell,j}.
\]
Therefore a principal relation $\sum n_iP_i$ forces $n_{j+1}=n_j$ for every $j$. Conversely equal coefficients give a power of $x$. We have proved the integral presentation
\begin{equation}\label{d3:eq:rootlattice}
\Cl(B)=\Z^N/\Z(1,\ldots,1).
\end{equation}

Since $B/\pi B$ is the domain $k[x,y,z]_{(x,y,z)}/(xy+z^N)$, inverting $\pi$ removes just the principal prime $(\pi)$. Thus \eqref{d3:eq:rootlattice} is also $\Pic(B[1/\pi])$. On the full generic surface
\[
Q_L=\Spec L[x,y,z]/\left(xy+\prod_i(z-a_i)\right),
\]
divisor localization at $x$ gives the same presentation: the units of $L[x,x^{-1},z]$ are $L^\times x^{\Z}$. Consequently
\begin{equation}\label{d3:eq:generalpiciso}
\Pic(Q_L)\xrightarrow{\sim}\Pic(B[1/\pi]).
\end{equation}

\subsubsection{Betti classes and finite localizations}

Assume that $L$ embeds in $\C$. The complex surface $Q_L(\C)$ is the blowup of $\A^2_{\C}$ at the $N$ points $(0,a_i)$ with the strict transform of the line $x=0$ removed. Indeed, blow up the ideal $(x,\prod(z-a_i))$; its $x$-chart is precisely $Q_L$, and the complement is that strict transform.

The blowup has $H^2=\Z^N$ on the exceptional classes $E_i$, and no other positive-degree cohomology. The removed line is isomorphic to $\A^1$, with class $-\sum E_i$, since its sum with the exceptional divisors is the principal divisor of $x$. The Gysin sequence therefore gives
\[
H^2(Q_L(\C),\Z)=\Z^N/\Z(1,\ldots,1),
\]
and first Chern class identifies this with $\Pic(Q_L)$. All the root divisors are defined over $L$, so the absolute Galois action on the geometric Picard lattice is trivial.

For any denominator allowed in the localization $B[1/\pi]$, every removed $L$-prime divisor has zero Picard class by \eqref{d3:eq:generalpiciso}. Its geometric components are a single Galois orbit; they have equal classes in a torsion-free lattice, and their sum is zero. Each component therefore has zero class. The degree-two support argument of Section~\ref{d3:sec:persistence} proves
\[
H^2(Q_L(\C),\R)\hookrightarrow H^2((Q_L)_s(\C),\R)
\]
for every such finite denominator $s$. The analytic Deligne injection in every twist greater than two follows exactly as in Section~\ref{d3:sec:Delignedetection}.

\subsubsection{The arithmetic rank calculation}

We now prove Theorem~\ref{d3:thm:eisenstein}. Eisenstein's condition makes $W$ a DVR with uniformizer $\theta$ and residue field $\F_p$. It also makes the defining equation of $A_f$ have nonzero linear term $p$ in its regular four-dimensional ambient local ring. Thus $A_f$ is regular local of dimension three, and
\[
A_f/(x)=W[y]_{(\theta,y)}.
\]
Solving for $y$ after inverting $x$ gives $\Frac A_f=\Q(x,z)$. Generic vanishing of $g_n$ follows from localization.

Choose a splitting field $L$ of $f$ and a prime above $p$, and let $O$ be the corresponding coefficient DVR. Every root $a_i$ has positive valuation. The map $A_f\to B$ used above is flat: $O$ is torsion-free over $\Z_{(p)}$, followed by localization. There is no need for $O$ itself to be a finite $\Z_{(p)}$-module. After inverting $\pi$, flat base change expresses the Gysin image as
\begin{equation}\label{d3:eq:generalbasechange}
g_n(\beta)|_{B[1/\pi]}=
\sum_{i=1}^N \sigma_i(\beta)\,\eta_i,
\qquad \eta_i=[\OO_{(x,z-a_i)}],
\end{equation}
where $\sigma_i:K\hookrightarrow L$ are the embeddings indexed by the roots.

Write $n=2m-1$. At a fixed complex embedding of $L$, the regulator in weight $m+1$ sends \eqref{d3:eq:generalbasechange} to the vector of Borel coefficients
\[
\bigl(r_m(\sigma_1\beta),\ldots,r_m(\sigma_N\beta)\bigr)
\quad\text{modulo the diagonal line},
\]
with the common nonzero Tate normalization coming from first Chern class. Its target is
\[
\HD^3((Q_L)_s,\R(m+1))
=H^2((Q_L)_s,\C)/H^2((Q_L)_s,\R(m+1)),
\]
because $m+1\geq3$. The preceding finite-localization injection ensures that a nonzero vector modulo the diagonal remains nonzero at every finite stage, hence in $K_n(B[1/\pi])_{\Q}$.

If $m$ is even, the two coefficients at conjugate embeddings have opposite signs, and a real embedding contributes zero. Borel's theorem \cite[Proposition~12.2]{Borel} identifies the real scalar extension of $K_n(K)_{\Q}$ with this $r_2$-dimensional coefficient space. A vector in it is diagonal only when it is zero. Thus $g_n$ is injective for $n\equiv3\pmod4$.

If $m$ is odd, conjugate coefficients agree, and the coefficient space has dimension $r_1+r_2$. Its diagonal line is exactly the real span of restriction from $K_n(\Q)_{\Q}$, which is one-dimensional. Since this line is generated by a rational class and Borel's real regulator is injective \cite[Section~6]{BunkeTamme}, its intersection with $K_n(K)_{\Q}$ is exactly that rational restriction line. Finally, that line really is killed by $g_n$: its classes lift from $K_n(\Z_{(p)})_{\Q}$, and the projection formula multiplies them by
\[
[\OO_{V(x)}]=[\OO]-[\OO]=0\quad\text{in }K_0(A_f).
\]
This proves both the exact kernel and the stated image dimensions.

\begin{corollary}
For every prime $p$ and integer $r\geq1$, the three-dimensional regular local ring
\[
\left(\Z_{(p)}[x,y,z]/(xy+z^{2r}+p)\right)_{(x,y,z)}
\]
has a rational degree-three generic kernel of dimension at least $r$. Thus these kernel dimensions are unbounded in fixed local dimension three. Moreover the single ring with $xy+z^3+p$ has a nonzero rational generic kernel in every odd degree $n\geq3$: its number field has signature $(1,1)$.
\end{corollary}

\section{Explicit triples of units at every residue prime}\label{d3:sec:explicitallprimes}

The preceding geometry also detects explicit symbols. The following formulas use the Dennis--Stein convention of Section~\ref{d3:sec:explicit}; unlike \eqref{d3:eq:explicitfront}, they require no prefactor.

\begin{theorem}\label{d3:thm:explicitallprimes}
For an odd prime $p$, put
\[
 A_p=\left(\Z_{(p)}[x,y,z]/(xy+\Phi_{2p}(z-1))\right)_{(x,y,z)}.
\]
Then
\begin{equation}\label{d3:eq:explicitoddprime}
 c_p=\left\{2-z,\ 1-x,\
 1+(1-x)yz\bigl((z-1)^p-1\bigr)\right\}\in K_3(A_p)
\end{equation}
has infinite order and maps to zero in $K_3(A_p[1/x])$ integrally.
For $p=2$, put
\[
 A_2=\left(\Z_{(2)}[x,y,t]/(xy+t^2+1)\right)_{(2,x,y,t-1)},
 \qquad H(t)=t^{10}-t^8+t^6-t^4+t^2-1.
\]
The same assertions hold for
\begin{equation}\label{d3:eq:explicittwoprime}
 c_2=\left\{1+t+t^2,\ 1-x,\ 1+(1-x)yH(t)\right\}\in K_3(A_2).
\end{equation}
All entries in these formulas are units. Each $A_p$ is regular local of dimension three, with residue field $\F_p$, so the displayed classes give nonzero rational generic kernels.
\end{theorem}

\subsection{The integral identities}

We first record the common calculation. Suppose $1-xb=u^n$, and put
\[
 \delta=\langle x,b\rangle,\qquad \lambda=1-x,\qquad
 \mu=1+\lambda b.
\]
Assume $u,g,\lambda,\mu$ are units and $\{g,u^n\}=0$ in $K_2$.
Relation \eqref{d3:eq:DS2}, together with $\langle x,1\rangle=0$, gives
\[
 \delta=\langle x,b+1-xb\rangle=\langle x,\mu\rangle
       =\{\lambda u^n,\mu\},
 \qquad 1-x\mu=\lambda u^n.
\]
Consequently
\begin{equation}\label{d3:eq:allprimesDS}
 [g]\delta=\{g,\lambda,\mu\},\qquad
 [g]\delta|_{D(x)}=\{g,x,u^n\}=-[x]\{g,u^n\}=0.
\end{equation}
These are integral identities in Quillen $K$-theory.

For odd $p$, write $a=z-1$ and set
\[
 H(a)=(a+1)(a^p-1),\qquad b=yH(a),\qquad g=1-a,\qquad n=2p.
\]
The polynomial identity
\[
 a^{2p}-1=H(a)\Phi_{2p}(a)=-xyH(a)
\]
gives $1-xb=a^{2p}$. At the closed point,
$a=-1$, $g=2$, and $\lambda=\mu=1$ in $\F_p$; these elements are units.
Steinberg gives $\{1-a,a^{2p}\}=2p\{1-a,a\}=0$.
Thus \eqref{d3:eq:allprimesDS} proves \eqref{d3:eq:explicitoddprime} and its integral vanishing after inverting $x$.
Moreover,
\[
 \Phi_{2p}(z-1)
 =\sum_{k=0}^{p-1}(1-z)^k
 =\sum_{j=1}^{p}(-1)^{j+1}\binom pj z^{j-1}
\]
is monic Eisenstein of degree $p-1$, with constant term $p$. The regularity assertion follows as in Theorem~\ref{d3:thm:eisenstein}.

At two, take $u=t$, $n=12$, $b=yH(t)$ and $g=1+t+t^2$.
Here
\[
 (t^2+1)H(t)=t^{12}-1,\qquad 1-xb=t^{12}.
\]
All required units have nonzero residue at $t=1$. To prove
$\{g,t^{12}\}=0$ integrally, use the finite \'etale local extension
$A_2'=A_2[\zeta]$, where $\zeta^2+\zeta+1=0$.
The factor $1-\zeta t$ is a unit, since
\[
 (1-\zeta t)(1-\zeta^2t)=g,
 \qquad (\zeta t)^{12}=t^{12}.
\]
Steinberg and the transfer projection formula therefore give
\[
 0=\operatorname{Tr}_{A_2'/A_2}
       \{1-\zeta t,t^{12}\}
   =\{\operatorname{Norm}_{A_2'/A_2}(1-\zeta t),t^{12}\}
   =\{g,t^{12}\}.
\]
There is no factor of two in this formula. Indeed, multiplication by
$1-\zeta t$ in the basis $1,\zeta$ has matrix
$\left(\begin{smallmatrix}1&t\\-t&1+t\end{smallmatrix}\right)$,
elementary-equivalent to $\operatorname{diag}(1,g)$; its $K_1$ transfer is exactly $[g]$.
See \cite[Chapter~III, Definition~5.11 and relations (D1)--(D3)]{Weibel}.
Equation \eqref{d3:eq:allprimesDS} now proves \eqref{d3:eq:explicittwoprime}.
Finally, with $v=t-1$, the equation becomes
$xy+v^2+2v+2=0$, which proves regularity at the specified maximal ideal.

\subsection{A supported lift on a finite surface}

For odd $p$ let $E=\Q(\zeta_{2p})$, and at two let $E=\Q(i)$.
Use the following smooth split surfaces and finite opens:
\[
\begin{array}{c|c|c}
 & Q & U\\ \hline
p\text{ odd}
 & \Spec E[x,y,a]/(xy+\Phi_{2p}(a))
 & D(a(1-a))\\
p=2
 & \Spec E[x,y,t]/(xy+t^2+1)
 & D(t(1+t+t^2)).
\end{array}
\]
The Dennis--Stein product $\alpha_p=[g]\delta$ is defined on $U$
and vanishes on $U[1/x]$ by the preceding calculation.
At two the same norm argument uses $E(\zeta)/E$.
The entire divisor $D_x=V(x)\subset Q$ lies in $U$:
its coordinates are primitive $2p$-th roots in the odd case, and
$t=\pm i$, $g(\pm i)=\pm i$, at two.
Localization and d\'evissage supply $\xi_p\in K_3(D_x)$ such that its
pushforward to $U$ is $\alpha_p$. Pushing the same class into $Q$ gives
\begin{equation}\label{d3:eq:allprimeslift}
 \widetilde c_p=i_*\xi_p\in K_3(Q),
 \qquad \widetilde c_p|_U=\alpha_p.
\end{equation}
Only this supported lift will be needed; its coefficients need not be specified.

We calculate the regulator of $\alpha_p$ using its Dennis--Stein expression.
The entries $\lambda,\mu$ need only be units in the arithmetic local ring;
they are not needed on the analytic cycle below.
On $D(x)\subset U$ one has $\delta=\{x,u^n\}$, and its holomorphic
weight-two curvature extends to $U$ as
\[
 \omega=n\,\frac{dx}{x}\wedge\frac{du}{u}
        =-\frac{dx\wedge db}{1-xb}.
\]
Whenever $g=e^L$ on an analytic neighborhood, the interval product
calculation \eqref{d3:eq:Lomega} gives
$r_3(\alpha_p)=[-L\omega]$ there.
This uses the multiplicativity and unit normalization of
\cite[Proposition~2.8, Theorem~2.31 and Lemma~4.5]{BunkeTamme}.

\subsection{Two nonzero matching-sphere periods}

We use a single endpoint formula. Let $\gamma:[-1,1]\to\C$ be a smooth
embedded arc joining two distinct simple roots of the relevant polynomial
$P(u)$, avoiding the other roots. On $xy+P(u)=0$, take
\[
 u=\gamma(\tau),\qquad
 x=\sqrt{|P(\gamma(\tau))|}\,e^{i\vartheta},\qquad
 y=-P(\gamma(\tau))/x.
\]
Adding $x=y=0$ at both ends gives a smooth sphere $\Sigma_\gamma$,
oriented by $d\tau\wedge d\vartheta$. Smoothness at an end follows by
using $|x|^2$ as a radial coordinate and writing $y$ as a smooth
nonzero phase times $\overline x$.

Suppose $L$ and $F$ are holomorphic near the arc, with
$e^L=g$ and $dF=-L\,du/u$. On the sphere away from its ends,
\[
 L\omega=n\,d\left(F\,\frac{dx}{x}\right).
\]
Stokes' theorem on a truncated cylinder, followed by shrinking its
two missing caps, gives
\begin{equation}\label{d3:eq:cyclotomicendpoint}
 \int_{\Sigma_\gamma}r_3(\alpha_p)
 =-2\pi i n\bigl(F(\gamma(1))-F(\gamma(-1))\bigr)
 \pmod{\R(3)}.
\end{equation}
The caps contribute zero because $L\omega$ extends smoothly.
This calculation uses $dx/x$ on the cylinder and never requires a
logarithm of $x$. Since $\R(3)=i\R$, a nonzero real period detects
the regulator.

For odd $p$, put $\theta=\pi/p$, $r=e^{i\theta}$, and choose
\[
 \gamma(\tau)=\cos\theta+i\sin\theta\,\tau .
\]
This chord joins $\overline r$ to $r$, lies in $0<\Re a<1$, and its
interior lies strictly inside the unit circle. It therefore avoids
$0,1$ and all other roots. On one neighborhood of the whole arc,
take the principal branches
\[
 L=\operatorname{Log}(1-a),\qquad F=\operatorname{Li}_2(a).
\]
Here $\operatorname{Li}_2$ is the classical dilogarithm, initially
given by $\sum_{r\geq1}a^r/r^2$ on the unit disc. The chosen
branches satisfy the required differential equation.
Writing $D$ for the Bloch--Wigner dilogarithm,
$F(r)-F(\overline r)=2iD(r)$, so \eqref{d3:eq:cyclotomicendpoint} yields
\begin{equation}\label{d3:eq:oddcyclotomicperiod}
 \int_{\Sigma_\gamma}\Re r_3(\alpha_p)
 =8\pi p\,D(e^{i\pi/p})>0.
\end{equation}
For example, positivity follows from
$D(e^{i\theta})=-\int_0^\theta\log(2\sin(v/2))\,dv$
for $0<\theta\leq\pi/3$.

At two, the straight segment from $-i$ to $i$ would meet $t=0$.
Instead use
\[
 \gamma(\tau)=\tfrac12(1-\tau^2)+i\tau .
\]
Its interior has positive real part and modulus less than one.
Thus the arc avoids $0$, the roots of $g$, and the other twelfth
roots of unity. With $\zeta=e^{2\pi i/3}$, both
$1-\zeta\gamma(\tau)$ and $1-\zeta^2\gamma(\tau)$ have positive
real part, including at the endpoints. Hence one neighborhood
supports the principal branches
\[
 L=\operatorname{Log}(1-\zeta t)+\operatorname{Log}(1-\zeta^2t),
 \qquad
 F=\operatorname{Li}_2(\zeta t)+\operatorname{Li}_2(\zeta^2t).
\]
Here $e^L=g$, $dF=-L\,dt/t$, and
$F(-i)=\overline{F(i)}$.
The absolutely convergent dilogarithm series on the unit circle gives
\[
 \operatorname{Li}_2(w)+\operatorname{Li}_2(\zeta w)
 +\operatorname{Li}_2(\zeta^2w)=\tfrac13\operatorname{Li}_2(w^3).
\]
Taking imaginary parts at $w=i$ therefore gives
$D(\zeta i)+D(\zeta^2i)=-\tfrac43D(i)$, and
\begin{equation}\label{d3:eq:twocyclotomicperiod}
 \int_{\Sigma_\gamma}\Re r_3(\alpha_2)
 =-64\pi D(i)\ne0.
\end{equation}
The number $D(i)=\sum_{k\geq0}(-1)^k/(2k+1)^2$ is positive.

\subsection{Persistence in the arithmetic local ring}

We finish the proof of Theorem~\ref{d3:thm:explicitallprimes}.
For odd $p$, take the coefficient DVR
$\OO=\Z_{(p)}[\zeta_{2p}]$ and use the original coordinate $z=a+1$.
All roots $1+r_i$ of $\Phi_{2p}(z-1)$ lie in its maximal ideal.
At two take $\OO=\Z_{(2)}[i]$ and the coordinate $v=t-1$; its
two roots $\pm i-1$ again lie in the maximal ideal.
The split local rings therefore satisfy \eqref{d3:eq:rootlattice}
and \eqref{d3:eq:generalpiciso}. The Betti and analytic Deligne
persistence argument of Section~\ref{d3:sec:eisenstein} applies to
every permitted finite localization of $Q$.

The open $U$ is one of those permitted localizations:
$a(1-a)$ has nonzero residue $-2$ for odd $p$, and $tg$ has
nonzero residue $1$ at two. On the generic fiber $T$ of the split
arithmetic local ring, the unit triple $c_p$, the Dennis--Stein
product $\alpha_p$, and the lift $\widetilde c_p$ in
\eqref{d3:eq:allprimeslift} have the same image.
Equations \eqref{d3:eq:oddcyclotomicperiod} and
\eqref{d3:eq:twocyclotomicperiod} show that the regulator of
$\widetilde c_p$ is nonzero already on $Q$.
The cited persistence makes it nonzero on every further permitted
finite open. A torsion relation on $c_p$ would, after passage to
$T$, descend to one such finite open by \eqref{d3:eq:continuity},
contradicting that real regulator.
Thus $c_p$ has infinite order, as asserted.


\section{Completion and arithmetic \texorpdfstring{$p$}{p}-adic detection}\label{d3:sec:completed}

There is a separate route to completed local rings. It uses the arithmetic comparison theorem for dagger varieties of Colmez--Nizio\l, rather than an interchange of a localization limit with a $p$-adic inverse limit. We spell out the map and the order of limits, and distinguish actual $K_3$-classes from classes in a completed spectrum.

\subsection{A ruling-cup injection on the completed node}

\begin{lemma}\label{d3:lem:completedcup}
Let $E/\Q_p$ be finite, let $W=\OO_E$, and let $0\ne d\in\mathfrak m_W$. Put
\[
\widehat B_d=W[[x,y,w]]/(xy+w(w+d)),\qquad
T_d=\Spec\widehat B_d[1/d].
\]
The ideal $I=(x,w)$ is invertible on $T_d$, and cup product induces an injection
\[
H^1(E,\Q_p(2))\longrightarrow
H^3_{\mathrm{cont},\mathrm{\acute et}}(T_d,\Q_p(3)),
\qquad b\longmapsto b\smile c_1(I).
\]
Here continuous cohomology is formed by taking the derived inverse limit of finite-coefficient cohomology and then inverting $p$.
\end{lemma}
\begin{proof}
At a prime containing $x,w$, the element $w+d$ is a unit, and the equation gives $I=(x)$. Outside $V(I)$ the ideal is the unit ideal. It is therefore an invertible ideal, generated globally by two elements, and defines a morphism
\[
f_T:T_d\longrightarrow\PP^1_E,\qquad f_T^*\OO(1)=I.
\]

Consider the smooth dagger affinoid
\[
U^\dagger=\operatorname{Sp}^\dagger
E\langle X,Y,Z\rangle^\dagger/(XY+Z(Z+1)).
\]
The two roots $0,-1$ are distinct in every residue characteristic. Fix
\[
\rho=|d|^{-1/2}\in\sqrt{|E^\times|},\qquad 1<\rho<|d|^{-1},
\]
and let $U_\rho$ be the closed radius-$\rho$ domain of the same algebraic quadric. The substitutions
\begin{equation}\label{d3:eq:fixedradius}
x=dX,\qquad y=dY,\qquad w=dZ
\end{equation}
define a homomorphism $\widehat B_d[1/d]\to\OO(U_\rho)$. Indeed every formal series with coefficients in $W$ converges on this \emph{same fixed domain}: terms of total degree $n$ have weighted norm at most $(|d|\rho)^n$, which tends to zero. A fixed power of $d^{-1}$ does not change the convergence radius.

For completeness, a noetherian integral model of this fixed affinoid is
\[
 A_\rho=W\langle a,b,c,r,s,t\rangle/
 (a^2-dr,\ b^2-ds,\ c^2-dt,\ ab+c(c+d)).
\]
Its generic fiber is $U_\rho$ under $X=a/d$, $Y=b/d$, $Z=c/d$:
the last three auxiliary variables impose
$|dX^2|,|dY^2|,|dZ^2|\leq1$.
The elements $a,b,c$ are topologically nilpotent in $A_\rho$,
since their squares lie in $dA_\rho$. Substitution therefore defines
$\widehat B_d\to A_\rho$ before inverting $d$.
The pair $(A_\rho,(p))$ is a noetherian henselian pair; removing its
$W$-torsion, if desired, preserves the generic fiber.

Pull scheme cohomology to $\Spec\OO(U_\rho)$, then to the affinoid \'etale topos. The latter map exists and is compatible with all finite coefficients, including $p$-power coefficients; see the Gabber--Fujiwara comparison in \cite[Theorem~6.2.1]{Achinger}, applied to a complete noetherian integral model of this one affinoid. At this fixed $\rho$, take the derived inverse limit over $p^\nu$ and invert $p$. Continuous \'etale and pro-\'etale cohomology agree for the affinoid, as recalled in \cite[Section~3.2.3]{ColmezNiziol}. Finally map into the homotopy colimit over shrinking strict neighborhoods, which is the dagger pro-\'etale complex in that reference. This constructs
\begin{equation}\label{d3:eq:schemetodagger}
H^3_{\mathrm{cont},\mathrm{\acute et}}(T_d,\Q_p(3))
\longrightarrow H^3_{\mathrm{pro\acute et}}(U^\dagger,\Q_p(3)).
\end{equation}
It respects pullback from $\PP^1_E$ and first Chern classes. We have not exchanged the inverse limit over coefficients with a direct limit over opens. The reverse comparison from dagger \'etale to dagger pro-\'etale cohomology is not assumed.

The ruling morphism $f:U^\dagger\to\PP^1_E$ obtained from $(X,Z)$ has an explicit two-chart description. Above the two standard closed discs of $\PP^1$ its charts are dagger bidiscs
\begin{align*}
U_0: &(X,t),&Z&=Xt,&Y&=-t(1+Xt),\\
U_\infty: &(s,Y),&Z&=-1-sY,&X&=-s(1+sY).
\end{align*}
Their intersection is a dagger unit annulus times a disc, with
\[
s=t^{-1},\qquad Y=-t-Xt^2.
\]
The bidiscs have no positive-degree de Rham cohomology, and the intersection has $H^1=E\,d\log t$. These facts follow by coefficientwise integration of overconvergent power and Laurent series: division by integers costs only a slight reduction of the overconvergence radius, and the $t^{-1}dt$ term remains. Coherent acyclicity for dagger affinoids \cite[Proposition~3.1]{GrosseKlonne} and Mayer--Vietoris give
\[
H^2_{\mathrm{dR}}(U^\dagger)=E[d\log t],\qquad H^3_{\mathrm{dR}}(U^\dagger)=0,
f^*:H^2_{\mathrm{dR}}(\PP^1_E)\xrightarrow{\sim}H^2_{\mathrm{dR}}(U^\dagger).
\]

The arithmetic Theorem~1.3(2), or Theorem~7.9(2) together with Proposition~5.9(2), of \cite{ColmezNiziol} gives a natural boundary
\[
\partial_3:\widetilde H^2(R\Gamma_{\mathrm{dR}}(U^\dagger)/F^3)
\longrightarrow\widetilde H^3_{\mathrm{pro\acute et}}(U^\dagger,\Q_p(3)).
\]
We verify that it is an isomorphism, so no passage from a monomorphism in a topological heart to ordinary cohomology is needed. The space $U^\dagger$ has the smooth weak formal model
$\OO_E\langle X,Y,Z\rangle^\dagger/(XY+Z(Z+1))$; the two factors $Z,Z+1$ generate the unit ideal. The ideal $(X,Z)$ is invertible on this model and extends the ruling morphism to the weak formal projective line over $\OO_E$. Let $F_0=\Frac W(k_E)$. The functorial Hyodo--Kato to de Rham comparison of \cite[equation~(5.4) and Propositions~5.5--5.6]{ColmezNiziol} becomes a strict quasi-isomorphism after tensoring from $F_0$ to $E$. Our de Rham calculation, and its vanishing in degree three, therefore give
\[
f^*:H^2_{\mathrm{HK}}(\PP^1_E)\xrightarrow{\sim}H^2_{\mathrm{HK}}(U^\dagger),
\qquad H^3_{\mathrm{HK}}(U^\dagger)=0.
\]
Indeed scalar extension to the field $E$ is faithful. Both degree-two groups are the rank-one slope-one isocrystal $F_0h$ with $\varphi=p\sigma$ and $N=0$. Its $\varphi=p^2$ eigenspace is zero by valuation. The second exact sequence of Theorem~1.3(2) now makes the term after $\widetilde H^3_{\mathrm{pro\acute et}}$ vanish, on both $U^\dagger$ and $\PP^1_E$. Since $F^3=0$ on this surface, the first exact sequence gives isomorphisms
\begin{align*}
H^2_{\mathrm{dR}}(U^\dagger)&\xrightarrow{\sim}\widetilde H^3_{\mathrm{pro\acute et}}(U^\dagger,\Q_p(3)),\\
H^2_{\mathrm{dR}}(\PP^1_E)&\xrightarrow{\sim}\widetilde H^3_{\mathrm{pro\acute et}}(\PP^1_E,\Q_p(3)).
\end{align*}
Here the sources are classical. Applying the classical-part functor of \cite[Section~3.1.1]{ColmezNiziol} to these isomorphisms gives the corresponding isomorphisms on ordinary cohomology. Naturality and the de Rham ruling isomorphism show that
\[
f^*:H^3_{\mathrm{\acute et}}(\PP^1_E,\Q_p(3))
\longrightarrow H^3_{\mathrm{pro\acute et}}(U^\dagger,\Q_p(3))
\]
is an isomorphism. Proper pro-\'etale and analytic \'etale cohomology
are identified by \cite[Section~7.2.4]{ColmezNiziol}. For the scheme
$\PP^1_E$, the torsion algebraic--analytic comparison of
\cite[Theorem~6.1.1]{BerkovichEtale} gives the same identification,
first for every $p^\nu$ and then by derived inverse limit.
These comparisons preserve pullbacks and Kummer Chern classes.
For finite coefficients, put $\Lambda=\Z/m$, where $m>0$ is
invertible in $E$, and write
$\mathrm{pr}\colon\PP^1_E\to\Spec E$.
The Kummer boundary defines
$\xi_m=c_1(\OO(1))\in H^2_{\mathrm{\acute et}}(\PP^1_E,\Lambda(1))$
\cite[Example~7.9(a)]{MilneEtale}.
For every integer $r$, pullback and cup product give a morphism
of \'etale complexes
\begin{equation}\label{d3:eq:finitePonebundle}
 \Lambda(r)\oplus\Lambda(r-1)[-2]
 \xrightarrow{(\mathrm{pr}^*,\,\mathrm{pr}^*(-)\smile\xi_m)}
 R\mathrm{pr}_*\Lambda(r).
\end{equation}
To check that it is a quasi-isomorphism, take a geometric stalk.
Proper base change identifies the target with the cohomology
complex of a geometric projective line, whose only nonzero
cohomology is in degrees zero and two
\cite[Proposition~14.2 and Theorems~17.7 and~17.10]{MilneEtale}.
The degree-zero map is the identity.
In degree two, the Kummer calculation sends the generator
$\OO(1)$ of $\Pic(\PP^1)=\Z$ to a generator modulo $m$;
this is also the Chern normalization in
\cite[Theorems~23.2 and~23.3(b)]{MilneEtale}.
Thus the map is a quasi-isomorphism on geometric stalks and
hence on $\Spec E$. This argument works over any field in
which $m$ is invertible.

The Kummer classes make these maps compatible for $m=p^\nu$.
Assume $p\ne\operatorname{char}(E)$.
Applying $R\Gamma(E,-)$, then $R\varprojlim_\nu$, and then
inverting $p$ preserves this finite direct sum; all
$\Q_p$-cohomology groups here, including $H^1(E,\Q_p(2))$,
are understood as continuous cohomology obtained by this
construction. Taking $r=3$ gives a split injection
\[
 H^1(E,\Q_p(2))
 \xrightarrow{\ \smile c_1(\OO(1))\ }
 H^3_{\mathrm{cont},\mathrm{\acute et}}(\PP^1_E,\Q_p(3)).
\]
Composing the ruling-cup map on $T_d$ with
\eqref{d3:eq:schemetodagger} is this split injection followed
by the already proved isomorphism $f^*$.
Hence the ruling-cup map on $T_d$ is injective.
\end{proof}

\subsection{Simultaneous ruling classes}
A map from a completed quadratic node isolates any chosen pair of roots.
The resulting detector records their coefficients modulo the diagonal.

\begin{lemma}\label{d3:lem:completedrootcup}\label{d3:lem:completedtamecup}
Let $E/\Q_p$ be finite and let $a_1,\ldots,a_N$ be distinct elements of $\mathfrak m_{\OO_E}$, with $N\geq2$. Put
\[
\widehat B=\OO_E[[x,y,z]]/
\left(xy+\prod_{i=1}^N(z-a_i)\right),
\qquad T=\Spec\widehat B[1/p],
\qquad J_i=(x,z-a_i).
\]
The $J_i$ are invertible ideals on $T$, their product is $(x)$, and the map
\begin{equation}\label{d3:eq:completedtamecup}
\frac{H^1(E,\Q_p(2))^N}
{\operatorname{diag}H^1(E,\Q_p(2))}
\longrightarrow H^3_{\mathrm{cont},\mathrm{\acute et}}(T,\Q_p(3)),
\qquad
[(\alpha_i)]\longmapsto\sum_i\alpha_i\smile c_1(J_i)
\end{equation}
is injective.
\end{lemma}

\begin{proof}
At a point with $x=z-a_i=0$, every other root factor is a unit, since each nonzero $a_i-a_k$ is invertible after inverting $p$. The equation then gives $J_i=(x)$ locally. Away from its vanishing locus $J_i$ is the unit ideal. This proves invertibility and, by the same local calculation, $\prod_iJ_i=(x)$.

Fix distinct indices $i,j$, and set
\[
d=a_j-a_i,\qquad H_{ij}(z)=\prod_{k\ne i,j}(z-a_k),
\]
where the empty product is $1$. Consider the completed quadratic node
\[
C_{ij}=\OO_E[[u,v,w]]/(uv+w(w-d)),\qquad
T_{ij}=\Spec C_{ij}[1/p].
\]
There is a continuous homomorphism of complete local $\OO_E$-algebras
\begin{equation}\label{d3:eq:pairwiseformalmap}
\widehat B\longrightarrow C_{ij},\qquad
x\longmapsto u,\quad
y\longmapsto H_{ij}(a_i+w)v,\quad
z\longmapsto a_i+w.
\end{equation}
All three images lie in the maximal ideal, so substitution defines every formal series. The defining relation maps to
\[
H_{ij}(a_i+w)\bigl(uv+w(w-d)\bigr)=0.
\]
In particular no inverse of the cofactor $H_{ij}$ is used. Inverting $p$ gives a morphism $q_{ij}:T_{ij}\to T$.

The generated ideals obtained from $J_i,J_j$ are
\[
I_i=(u,w),\qquad I_j=(u,w-d).
\]
For $k\ne i,j$, put $b=a_k-a_i$, so $b\ne0,d$. In $C_{ij}[1/p]$ the identity
\[
b(b-d)=-uv-(w-b)(w+b-d)
\]
shows that $(u,w-b)$ is the unit ideal, since $b(b-d)\in E^\times$.

These computations identify the pullbacks of the line bundles even though $q_{ij}$ need not be flat. The pullback of each invertible $J_k$ surjects onto its generated ideal. For $k=i,j$ that image is an invertible ideal by the quadratic-node calculation, and for every other $k$ it is the unit ideal. A surjection between invertible modules is an isomorphism. Hence
\[
q_{ij}^*J_i\cong I_i,\qquad
q_{ij}^*J_j\cong I_j,\qquad
q_{ij}^*J_k\cong\OO\quad(k\ne i,j).
\]
Also $I_iI_j=(u)$, so $c_1(I_j)=-c_1(I_i)$.

Pulling the cup sum to $T_{ij}$ therefore gives
\[
q_{ij}^*\left(\sum_k\alpha_k\smile c_1(J_k)\right)
=(\alpha_i-\alpha_j)\smile c_1(I_i).
\]
Lemma~\ref{d3:lem:completedcup}, with parameter $-d$, makes the right-hand cup map injective. Its prescribed inversion is the same as inversion of $p$, because $d$ is a nonzero element of the maximal ideal of the DVR $\OO_E$. Thus vanishing of the original sum forces $\alpha_i=\alpha_j$. Taking the pairs $(1,j)$ proves that its kernel is the diagonal. The diagonal already vanishes because $\prod_iJ_i=(x)$, proving \eqref{d3:eq:completedtamecup}.
\end{proof}

\subsection{Actual local coefficients and the line-bundle formula}

\begin{lemma}\label{d3:lem:continuouschern}
Let $E/\Q_p$ be finite and let $R$ be an $E$-algebra. The higher Chern maps used below have natural values in continuous \'etale cohomology, and their finite-coefficient line-bundle product identity holds in that continuous target. No vanishing of $\varprojlim^1$ for $\Spec R$ is required.
\end{lemma}
\begin{proof}
Write $R=\colim R_\lambda$ over its finitely generated $E$-subalgebras. For $X_\lambda=\Spec R_\lambda$, geometric \'etale finiteness and finiteness of local Galois cohomology, through Hochschild--Serre, show that every
$H^q_{\mathrm{\acute et}}(X_\lambda,\mu_{p^\nu}^{\otimes i})$ is finite. See \cite[Theorem~19.1(b)]{MilneEtale} and \cite[Chapter~I, Corollary~2.3]{MilneDuality}; the relevant field has characteristic zero, so this includes $p$-power coefficients and $p=2$. Jannsen's exact sequence \cite[equation~(0.2)]{Jannsen} consequently identifies continuous cohomology with the ordinary inverse limit at each such finite-type stage. The compatible finite-coefficient Chern classes therefore have unique continuous lifts there, and their product identities hold there as well.

Now use $K_3(R)=\colim K_3(R_\lambda)$ and pull those continuous classes to $\Spec R$. An element, an equality of elements, and any finite collection of projective modules descend to one common stage. This proves independence of the stage, naturality, and the product identity. It is the finite-type descent construction of \cite[Section~5.1]{LevineIndK3}, applied here to affine $E$-algebras. In particular it applies to the completed generic rings and the affinoid algebras in this section. We do not identify $H^3_{\mathrm{cont}}(\Spec R,\Z_p(3))$ with $\varprojlim H^3_{\mathrm{\acute et}}(\Spec R,\mu_{p^\nu}^{\otimes3})$, and do not interchange these cohomological limits with the filtered limit of rings.
\end{proof}

\begin{lemma}\label{d3:lem:actuallocal}
For a finite extension $E/\Q_p$, the images of actual elements of $K_3(\OO_E)$ under the rational \'etale Chern map
\[
c_{2,1}:K_3(\OO_E)\longrightarrow H^1(E,\Q_p(2))
\]
span the target over $\Q_p$. If $E/F$ is quadratic, actual integral anti-invariant elements span its anti-invariant subspace, whose $\Q_p$-dimension is $[F:\Q_p]$.
\end{lemma}
\begin{proof}
Levine's finite-coefficient theorem \cite[Theorem~4.12]{LevineIndK3}, also stated as the exact sequence in \cite[Section~5]{WeibelChern}, gives
\[
K_3^M(E)\longrightarrow K_3(E;\Z/p^\nu)
\xrightarrow{c_{2,1}}H^1(E,\mu_{p^\nu}^{\otimes2})\longrightarrow0,
\]
with the induced map on the indecomposable quotient an isomorphism, including $p=2$. Milnor $K_3$ of a local field is divisible \cite[Chapter~III, Exercise~7.4]{Weibel}; hence $K_3^M(E)/p^\nu=0$. The first image is therefore zero, and the displayed Chern map is a natural isomorphism.

The coefficient sequence for $G=K_3(E)$ is
\cite[IV, Section~2]{Weibel}
\[
0\longrightarrow G/p^\nu G\longrightarrow K_3(E;\Z/p^\nu)
\longrightarrow K_2(E)[p^\nu]\longrightarrow0.
\]
The transition on the right is multiplication by $p$. By \cite[III, Theorem~6.2.4]{Weibel}, with the characteristic-zero torsion assertion attributed there to \cite{MerkurjevTorsion}, $K_2(E)$ is the direct sum of a uniquely divisible group and a finite cyclic group. Hence that right-hand inverse limit is zero. Each $G/p^\nu G$ is finite, since it injects into the finite middle group, and this inverse system has vanishing $\varprojlim^1$. The finite groups $H^0(E,\mu_{p^\nu}^{\otimes2})$ also have vanishing $\varprojlim^1$, so the inverse limit of the $H^1$ groups is the continuous group $H^1(E,\Z_p(2))$. Taking inverse limits therefore gives the natural identification
\[
\varprojlim_\nu G/p^\nu G\cong H^1(E,\Z_p(2)).
\]
The actual group $G$ is dense in its completion, so its rational Chern image spans $H^1(E,\Q_p(2))$. There is no Tate-module contribution from $K_2$ hidden in this argument.

Localization lifts every element of $K_3(E)$ to $K_3(\OO_E)$, since $K_2(k_E)=0$. The natural Bloch--Kato isomorphism
\[
\exp_E:E\xrightarrow{\sim}H^1(E,\Q_p(2))
\]
identifies the anti-invariant subspace for a quadratic $E/F$ with the trace-zero subspace of $E$, of dimension $[F:\Q_p]$. Apply $1-\sigma$ to actual lifts whose images span that space. This gives integral classes $\beta$ with $\sigma\beta=-\beta$, without dividing by two. Naturality of the Chern maps and the exponential ensures all these assertions are equivariant. For the exponential and its normalization, see \cite[Section~1.3]{HuberKings}.
\end{proof}

We also need the precise line-bundle product. If $L$ is a line bundle and $\beta\in K_3$, then
\begin{equation}\label{d3:eq:etaleproduct}
c_{3,3}\bigl(([L]-1)\beta\bigr)
=-2c_1(L)\smile c_{2,1}(\beta).
\end{equation}
This is \cite[Theorem~3.11(i)]{WeibelChern}. The $K$-degree is three, so the exceptional degree-two correction in that paper does not apply, even at $p=2$. Lemma~\ref{d3:lem:continuouschern} puts this formula in continuous cohomology before passing to $\Q_p$. The nonzero factor $-2$ is harmless for rational detection.

\subsection{Complete quadratic examples over every $p$-adic base}

\begin{theorem}\label{d3:thm:completed}
Let $V=\OO_F$ for a finite extension $F/\Q_p$, and let $f(z)\in V[z]$ be an Eisenstein quadratic polynomial. Then
\[
\widehat A_f=V[[x,y,z]]/(xy+f(z))
\]
is a regular local ring of dimension three, and
\[
\dim_{\Q}\ker\bigl(K_3(\widehat A_f)_{\Q}
\longrightarrow K_3(\Frac\widehat A_f)_{\Q}\bigr)\geq[F:\Q_p].
\]
Taking $V=\Z_p$, this applies to $f=z^2+p$ for every prime $p$, including $p=2$.
\end{theorem}
\begin{proof}
Write $E=F(\theta)$ for a root of $f$, $W=V[\theta]=\OO_E$, and $\sigma$ for the quadratic involution. Eisenstein's condition proves regularity exactly as before, and the divisor $x=0$ has regular coordinate ring $W[[y]]$. Choose actual anti-invariant classes $\beta_j\in K_3(W)$ whose \'etale Chern classes are $\Q_p$-independent; Lemma~\ref{d3:lem:actuallocal} supplies $[F:\Q_p]$ such classes. Push them forward from $x=0$ to obtain $c_j\in K_3(\widehat A_f)$. Each generic image is zero by localization.

After the finite free base change from $V$ to $W$, put $w=z-\theta$ and $d=\theta-\sigma\theta$. Since both roots are topologically nilpotent and distinct, this is a valid formal change of variables and $0\ne d\in\mathfrak m_W$. The equation becomes
\[
xy+w(w+d)=0.
\]
After inverting $d$, the two components of $x=0$ are disjoint. With $I=(x,w)$ and $\eta_+=[\OO_{(x,w)}]=1-[I]$, flat base change gives
\[
c_j|_{T_d}=\beta_j\eta_++\sigma(\beta_j)\eta_-
=2\beta_j\eta_+,
\qquad \eta_++\eta_-=0.
\]
Applying \eqref{d3:eq:etaleproduct} gives
\[
c_{3,3}(c_j|_{T_d})=4c_1(I)\smile c_{2,1}(\beta_j).
\]
These classes are $\Q_p$-independent by Lemma~\ref{d3:lem:completedcup}. Thus the integral classes $c_j$ are $\Q$-independent after rationalization, proving the assertion.
\end{proof}

\begin{remark}
The simplest residue-two complete example in this theorem is $\Z_2[[x,y,z]]/(2+xy+z^2)$. It is different from the noncomplete residue-two example of Section~\ref{d3:sec:two}, which uses the imaginary quadratic number field $\Q(i)$ to obtain an archimedean source. The completed theorem uses local \'etale regulator classes and requires no imaginary number field.
\end{remark}

\subsection{The explicit class also survives completion at three}\label{d3:sec:explicitcomplete}

\begin{proposition}\label{d3:prop:explicitcomplete}
For
\[
\widehat A=\Z_3[[x,y,z]]/(3+xy+z^2),
\]
the explicit class $c_{\mathrm{exp}}$ of \eqref{d3:eq:explicitfront} has infinite order and zero generic image. Consequently the same explicit class has infinite order in the henselization of the original ring $A$ as well.
\end{proposition}
\begin{proof}
Let $E_0=\Q(u)$, $u^2=-3$, and retain the supported lift $\widetilde c\in K_3(Q)$ constructed in Section~\ref{d3:sec:explicit}. Since $D_+$ and $D_-$ are affine lines over $E_0$, homotopy invariance and their pushforwards give
\[
\widetilde c=\gamma_+\eta_++\gamma_-\eta_-
=(\gamma_+-\gamma_-)\eta_+=\gamma\eta_+,
\qquad \gamma\in K_3(E_0).
\]
The positive complex period in \eqref{d3:eq:positiveperiod} implies $\gamma\ne0$ in $K_3(E_0)_{\Q}$. Borel's theorem \cite[Proposition~12.2]{Borel} gives $\dim_{\Q}K_3(E_0)_{\Q}=1$.

Put $\zeta=(1+u)/2$, so $1-\zeta=\zeta^{-1}$. The special-unit Bloch cycle $[\zeta]_2$ gives a rational class $\beta_{\mathrm{cyc}}\in K_3(E_0)_{\Q}$, compatible with its integral model at three and with the embedding $E_0\hookrightarrow E=\Q_3(u)$. Besser--de Jeu's Theorem~1.10(2), in its unconditional weight-two case, computes its local syntomic regulator as
\[
\pm L_{\mathrm{mod},2}(\zeta)=\pm\operatorname{Li}_2(\zeta),
\]
because $\log\zeta=\log(1-\zeta)=0$ \cite{BesserDeJeu}. We verify directly that this value is nonzero.

Normalize $v_3(3)=1$ and put
\[
 \epsilon=1+\zeta=\frac{3+u}{2}.
\]
Then $v_3(\epsilon)=1/2$ and $\epsilon^3=3u$.
Coleman's reflection and inversion identities, and the duplication
formula, give
\[
 \operatorname{Li}_2(\epsilon)
 =-\operatorname{Li}_2(-\zeta),\qquad
 \operatorname{Li}_2(\zeta)+\operatorname{Li}_2(-\zeta)
 =\frac12\operatorname{Li}_2(\zeta^2)
 =-\frac12\operatorname{Li}_2(-\zeta).
\]
Here $1-\epsilon=-\zeta$, $( -\zeta)^{-1}=\zeta^2$,
and all logarithms of roots of unity vanish. See
\cite[equation~(1.2) and (2.4)]{BesserDeJeu}
and \cite[p.~377]{BesserDeJeuCurves}.
Consequently
\[
 \frac{\operatorname{Li}_2(\zeta)}u
 =\frac{3}{2u}\operatorname{Li}_2(\epsilon)
 =\frac{3}{2u}\sum_{n\geq1}\frac{\epsilon^n}{n^2}.
\]
The term $n=3$ is exactly $1/2$. Every other term has valuation
\[
 \frac{n+1}{2}-2v_3(n)\geq1.
\]
For $v_3(n)=0$ this is immediate; for $v_3(n)=1$ and $n\ne3$
use $n\geq6$; for $r=v_3(n)\geq2$ use
$(3^r+1)/2-2r\geq1$. The valuations tend to infinity.
Galois equivariance and inversion show that the quotient
$\operatorname{Li}_2(\zeta)/u$ lies in $\Q_3$.
Intersecting the resulting congruence modulo $3\OO_E$ with
$\Q_3$ therefore gives
\begin{equation}\label{d3:eq:localdilog}
\boxed{\quad
\frac{\operatorname{Li}_2(\zeta)}u\in\frac12+3\Z_3,
\qquad v_3(\operatorname{Li}_2(\zeta))=\frac12.
\quad}
\end{equation}

The syntomic-to-\'etale map for $\Spec\OO_E$ is compatible with Chern classes, and under the identification
\[ H^1_{\mathrm{syn}}(\OO_E,2)\cong E \]
it is the Bloch--Kato exponential. Precise references are \cite[Example~2.2.4 and Propositions~2.2.7, 2.2.9, 2.3.4]{HuberKings}; see also the proof of \cite[Corollary~5.19]{Tamme}. These statements allow ramified $E/\Q_3$. Since the exponential for $\Q_3(2)$ is an isomorphism, \eqref{d3:eq:localdilog} gives
\[
c_{2,1}(\beta_{\mathrm{cyc}}|_E)\ne0.
\]
In particular $\beta_{\mathrm{cyc}}$ is nonzero globally. The one-dimensionality of $K_3(E_0)_{\Q}$ now gives $\gamma=q\beta_{\mathrm{cyc}}$ in $K_3(E_0)_{\Q}$ for some $q\in\Q^\times$. No exact comparison scalar between a Dennis--Stein presentation and a Bloch lift has been assumed or is needed.

On the completed split generic ring $T$ the image of $c_{\mathrm{exp}}$ equals the image of $\widetilde c$: indeed $1-z^2$ is a unit in $\widehat A$, and the equality $\widetilde c|_U=2\alpha$ from Section~\ref{d3:sec:explicit} pulls back there. Thus
\[
c_{\mathrm{exp}}|_T=q\beta_{\mathrm{cyc}}\eta_+\qquad\text{in }K_3(T)_{\Q}.
\]
Its \'etale Chern class is nonzero by \eqref{d3:eq:etaleproduct} and Lemma~\ref{d3:lem:completedcup}, with $w=z-u$, $d=2u$. Therefore $c_{\mathrm{exp}}$ has infinite order in $K_3(\widehat A)$. Its generic image is zero by the integral identity \eqref{d3:eq:explicitgeneric}, which holds after completion. The map from the original local ring to its completion factors through its henselization, so infinite order after completion also proves infinite order there.
\end{proof}

\subsection{The explicit class also survives completion at two}
\label{d3:sec:explicittwocomplete}

\begin{lemma}\label{d3:lem:localdilogtwo}
For the Coleman dilogarithm on $E=\Q_2(i)$, $i^2=-1$, normalized
by its power series at zero and with $\log2=0$, one has
\begin{equation}\label{d3:eq:localdilogtwo}
 \frac{\operatorname{Li}_2(i)}{i}\in\frac12+\Z_2,
 \qquad v_2(\operatorname{Li}_2(i))=-1,
 \qquad v_2(2)=1.
\end{equation}
\end{lemma}
\begin{proof}
The modified dilogarithm
$D(z)=\operatorname{Li}_2(z)+\frac12\log z\log(1-z)$
satisfies $D(1-z)=-D(z)$ and $D(z^{-1})=-D(z)$;
see \cite[equation~(1.8), p.~307, and p.~377]{BesserDeJeuCurves}.
Put $q=1-i$, so $v_2(q)=1/2$.
Since $\log i=0$, reflection gives the convergent expansion
\[
 \operatorname{Li}_2(i)
 =-\operatorname{Li}_2(q)
 =-\sum_{n\geq1}\frac{q^n}{n^2}.
\]
Now $q^2=-2i$, $q^4=-4$, and $q^8=16$. Thus the fourth and eighth
terms cancel exactly, while the second term is $-i/2$.
For $n\notin\{2,4,8\}$,
\[
 v_2(q^n/n^2)=n/2-2v_2(n)\geq0.
\]
Indeed, writing $n=2^rk$ with $k$ odd proves this immediately
for $r\leq3$; for $r\geq4$ use
$2^{r-1}k-2r\geq2^{r-1}-2r\geq0$.
These valuations tend to infinity. Hence
$\operatorname{Li}_2(i)-i/2\in\OO_E$, proving the asserted valuation.
Galois equivariance and inversion give
$\sigma(\operatorname{Li}_2(i))=-\operatorname{Li}_2(i)$
for $\sigma(i)=-i$; see
\cite[Remark~2.3 and equation~(1.2)]{BesserDeJeu}.
Thus $\operatorname{Li}_2(i)/i\in\Q_2$, and intersecting the integral
congruence with $\Q_2$ proves \eqref{d3:eq:localdilogtwo}.
\end{proof}

\begin{proposition}\label{d3:prop:explicittwocomplete}
The explicit triple $c_2$ of \eqref{d3:eq:explicittwoprime} has
infinite order and zero generic image in
\[
 \widehat A_2
 =\Z_2[[x,y,w]]/(xy+w^2+2w+2),\qquad t=1+w.
\]
The same holds in the henselization of its original local ring.
\end{proposition}
\begin{proof}
Let $E_0=\Q(i)$ and let $\widetilde c_2$ be the supported lift
of \eqref{d3:eq:allprimeslift}. Homotopy invariance on the two
affine-line divisors and $\eta_-=-\eta_+$ give
\[
 \widetilde c_2=\gamma\eta_+,\qquad
 \gamma\in K_3(E_0)_{\Q},\qquad
 \eta_+=1-[(x,t-i)].
\]
Its nonzero period \eqref{d3:eq:twocyclotomicperiod} implies
$\gamma\ne0$.

The rational Bloch cycle $[i]_2$ defines
$\beta_i\in K_3(E_0)_{\Q}$. Interpret it over
$\Z[i]_{(1-i)}$ by the rational localization isomorphism.
Theorem~1.12 of \cite{BesserDeJeu} applies to every root of unity,
including $p$-power order, and gives its local syntomic regulator
as $\pm\operatorname{Li}_2(i)$ in $E=\Q_2(i)$.
The theorem has no odd-prime or $p>$weight restriction.
The Chern-compatible syntomic-to-\'etale map at the point
$\Spec\Z_2[i]$ is the Bloch--Kato exponential, also at $p=2$;
see \cite[Example~2.2.4 and Propositions~2.2.7, 2.3.4]{HuberKings}
and \cite[Corollary~5.19 and its proof]{Tamme}.
This exponential is an isomorphism in twist two, so
Lemma~\ref{d3:lem:localdilogtwo} proves
$c_{2,1}(\beta_i|_E)\ne0$.
Borel's equality $\dim_{\Q}K_3(\Q(i))_{\Q}=1$
\cite[Proposition~12.2]{Borel} now gives
$\gamma=a\beta_i$ for some $a\in\Q^\times$.

After finite free coefficient extension to $W=\Z_2[i]$, put
\[
 s=t-i=w-(i-1),\qquad d=2i.
\]
Because $i-1$ is topologically nilpotent, this is a valid formal
coordinate change, and the split completed ring is
\[
 \widehat B=W[[x,y,s]]/(xy+s(s+d)).
\]
Let $T=\Spec\widehat B[1/2]$ and $I=(x,s)$ on $T$.
The open used for the supported lift is $D(tg)$,
where $g=1+t+t^2$. Both $t=1+w$ and
$g=3+3w+w^2$ are units after completion. Thus
\eqref{d3:eq:allprimeslift} pulls back to
$c_2|_T=a\beta_i(1-[I])$.
Equation~\eqref{d3:eq:etaleproduct} gives
\[
 c_{3,3}(c_2|_T)
 =2a\,c_1(I)\smile c_{2,1}(\beta_i|_E)\ne0
\]
by Lemma~\ref{d3:lem:completedcup}, with parameter $d=2i$.
Its scaled model is $XY+Z(Z+1)=0$, obtained from
$x=dX$, $y=dY$, $s=dZ$; hence that lemma applies in residue
characteristic two.
This proves infinite order in $K_3(\widehat A_2)$.
The integral generic-zero identity for $c_2$ remains valid after
completion. Finally the completion map factors through the
henselization, giving the same conclusion there.
\end{proof}


\subsection{Cyclotomic local values and a common normalization}
\label{d3:sec:cyclotomicnormalization}


\begin{lemma}\label{d3:lem:cyclotomicdilognonzero}
Let $p$ be odd, $r\geq1$, and let $\rho$ be a primitive $p^r$-th
root of unity in $E=\Q_p(\rho)$. For the Coleman dilogarithm
normalized by its series at zero and $\log p=0$, one has
\[
 v_p(\operatorname{Li}_2(\rho))
 =\frac p{p-1}-2r,\qquad
 \operatorname{Li}_2(-\rho)\ne0.
\]
\end{lemma}
\begin{proof}
Put $e=p^{r-1}(p-1)$ and $q=1-\rho$, so $v_p(q)=1/e$.
Coleman reflection for
$D(z)=\operatorname{Li}_2(z)+\frac12\log z\log(1-z)$
and $\log\rho=0$ give
\[
 \operatorname{Li}_2(\rho)
 =-\sum_{n\geq1}\frac{q^n}{n^2};
\]
see \cite[p.~377]{BesserDeJeuCurves}.
Writing $n=p^sk$ with $p\nmid k$, its term valuation is
$p^sk/e-2s$. For $k=1$, the consecutive differences are
$p^{s-r+1}-2$, negative for $s<r$ and positive for $s\geq r$.
Thus the unique least term is $n=p^r$, with the asserted
valuation; every other term has valuation at least one more,
and the valuations tend to infinity.

Let $\sigma_2(\rho)=\rho^2$. Galois equivariance and duplication
\cite[Remark~2.3 and equation~(2.4)]{BesserDeJeu} give
\[
 \operatorname{Li}_2(-\rho)
 =\tfrac12\sigma_2(\operatorname{Li}_2(\rho))
       -\operatorname{Li}_2(\rho).
\]
If this vanished, iteration through the finite order $d$ of
$\sigma_2$ would give
$(2^d-1)\operatorname{Li}_2(\rho)=0$, a contradiction.
\end{proof}

\begin{lemma}\label{d3:lem:cyclotomicnormalization}
Fix the Besser--de Jeu weight-two symbol map with one natural
sign convention. For a root of unity $a$ of order $N$ in a
number field $F$, with $a,1-a\in\OO_F^\times$, let
$\beta_{\mathrm{BDJ}}(a)$ denote its rational symbol class and set
\[
 \beta_{\mathrm{BT}}(a)=\frac1N\operatorname{bl}(N[a])
       \in K_3(F)_{\Q}
\]
using \cite[Theorem~4.9]{BunkeTamme}.
There is one $s\in\Q^\times$, independent of $F$ and $a$, such that
\[
 \beta_{\mathrm{BDJ}}(a)=s\,\beta_{\mathrm{BT}}(a).
\]
\end{lemma}
\begin{proof}
The multiple $N[a]$ has zero integral wedge boundary because
$a^N=1$, including with the antisymmetric-tensor convention.
The BT class is unique rationally by Borel injectivity
\cite[Section~6]{BunkeTamme} and has
regulator $(-iD_{\mathrm{BW}}(\sigma a))_\sigma$.

The field map in \cite[equation~(1.1), Theorem~1.10(1)]{BesserDeJeu}
is de Jeu's map. Its complex regulator is computed in
\cite[Proposition~4.1 and Remark~5.2]{DeJeuWedge}, also restated in \cite[Theorem~2.3]{DeJeuCurves}, simultaneously
at all embeddings, by $P_2=D_{\mathrm{BW}}$.
Transporting its degree-dependent target identifications to the
BT target therefore gives one nonzero real constant $C$ with
\[
 r_{\mathrm{BT}}(\beta_{\mathrm{BDJ}}(a))_\sigma
       =C\,iD_{\mathrm{BW}}(\sigma a).
\]
The sign choice is made once in the fixed weight and degree,
as in \cite[Remark~1.7]{BesserDeJeu}.

Calibrate at $\zeta=(1+\sqrt{-3})/2$ in $\Q(\sqrt{-3})$.
Both classes have nonzero regulator and lie in its rank-one
rational $K_3$ \cite[Proposition~12.2]{Borel}, so
$\beta_{\mathrm{BDJ}}(\zeta)=s\beta_{\mathrm{BT}}(\zeta)$
for some $s\in\Q^\times$.
Their regulator values imply $C=-s$.
For every $F,a$, the two classes in the statement now have
the same regulator at all embeddings. Borel injectivity proves
their equality.
\end{proof}


\subsection{The explicit odd-prime triples survive completion}
\label{d3:sec:explicitoddcomplete}

We identify all supported coefficients before applying the local
detector. This avoids using a single complex embedding to identify
a class in a higher-rank number-field $K$-group.

\begin{lemma}\label{d3:lem:explicitoddcoefficients}
Let $p$ be odd, let $E_0=\Q(\mu_{2p})$, and enumerate its primitive
$2p$-th roots as $a_1,\ldots,a_{p-1}$. On
\[
 Q=\Spec E_0[x,y,a]/(xy+\Phi_{2p}(a))
\]
put $D_i=(x,a-a_i)$, $J_i=\OO_Q(-D_i)$ and
$\eta_i=[\OO_{D_i}]=1-[J_i]$.
Normalize the rational Bloch classes $\beta_{\rm BT}(a_i)$ by
\[
 r_2(\beta_{\rm BT}(a_i))_\sigma
   =-iD_{\rm BW}(\sigma(a_i))
\]
in the Bunke--Tamme cone coordinate, at every complex embedding
$\sigma$ of $E_0$.
Then every supported lift $\widetilde c_p$ in
\eqref{d3:eq:allprimeslift} satisfies
\begin{equation}\label{d3:eq:explicitoddcoefficient}
 \widetilde c_p=-2p\sum_{i=1}^{p-1}\beta_{\rm BT}(a_i)\eta_i
 \qquad\text{in }K_3(Q)_{\Q}.
\end{equation}
\end{lemma}

\begin{proof}
The stated normalization defines actual rational classes:
$a_i,1-a_i$ are global units, because $\Phi_{2p}(1)=1$,
and $2p[a_i]$ lies in the kernel of
$[a]\mapsto a\wedge(1-a)$.
One may therefore take
\[
 \beta_{\rm BT}(a_i)=\frac1{2p}\operatorname{bl}(2p[a_i])|_{E_0}
\]
using \cite[Theorem~4.9 and identification~(4.8)]{BunkeTamme}.
Borel injectivity \cite[Section~6]{BunkeTamme} makes this rational class independent of the
integral choice and equivariant under field embeddings.

Homotopy invariance on the affine-line components of $D_x$ and
the projection formula give
$\widetilde c_p=\sum_i\gamma_i\eta_i$, with
$\gamma_i\in K_3(E_0)$.
Fix distinct $i,j$ and an embedding $\sigma:E_0\hookrightarrow\C$.
The chord from $\sigma(a_i)$ to $\sigma(a_j)$ avoids zero:
primitive $2p$-th roots cannot be antipodal when $p$ is odd.
Its interior lies inside the unit circle and meets no other root,
and the whole chord has $\Re a<1$. Thus the principal
$\operatorname{Log}(1-a)$ and $\operatorname{Li}_2(a)$ are defined
on one neighborhood of it.

Orient the associated matching sphere from the $i$th root to
the $j$th root as in \eqref{d3:eq:cyclotomicendpoint}.
That formula gives the real period
\[
 2\pi(2p)\bigl(D_{\rm BW}(\sigma(a_j))
                  -D_{\rm BW}(\sigma(a_i))\bigr).
\]
The sphere meets $D_i$ with oriented intersection $+1$,
$D_j$ with intersection $-1$, and misses the other divisors.
Indeed, near its starting pole $\tau+1$ is a positive smooth
multiple of $|x|^2$, whereas near its ending pole $1-\tau$ is.
Consequently the periods of $\operatorname{ch}_1(\eta_i)$
and $\operatorname{ch}_1(\eta_j)$ are $+2\pi i$ and $-2\pi i$.

Write $r_2(\gamma_k)_\sigma=iq_k(\sigma)$ with $q_k(\sigma)$ real.
Multiplicativity therefore gives
\[
 q_i(\sigma)-q_j(\sigma)
 =2p\bigl(D_{\rm BW}(\sigma(a_i))
                  -D_{\rm BW}(\sigma(a_j))\bigr).
\]
Since this holds at every complex embedding, Borel injectivity
implies
\[
 \gamma_i-\gamma_j
 =-2p\bigl(\beta_{\rm BT}(a_i)-\beta_{\rm BT}(a_j)\bigr)
 \quad\text{in }K_3(E_0)_{\Q}.
\]
Thus $\gamma_i+2p\beta_{\rm BT}(a_i)$ is one common coefficient.
Its contribution is zero because
$\sum_i\eta_i=[\OO_{V(x)}]=0$, proving
\eqref{d3:eq:explicitoddcoefficient}.
\end{proof}

\begin{proposition}\label{d3:prop:explicitoddcomplete}
For every odd prime $p$, the explicit triple $c_p$ in
\eqref{d3:eq:explicitoddprime} has infinite order and zero generic
image in
\[
 \widehat A_p=\Z_p[[x,y,z]]/(xy+\Phi_{2p}(z-1)).
\]
The same assertions hold in the henselization of its original
arithmetic local ring.
\end{proposition}

\begin{proof}
Put $E=\Q_p(\mu_{2p})$ and $W=\OO_E$.
After the finite free coefficient extension to $W$, the completed
ring has the split form
\[
 \widehat B=
 W[[x,y,z]]\Big/\left(xy+\prod_i(z-(1+a_i))\right).
\]
All $1+a_i$ lie in the maximal ideal of $W$.
Let $T=\Spec\widehat B[1/p]$ and still write
$J_i=(x,z-(1+a_i))$ on $T$.
The finite open $U=D(a(1-a))$, with $a=z-1$, contains this generic
scheme, since $a$ and $1-a$ are units in $\widehat B$.
Thus \eqref{d3:eq:explicitoddcoefficient} pulls back to $c_p|_T$.

By Lemma~\ref{d3:lem:cyclotomicnormalization}, the classes
$\beta_{\rm BT}(a_i)$ are one common nonzero rational multiple of
the de Jeu classes represented by $[a_i]_2$.
Theorem~1.12 of \cite{BesserDeJeu}, followed by the point
syntomic-to-\'etale comparison and the Bloch--Kato isomorphism,
therefore gives
\[
 c_{2,1}(\beta_{\rm BT}(a_i)|_E)
   =\kappa\,\exp_E(\operatorname{Li}_2(a_i))
\]
with one $\kappa\in\Q_p^\times$, independent of $i$.
Lemma~\ref{d3:lem:cyclotomicdilognonzero} proves
$\operatorname{Li}_2(a_i)\ne0$, and inversion gives
$\operatorname{Li}_2(a_i^{-1})=-\operatorname{Li}_2(a_i)$.
The displayed coefficient vector is consequently non-diagonal.

The line-bundle formula \eqref{d3:eq:etaleproduct} gives
\[
 c_{3,3}(c_p|_T)
 =-4p\sum_i c_1(J_i)\smile
                c_{2,1}(\beta_{\rm BT}(a_i)|_E)\ne0
\]
by Lemma~\ref{d3:lem:completedrootcup}.
This proves infinite order after completion.
The integral vanishing identity \eqref{d3:eq:allprimesDS} remains
valid after every base change, so its generic image is zero.
Finally, the completion map factors through the henselization,
which proves both assertions there as well.
\end{proof}

\section{Completed Eisenstein and wild cyclotomic families}\label{d3:sec:completedfamilies}
Pairwise ruling maps detect all root coefficients modulo the diagonal.
We apply this detector first to arbitrary Eisenstein polynomials and
then to explicit cyclotomic coefficients with independent wild conjugates.

\subsection{Eisenstein polynomials and unbounded rank}
\label{d3:sec:completedarbitraryrank}\label{d3:sec:completedtamerank}

Lemma~\ref{d3:lem:completedrootcup} applies to every Eisenstein polynomial, including wild degrees. We combine it with actual local coefficients to obtain the following rank bound.


\begin{theorem}
\label{d3:thm:completedeisensteinrank}\label{d3:thm:completedtamerank}
Let $F/\Q_p$ be finite, let $V=\OO_F$, and let $f(z)\in V[z]$ be an Eisenstein polynomial of degree $N\geq2$. Then
\[
A_f=V[[x,y,z]]/(xy+f(z))
\]
is a complete regular local ring of dimension three with residue field $k_F$, and
\[
\dim_{\Q}\ker\left(K_3(A_f)_{\Q}
\longrightarrow K_3(\Frac A_f)_{\Q}\right)
\geq [F:\Q_p](N-1).
\]
The lower bound is witnessed by actual integral classes whose generic images vanish integrally.
\end{theorem}

\begin{proof}
In the regular local ring $V[[x,y,z]]$, the equation $xy+f(z)$ has a nonzero linear term in a uniformizer of $V$. Indeed the constant coefficient of $f$ has valuation one, while all other terms lie in the square of the maximal ideal. Its quotient is therefore regular of dimension three, with maximal ideal $(x,y,z)$ and residue field $k_F$.

Choose a root $\theta$ of $f$ and put
\[
K=F(\theta),\qquad W=V[\theta]=\OO_K.
\]
Eisenstein's criterion gives $[K:F]=N$, and the divisor $x=0$ has regular coordinate ring
\[
A_f/(x)\cong W[[y]].
\]
For an actual $\beta\in K_3(W)$ define
\[
c(\beta)=i_*(\beta|_{W[[y]]})\in K_3(A_f).
\]
Quillen localization \cite{Quillen} makes its image in $K_3(A_f[1/x])$, hence in $K_3(\Frac A_f)$, zero.

Let $E/F$ be a splitting field of $f$, let $\sigma_1,\ldots,\sigma_N:K\hookrightarrow E$ be its $F$-embeddings, and write $a_i=\sigma_i(\theta)$. Characteristic zero makes the roots distinct; Eisenstein's condition puts every $a_i$ in $\mathfrak m_{\OO_E}$. Finite free coefficient extension gives
\[
A_f\otimes_V\OO_E
\cong
\widehat B=\OO_E[[x,y,z]]/
\left(xy+\prod_i(z-a_i)\right).
\]
Set $T=\Spec\widehat B[1/p]$. This scheme is regular: it is the finite \'etale base change of $\Spec A_f[1/p]$ along the separable field extension $E/F$.

After this flat base change and inversion of $p$, the divisor $W[[y]]$ splits into the disjoint regular Cartier divisors
\[
D_i=V(x,z-a_i)\subset T,\qquad
D_i\cong\Spec\OO_E[[y]][1/p].
\]
Let $J_i=(x,z-a_i)$ and $\eta_i=[\OO_{D_i}]=1-[J_i]$. Flat base change and the projection formula give
\begin{equation}\label{d3:eq:tamegysin}
c(\beta)|_T=\sum_{i=1}^N\sigma_i(\beta)\eta_i.
\end{equation}
Here $\sigma_i(\beta)$ denotes restriction to $K_3(K)$ followed by extension along $\sigma_i$, and then the coefficient pullback from $K_3(E)$ to $K_3(T)$. The displayed formal divisor rings are localized power-series rings; they are not being replaced by $E[[y]]$.

Put
\[
\alpha=c_{2,1}(\beta|_K)\in H^1(K,\Q_p(2)),
\qquad a=\exp_K^{-1}(\alpha)\in K.
\]
Naturality of the Bloch--Kato exponential gives
\begin{equation}\label{d3:eq:tamecoefficientvector}
\bigl(c_{2,1}(\sigma_i(\beta))\bigr)_i
=\bigl(\exp_E(\sigma_i(a))\bigr)_i.
\end{equation}
The continuous Chern maps of Lemma~\ref{d3:lem:continuouschern} and the product formula \eqref{d3:eq:etaleproduct} turn \eqref{d3:eq:tamegysin} into
\[
c_{3,3}(c(\beta)|_T)
=2\sum_i c_{2,1}(\sigma_i(\beta))\smile c_1(J_i).
\]
The factor $2$ remains nonzero over $\Q_p$ for every $p$.

By Lemma~\ref{d3:lem:completedrootcup}, the cup detector kills exactly the diagonal coefficient vectors. By \eqref{d3:eq:tamecoefficientvector}, this means that all the $\sigma_i(a)$ are equal. This happens exactly when $a\in F$: every element of $\Gal(E/F)$ restricts to one of the embeddings $\sigma_i$, so equality of all conjugates makes $a$ fixed by $\Gal(E/F)$. Consequently the map
\[
g:K\longrightarrow E^N/E(1,\ldots,1),
\qquad a\longmapsto[(\sigma_i(a))_i]
\]
has kernel $F$. Its image has $\Q_p$-dimension
$[K:\Q_p]-[F:\Q_p]=[F:\Q_p](N-1)$.

By Lemma~\ref{d3:lem:actuallocal}, the values
$\exp_K^{-1}c_{2,1}(\beta)$ of actual $\beta\in K_3(W)$ span $K$ over $\Q_p$. Choose $[F:\Q_p](N-1)$ actual integral coefficient classes whose images under $g$ are $\Q_p$-independent. The displayed detector makes their pushforwards $\Q$-independent after rationalization: a rational relation would give a $\Q_p$-linear relation among those detected coefficient vectors. Each generic image already vanishes integrally by localization, proving the theorem.
\end{proof}

In particular, for every fixed prime $p$ and every integer $N\geq2$, with no restriction on divisibility by $p$,
\[
A_N=\Z_p[[x,y,z]]/(p+xy+z^N)
\]
has residue field $\F_p$, dimension three, and a rational generic kernel in $K_3$ of dimension at least $N-1$. Thus these dimensions are unbounded with the prime, residue field, and local dimension all fixed.


\subsection{Explicit wild-orbit families at every residue prime}
\label{d3:sec:explicitwildorbit}

We prove a local dilogarithm independence statement before applying the
geometric detector. It gives the whole minus space at the primes two
and three. For larger primes the assertion concerns one wild orbit;
no nonvanishing assertion for the other odd tame characters is needed.
Throughout, the Coleman dilogarithm has its usual series at zero and
the logarithm is chosen with $\log p=0$.

\begin{lemma}\label{d3:lem:wilddepletion}
Let $\rho_r$ be a primitive $p^r$-th root with
$\rho_r^p=\rho_{r-1}$, put $R_r=\operatorname{Li}_2(\rho_r)$,
and normalize $v_p(p)=1$.
For odd $p$ and $r\geq2$, write $e=p^{r-1}(p-1)$. Then
\[
 v_p(R_r-p^{-2}R_{r-1})=
 \begin{cases}
 -3/e,&p=3,\\
 -(p-2)/e,&p\geq5.
 \end{cases}
\]
For $p=2$ and $r\geq3$ one has
\[
 v_2(R_r-\tfrac14R_{r-1})=-2^{-(r-2)}.
\]
\end{lemma}
\begin{proof}
Put $F_p(z)=\operatorname{Li}_2(z)-p^{-2}\operatorname{Li}_2(z^p)$.
Theorem~2.7(3) and equation~(3.3) of
\cite{BesserDeJeuAlgorithm} identify it with the analytic continuation
of $\sum_{p\nmid n}z^n/n^2$ to $|1-z|>p^{-1/(p-1)}$.
These statements apply to every prime, including two.
For $j\geq0$ put
\[
 S_j(t)=\sum_{m\geq0}m^jt^m=\frac{P_j(t)}{(1-t)^{j+1}},
 \quad P_j\in\Z[t],\quad P_0=1,
 \qquad
 A_j(z)=\sum_{a=1}^{p-1}\frac{z^a}{a^{j+2}},
\]
with $m^0=1$ also at $m=0$.
Expanding $(a+pm)^{-2}$ first on $|z|<1$ gives
\begin{equation}\label{d3:eq:wilddepletedexpansion}
 F_p(z)=\sum_{j\geq0}(-1)^j(j+1)p^j A_j(z)S_j(z^p).
\end{equation}
For any rational $0<a<1$, the right side converges uniformly on
\[
 |z|\leq1,\qquad v_p(1-z^p)\leq a,
\]
since its $j$th summand has valuation at least $j-(j+1)a$.
This domain is the connected annulus
$p^{-a/p}\leq|1-z|\leq1$, and lies in the quoted domain of $F_p$.
The rigid analytic identity principle proves
\eqref{d3:eq:wilddepletedexpansion} there. For all the roots under
consideration one may take $a=2/3$.

At $p=2$ one has $A_j(z)=z$. Write
$d=v_2(1-\rho_r^2)=2^{-(r-2)}\leq1/2$.
The $j=0$ term is $\rho_r/(1-\rho_r^2)$, of valuation $-d$;
every $j\geq1$ term has valuation at least
$j-(j+1)d=-d+j(1-d)>-d$. This proves the assertion at two.

At $p=3$ the numerator of the $j=0$ term is
$\rho_r+\rho_r^2/4$, which reduces to $2$.
With $d=v_3(1-\rho_r^3)=3/e\leq1/2$, the same argument gives $-3/e$.

Suppose now that $p\geq5$. For $A=A_0$, reduction modulo $p$ gives
\[
 A(1)=A'(1)=0,\qquad A''(1)/2=-1/2\quad\text{in }\F_p.
\]
Indeed these are the usual sums of inverse squares, inverses,
and $(a-1)/(2a)$ over $a\in\F_p^\times$.
Since $q=1-\rho_r$ has valuation $1/e$ and $2/e<1$,
the unique least term of $A(1-q)$ is its quadratic term.
Thus $v_p(A(\rho_r))=2/e$ and the $j=0$ term has valuation $(2-p)/e$.
For $j\geq1$ the difference between its lower valuation bound
$j-(j+1)p/e$ and $(2-p)/e$ is
\[
 j(1-p/e)-2/e\geq1-(p+2)/e>0,
\]
because $e\geq p(p-1)>p+2$.
Again the $j=0$ term is uniquely minimal.
\end{proof}

\begin{proposition}\label{d3:prop:wildorbit}
For odd $p$, let $E_r=\Q_p(\rho_r)$ and
$\Gamma_r=\Gal(E_r/E_1)$.
The $p^{r-1}$ elements
\[
 \{\sigma\operatorname{Li}_2(-\rho_r):\sigma\in\Gamma_r\}
\]
are linearly independent over $\Q_p$.
For $p=3$ they form a basis of $E_r^-$.

At $p=2$, let $r\geq2$, $E_r=\Q_2(\rho_r)$, and
$\Gamma_r=\Gal(E_r/\Q_2(i))$.
Put
\[
 \Psi_2(\rho_r)=\tfrac13\operatorname{Li}_2(\rho_r^3)
                              -\operatorname{Li}_2(\rho_r).
\]
Its $2^{r-2}$ conjugates under $\Gamma_r$ form a $\Q_2$-basis of
$E_r^-$.
\end{proposition}
\begin{proof}
For odd $p$, write
$\Gal(E_r/\Q_p)=\Delta_r\times\Gamma_r$, with $|\Delta_r|=p-1$,
and let $\omega:\Delta_r\to\Q_p^\times$ be the Teichm\"uller
character: $\omega(\sigma_a)\equiv a\pmod p$ when
$\sigma_a(\rho_r)=\rho_r^a$.
Use its idempotent
\[
 e_\omega=\frac1{p-1}\sum_{\sigma\in\Delta_r}
                        \omega(\sigma)^{-1}\sigma.
\]
If $y\ne0$ has valuation $m/e$ with $m\equiv1\pmod{p-1}$,
then $e_\omega y$ has the same valuation as $y$.
In fact, write $y=q^m u$ with $u$ a unit. Since the extension is
totally ramified, $\sigma_a(u)/u$ reduces to $1$, while
$\sigma_a(q)/q$ reduces to $a$. Each summand of $e_\omega y/y$ reduces to $1$, as does their average.

The integer numerators $-3$ and $-(p-2)$ in
Lemma~\ref{d3:lem:wilddepletion} are congruent to $1$ modulo $p-1$.
Consequently the $\omega$ projection of every depleted value is nonzero.
The base value satisfies
$v_p(R_1)=-(p-2)/(p-1)$ by
Lemma~\ref{d3:lem:cyclotomicdilognonzero}, and its $\omega$ projection
is nonzero for the same reason.

Put $C_r=p^{2r-1}R_r$.
Distribution \cite[Proposition~2.10(1)]{BesserDeJeuAlgorithm} gives
\[
 \operatorname{Tr}_{E_r/E_{r-1}}R_r=p^{-1}R_{r-1},
 \qquad p^{-1}\operatorname{Tr}_{E_r/E_{r-1}}C_r=C_{r-1}.
\]
Hence the successive rational isotypic projections of
$e_\omega C_r$ under the cyclic group $\Gamma_r$ are
\[
 e_\omega C_1,\quad e_\omega(C_2-C_1),\quad\ldots,
 \quad e_\omega(C_r-C_{r-1}).
\]
They are all nonzero. The normal basis theorem identifies
$e_\omega E_r$ with the regular $\Q_p[\Gamma_r]$-module
\cite[Chapter~I, p.~34]{MilneDuality}.
Its irreducible constituents have minimal polynomials
$T-1,\Phi_p(T),\ldots,\Phi_{p^{r-1}}(T)$, each once;
these are irreducible over $\Q_p$.
Thus $e_\omega R_r$ is cyclic and its orbit is a basis.
Projecting any relation shows that the orbit of $R_r$ itself
is independent. Duplication gives
\[
 \operatorname{Li}_2(-\rho_r)=(\tfrac12\sigma_2-1)R_r.
\]
The displayed operator is invertible, commutes with $\Gamma_r$,
and preserves independence: $\sigma_2$ has finite order whereas
$2$ is not a root of unity.
Inversion makes every value anti-invariant. At $p=3$,
$\Delta_r=\{1,\text{inversion}\}$, so this is the whole minus space.

At $p=2$ use the base level $r=2$ and
$\Gal(E_r/\Q_2)=\{1,\text{inversion}\}\times\Gamma_r$.
The group $\Gamma_r$ is cyclic of order $2^{r-2}$, generated by
$\sigma_5$. Its minus space is its regular representation
\cite[Chapter~I, p.~34]{MilneDuality}.
Set again $C_r=2^{2r-1}R_r$; distribution gives average trace
$C_r\mapsto C_{r-1}$. The base $C_2$ is nonzero by
Lemma~\ref{d3:lem:localdilogtwo}, and every
$C_j-C_{j-1}$ for $j\geq3$ is nonzero by
Lemma~\ref{d3:lem:wilddepletion}.
The same multiplicity-one argument, now using
$T-1,\Phi_2(T),\ldots,\Phi_{2^{r-2}}(T)$ over $\Q_2$,
shows that $R_r$ is a cyclic vector in $E_r^-$.
Finally $\Psi_2(\rho_r)=(\sigma_3/3-1)R_r$, an invertible
operator commuting with $\Gamma_r$, so it too is cyclic.
\end{proof}

\begin{theorem}
\label{d3:thm:explicitwildtriples}
For odd $p$ and $r\geq1$ put
\[
 A_{p,r}=\left(\Z_{(p)}[x,y,z]/
       (xy+\Phi_{2p^r}(z-1))\right)_{(p,x,y,z)}.
\]
Write $a=z-1$, $n=2p^r$, and
\[
 H=\frac{a^n-1}{\Phi_{2p^r}(a)},\qquad
 \lambda=1-x,\qquad \mu=1+\lambda yH.
\]
For $0\leq j<p^{r-1}$, let $k_j$ be the unique odd integer with $1\leq k_j<2p^r$ and
$k_j\equiv(1+p)^j\pmod{p^r}$.
Then the explicit classes
\[
 c_j=k_j\{1-a^{k_j},\lambda,\mu\}\in K_3(A_{p,r})
\]
are linearly independent over $\Q$ after tensoring with $\Q$,
and each has zero generic image integrally.
Their images retain these properties in the henselization and
completion of $A_{p,r}$.

At $p=2$, for $r\geq2$ put
\[
 A_{2,r}=\left(\Z_{(2)}[x,y,t]/
                 (xy+\Phi_{2^r}(t))\right)_{(2,x,y,t-1)},
 \qquad n=3\cdot2^r,
\]
and set $H=(t^n-1)/\Phi_{2^r}(t)$,
$\lambda=1-x$, $\mu=1+\lambda yH$.
Let $k_j$ be the least positive representative of $5^j$ modulo
$2^r$, for $0\leq j<2^{r-2}$.
The same assertions hold for the $2^{r-2}$ classes
\[
 c_j=k_j\{1+t^{k_j}+t^{2k_j},\lambda,\mu\}.
\]
All the displayed local rings are regular of dimension three.
\end{theorem}
\begin{proof}
The polynomials $\Phi_{2p^r}(z-1)$ for odd $p$ and
$\Phi_{2^r}(1+z)$ at two are Eisenstein.
Their constant terms are respectively $p$ and $2$; modulo that
prime they are powers of $z$. This proves regularity and dimension
as in Theorem~\ref{d3:thm:eisenstein}.
Every entry of each displayed triple is a unit: at the closed
point $a=-1$ for odd $p$, and $t=1$ at two.

For either prime let $u$ denote $a$ or $t$, and put $b=yH$.
Then $1-xb=u^n$ and the Dennis--Stein identity
\eqref{d3:eq:allprimesDS} applies to
$\delta=\langle x,b\rangle$.
In the odd case, for any allowed $k$,
\[
 k\{1-a^k,a^n\}=\{1-a^k,(a^k)^n\}=0.
\]
At two, use the finite \'etale quadratic extension adjoining
$\zeta_3$. Steinberg and the transfer projection formula give
\[
 \begin{split}
 0&=\operatorname{Tr}\{1-\zeta_3t^k,(\zeta_3t^k)^n\}\\
  &=\{1+t^k+t^{2k},t^{kn}\}
   =k\{1+t^k+t^{2k},t^n\}.
 \end{split}
\]
Here $3\mid n$, and the norm of $1-\zeta_3t^k$ is exactly
$1+t^k+t^{2k}$, with no extra factor of two.
Only oddness of $k$ is needed; $k$ may be divisible by $3$.
If $g_k$ denotes the first entry of the corresponding triple, then
\[
 k[g_k]\delta=k\{g_k,\lambda,\mu\},\qquad
 k[g_k]\delta|_{D(x)}=0.
\]
This proves integral generic vanishing before and after every base change.

We specify the coefficient calculation needed for independence.
On the split finite surface
$Q=\Spec F[x,y,u]/(xy+P(u))$, with $F=\Q(\mu_{2p^r})$ for
odd $p$ and $F=\Q(\mu_{2^r})$ at two, write the roots as $u_i$,
and put $\eta_i=1-[(x,u-u_i)]$.
All roots lie in the finite open $D(ug_k)$.
Localization and d\'evissage give a supported lift of
$k[g_k]\delta$ to $K_3(Q)$.
The matching-sphere calculation \eqref{d3:eq:cyclotomicendpoint}
identifies any such lift, rationally, as
\begin{equation}\label{d3:eq:wildsupportedcoefficients}
 \widetilde c_k=-n\sum_i b_k(u_i)\eta_i,
\end{equation}
where
\[
 b_k(u)=
 \begin{cases}
 \beta_{\rm BT}(u^k),&p\text{ odd},\\
 \operatorname{Tr}^{K}_{F(\zeta_3)/F}
       \beta_{\rm BT}(\zeta_3u^k),&p=2.
 \end{cases}
\]
For completeness, choose an arc between each pair of roots whose
interior lies in the punctured open unit disc. There
$\operatorname{Log}(1-u^k)$, or the sum
$\operatorname{Log}(1-\zeta_3u^k)+
 \operatorname{Log}(1-\zeta_3^2u^k)$, has a holomorphic branch.
Its primitive for $-L\,du/u$ is respectively
\[
 \frac1k\operatorname{Li}_2(u^k),\qquad
 \frac1k\bigl(\operatorname{Li}_2(\zeta_3u^k)
                  +\operatorname{Li}_2(\zeta_3^2u^k)\bigr).
\]
The factor $k$ in $c_k$ cancels this denominator.
The endpoint periods therefore have the common factor $n$,
independent of $k$. Checking the periods at every complex embedding,
and then applying Borel injectivity as in
Lemma~\ref{d3:lem:explicitoddcoefficients}, determines all coefficient
differences. The remaining diagonal coefficient vanishes because
$\sum_i\eta_i=0$. This proves
\eqref{d3:eq:wildsupportedcoefficients}.

The common normalization of
Lemma~\ref{d3:lem:cyclotomicnormalization}, the root-of-unity formula
\cite[Theorem~1.12]{BesserDeJeu}, and the point comparison give,
with one nonzero scalar for the whole vector, the local coefficient
values
\[
 c_{2,1}(b_k(u_i))=
 \begin{cases}
 \kappa\exp_E(\operatorname{Li}_2(u_i^k)),&p\text{ odd},\\
 \kappa\exp_E(\Psi_2(u_i^k)),&p=2.
 \end{cases}
\]
At two this follows by restriction-transfer and distribution:
$\operatorname{Li}_2(\zeta_3u^k)+
\operatorname{Li}_2(\zeta_3^2u^k)=\Psi_2(u^k)$.
All these roots and their complements used in $\beta_{\rm BT}$
are global units. No root-dependent normalization is introduced.

After completion and the finite free extension to
$\OO_E$, every shifted root lies in its maximal ideal, so
Lemma~\ref{d3:lem:completedrootcup} and the line-bundle Chern formula
apply to \eqref{d3:eq:wildsupportedcoefficients}.
They kill exactly the diagonal coefficient vectors.
Every vector displayed above changes sign when $u_i$ is replaced
by $u_i^{-1}$. A diagonal linear combination must therefore be zero.
Looking at a single fixed primitive root now turns a relation
among the $c_j$ into a relation among the wild-orbit values in
Proposition~\ref{d3:prop:wildorbit}. Its independence forces every
coefficient to vanish, even over $\Q_p$ at the regulator level.
Thus the actual $K_3$ classes are $\Q$-linearly independent after
completion. The arithmetic local ring and its henselization map
to that completion, proving the corresponding assertions there.
\end{proof}

\begin{remark}
The explicit lower bounds are $p^{r-1}$ for odd $p$ and
$2^{r-2}$ at two. They equal the cyclotomic minus dimension at
$p=3$ and $p=2$. For $p\geq5$ this argument proves independence
through the Teichm\"uller component; it does not assert full minus
rank $(p-1)p^{r-1}/2$.
\end{remark}


\section{Henselization and infinite rank}\label{d3:sec:henselian}
Actual local coefficients descend to henselian arithmetic valuation rings.
Their finite families give the same rank bounds before completion and
produce infinite rank after strict henselization.

\subsection{Actual coefficients over a henselian valuation ring}
\label{d3:sec:henselianeisenstein}

The coefficient classes needed by the completed argument can already be
chosen over henselian number-field valuation rings.  The finite-quotient
assertion in Levine's theorem gives this directly.

\begin{lemma}
\label{d3:lem:henseliancoefficients}
Let $W_0$ be a discrete valuation ring whose fraction field is a number
field and whose residue field has characteristic $p$.  Put
\[
H=\Frac(W_0^h),\qquad K=\Frac(\widehat W_0).
\]
For every $\nu\geq1$, the natural map
\[
K_3(W_0^h)\longrightarrow K_3(K)/p^\nu
\]
is surjective.  In particular the rational \'etale Chern classes of
actual elements of $K_3(W_0^h)$ span $H^1(K,\Q_p(2))$ over $\Q_p$.
\end{lemma}

\begin{proof}
The henselization $W_0^h$ is a discrete valuation ring with the same
residue field and completion as $W_0$; see
\cite[Section~15.46, Lemmas~15.46.1, 15.46.3 and 15.46.11]{StacksHenselization}.
Its fraction field embeds densely in $K$, and
\[
H=\{a\in K:a\text{ is algebraic over }\Q\}.
\]
Indeed $H$ is algebraic over $\Frac(W_0)$.  Conversely a henselian
rank-one valued field is separably algebraically closed in its completion:
approximate a separable algebraic root by an element of the field closely
enough to apply Hensel's lemma, and isolate that root from the finitely
many other roots.  The Henselian root then equals the original root.
Characteristic zero makes every algebraic extension separable.

The second assertion of \cite[Theorem~4.13, p.~336]{LevineIndK3},
applied to $K$ and its field of constants $H$, gives
\[
K_3(H)^{\mathrm{ind}}/p^\nu
\xrightarrow{\sim}K_3(K)^{\mathrm{ind}}/p^\nu.
\]
Milnor $K_3(K)$ is $p$-divisible, as used in
Lemma~\ref{d3:lem:actuallocal}.  Its image in $K_3(K)$ is therefore
contained in $p^\nu K_3(K)$, so
$K_3(K)^{\mathrm{ind}}/p^\nu=K_3(K)/p^\nu$.
Consequently $K_3(H)\to K_3(K)/p^\nu$ is surjective.
Quillen localization lifts every element of $K_3(H)$ to $K_3(W_0^h)$,
because $K_2(k_{W_0})=0$.

The proof of Lemma~\ref{d3:lem:actuallocal} identifies
$\varprojlim_\nu K_3(K)/p^\nu$ with $H^1(K,\Z_p(2))$.
The surjectivity just proved makes the image of the actual group
$K_3(W_0^h)$ dense in this inverse limit.  Its rational Chern image
therefore spans $H^1(K,\Q_p(2))$.
\end{proof}

\begin{theorem}
\label{d3:thm:henselianeisensteinrank}
Let $F_0$ be a number field, let $V_0=(\OO_{F_0})_{\mathfrak p}$ for a
prime above $p$, and let $F$ be its completion.  Let $f(z)\in V_0[z]$
be Eisenstein of degree $N\geq2$, and put
\[
R_0=\left(V_0[x,y,z]/(xy+f(z))\right)_{(\mathfrak p,x,y,z)},
\qquad R=R_0^h.
\]
Then $R$ is a henselian regular local ring of dimension three with
residue field $k_F$, and
\[
\dim_{\Q}\ker\left(K_3(R)_{\Q}\longrightarrow
K_3(\Frac R)_{\Q}\right)\geq [F:\Q_p](N-1).
\]
The lower bound is witnessed by actual integral classes whose generic
images vanish integrally.
\end{theorem}

\begin{proof}
Choose a root $\theta$ of $f$ and put
\[
W_0=V_0[\theta],\qquad K=F(\theta),\qquad W=\OO_K.
\]
Eisenstein's condition makes $W_0$ a discrete valuation ring with
uniformizer $\theta$, residue field $k_F$, and completion $W$.
It also makes $R_0$ regular of dimension three, by the same linear-term
calculation as in Theorem~\ref{d3:thm:completedeisensteinrank}.
Regularity, dimension and the residue field are preserved by
henselization, and
\[
\widehat R=\OO_F[[x,y,z]]/(xy+f(z)).
\]
These assertions, including the faithful flatness of $R\to\widehat R$,
follow from \cite[Section~15.46]{StacksHenselization}.

Henselization commutes with the quotient by $x$, so the regular divisor
$D=\Spec(R/(x))$ has coordinate ring
\[
R/(x)\cong \left(W_0[y]_{(\theta,y)}\right)^h.
\]
Here we use \cite[Section~10.156, Lemma~10.156.2]{StacksHenselization}.
The local map $W_0\to R/(x)$ extends to $W_0^h\to R/(x)$.
On completion this map is compatible with
$W\to W[[y]]=\widehat R/(x)$, by the universal property of
henselization.

For $\beta\in K_3(W_0^h)$ define
\[
c^h(\beta)=i_*(\beta|_D)\in K_3(R),
\qquad i:D\hookrightarrow\Spec R.
\]
Localization makes its image in $K_3(R[1/x])$, hence in
$K_3(\Frac R)$, zero integrally.  Flat base change along
$R\to\widehat R$ sends $c^h(\beta)$ to exactly the class $c(\widehat
\beta)$ constructed in Theorem~\ref{d3:thm:completedeisensteinrank},
where $\widehat\beta$ is the image in $K_3(W)$.

By Lemma~\ref{d3:lem:henseliancoefficients}, the values
\[
a(\beta)=\exp_K^{-1}c_{2,1}(\beta|_K)
\]
of actual $\beta\in K_3(W_0^h)$ span $K$ over $\Q_p$.
The completed detector in
Theorem~\ref{d3:thm:completedeisensteinrank} has kernel precisely $F$
on this coefficient space.  Choose $[F:\Q_p](N-1)$ actual coefficient
classes whose values are independent in $K/F$ over $\Q_p$.
Their completed pushforwards have independent detected classes.
Any rational relation among their henselian pushforwards would remain
a relation after completion and would contradict that independence.
This proves the asserted lower bound.
\end{proof}

In particular, for every prime $p$ and every $N\geq2$, the henselization
\[
\left(\left(\Z_{(p)}[x,y,z]/(p+xy+z^N)\right)_{(x,y,z)}\right)^h
\]
has residue field $\F_p$, dimension three, and rational generic
$K_3$-kernel of dimension at least $N-1$.


\subsection{One complete strictly henselian ring with infinite rational rank}

\begin{lemma}\label{d3:lem:completedcupperfect}
Lemma~\ref{d3:lem:completedcup} remains valid when $E$ is a complete
discretely valued field of mixed characteristic $(0,p)$ with perfect
residue field, containing a fixed finite extension of $\Q_p$ as a valued
subfield.
\end{lemma}
\begin{proof}
We check the points in that proof where the base could matter.
The valuation ring $\OO_E$ is still a noetherian complete DVR.
The fixed rational-radius models used in \eqref{d3:eq:fixedradius}
are therefore noetherian henselian models to which
\cite[Theorem~6.2.1]{Achinger} applies, including for $p$-power coefficients.
The order of operations in \eqref{d3:eq:schemetodagger} is unchanged:
finite coefficients at one fixed radius, derived inverse limit,
inversion of $p$, and finally the dagger homotopy colimit.
The two-chart calculation by integration of overconvergent series gives
the same isomorphism
\[
 f^*:H^2_{\mathrm{dR}}(\PP^1_E)
       \xrightarrow{\sim}H^2_{\mathrm{dR}}(U^\dagger).
\]

The standing hypotheses of \cite{ColmezNiziol}, stated at the beginning
of its introduction, allow any complete mixed-characteristic DVR with
perfect residue field. Thus its Theorem~1.3(2) still gives the injective
boundary $\partial_3$ on $U^\dagger$.
For $\PP^1_E$ the boundary is an isomorphism:
$H^2_{\mathrm{HK}}(\PP^1)=F_0h$ with
$\varphi(ah)=p\sigma(a)h$, where
$F_0=\Frac W(k_E)$, while $H^3_{\mathrm{HK}}(\PP^1)=0$.
The eigenspace $\varphi=p^2$ is zero, since
$p\sigma(a)=p^2a$ for $a\ne0$ contradicts
$v_p(\sigma(a))=v_p(a)$.
This argument does not require finite order of $\sigma$.
Proper \'etale and pro-\'etale cohomology agree by
\cite[Section~7.2.4]{ColmezNiziol}.

We also remove the distinction between topological and ordinary
cohomology in the target. The unit quadric has the smooth weak formal
model
$\operatorname{Spwf}\OO_E\langle X,Y,Z\rangle^\dagger/
(XY+Z(Z+1))$, and its ruling extends to the weak formal projective
line: the ideal $(X,Z)$ is already invertible on the integral model.
The natural Hyodo--Kato comparison in
\cite[equation~(5.4) and the following strict quasi-isomorphism,
Propositions~5.5--5.6]{ColmezNiziol}, together with faithful extension
from $F_0$ to $E$, transports the de Rham degree-two isomorphism to
\[
 H^2_{\mathrm{HK}}(\PP^1_E)
       \xrightarrow{\sim}H^2_{\mathrm{HK}}(U^\dagger);
 \qquad H^3_{\mathrm{HK}}(U^\dagger)=0.
\]
Thus the same Hyodo--Kato term vanishes for $U^\dagger$, and
$\partial_{3,U}$ is an isomorphism.
Its de Rham source is classical by
\cite[Section~5.1]{ColmezNiziol}, even if $E$ has infinite dimension
over $\Q_p$. Hence its target is classical too.
The classical-part functor of \cite[Section~3.1.1]{ColmezNiziol}
sends $\widetilde H^i$ to ordinary $H^i$; applying it to this
isomorphism gives an isomorphism on ordinary cohomology.
Naturality now proves injectivity of pullback from
$H^3_{\mathrm{cont},\mathrm{\acute et}}(\PP^1_E,\Q_p(3))$ on the
ordinary cohomology used by the Chern maps.

For the last step, the splitting \eqref{d3:eq:finitePonebundle}
with $r=3$ gives the split injection
\[
 H^1(E,\Q_p(2))
 \xrightarrow{\ \smile c_1(\OO(1))\ }
 H^3_{\mathrm{cont},\mathrm{\acute et}}(\PP^1_E,\Q_p(3)).
\]
Its finite-coefficient splitting survives derived inverse limits and
inversion of $p$. No vanishing assertion about
$H^3(E,\Q_p(3))$ is required.
Its pullback is the ruling-cup map, which is therefore
injective.
\end{proof}

\begin{theorem}
\label{d3:thm:stricthenselianinfinite}
Set
\[
 V_\infty=W(\overline{\mathbf F}_3),\qquad
 A_\infty=V_\infty[[x,y,z]]/(3+xy+z^2).
\]
Then $A_\infty$ is a complete strictly henselian regular local ring of
dimension three, and
\[
 \dim_{\Q}\ker\bigl(
   K_3(A_\infty)_{\Q}
       \longrightarrow K_3(\Frac A_\infty)_{\Q}\bigr)=\infty.
\]
For every $m\geq1$ this kernel contains $m$ independent classes
represented by actual integral elements of $K_3(A_\infty)$.
\end{theorem}
\begin{proof}
The Witt ring $V_\infty$ is a complete DVR with uniformizer $3$.
The equation has nonzero linear term $3$ in the cotangent space of
$V_\infty[[x,y,z]]$, so its quotient is regular local of dimension three,
with maximal ideal $(x,y,z)$. Completeness and the algebraically closed
residue field imply strict henselianity.

For each positive integer $m$, put
\[
 V_m=W(\mathbf F_{3^m}),\quad F_m=\Frac V_m,\quad
 A_m=V_m[[x,y,z]]/(3+xy+z^2).
\]
The canonical local map $A_m\to A_\infty$ is flat by miracle flatness:
both rings are regular of dimension three and the closed fiber has
dimension zero \cite[Tag~00R4]{StacksFlatness}.

Choose $u^2=-3$ and write
\[
 E_m=F_m(u),\quad W_m=V_m[u],\qquad
 E_\infty=(\Frac V_\infty)(u),\quad W_\infty=V_\infty[u].
\]
The field $E_\infty$ is the completion of $E_m^{\mathrm{ur}}$ for every
$m$. Indeed its Witt coefficients can be approximated to every finite
precision in finite residue subfields. Completion of a henselian
discretely valued field preserves its category of finite separable
extensions \cite[Lemma~4.2.1]{SaitoCompletion}; consequently
\[
 G_{E_\infty}\simeq I_{E_m}.
\]
Restriction therefore induces an injection
\begin{equation}\label{d3:eq:inertiarestriction}
 H^1(E_m,\Q_3(2))\hookrightarrow H^1(E_\infty,\Q_3(2)).
\end{equation}
Here is a direct verification. For $G=G_{E_m}$ and $I=I_{E_m}$,
$\Q_3(2)^I=0$ because the cyclotomic character squared is nontrivial
on inertia. A cocycle trivial in $H^1(I,\Q_3(2))$ can be changed by a
coboundary to vanish on $I$. For such a cocycle $b$,
$i b(g)=b(ig)=b(g(g^{-1}ig))=b(g)$, so every $b(g)$ belongs to
$\Q_3(2)^I$ and is zero.

By Lemma~\ref{d3:lem:actuallocal}, choose actual integral
anti-invariant classes
$\beta_{m,1},\ldots,\beta_{m,m}\in K_3(W_m)$ whose Chern classes
\[
 \alpha_{m,j}=c_{2,1}(\beta_{m,j})
       \in H^1(E_m,\Q_3(2))
\]
are $\Q_3$-independent. Pull these classes to the regular divisor
$A_m/(x)\simeq W_m[[y]]$ and push forward to obtain
$c_{m,j}\in K_3(A_m)$. Localization makes them zero after inverting
$x$. Their images $c^\infty_{m,j}\in K_3(A_\infty)$ are therefore
also generically zero.

Make the finite free base change $V_\infty\to W_\infty$ and put
\[
 T_\infty=\Spec
 \bigl(W_\infty[[x,y,z]]/(xy+z^2-u^2)\bigr)[1/u],
 \qquad I_+=(x,z-u),\qquad \eta_+=1-[I_+].
\]
Flat base change and the projection formula give
\[
 c^\infty_{m,j}|_{T_\infty}
      =2\,\beta_{m,j}|_{E_\infty}\,\eta_+.
\]
All rings here are algebras over the fixed finite extension
$\Q_3(u)$, so Lemma~\ref{d3:lem:continuouschern} applies without any
new finiteness assumption over $E_\infty$.
Using \eqref{d3:eq:etaleproduct} gives
\[
 c_{3,3}(c^\infty_{m,j}|_{T_\infty})
   =4\,c_1(I_+)\smile
       \operatorname{res}_{E_\infty/E_m}(\alpha_{m,j}).
\]
The coefficient classes remain $\Q_3$-independent by
\eqref{d3:eq:inertiarestriction}. Lemma~\ref{d3:lem:completedcupperfect},
with $w=z-u$ and $d=2u$, proves independence of these Chern images.
Hence the actual integral classes $c^\infty_{m,1},\ldots,
c^\infty_{m,m}$ are $\Q$-independent after rationalization.
Since $m$ is arbitrary and $A_\infty$ is fixed, the asserted kernel
is infinite dimensional.
\end{proof}

\begin{remark}
This proof uses arbitrarily large finite families from finite
unramified coefficient fields. It neither requires that actual
regulators exhaust $E_\infty$ nor identifies $A_\infty$ with a
filtered colimit of the rings $A_m$.
\end{remark}


\subsection{Infinite rank already in an arithmetic strict henselization}
\label{d3:sec:stricthenseldescent}

\begin{lemma}
\label{d3:lem:stricthenselcompletion}
For a prime $p$, set
\[
 A_{0,p}=\left(\Z_{(p)}[x,y,z]/(p+xy+z^2)\right)_{(x,y,z)},
 \qquad S_p=A_{0,p}^{sh},
\]
using the geometric residue field $\overline{\F}_p$.
There is a canonical compatible local map
\[
 S_p\longrightarrow
 A_{\infty,p}\coloneqq W(\overline{\F}_p)[[x,y,z]]/(p+xy+z^2),
\]
and it identifies $A_{\infty,p}$ with the completion of $S_p$.
In particular $S_p$ is noetherian regular local of dimension three
with residue field $\overline{\F}_p$.
\end{lemma}
\begin{proof}
The target is a complete local ring with algebraically closed residue
field, hence strictly henselian.  The local map from $A_{0,p}$ extends
to $S_p$ by the universal property with the prescribed geometric
residue field.  The assertions about $S_p$ follow from
\cite[Section~15.46]{StacksHenselization}.

Write $\m=(x,y,z)$.  Fix $r\geq1$ and choose $s$ with $2s\geq r$.
Since $p\in\m^2$, the Artinian local ring
$R_r=A_{0,p}/\m^r$ is a $\Z/p^s$-algebra, and
\[
 A_{\infty,p}/\m^r
 \cong R_r\otimes_{\Z/p^s}W_s(\overline{\F}_p).
\]
Here $W_s$ denotes Witt vectors of length $s$.  The equality
\[
 W_s(\overline{\F}_p)=\colim_m W_s(\F_{p^m}),
\]
with $m$ ordered by divisibility, follows because a finite Witt vector
has only finitely many residue-field coordinates.  Each
$R_r\otimes_{\Z/p^s}W_s(\F_{p^m})$ is the finite \'etale local extension
of $R_r$ with residue field $\F_{p^m}$.
Finite \'etale extensions of a henselian local ring are determined by
their residue extensions \cite[Tag~04GK]{StacksHenselization}.
Thus $A_{\infty,p}/\m^r$ is the strict henselization of $R_r$.
By quotient compatibility
\cite[Section~10.156, Lemma~10.156.4]{StacksHenselization},
\[
 A_{\infty,p}/\m^r\cong S_p/\m^rS_p.
\]
These identifications respect the transition maps and the chosen
residue field.  Taking inverse limits proves
$\widehat S_p\cong A_{\infty,p}$.
\end{proof}

\begin{corollary}
\label{d3:cor:stricthenselizationinfinite}
Let
\[
 S=\left(\left(\Z_{(3)}[x,y,z]/(3+xy+z^2)\right)_{(x,y,z)}\right)^{sh}.
\]
Then
\[
 \dim_{\Q}\ker\left(K_3(S)_{\Q}
      \longrightarrow K_3(\Frac S)_{\Q}\right)=\infty.
\]
For every $m\geq1$ there are $m$ independent classes represented by
actual integral elements of $K_3(S)$ whose generic images vanish
integrally.
\end{corollary}
\begin{proof}
For each $m$, choose a monic polynomial $g_m\in\Z[T]$ whose reduction
modulo $3$ is irreducible of degree $m$, and choose a residue root in
$\F_{3^m}\subset\overline{\F}_3$.
The ring
\[
 V_{m,0}=\Z_{(3)}[T]/(g_m)
\]
is a finite \'etale local discrete valuation ring with number-field
fraction field, residue field $\F_{3^m}$, and completion
$V_m=W(\F_{3^m})$.  The chosen residue root gives compatible embeddings
of $V_{m,0}$ into $\Z_{(3)}^{sh}$ and into $V_m$.
Put
\[
 R_{m,0}=A_{0,3}\otimes_{\Z_{(3)}}V_{m,0},\qquad R_m=R_{m,0}^h.
\]
The finite \'etale algebra $R_{m,0}$ is local, because its closed fiber
is the field $\F_{3^m}$.  There are compatible local maps
\[
 R_m\longrightarrow S\longrightarrow A_{\infty,3},
 \qquad
 \widehat R_m=A_m=V_m[[x,y,z]]/(3+xy+z^2)
       \longrightarrow A_{\infty,3}.
\]
The two maps from $R_m$ to $A_{\infty,3}$ agree by the universal
property of henselization.  The second row is the usual inclusion of
Witt coefficients after completion.

Set $u^2=-3$ and $W_{m,0}=V_{m,0}[u]$.
Lemma~\ref{d3:lem:henseliancoefficients}, followed by $1-\sigma$ for
$\sigma(u)=-u$, supplies actual anti-invariant classes
\[
 \beta_{m,1}^h,\ldots,\beta_{m,m}^h\in K_3(W_{m,0}^h)
\]
whose Chern classes in $H^1(\Frac(V_m)(u),\Q_3(2))$ are
$\Q_3$-independent.
Pull them to $R_m/(x)$ and push forward along $x=0$ as in
Theorem~\ref{d3:thm:henselianeisensteinrank}.  Their resulting classes in
$K_3(R_m)$ vanish after inverting $x$, and so do their images in
$K_3(S)$.

After mapping to $A_{\infty,3}$ these are exactly the completed
finite-level classes detected in
Theorem~\ref{d3:thm:stricthenselianinfinite}.  Its proof applies to
arbitrary independent anti-invariant coefficient classes from each
finite unramified level.  It therefore proves that these $m$ images
are $\Q$-independent.  Any rational relation in $K_3(S)$ would remain
a relation in $K_3(A_{\infty,3})$, so the classes in $K_3(S)$ are
independent as well.  Since $m$ is arbitrary, the result follows.
\end{proof}


\section{Cyclotomic coefficients in higher odd degrees}
\label{d3:sec:completedhigherodd}

The following construction supplies actual higher $K$-classes without
requiring a general surjectivity assertion from actual $K$-theory to
its $p$-completed realization. Fix a weight $n\geq2$. We write
$c_{n,1}:K_{2n-1}\to H^1(-,\Q_p(n))$ and
$c_{n+1,3}:K_{2n-1}\to H^3(-,\Q_p(n+1))$; the second subscript
denotes the cohomological degree.

\begin{lemma}
\label{d3:lem:highercoefficients}
Let $E/\Q_p$ be finite and $n\geq2$.
For a root of unity $\rho\in E$, the number-field cyclotomic class
$b_n(\rho)\in K_{2n-1}(\OO_E)_{\Q}$ has
\begin{equation}\label{d3:eq:highercyclotomicregulator}
 c_{n,1}(b_n(\rho))
 =\exp_{E,n}\bigl(\kappa_n\operatorname{Li}_n(\rho)\bigr),
 \qquad \kappa_n\in\Q^\times,
\end{equation}
where $\kappa_n$ is common to all roots and embeddings under one
fixed symbol convention. The exponential is an isomorphism for
$n\geq2$. For an invertible module $L$ and an actual class $\beta$
of degree $2n-1$, one has
\begin{equation}\label{d3:eq:higherlineproduct}
 c_{n+1,3}\bigl(\beta(1-[L])\bigr)
 =n\,c_{n,1}(\beta)\cup c_1(L).
\end{equation}
Finally, Lemmas~\ref{d3:lem:completedcup} and
\ref{d3:lem:completedrootcup} hold with coefficient group
$H^1(E,\Q_p(n))$ and target twist $n+1$.
\end{lemma}
\begin{proof}
Theorem~1.12 of \cite{BesserDeJeu} constructs the actual rational
cyclotomic class and includes roots of $p$-power order.
Its normalized syntomic value is
$\pm(n-1)!\operatorname{Li}_n(\rho)$.
To compare the normalizations, put $q=p^f=|k_E|$, let $E_0$ be the
maximal unramified subfield of $E$, and write $\sigma$ for its Frobenius.
Set
\[
 A=1-p^{-n}\sigma,\qquad
 N=\sum_{j=0}^{f-1}p^{-nj}\sigma^j,\qquad P=1-q^{-n}.
\]
Then $NA=AN=P$.  The point syntomic complex of
\cite[Example~2.2.4]{HuberKings} is
$[E_0\longrightarrow E\oplus E_0]$, with
$c\mapsto(c,Ac)$, and its canonical $E$-coordinate is
$\eta([a,b])=a-A^{-1}b$.
The comparison of \cite[Proposition~8.6(2)--(3), its proof and equation~(8.1)]{BesserRegulators}
uses $N$ to pass from the semilinear $p$-Frobenius to the linear
$q$-Frobenius.  On the point the latter acts as the identity on $E$,
and the raw modified-cone coordinate is
\[
 \epsilon=Pa-Nb=P\eta([a,b]).
\]
Definition~4.6 of \cite{BesserDeJeu} sends this coordinate to
\begin{equation}\label{d3:eq:pointnormalization}
 (1-\varphi_q^*/q^n)^{-1}\epsilon
   =P^{-1}\epsilon=\eta([a,b]).
\end{equation}
Thus \cite[Corollary~9.10 and Proposition~9.11]{BesserRegulators},
or equivalently \cite[Proposition~2.3.4]{HuberKings}, identifies the
normalized regulator with the Chern-compatible Bloch--Kato
exponential.  This proves \eqref{d3:eq:highercyclotomicregulator};
the only remaining constant is the fixed rational symbol normalization
in Theorem~1.12 and Remark~1.7 of \cite{BesserDeJeu}.
The calculation applies to arbitrary ramification and to $p=2$.

Formula~(3.10) and Theorem~3.11(i) of \cite{WeibelChern}, applied
to $[L]-1$ in degree zero and $\beta$ in degree $2n-1$, have only
the term $i_1=1,i_2=n$ in weight $n+1$.
Its coefficient is $-n$; negating $[L]-1$ gives
\eqref{d3:eq:higherlineproduct}. The exceptional $K$-degree two
case does not occur. The finite-type descent proof of
Lemma~\ref{d3:lem:continuouschern} applies unchanged in degree
$2n-1\geq3$, so the identity holds in continuous cohomology.

For the cup assertion use the same fixed-radius map
\eqref{d3:eq:schemetodagger} and the same ruling morphism
$f:U^\dagger\to\PP^1_E$. The boundary with twist $n+1$ in
\cite[Theorem~1.3]{ColmezNiziol} is injective in degree three
when $n=2$ and is an isomorphism when $n\geq3$, since
$3\leq(n+1)-1$. On $\PP^1_E$ the boundary is an isomorphism
also for $n=2$, as in Lemma~\ref{d3:lem:completedcup}.
Its naturality and the de Rham ruling isomorphism make
$f^*$ injective on $H^3_{\mathrm{pro\acute et}}(-,\Q_p(n+1))$.
The splitting \eqref{d3:eq:finitePonebundle}, with $r=n+1$,
identifies cup product with $c_1(\OO(1))$ as a split injection
from $H^1(E,\Q_p(n))$ into its source.
Composing with the injective $f^*$ proves the quadratic assertion.
The pairwise formal maps in Lemma~\ref{d3:lem:completedrootcup}
then give the simultaneous assertion with the same diagonal kernel.
\end{proof}

\begin{lemma}
\label{d3:lem:evenweightibase}
For every even integer $n\geq2$ and $i^2=-1$,
\[
 \frac{\operatorname{Li}_n(i)}{i}\in\frac12+\Z_2.
\]
In particular $v_2(\operatorname{Li}_n(i))=-1$.
\end{lemma}
\begin{proof}
Choose a primitive cube root $\zeta$ in an unramified quadratic
extension and put
$D_n=\operatorname{Li}_n(\zeta i)+\operatorname{Li}_n(\zeta^2i)$.
Write $F_{2,n}(z)=\operatorname{Li}_n(z)-2^{-n}\operatorname{Li}_n(z^2)$.
Distribution and inversion give
\[
 D_n=-(1+3^{1-n})\operatorname{Li}_n(i),\qquad
 \operatorname{Li}_n(-\zeta)+\operatorname{Li}_n(-\zeta^2)=0.
\]
Thus $D_n=F_{2,n}(\zeta i)+F_{2,n}(\zeta^2i)$.
For $S_j(t)=\sum_{m\geq0}m^jt^m
=P_j(t)/(1-t)^{j+1}$, expand the odd denominators by the binomial
series first on $|z|<1$. The terms converge uniformly when
$|z|\leq1$ and $v_2(1-z^2)\leq2/3$, with valuation at least
$j-(j+1)2/3$. The rigid identity principle and
\cite[Theorem~2.7(3)]{BesserDeJeuAlgorithm} therefore give
\[
 \frac{D_n}{i}
 =\sum_{j\geq0}(-1)^j\binom{n+j-1}{j}2^j A_j,
 \qquad A_j=\sum_{a=1}^2\zeta^a S_j(-\zeta^{2a}).
\]
It applies because $1+\zeta^{2a}$ is a unit.
Every $A_j$ belongs to $\Z[\zeta]$ and is fixed by
$\zeta\mapsto\zeta^2$, hence lies in $\Z$.
Directly, $A_0=-2$ and $A_1=1$.
Consequently
\[
 D_n/i\in-2-2n+4\Z_2=-2+4\Z_2.
\]
Since $n-1$ is odd, $3^{n-1}\equiv3\pmod8$ and
$v_2(1+3^{1-n})=2$.
The first displayed identity now gives
$\operatorname{Li}_n(i)/i\in\Q_2$ of valuation $-1$,
which is exactly the asserted coset.
\end{proof}

\begin{proposition}
\label{d3:prop:higherwildaugmentation}
Let $r\geq3$, let $\rho_r$ be a primitive $2^r$-th root with
$\rho_r^2=\rho_{r-1}$, and put
$E_r=\Q_2(\rho_r)$ and $D=2^{r-2}$.
For every $n\geq2$,
\begin{equation}\label{d3:eq:higherwilddepletion}
 v_2\left(\operatorname{Li}_n(\rho_r)
       -2^{-n}\operatorname{Li}_n(\rho_{r-1})\right)
 =-2^{-(r-2)}.
\end{equation}
The $D-1$ elements
\[
 \operatorname{Li}_n(\rho_r^{5^{j+1}})
 -\operatorname{Li}_n(\rho_r^{5^j}),
 \qquad 0\leq j\leq D-2,
\]
are $\Q_2$-linearly independent, and their span intersects
$\Q_2\subset E_r$ only in zero.
If $n$ is even, the $D$ elements
\[
 \operatorname{Li}_n(\rho_r^{5^j}),\qquad 0\leq j\leq D-1,
\]
form a $\Q_2$-basis of the inversion-odd subspace $E_r^-$.
\end{proposition}
\begin{proof}
Put $F_{2,n}(z)=\operatorname{Li}_n(z)-2^{-n}\operatorname{Li}_n(z^2)$.
With $S_j(t)=P_j(t)/(1-t)^{j+1}$ as above, expansion on $|z|<1$ gives
\[
 F_{2,n}(z)=z\sum_{j\geq0}(-1)^j
       \binom{n+j-1}{j}2^j S_j(z^2).
\]
The sum converges uniformly on
$|z|\leq1$, $v_2(1-z^2)\leq2/3$, since its $j$th term has
valuation at least $j-(j+1)2/3\to\infty$.
This is a connected annulus inside $|1-z|>1/2$;
\cite[Theorem~2.7(3)]{BesserDeJeuAlgorithm} and the rigid identity
principle extend the equality there. At $\rho_r$ put
$d=v_2(1-\rho_r^2)=2^{-(r-2)}\leq1/2$.
The $j=0$ term has valuation $-d$, while every $j\geq1$ term
has valuation at least $j-(j+1)d>-d$.
This proves \eqref{d3:eq:higherwilddepletion}.

Write $R_s=\operatorname{Li}_n(\rho_s)$ and
$\Gamma_r=\Gal(E_r/\Q_2(i))=\langle\gamma\rangle$, with
$\gamma(\rho_r)=\rho_r^5$ and $|\Gamma_r|=D$.
Inversion acts on $R_r$ by $\epsilon=(-1)^{n+1}$.
The normal-basis theorem makes $E_r^\epsilon$ the regular
$\Q_2[\Gamma_r]$-module \cite[Chapter~I, p.~34]{MilneDuality}. Distribution gives average trace
$R_s\mapsto2^{-n}R_{s-1}$.
For every $3\leq s\leq r$, the rational conductor
$2^{s-2}$ projection of $R_r$ is therefore
\[
 2^{-n(r-s)}(R_s-2^{-n}R_{s-1})\ne0.
\]
Each nontrivial rational irreducible of the cyclic group occurs
once, and its polynomial $\Phi_{2^a}$ is irreducible over $\Q_2$.
Hence the augmentation ideal acts injectively on $R_r$.
Its basis $\gamma^{j+1}-\gamma^j$, $0\leq j\leq D-2$, gives
the asserted independent values. Every such value has
$\Gamma_r$-average zero, so their span meets $\Q_2$ only in zero.
If $n$ is even, the trivial $\Gamma_r$ projection is
$2^{-n(r-2)}\operatorname{Li}_n(i)\ne0$ by
Lemma~\ref{d3:lem:evenweightibase}. All rational blocks then occur,
so the whole orbit is a basis of $E_r^-$.
\end{proof}

\begin{theorem}
\label{d3:thm:completedhigherodd}
For $r\geq3$ put $D=2^{r-2}$ and
\[
 A_r=\Z_2[[x,y,z]]/
 \bigl(xy+\Phi_{2^r}(1+z)\bigr).
\]
For every $n\geq2$,
\[
 \dim_{\Q}\ker\left(K_{2n-1}(A_r)_{\Q}
       \longrightarrow K_{2n-1}(\Frac A_r)_{\Q}\right)
 \geq
 \begin{cases}
 D,& n\text{ even},\\
 D-1,& n\text{ odd}.
 \end{cases}
\]
More precisely, let $W_r=\Z_2[\rho_r]$, identify
$A_r/(x)=W_r[[y]]$, and let $i$ be this closed immersion.
For even $n$, use the $D$ classes
\[
 i_*b_n(\rho_r^{5^j}),\qquad 0\leq j\leq D-1.
\]
For odd $n$, use the $D-1$ classes
\[
 i_*\bigl(b_n(\rho_r^{5^{j+1}})-b_n(\rho_r^{5^j})\bigr),
 \qquad 0\leq j\leq D-2.
\]
The coefficients are first pulled to $W_r[[y]]$.
The indicated classes are independent in the displayed kernel.
\end{theorem}
\begin{proof}
The polynomial $\Phi_{2^r}(1+z)$ is Eisenstein at two.
Its linear uniformizer term makes $A_r$ regular local of
dimension three, with residue field $\F_2$.
Localization makes each supported class generically zero.
After the finite free coefficient extension from $\Z_2$ to $W_r$,
the roots are $a_\sigma=\sigma(\rho_r)-1$.
On the regular generic open, flat base change and projection give
\[
 i_*(\beta)|_T
 =\sum_{\sigma\in\Gal(E_r/\Q_2)}
      \sigma(\beta)(1-[J_\sigma]).
\]
Base change may be performed in $G$-theory before restricting
to this regular open, so regularity of the split integral model
is not required. By Lemma~\ref{d3:lem:highercoefficients}, its Chern
class is
\[
 n\sum_\sigma
 \exp_{E_r,n}(\kappa_n\sigma(R))\cup c_1(J_\sigma),
\]
where $R$ is the corresponding difference in
Proposition~\ref{d3:prop:higherwildaugmentation}, or the individual
value there when $n$ is even.
The simultaneous cup detector kills a coefficient vector exactly
when it is diagonal. Since the exponential is injective, a
relation among these detectors would put the corresponding
$\Q_2$-linear combination of the $R$'s in $\Q_2$.
Proposition~\ref{d3:prop:higherwildaugmentation} makes that combination
zero (using $E_r^-\cap\Q_2=0$ in even weight), and then makes
all its coefficients zero.
The detectors are $\Q_2$-independent, hence the actual rational
$K$-classes are $\Q$-independent. A common integer clears the
denominators of this finite family and gives integral classes
whose generic images are zero integrally.
\end{proof}

For odd weight the trivial wild-character block lies in the
constant field and is deliberately removed by taking differences.
No assertion about surjectivity of the actual local higher
$K$-regulator onto its completed target is needed.


\begin{corollary}
For the single regular local ring
\[
 R=\Z_2[[x,y,z]]/
 \bigl(xy+z^4+4z^3+6z^2+4z+2\bigr),
\]
the generic restriction has nonzero rational kernel in every
odd degree at least three. In degree $4j-1$ its kernel has
$\Q$-dimension at least two, and in degree $4j+1$, $j\geq1$,
at least one.
\end{corollary}
\begin{proof}
Take $r=3$ in Theorem~\ref{d3:thm:completedhigherodd} and use
$\Phi_8(1+z)=(1+z)^4+1$.
\end{proof}


\subsection{Higher polylogarithms on odd-prime wild orbits}
\label{d3:sec:oddprimehigherwild}

Fix an odd prime $p$, put $M=p-1$, and choose compatible primitive
$p^r$-th roots $\rho_r$. Write
\[
 E_r=\Q_p(\rho_r),\qquad e_r=Mp^{r-1},\qquad
 R_r=\operatorname{Li}_n(\rho_r),\qquad n\geq2.
\]
Let $\Delta$ be the tame subgroup and $\omega:\Delta\to\mu_M$
its Teichm\"uller character. Put $\chi_n=\omega^{n-1}$ and
$\Gamma_r=\Gal(E_r/E_1)$, generated by
$\gamma(\rho_r)=\rho_r^{1+p}$.

\begin{proposition}
\label{d3:prop:oddprimehigherwild}
For even $n$, let $\ell\in\{1,3,\ldots,p-2\}$ satisfy
$\ell\equiv n-1\pmod M$. Then
\begin{equation}\label{d3:eq:oddprimehigherbase}
 v_p(R_1)=\frac{\ell}{p-1}-1-v_p(n-1).
\end{equation}
For either parity of $n$, let $q\in\{0,\ldots,p-2\}$ satisfy
$q\equiv n\pmod M$. For $r\geq2$,
\begin{equation}\label{d3:eq:oddprimehighernew}
 v_p(R_r-p^{-n}R_{r-1})=\frac{q-p}{e_r}.
\end{equation}
The $\chi_n$-projection preserves the leading term in both formulas.
Writing $D=p^{r-1}$, the following families are $\Q_p$-linearly
independent:
\[
 \{\gamma^jR_r:0\leq j<D\}\quad(n\text{ even}),
\]
and
\[
 \{\gamma^{j+1}R_r-\gamma^jR_r:0\leq j<D-1\}
                                  \quad(n\text{ odd}).
\]
The span of either displayed family has zero intersection with
$\Q_p\subset E_r$.
\end{proposition}

\begin{proof}
Set $F_{p,n}(z)=\operatorname{Li}_n(z)-p^{-n}
\operatorname{Li}_n(z^p)$ and
$S_j(t)=P_j(t)/(1-t)^{j+1}$, $P_j\in\Z[t]$, where
$S_j(t)=\sum_{h\geq0}h^jt^h$ and $S_0(t)=1/(1-t)$.
Grouping the prime-to-$p$ denominators as $a+ph$ gives
\begin{equation}\label{d3:eq:oddprimehigherdepleted}
 F_{p,n}(z)=\sum_{j\geq0}(-1)^j\binom{n+j-1}{j}p^j
       \sum_{a=1}^{p-1}\frac{z^a}{a^{n+j}}S_j(z^p).
\end{equation}
On $|z|\leq1$, $v_p(1-z^p)\leq b<1$, the $j$th term has
valuation at least $j-(j+1)b$, so the series converges uniformly.
This domain is the annulus
$p^{-b/p}\leq|1-z|\leq1$ and lies in the continuation region
of \cite[Theorem~2.7(3)]{BesserDeJeuAlgorithm}.
Both sides agree on $|z|<1$, hence on the annulus. Taking
$b=2/3$ includes all points evaluated below.

Suppose first that $n$ is even, and put $\rho=\rho_1$.
Distribution and inversion
\cite[Proposition~2.10]{BesserDeJeuAlgorithm} give
\[
 \sum_{\xi\in\mu_M\setminus\{1\}}\operatorname{Li}_n(\xi\rho)
       =-(1+M^{1-n})\operatorname{Li}_n(\rho).
\]
Since $(\xi\rho)^p=\xi$ and
$\sum_{\xi\ne1}\operatorname{Li}_n(\xi)=0$ by inversion,
the sum on the left can be replaced by
$\sum_{\xi\ne1}F_{p,n}(\xi\rho)$.
All terms of \eqref{d3:eq:oddprimehigherdepleted} with $j\geq1$
then belong to $p\mathcal O_{E_1}$. Its $j=0$ sum is
\[
 P_n(\rho)=\sum_{a=1}^{p-1}\frac{C_a\rho^a}{a^n},\qquad
 C_a=\sum_{\xi\ne1}\frac{\xi^a}{1-\xi}=a-\frac p2.
\]
Indeed, pairing inverse roots gives
$\sum_{\xi\ne1}(1-\xi)^{-1}=(M-1)/2$; hence
$C_1=1-p/2$, and $C_{a+1}-C_a=1$ for $1\leq a<M$.

For any integer $b$, the coefficient of $(Z-1)^k$ in
$\sum_{a=1}^{p-1}a^bZ^a\in\F_p[Z]$ is
$\sum_a a^b\binom ak$.
If $q$ is the least nonnegative residue of $-b$ modulo $M$,
the first nonzero coefficient is at $k=q$ and equals $-1/q!$:
for $k<q$ no power occurring in the degree-$k$ polynomial
$a^b\binom ak$ has exponent divisible by $M$, while at $q$
only its leading coefficient contributes.
Applying this with $b=1-n$ shows that $\overline P_n$
has a zero of exact order $\ell$ at $1$. Thus
\[
 v_pP_n(\rho)=\ell/M<1.
\]
The lower Taylor coefficients are divisible by $p$, and the
higher terms have strictly greater valuation. They and the
$j\geq1$ remainder cannot cancel this term.
Since $n-1$ is odd, the elementary lifting-the-exponent identity gives
\[
 v_p(1+M^{1-n})=1+v_p(n-1),
\]
proving \eqref{d3:eq:oddprimehigherbase}.

For $r\geq2$ and either parity, the $j=0$ term in
\eqref{d3:eq:oddprimehigherdepleted} is
\[
 \frac{\sum_{a=1}^{p-1}a^{-n}\rho_r^a}{1-\rho_r^p}.
\]
The same Taylor calculation gives numerator valuation $q/e_r$,
including $q=0$. The denominator has valuation $p/e_r$.
For $j\geq1$ the valuation is at least
\[
 j-(j+1)p/e_r\ \geq\ j-(j+1)/(p-1)\ \geq\ 0,
\]
and this bound tends to infinity. The $j=0$ term has negative
valuation, so it is the unique leading contribution. This proves
\eqref{d3:eq:oddprimehighernew}.

It remains to verify the character assertion, rather than merely
nonvanishing of the total values. Put $\pi_r=\rho_r-1$.
If $\tau\in\Delta$ has $\overline{\omega(\tau)}=a$, then
$\tau(\pi_r)/\pi_r\equiv a\pmod{\pi_r}$.
For $0\ne x\in E_r$ with $v_{\pi_r}(x)=m$, the leading term
therefore transforms by $a^m$. The idempotent
\[
 e_{\chi_n}=\frac1M\sum_{\tau\in\Delta}\chi_n(\tau)^{-1}\tau
\]
preserves it if $m\equiv n-1\pmod M$.
For the base value this congruence follows from
$m=\ell-M(1+v_p(n-1))$; for the wild difference it follows
from $m=q-p$.

The normal basis theorem identifies $E_r^{\chi_n}$ with the
regular $\Q_p[\Gamma_r]$-module
\cite[Chapter~I, p.~34]{MilneDuality}. Its rational character blocks
occur once, with polynomials $\Phi_{p^j}$, $0\leq j\leq r-1$.
If $H_r=\Gal(E_r/E_{r-1})$, distribution gives
\[
 \frac1p\sum_{h\in H_r}hR_r=p^{-n}R_{r-1}.
\]
Thus \eqref{d3:eq:oddprimehighernew} supplies a nonzero projection
to the new block $\Phi_{p^{r-1}}$, and the averaging identity
preserves every previously nonzero block.
For even $n$, \eqref{d3:eq:oddprimehigherbase} starts this induction
with a nonzero trivial wild block, giving a cyclic vector of
the full regular module. For odd $n$, all nontrivial blocks
occur, so the augmentation ideal maps injectively.
Applying $e_{\chi_n}$ to a relation now proves the two asserted
independence statements for the unprojected values themselves.
Finally, for even $n$ inversion is odd, excluding $\Q_p$;
for odd $n$ every displayed difference has zero $\Gamma_r$-average,
which likewise excludes $\Q_p$.
\end{proof}

\begin{remark}
The tame idempotent is used only on the regulator target.
Actual coefficients may remain the root classes
$b_n(\rho_r^{(1+p)^j})$ or their adjacent differences; no
$\Q_p$-linear projector is applied inside actual rational
$K$-theory. With the common normalization
\eqref{d3:eq:highercyclotomicregulator}, their usual local Chern
values are the corresponding displayed values followed by
multiplication by $\kappa_n\ne0$ and the Bloch--Kato exponential.
Since these operations are $\Q_p$-linear isomorphisms, the
actual root family has $\Q$-span of dimension $D$ in even
weight, and the actual difference family has $\Q$-span of
dimension $D-1$ in odd weight: any rational relation would
give a forbidden $\Q_p$-linear relation of regulator values.
The proposition detects one tame character and
does not assert nonvanishing of all tame character components.
\end{remark}

\begin{corollary}
\label{d3:cor:oddprimecompletedhigher}
For an odd prime $p$ and $r\geq2$ put $D=p^{r-1}$ and
\[
 A_{p,r}=\Z_p[[x,y,z]]/(xy+\Phi_{p^r}(1+z)).
\]
For every $n\geq2$, the rational generic kernel in degree $2n-1$
has dimension at least $D$ when $n$ is even and at least $D-1$
when $n$ is odd. The classes are the pushforwards from $x=0$ of
$b_n(\rho_r^{(1+p)^j})$ in even weight, or their consecutive
differences in odd weight, after one common denominator is cleared.
They vanish after inverting $x$ integrally.
In particular, fixing $r=2$ gives one complete regular local ring
with residue field $\F_p$ for which rational generic injectivity
fails in every odd degree at least three.
\end{corollary}
\begin{proof}
The polynomial is Eisenstein and the equation has nonzero linear
term $p$, so the displayed ring is regular local of dimension three.
Apply the proof of Theorem~\ref{d3:thm:completedhigherodd}, with
Proposition~\ref{d3:prop:oddprimehigherwild} in place of
Proposition~\ref{d3:prop:higherwildaugmentation}.
The simultaneous ruling-cup detector in
Lemma~\ref{d3:lem:highercoefficients} kills exactly the diagonal
coefficient vectors. The regulator spans just computed meet the
constant field only in zero, so the same proof gives precisely
these ranks and the stated integral supported representatives.
\end{proof}

\section{Continuous cyclic characters and explicit families}\label{d3:sec:directcyclic}

This section supplies an independent proof for the explicit completed
three-adic class, and a direct endpoint criterion for all the explicit
prime-power families. The first proof uses an elementary convergent
series. The second uses Coleman functions only as analytic primitives.
The coefficient-cup argument at the end independently supplies actual
local coefficients for existence and rank constructions.

\subsection{A direct detector in continuous cyclic homology}
\label{d3:sec:directtcdetector}

For a spectrum $M$, write $M_{\mathbf Q_p}=M^{\wedge}_p[1/p]$.
Thus $K(R;\mathbf Q_p)$ denotes a rationalized completed spectrum;
it does not denote the tensor product of each actual $K$-group with
$\mathbf Q_p$. Products of units below are products of actual
Quillen $K_1$-classes.

Let $W=\mathcal O_E$, where $E/\mathbf Q_p$ is finite, and let $\pi$
be a uniformizer. If $\mathfrak X=\operatorname{Spf}B$ is smooth and
topologically of finite presentation over $\operatorname{Spf}W$, put
\[
 \Omega_B^j=\widehat\Omega^j_{B/W}[1/p].
\]
The hat denotes continuous differentials, completed before inverting
$p$. Cohomology is ordinary cohomology in $D(E)$, with differential images taken without closure.

\begin{lemma}
\label{d3:lem:relativelogproduct}
Suppose that $\mathfrak X=\operatorname{Spf}B$ is smooth, affine,
topologically of finite presentation, and of relative dimension at
most two over $\operatorname{Spf}W$. There is a natural relative map
\[
 \chi_B:K(B,\pi B)\longrightarrow
                  \Sigma HC^{\mathrm{cts}}(B/W;\mathbf Q_p).
\]
For $v\in1+\pi B$ and $b_1,b_2\in B^\times$, the top HKR component
of its degree-three value is
\begin{equation}\label{d3:eq:tcreviewlogproduct}
 \chi_B([v][b_1][b_2])_{\mathrm{top}}
   =\log(v)\,d\log b_1\wedge d\log b_2
       \quad\text{in }\Omega_B^2/d\Omega_B^1,
\end{equation}
up to one universal sign determined by the norm and suspension
conventions. Here $[v]$ has its relative lift furnished by its reduction
to the identity unit.
\end{lemma}

\begin{proof}
The fiber square of \cite[Proposition~4.10]{AMMN} has rows
\[
\begin{matrix}
 K^{\mathrm{cts}}(\mathfrak X;\mathbf Q_p)
       &\longrightarrow&K(B/\pi;\mathbf Q_p)\\
 \downarrow&&\downarrow\\
 TC^{\mathrm{cts}}(\mathfrak X;\mathbf Q_p)
       &\longrightarrow&TC(B/\pi;\mathbf Q_p)\\
 \downarrow&&\downarrow\\
 HC^{-,\mathrm{cts}}(B/W;\mathbf Q_p)
       &\longrightarrow&HP^{\mathrm{cts}}(B/W;\mathbf Q_p).
\end{matrix}
\]
The last horizontal fiber is $\Sigma HC^{\mathrm{cts}}(B/W;\mathbf Q_p)$,
not negative cyclic homology. Mapping the actual relative spectrum to
these fibers gives $\chi_B$. Consequently relative $K_3$ has target
$HC_2$, and its top component is $\Omega_B^2/d\Omega_B^1$.

We record why the fiber comparison respects multiplication by absolute
$K$-classes. Set $H=THH(B)$, a commutative algebra in cyclotomic spectra.
If $TC_{\mathrm{cyc}}$ denotes the functor on that category, then
$TC_{\mathrm{cyc}}(H)=TC(B)$; no second $THH$ is intended.
All objects $H\otimes D$ and maps $\mathrm{id}_H\otimes f$ in
\cite[Proposition~2.8, Corollary~2.10, Theorem~2.12]{AMMN} are
$H$-modules and module maps. The fixed-point and Tate functors, their
canonical maps, and the norm triangle preserve the induced $TC(B)$-actions
\cite[Corollary~I.4.3]{NikolausScholze}. Thus the inverses of the
equivalences used there are module inverses. The comparisons from
derived to ordinary reduction and then to $B/\pi$ are induced by
$B$-algebra maps. The projection $HH(B)\to HH(B/W)$ in the proof of
Proposition~4.10 also preserves the action. These inverses and fiber
equivalences are used after $p$-completion and rationalization.

The multiplicative cyclotomic trace
\cite[Theorem~7.4]{BGTTrace} now gives $K(B)$-linearity, with target
action through the canonical map $K(B)\to HC^{-,\mathrm{cts}}(B/W;\mathbf Q_p)$.
At the cyclic-complex level, the compatibility of the norm boundary
and its suspension with the $HC^-$-module action is the Connes-sequence
compatibility of \cite[Proposition~2.3]{BokstedtOttosen}.
The identification of the spectral norm with the Connes operator is \cite[Theorem~2.1, Lemma~2.2, Theorem~2.3]{HoyoisCircle};
this includes
the usual graded sign of a connecting map. In particular, no assumption
that an unspecified additive fiber equivalence is multiplicative is
being made.

To normalize a relative unit, first take $B=W\langle T\rangle$ and
$v=1+\pi T$. Its relative character is a function
$f\in HC_0^{\mathrm{cts}}(B/W;\mathbf Q_p)=E\langle T\rangle$.
Choose the norm convention in which its image in $HC^-_1$ is $df$.
The left vertical arrow in \cite[Theorem~2.12]{AMMN} is canonical;
its projection to Hochschild homology takes a unit to the Dennis trace
$v^{-1}\otimes v$. Hence $df=d\log v$. Evaluation at $T=0$ gives
$f(0)=0$. Continuous differentiation on $E\langle T\rangle$ has
only constants in its kernel, and therefore
\[
 f=\log(1+\pi T)=\sum_{n\ge1}(-1)^{n+1}\pi^nT^n/n.
\]
The series converges beyond the closed unit disc. Substitution
$T=(v-1)/\pi$ proves the assertion for every relative unit under
consideration.

Finally use the normalized cyclic complex
$C^+=E[t,t^{-1}]/tE[t]\otimes C_*$, with $|t|=-2$.
The $HC^-\times HC$ action is shuffle plus $t$ times cyclic shuffle
\cite[Definition~2.1 and Remark~2.2]{BokstedtOttosen}.
Starting with the column-zero representative $\log v$, all positive
columns and cyclic-shuffle corrections vanish in $C^+$. After each
of the two unit multiplications only the ordinary shuffle remains.
Under
\[
 b_0\otimes\cdots\otimes b_n\longmapsto
       b_0\,db_1\wedge\cdots\wedge db_n/n!,
\]
the two shuffles in degree two cancel the factor $2!$, giving
\eqref{d3:eq:tcreviewlogproduct}, with the suspension signs just specified.
This calculation involves only degrees through two and commutes with
the continuous, relative-to-$W$ projection. It does not assume a
multiplicative splitting of an infinite HKR filtration.
\end{proof}

\begin{lemma}
\label{d3:lem:tcreviewresidue}
Let $a_0,a_1\in W$ have unit difference, and put
\[
 B=W\langle X,Y,Z\rangle/
           \bigl(XY+(Z-a_0)(Z-a_1)\bigr).
\]
For $F\in E\langle Z\rangle$, the form
$\omega_F=(dX/X)\wedge dF$ is globally regular. If it is exact in
$\Omega_B^2/d\Omega_B^1$, then $F(a_0)=F(a_1)$.
\end{lemma}
\begin{proof}
The two root-deletion affine opens have completed coordinates
\[
 Z=a_i+Xt_i,\qquad
 Y=-t_i(a_i-a_{1-i}+Xt_i)\quad(i=0,1).
\]
They are the completions of $D(X)\cup D(Z-a_{1-i})$, cover the formal
scheme, and have intersection $W\langle X,X^{-1},Z\rangle$.
The forms
\[
 q_i=-(F(Z)-F(a_i))\,dX/X
\]
are regular on the respective bidiscs, satisfy $dq_i=\omega_F$, and
have difference $(F(a_1)-F(a_0))dX/X$.

A closed bidisc form $A(X,t)dX+B(X,t)dt$ has zero Laurent residue
after $t=c/X$, for every $|c|\le1$. Indeed its residue is the
convergent series
\[
 \sum_{m\ge0}(a_{m,m+1}-b_{m+1,m})c^{m+1}=0,
\]
because closedness equates $(m+1)a_{m,m+1}$ and
$(m+1)b_{m+1,m}$. Division by $m+1$ is used only in this individual
coefficient equality, not to construct a bounded primitive.
If $\omega_F=dq$ globally, apply this to $q_i-q$ after restricting
the overlap to $Z=0$, where $c=-a_i$. The difference then has residue
zero, which is the required endpoint equality. This argument includes
$c=0$ and works over the residue field $\mathbf F_2$.
\end{proof}

\begin{proposition}
\label{d3:prop:directtcoriginal}
Let $A=\mathbf Z_3[[x,y,z]]/(3+xy+z^2)$. The actual element
\[
 c=2\left\{\frac{1+z}{2},\,1+x,\,
                        1+\frac{(1+x)y}{4}\right\}\in K_3(A)
\]
has infinite order and has zero image in $K_3(A[1/x])$.
The ring $A$ is regular local of dimension three.
\end{proposition}
\begin{proof}
This is the completed case of Theorem~\ref{thm:explicit-three-dimensional}.
The proof there constructs the relative class on the smooth formal quadric,
uses the actual vanishing $K_4(C)=0$ to pass to absolute $K$-theory,
and detects the resulting class by an endpoint difference in the ordinary
continuous de Rham quotient.
\end{proof}

\paragraph{The finite continuous HKR calculation.}
For the three-adic quadric, put
\[
 B_0=W[X,Y,Z]/(XY+Z^2-1),\qquad B=\widehat{B_0}_3,
 \qquad \Omega_B^i=\widehat\Omega^i_{B/W}[1/3].
\]
The algebra $B_0$ is smooth over $W$ of relative dimension two.
The normalized HKR map on its Hochschild complex is
\[
 b_0\otimes\cdots\otimes b_n\longmapsto
 \frac{1}{n!}b_0\,db_1\wedge\cdots\wedge db_n
 \quad(0\le n\le2),
\]
and is zero in higher degrees.  It is defined over $W$, since $2$ is
a unit, and intertwines the Hochschild differential with zero and the
Connes operator with $d$.  Smoothness makes it a quasi-isomorphism of
mixed complexes: the Hochschild homology is $\Omega^n_{B_0/W}$ for
$0\le n\le2$ and vanishes above degree two.
The comparison with the continuous invariant is supplied by
\cite[Proposition~4.2]{AMMN}.  In the cyclic, negative cyclic, and
periodic total complexes of this finite mixed complex, each total
degree has only finitely many differential-form terms.  Their
reduction towers are surjective, so derived three-adic completion is
computed termwise.  After inverting $3$, this gives in particular
\[
 HC^{\mathrm{cts}}_0(B/W;\mathbf Q_3)=\Omega_B^0,
 \qquad
 HC^{\mathrm{cts}}_2(B/W;\mathbf Q_3)
   =\frac{\Omega_B^2}{d\Omega_B^1}
     \oplus\ker(d:\Omega_B^0\to\Omega_B^1).
\]
The norm map from the degree-zero cyclic term to the degree-one
negative cyclic term is $f\mapsto df$, with the chosen global
suspension sign.  In particular, the displayed quotient is the
ordinary quotient by the image of $d$; its image is not replaced by
its closure.  The same calculation in relative dimension one applies
to the universal unit in $W\langle T\rangle$.


\paragraph{The ordinary quotient is essential.}
The target in \cite[Constructions~4.6 and~4.11]{AMMN} uses the
ordinary derived category. For example, in the preceding quadric put
$\omega_0=(dX/X)\wedge dZ$. The globally regular primitives
$-(Z-Z^{3^j})dX/X$ satisfy
\[
 d\!\left(-(Z-Z^{3^j})\frac{dX}{X}\right)
   =\omega_0-3^jZ^{3^j-1}\omega_0\longrightarrow\omega_0.
\]
The ruling form is therefore in the closure of exact forms although
the residue lemma shows it is not exact. No assertion about
Hausdorff-reduced cohomology or the order at each finite Artin stage
is made here.

\subsubsection{The same computation for the explicit families}

For odd $p$ and $r\ge1$, put $M=p^r$, $n=2M$, and
\[
 A_{p,r}=\mathbf Z_p[[x,y,z]]/(xy+\Phi_{2M}(z-1)),
 \qquad u=z-1,\qquad P(u)=\Phi_{2M}(u).
\]
For $p=2$, take $r\ge2$, $M=2^r$, $n=3M$, and
\[
 A_{2,r}=\mathbf Z_2[[x,y,w]]/(xy+\Phi_M(1+w)),
 \qquad u=1+w,\qquad P(u)=\Phi_M(u).
\]
Let $H(u)=(u^n-1)/P(u)$, $b=yH(u)$,
$\lambda=1-x$, and $\mu=1+\lambda b$. For positive odd $k$ prime
to $p$, put
\[
 g_k=\begin{cases}1-u^k,&p\text{ odd},\\
                  1+u^k+u^{2k},&p=2,
       \end{cases}
 \qquad c_k=k\{g_k,\lambda,\mu\}\in K_3(A_{p,r}).
\]
Every displayed entry is a unit. We use Coleman $\operatorname{Li}_2$
with its series at zero and $\log p=0$, and at roots set
\[
 D_k(a)=\begin{cases}
       \operatorname{Li}_2(a^k),&p\text{ odd},\\
       \tfrac13\operatorname{Li}_2(a^{3k})-
                        \operatorname{Li}_2(a^k),&p=2.
       \end{cases}
\]

\begin{proposition}
\label{d3:prop:directtcendpoints}
The classes $c_k$ vanish integrally after inverting $x$.
If a rational relation $\sum_k q_kc_k=0$ holds in
$K_3(A_{p,r})\otimes\mathbf Q$, then
\[
       \sum_k q_kD_k(a)=0
\]
for every root $a$ of $P$.
\end{proposition}
\begin{proof}
Write $s=1-xb=u^n$. The Dennis--Stein identity gives
$\delta=\langle x,b\rangle=\{\lambda s,\mu\}$.
At odd $p$, Steinberg applied to $u^k$ gives $k\{g_k,s\}=0$.
At two, apply Steinberg to $\zeta_3u^k$ after the finite \'etale
extension adjoining $\zeta_3$ and transfer: since $3\mid n$, the
second entry comes from $u^{kn}$ and the norm of the first is exactly
$g_k$. Thus the same relation holds, without an extra factor two.
It follows that $c_k=k[g_k]\delta$, and the unit expression for
$\delta$ on $D(x)$ makes its generic image zero.

Fix distinct roots $a,a'$ in $E=\mathbf Q_p(\mu_M)$, set
$W=\mathcal O_E$, $d=a'-a$, and
$C(u)=P(u)/((u-a)(u-a'))$. There is an actual convergent map to the
smooth good reduction model
\[
 B=W\langle X,Y,Z\rangle/(XY+Z(Z+1)),\qquad
 u\mapsto a-dZ,\quad x\mapsto dX,\quad
 y\mapsto dC(a-dZ)Y.
\]
All roots reduce to $-1$ at odd $p$ and to $1$ at two, so $d$ and
the images of the original formal parameters are topologically
nilpotent. No inverse of $C$ is used. The special fiber is again an
affine-line torsor over $\mathbf P^1$, so its $K_3$ and $K_4$
vanish after $p$-completion and rationalization.

Set $m=p-1$ at odd $p$ and $m=1$ at two. Then $g_k^m\in1+\pi B$,
and $k[g_k^m][\lambda][\mu]$ is an actual relative lift of $mc_k$.
Let $L_k=\log g_k$ and choose $F_k$ with
$dF_k=L_k\,d\log s=nL_k\,du/u$. These functions converge beyond
the closed $Z$-disc. The differential identity
\[
 d\log\lambda\wedge d\log\mu
 =d\log x\wedge d\log s-d\log s\wedge d\log\mu
 =-dx\wedge db/\mu
\]
gives the following global representative of the relative character:
\[
 \tau_k=mk\frac{dX}{X}\wedge dF_k
                         -d(mkF_k\,d\log\mu).
\]

At odd $p$, $u^k$ lies in the nonsingular residue disc of $-1$;
the analytic differential equation gives
$F_k=-(n/k)\operatorname{Li}_2(u^k)+\mathrm{constant}$.
At two, use the rigid analytic sum
\[
 \operatorname{Li}_2(\zeta_3u^k)+
                   \operatorname{Li}_2(\zeta_3^2u^k).
\]
Its two arguments lie in nonsingular residue discs, its derivative is
$-kL_k\,du/u$, and it descends from the unramified quadratic
extension. Distribution gives its endpoint values $D_k(a)$.
These are only identities of analytic functions; the necessary
analyticity, distribution and inversion are
\cite[Theorem~2.7(1), Proposition~4.4, Proposition~2.10]{BesserDeJeuAlgorithm}.
No assertion of analyticity for the individual $\operatorname{Li}_2$
terms across the residue disc of $1$ is used at two.

Lemma~\ref{d3:lem:tcreviewresidue}, with roots $0,-1$, now computes the
endpoint obstruction of $mc_k$ as
\[
          \pm mn\bigl(D_k(a')-D_k(a)\bigr).
\]
The prefactor $k$ has canceled the $1/k$ in the primitive, leaving
one common scalar. A rational relation therefore makes
$\sum_kq_kD_k(a)$ independent of $a$. Inversion sends roots to
roots and changes this value's sign, since logarithms vanish at
roots of unity. The common value must consequently be zero.
\end{proof}

\begin{corollary}\label{d3:cor:directtcwildorbit}
Using the purely analytic wild-orbit independence statement of
Proposition~\ref{d3:prop:wildorbit}, choose odd representatives
$k_j\equiv(1+p)^j\pmod{p^r}$ for $0\le j<p^{r-1}$ at odd $p$,
and $k_j\equiv5^j\pmod{2^r}$ for $0\le j<2^{r-2}$ at two.
The corresponding actual classes $c_{k_j}$ are rationally independent
in $K_3(A_{p,r})$ and have integral generic image zero.
The same holds over the polynomial quotients with coefficient ring
$\mathbf Z_{(p)}$, localized at the ideal generated by $p$ and the
displayed formal parameters, and over their henselizations.
\end{corollary}
\begin{proof}
Apply Proposition~\ref{d3:prop:directtcendpoints} at one fixed primitive
root and then the stated analytic independence. A rational relation
in a preceding local ring or its henselization would map to one in
the completion. The generic-zero Dennis--Stein relation holds already
in those local rings.
\end{proof}


\section{Finite mixed HKR and classes in every higher degree}
\label{d3:sec:allhigherdegreelaurent}

\begin{lemma}
\label{d3:lem:finitehkrgeneral}
Let $W$ be a complete mixed-characteristic DVR with perfect residue
field, and let $B$ be the $p$-adic completion of a smooth affine
$W$-algebra $B_0$ of relative dimension $D<\infty$.
The continuous cyclic, negative-cyclic, and periodic complexes of
$B/W$, completed before inverting $p$, are computed by the finite
mixed complex of continuous forms $(\Omega_B^{\leq D},0,d)$.
In particular the highest form component of $HC_D^{\mathrm{cts}}$
is $\Omega_B^D/d\Omega_B^{D-1}$.
For $v\in1+\pi B$ and $b_1,\ldots,b_D\in B^\times$, the character
of Lemma~\ref{d3:lem:relativelogproduct} satisfies
\[
 \chi_B\bigl([v]_{\mathrm{rel}}[b_1]\cdots[b_D]\bigr)_{\mathrm{top}}
 =\pm\log(v)\,d\log b_1\wedge\cdots\wedge d\log b_D.
\]
The sign is the same norm and suspension convention as before.
\end{lemma}

\begin{proof}
An integral comparison suffices to control all completion operations.
In Hochschild degree $j\leq D$ use the map
\[
 a_0\otimes\cdots\otimes a_j\longmapsto
       \frac{D!}{j!}a_0\,da_1\wedge\cdots\wedge da_j,
\]
and use zero in higher degrees. It intertwines $b$ with zero and
the Connes operator with $d$. Smooth HKR identifies its map on every
nonzero Hochschild homology group with multiplication by $D!$.
The cone thus has homology in degrees $0,\ldots,D$, all killed by
$p^a$, where $a=v_p(D!)$; flatness of the form modules is used here.

This error can be killed before taking any circle construction.
Its Postnikov tower lies in the category
$\operatorname{Fun}(BS^1,D(W))$ of complexes with circle action.
Each graded piece has trivial circle action: the automorphism
space of a module in the heart is discrete, while $BS^1$ is simply
connected. Hence $p^a$ kills that piece in the equivariant category.
In a triangle $A\to M\to C$, if $p^r$ kills $A$ and $p^s$ kills
$C$, then $p^s\mathrm{id}_M$ factors through $A$, so $p^{r+s}$
kills $M$. Induction on the finite Postnikov tower shows that
$p^{a(D+1)}$ kills the whole equivariant cone.
The exact orbit, fixed-point, Tate, and derived completion functors
preserve this null map. Their comparison errors therefore vanish
after inverting $p$.

In every total degree the cyclic complexes of the finite form
model contain only finitely many form modules. These modules are
$W$-flat and their reduction towers are surjective. Derived
completion is consequently termwise completion, also in the
unbounded total complexes: the cone-of-$1-\mathrm{shift}$ model
for the inverse limit is degreewise exact for these surjections.
Proposition~4.2 of \cite{AMMN}, including its relative version,
identifies this with the asserted continuous invariant. No exchange
of rationalization with an uncontrolled infinite product is used.

The module comparison and the one-variable normalization of a
relative unit in Lemma~\ref{d3:lem:relativelogproduct} do not depend
on the relative dimension. Its ordinary cyclic representative
$\log v$ starts in column zero. At every unit multiplication the
positive columns and the extra $t$ in the cyclic-shuffle correction
vanish in $C^+$. Only the ordinary shuffles of the lowest Dennis
unit traces remain. Their $D!$ permutations cancel the $1/D!$
normalized HKR factor, proving the formula.
\end{proof}

\subsection{The comparison for completed essentially smooth models}

\begin{lemma}
\label{d3:lem:essentiallysmoothAMMN}
Let $W_0=W(\overline{\mathbf F}_p)$ and let $W$
be a finite totally ramified extension, with uniformizer $\pi$
and fraction field $E$. Let $B_0$ be a localization of a smooth
$W$-algebra, with $\Omega^1_{B_0/W}$ finite projective of constant
rank $D$, and put $B=\widehat{B_0}_p$.
There is a natural actual relative character
\[
 \chi_B:K(B,\pi B)\longrightarrow
                 \Sigma HC^{\mathrm{cts}}(B/W;\mathbf Q_p),
\]
which factors through an equivalence from
$K(B,\pi B)^{\wedge}_p[1/p]$. It is $K(B)$-linear through the negative-cyclic
character, and the target has the finite continuous mixed-HKR
model, with ordinary de Rham quotients.
\end{lemma}

\begin{proof}
We justify the extension beyond topological finite presentation.
For each $r$, the ring $W_0/p^r$ is the filtered union of the
finite \emph{\'{e}tale} $\mathbf Z/p^r$-algebras
$W_r(\mathbf F_{p^m})$. Flat base change for cotangent complexes
therefore gives $\widehat L_{W_0/\mathbf Z,p}=0$.
Write $W=W_0[\pi]$ and let $f$ be the Eisenstein polynomial of
$\pi$. Then
\[
 \widehat L_{W/\mathbf Z,p}
   \simeq[W\xrightarrow{f'(\pi)}W]
\]
in cohomological degrees $-1,0$. Choose $a$ with
$p^a\in f'(\pi)W$; multiplication by $p^a$ on this complex is
null-homotopic.

Cotangent transitivity, after derived $p$-completion, gives
\[
 M\longrightarrow\widehat L_{B_0/\mathbf Z,p}
       \longrightarrow P,
 \qquad M=[B\xrightarrow{f'(\pi)}B],\quad
 P=\widehat\Omega^1_{B_0/W}.
\]
The target $P$ is finite projective in degree zero, so this
triangle splits as a triangle of underlying $B$-module
complexes. The error in the $i$th derived exterior power of its
projection is a sum of terms
$(\mathbf L\Lambda^rM)\otimes\Lambda^{i-r}P$, $1\le r\le i$.
The $r$th such term is killed by $p^{ar}$ as an object: derived
exterior power sends the null-homotopy of $p^a\operatorname{id}_M$
to one for $p^{ar}$.

The connective HKR filtration on Hochschild homology has graded
terms $(\mathbf L\Lambda^iL)[i]$. Any fixed homotopy range
involves only finitely many of these terms, even after completion.
Indeed derived $p$-completion preserves connectivity: for a
$c$-connective spectrum $X$, its only possible new bottom group
is $\lim^1_r(\pi_cX/p^r)=0$, because the quotient tower is
surjective. Thus the increasingly connective HKR tails remain
increasingly connective. The preceding bounds
and a finite extension argument therefore make
\[
 HH(B_0;\mathbf Z_p)\longrightarrow HH(B_0/W;\mathbf Z_p)
\]
a quasi-isogeny in each finite range. Its rationalization is an
equivalence. The bounds use the base different and the range;
they are unaffected by the localization set. No interchange of
an infinite localization colimit with completion is used.

To pass to cyclic homology, we control the comparison cone in the
equivariant derived category. Let $C$ be the cone of the map of completed Hochschild
complexes. For every integer $N$, the preceding argument gives an
integer $a_N$ killing each of the finitely many nonzero homotopy
modules of $\tau_{\leq N}C$. The pointwise $t$-structure on
$\operatorname{Fun}(BS^1,D(\mathbf Z_p))$ gives its equivariant
Postnikov tower. Each heart object has trivial circle action:
the automorphism space of a module in the heart is discrete and
$BS^1$ is simply connected. A finite extension argument therefore
gives an integer $b_N$ such that $p^{b_N}$ annihilates
$\tau_{\leq N}C$ as an equivariant object, not merely on its
underlying homotopy groups.

Taking homotopy orbits preserves this scalar null-homotopy.
The omitted Postnikov tail remains $(N+1)$-connective after
homotopy orbits and derived $p$-completion. Indeed, homotopy
orbits of a connective spectrum are connective, and in derived
$p$-completion the potentially new bottom $\lim^1$ term is zero
because the bottom coefficient tower is the surjective tower
$\pi_{N+1}(-)/p^r$. Thus the cone of the corresponding map of
completed cyclic complexes has bounded $p$-torsion in every fixed
degree, and vanishes after inverting $p$.

Equivalently, one can take a sufficiently high finite bar skeleton
for the circle-orbit construction in each degree. Its building blocks
are finite colimits of $(S^1)^j_+\wedge X$, so the comparison commutes
with derived completion. This argument requires no uniform bound
$b_N$ as $N$ varies and no interchange of rationalization with an
infinite product. The norm fiber sequence then makes the square
between absolute and relative $HC^-\to HP$ cartesian. It does not
require that the two vertical maps of that square be equivalences
individually.

By \cite[Proposition~4.2]{AMMN}, including the relative variant
stated just before Proposition~4.10, both $B_0$ and $B$ compute
the same continuous $p$-completed HH and cyclic theories.
The ordinary ring $B$ is henselian along $(p)$, so
\cite[Theorem~A]{AMMN} applies. The map $B/p\to B/\pi$ has
nilpotent kernel, and \cite[Corollary~3.8]{AMMN} identifies its
K-theory and TC after rationalized $p$-completion. Composing the
actual-to-completed map on relative K-theory with the resulting
fiber equivalence gives $\chi_B$.

The module argument is unchanged: for $H=THH(B)$, the objects
$H\otimes D$ and maps $\operatorname{id}_H\otimes f$ in
\cite[Proposition~2.8, Corollary~2.10, and Theorem~2.12]{AMMN}
form a diagram of $H$-modules in cyclotomic spectra.
The lax monoidal TC and circle functors give module diagrams,
their equivalences can be inverted there, and the relative HH
projection preserves the action. Together with the
multiplicative cyclotomic trace this gives the asserted actual
$K(B)$-linearity. This construction has no finite presentation
requirement.

Finally, the integral mixed HKR map in degree $j\le D$ has
coefficient $D!/j!$ and induces multiplication by $D!$ on
$HH_j(B_0/W)=\Omega^j_{B_0/W}$. Its equivariant cone has a
Postnikov tower of length $D+1$ whose graded objects are killed
by $p^{v_p(D!)}$. These heart objects have trivial circle action;
the extension argument kills the entire equivariant cone by
$p^{(D+1)v_p(D!)}$. Thus circle orbits, fixed points, Tate
constructions, and derived completion retain only bounded
torsion error. The rational mixed-HKR comparison follows.
Each total degree has finitely many flat form modules, with
surjective reduction towers, so completion is termwise and
cohomology uses the ordinary image of $d$.
\end{proof}

\subsection{Continuous residues and higher unit products}

\begin{lemma}
\label{d3:lem:continuouslaurentresidue}
Let $W$ be a complete DVR of mixed characteristic with uniformizer
$\pi$ and fraction field $E$. Let $B$ be a smooth formal $W$-algebra
of topologically finite type. Take a nonzero
$g\in W[T_1,\ldots,T_s]$ whose nonzero coefficients are units, and set
\[
 B_g=\bigl(B[T_1^{\pm1},\ldots,T_s^{\pm1},1/g]\bigr)^{\wedge}_{\pi}.
\]
There are $E$-linear maps on continuous forms, with base forms
placed before the logarithmic differentials,
\[
 \operatorname{Res}_s:
 \Omega^{q+s}_{B_g[1/p]/E,\mathrm{cts}}
       \longrightarrow\Omega^q_{B[1/p]/E,\mathrm{cts}},
\]
such that
\[
 \operatorname{Res}_s(d\eta)=d\operatorname{Res}_s(\eta),\qquad
 \operatorname{Res}_s\left(
 \omega\wedge d\log T_1\wedge\cdots\wedge d\log T_s\right)=\omega.
\]
Consequently wedging with these logarithmic differentials gives
a split injection on ordinary continuous de Rham cohomology.
There is also a degree-zero projection $P_0$ which takes the
constant coefficient of the summand containing no $d\log T_i$.
It commutes with $d$ and is the identity on continuous base forms.
\end{lemma}

\begin{proof}
Put $R_0=W[T_1^{\pm1},\ldots,T_s^{\pm1},1/g]$.
There is an integral formal expansion
\[
 R_0\longrightarrow W((T_1))\cdots((T_s)).
\]
To construct it, write $g=T_s^a(g_a+T_sg_{a+1}+\cdots)$ with
$a$ minimal. By induction $g_a$ is a unit in the preceding
iterated Laurent ring: the recursion ends at one of the unit
coefficients of $g$. Expand the remaining inverse as a formal
geometric series in $T_s$. This construction divides by no power
of $\pi$.
Taking the coefficient of $T_s^0$, then $T_{s-1}^0$, and so on,
defines a $W$-linear map
\[
 \operatorname{CT}:R_0\longrightarrow W,
 \qquad \operatorname{CT}(1)=1,\quad
 \operatorname{CT}(T_i\partial_{T_i}h)=0.
\]
It is uniformly $\pi$-adically continuous, since
$\operatorname{CT}(\pi^aR_0)\subseteq\pi^aW$ for every $a$.
The $B$-linear extension therefore extends uniquely to $B_g\to B$.
The same construction works for every finite projective module
of continuous base forms, by taking a direct summand of a finite
free module.

Decompose forms into base forms and the ordered $d\log T_i$.
Retain only the summand containing all $s$ logarithmic factors,
and apply this extended constant-term map to its coefficient.
This defines $\operatorname{Res}_s$ integrally before inverting $p$.
The base differential commutes with constant-term extraction;
every term from a Laurent-variable derivative has zero constant
term. The two displayed identities follow, first on the dense
uncompleted modules and then by continuity. They show directly
that boundaries map to boundaries and that the stated composite
is the identity on forms and cohomology.
For $P_0$, a Laurent-variable derivative always introduces a
logarithmic factor and is discarded, proving its chain identity
by the same calculation.

No embedding of the completed ring into the ordinary Laurent ring
is asserted. For example $\sum_{a\geq0}\pi^aT_s^{-a}$ has no lower
bound on its Laurent exponents. If a completed series target is
desired, it is
$\varprojlim_a(B/\pi^a)((T_1))\cdots((T_s))$.
Nor need the formal expansion of $1/g$ converge analytically
when every $T_i$ has norm one.
\end{proof}


\begin{lemma}
\label{d3:lem:laurentspecialvanishing}
Let $X$ be a smooth variety of dimension $d$ over a perfect field
of characteristic $p$. For every $j>d$,
\[
 \pi_jK(X)^{\wedge}_p=\pi_jK(X;\Q_p)=0.
\]
In particular, if $s\geq1$, $0\ne\overline g\in
\F_3[T_1^{\pm1},\ldots,T_s^{\pm1}]$, and
\[
 C_s=\Spec\!\left(
 \F_3[X,Y,Z,T_1^{\pm1},\ldots,T_s^{\pm1},1/\overline g]
                         /(XY+Z^2-1)\right),
\]
then
\[
 \pi_{s+3}K(C_s;\Q_3)=\pi_{s+4}K(C_s;\Q_3)=0.
\]
\end{lemma}
\begin{proof}
Geisser--Levine prove
\[
 K_j(X;\Z/p^a)=0\qquad(a\geq1,\ j>d)
\]
in \cite[Theorem~8.4, p.~491]{GeisserLevineCharP}.
These are the finite-coefficient $K$-groups, with the usual
coefficient long exact sequence; see their Section~2, p.~464.
In particular the assertion includes the torsion contribution
from the preceding actual $K$-group.
The homotopy-limit exact sequence is
\[
 0\longrightarrow\varprojlim_a^{1}K_{j+1}(X;\Z/p^a)
 \longrightarrow\pi_jK(X)^{\wedge}_p
 \longrightarrow\varprojlim_a K_j(X;\Z/p^a)
 \longrightarrow0.
\]
Both outside terms vanish when $j>d$. Inverting $p$ proves the
first assertion. The displayed $C_s$ is smooth of dimension
$s+2$, so its two asserted groups lie in this range.
\end{proof}


\begin{theorem}
\label{d3:thm:everydegreelaurent}
Put
\[
 A_0=\left(\Z_{(3)}[x,y,z]/(3+xy+z^2)\right)_{(x,y,z)},
 \qquad \mathfrak m_0=(x,y,z)A_0.
\]
For $s\geq0$, let
\[
 R_s=A_0[T_1,\ldots,T_s]_{\mathfrak m_0A_0[T_1,\ldots,T_s]}.
\]
This is a regular local ring of dimension three and residue field
$\F_3(T_1,\ldots,T_s)$. The actual class
\[
 c_s=2\left\{\frac{1+z}{2},\,1+x,\,
       1+\frac{(1+x)y}{4},\,T_1,\ldots,T_s\right\}
                    \in K_{s+3}(R_s)
\]
has infinite order and restricts to zero in $K_{s+3}(R_s[1/x])$
integrally. Consequently rational Gersten injectivity fails in
each degree $j\geq3$ on an explicitly specified ring of dimension
three.
\end{theorem}

\begin{proof}
Polynomial extension and localization preserve regularity; the
chosen prime has height three and the displayed residue field.
Every $T_i$ is a unit. The Dennis--Stein calculation of
Proposition~\ref{d3:prop:directtcoriginal}, followed by multiplication
by these units, proves integral vanishing after inverting $x$.
The case $s=0$ has already been proved.

Suppose that $N c_s=0$ for some positive integer $N$ and $s\geq1$.
Continuity of actual $K$-theory makes this relation hold over
$A_0[T_i^{\pm1},1/f]$ for one finite denominator with
$\bar f\ne0$ in $\F_3[T_1,\ldots,T_s]$.
Use $W=\Z_3[u]/(u^2+3)$ and the original good quadric
$B=W\langle X,Y,Z\rangle/(XY+Z^2-1)$.
Choose $g\in W[T_1,\ldots,T_s]$ by lifting exactly the nonzero
coefficients of $\bar f$ to units, and form $B_g$ as in
Lemma~\ref{d3:lem:continuouslaurentresidue}.
Under $(x,y,z)=(uX,uY,uZ)$, the image of $f$ is congruent to $g$
modulo $u$. Thus $f/g\in1+uB_g$ is a unit, and the finite
relation has an actual image in $K_{s+3}(B_g)$.

The specified image has the actual relative lift
\[
 \xi_s=[a^2]_{\mathrm{rel}}[\lambda][\mu]
                                  [T_1]\cdots[T_s].
\]
The formal algebra $B_g$ has relative dimension $D=s+2$ over $W$.
Lemma~\ref{d3:lem:finitehkrgeneral} gives its top character as
\[
 \chi_{B_g}(\xi_s)_{\mathrm{top}}
       =\pm\tau_c\wedge d\log T_1\wedge\cdots\wedge d\log T_s,
\]
where $[\tau_c]\ne0$ was calculated in
Proposition~\ref{d3:prop:directtcoriginal}.
The residue of Lemma~\ref{d3:lem:continuouslaurentresidue} sends this
class back to $\pm[\tau_c]$, so it is nonzero in the ordinary
quotient by exact forms.

By Lemma~\ref{d3:lem:laurentspecialvanishing},
$\pi_{s+4}K(B_g/u;\Q_3)=0$. The long exact sequence of the
completed relative spectrum therefore makes
\[
 \pi_{s+3}K(B_g,uB_g;\Q_3)
       \longrightarrow\pi_{s+3}K(B_g;\Q_3)
\]
injective. The nonzero character of $\xi_s$ shows that its absolute
image has infinite order, contrary to the image of $Nc_s=0$.
This proves the assertion for the actual class in $R_s$.
Only the auxiliary special fiber is subject to the finite-coefficient
vanishing theorem; its perfect ground field is $\F_3$ even though
the source ring has the stated imperfect residue field.
\end{proof}


\section{Coefficient products and infinite rank}
The continuous cyclic character also detects products of actual point
classes with a ruling class. This gives an independent construction
of the coefficient classes used in the rank arguments.

\subsection{Coefficient products in cyclic homology}
\label{d3:sec:directtccoefficients}

Write $K(R;\mathbf Q_p)=K(R)^{\wedge}_p[1/p]$; this notation denotes
a rationalized completed spectrum. Let $E/\mathbf Q_p$ be finite, let
$W=\mathcal O_E$, and let $k$ be its residue field. By
\cite[Proposition~4.10 and Example~4.3]{AMMN}, the point has a natural
relative character
\[
 r_W:K_3(W)\longrightarrow\pi_3K(W;\mathbf Q_p)
 \xrightarrow{\sim}HC_2^{\mathrm{cts}}(W/W;\mathbf Q_p)=E.
\]
The isomorphism uses
$\pi_3K(k;\mathbf Q_p)=\pi_4K(k;\mathbf Q_p)=0$.

\begin{lemma}
\label{d3:lem:cherncompletionfactor}
Let $E/\Q_p$ be finite, let $W=\mathcal O_E$, and let
$n\geq2$ and $s=2n-1$. Write
\[
 K(R)^{\wedge}_p=\operatorname*{holim}_{\nu}K(R)/p^\nu,
 \qquad K(R;\Q_p)=K(R)^{\wedge}_p[1/p].
\]
There is a natural $\Q_p$-linear map
\[
 \widetilde c_{n,E}:\pi_sK(W;\Q_p)
             \longrightarrow H^1(E,\Q_p(n))
\]
such that the usual \'etale Chern class on actual $K$-theory is
the composite
\[
 K_s(W)_{\Q}\longrightarrow\pi_sK(W;\Q_p)
       \xrightarrow{\widetilde c_{n,E}}H^1(E,\Q_p(n)).
\]
These maps are natural under finite extensions of $p$-adic fields.
Consequently, actual classes with $\Q_p$-linearly independent
\'etale Chern values have $\Q_p$-linearly independent images in
$\pi_sK(W;\Q_p)$.
\end{lemma}

\begin{proof}
First construct the factor over the field $E$, where $p$ is invertible.
For every $\nu\geq1$ the finite-coefficient Chern class is an
additive map
\[
 c^{(\nu)}_{n,1}:K_s(E;\Z/p^\nu)
       \longrightarrow H^1(E,\mu_{p^\nu}^{\otimes n}).
\]
This includes $p=2$: the exceptional $K$-degree in
\cite[Proposition~2.4, p.~10]{WeibelChern} is two, whereas $s\geq3$.
The maps are compatible with $p^{\nu+1}\to p^\nu$ reduction.
Indeed, take $q=p$ and $r=p^\nu$ in the middle square of
\cite[Proposition~2.8, p.~12]{WeibelChern}:
\[
\begin{array}{ccc}
 K_s(E;\Z/p^{\nu+1})&\longrightarrow&K_s(E;\Z/p^\nu)\\
 {\scriptstyle c^{(\nu+1)}_{n,1}}\downarrow&&
             \downarrow{\scriptstyle c^{(\nu)}_{n,1}}\\
 H^1(E,\mu_{p^{\nu+1}}^{\otimes n})&\longrightarrow&
                         H^1(E,\mu_{p^\nu}^{\otimes n}).
\end{array}
\]
The homotopy-limit projections therefore define a homomorphism
\begin{equation}\label{d3:eq:cherncompletionintegral}
 \pi_sK(E)^{\wedge}_p
 \longrightarrow\varprojlim_{\nu} K_s(E;\Z/p^\nu)
 \xrightarrow{\varprojlim c^{(\nu)}_{n,1}}
 \varprojlim_{\nu}H^1(E,\mu_{p^\nu}^{\otimes n}).
\end{equation}
No assertion about a source $\varprojlim^{1}$ term is required:
\eqref{d3:eq:cherncompletionintegral} uses only the canonical projections.
On the target, the groups $H^0(E,\mu_{p^\nu}^{\otimes n})$ are
finite. Their inverse system satisfies the Mittag--Leffler condition,
so \cite[equation~(0.2), p.~208]{Jannsen} gives
\[
 \varprojlim_{\nu}H^1(E,\mu_{p^\nu}^{\otimes n})
       \cong H^1_{\mathrm{cont}}(E,\Z_p(n)).
\]

The resulting map is $\Z_p$-linear. For clarity, the source homotopy
group has its natural $\Z_p=\pi_0\mathbb S^{\wedge}_p$ action.
Each composite to the $\nu$th cohomology group is additive and has
target annihilated by $p^\nu$. If $\lambda\in\Z_p$ and
$\lambda=a+p^\nu b$ with $a\in\Z$ and $b\in\Z_p$, that composite
sends $\lambda x$ to $a$ times the image of $x$, which is also
$\lambda$ times that image. Passing to the inverse limit proves
linearity. Inverting $p$ and using the definition of rational
continuous cohomology \cite[equation~(0.4)]{Jannsen} gives
\[
 \overline c_{n,E}:\pi_sK(E;\Q_p)
                 \longrightarrow H^1(E,\Q_p(n)).
\]
Define $\widetilde c_{n,E}$ by preceding this map with the natural
map $\pi_sK(W;\Q_p)\to\pi_sK(E;\Q_p)$.

By \cite[Lemma~2.3, p.~10]{WeibelChern}, at each finite level the
usual Chern class on $K_s(E)$ is the finite-coefficient class after
reduction. The target inverse-limit identification then gives the
asserted equality on actual classes. Naturality follows from
\cite[Proposition~2.1.1]{WeibelChern} and the naturality of every
map above. Finally, applying the $\Q_p$-linear map
$\widetilde c_{n,E}$ to a relation proves the independence assertion.
\end{proof}

\begin{remark}
This lemma does not assert that actual $K_s(W)$ surjects onto the
homotopy group of its completed spectrum. It also does not identify
any cyclic or syntomic coordinate with a polylogarithm. If an
independently constructed $\Q_p$-linear isomorphism
$\pi_sK(W;\Q_p)\cong E$ is available, the lemma only transports
independence to that coordinate.
\end{remark}

\begin{lemma}\label{d3:lem:directtcactualcoefficients}
The values of $r_W$ on actual classes span $E$ over $\mathbf Q_p$.
For a quadratic extension $E/F$, actual integral anti-invariant
classes span $E^{-}$ under this map.
\end{lemma}
\begin{proof}
Let $d=[E:\mathbf Q_p]$ and put
$V=\pi_3K(\mathcal O_E;\mathbf Q_p)$.
The point comparison of \cite[Proposition~4.10]{AMMN} identifies $V$ with $E$, so
$\dim_{\mathbf Q_p}V=d$. Independently, the Bloch--Kato exponential
identifies $H^1(E,\mathbf Q_p(2))$ with $E$, also of dimension $d$.
The finite-coefficient Chern construction gives a natural linear map
\[
 \widetilde c_{2,E}:V\longrightarrow H^1(E,\mathbf Q_p(2))
\]
through which the Chern map on actual $K_3(\mathcal O_E)$ factors.
This is the case $n=2$ of Lemma~\ref{d3:lem:cherncompletionfactor};
its proof uses only finite-coefficient Chern additivity,
coefficient compatibility, and the canonical map from the homotopy
limit to the limit of homotopy groups.

By Lemma~\ref{d3:lem:actuallocal}, whose proof uses Levine's
finite-coefficient theorem, Milnor $K_3$-divisibility, the
Moore--Merkurjev description of $K_2(E)$, and localization,
the actual Chern values span $H^1(E,\mathbf Q_p(2))$.
Let $V_{\rm act}$ be the span in $V$ of the images of actual classes.
Then $\widetilde c_{2,E}(V_{\rm act})=H^1(E,\mathbf Q_p(2))$.
Both ambient vector spaces have dimension $d$, whence
$V_{\rm act}=V$ and $\widetilde c_{2,E}$ is an isomorphism.
Thus the actual AMMN point values span $E$.

All maps are natural for finite extensions of local fields.
For a quadratic $E/F$, apply $1-\sigma$ to actual integral classes.
Their images span $E^-$ since $(1-\sigma)E=E^-$ as
$\mathbf Q_p$-vector spaces. This argument divides by no integer
inside actual integral $K$-theory, and applies equally at $p=2$.
\end{proof}

\begin{proposition}\label{d3:prop:directtccoefficientcup}
Let $0\ne d\in\mathfrak m_W$, and set
\[
 T=\bigl(W[[x,y,w]]/(xy+w(w+d))\bigr)[1/p],\qquad J=(x,w)\mathcal O_T.
\]
If $\beta\in K_3(W)$ satisfies $r_W(\beta)\ne0$, then
$\beta(1-[J])$ has infinite order in $K_3(T)$. Actual coefficients
with $\mathbf Q_p$-independent $r_W$-values give rationally independent
products.
\end{proposition}
\begin{proof}
Put
\[
 B=W\langle X,Y,Z\rangle/(XY+Z(Z+1)),\qquad I=(X,Z),\qquad C=B/\pi.
\]
The substitution $(x,y,w)=(dX,dY,dZ)$ converges for every formal
series and defines a ring map to $B$. On the generic fiber it takes
$J$ to a line bundle isomorphic to $I$. The special fiber $C$ is an
affine-line torsor over $\mathbf P^1_k$, so
$K(C)\simeq K(k)\oplus K(k)$ and
$\pi_3K(C;\mathbf Q_p)=\pi_4K(C;\mathbf Q_p)=0$.

Let $b=r_W(\beta)$ and $\eta=1-[I]$. The module compatibility of the
AMMN fiber character, verified in Lemma~\ref{d3:lem:relativelogproduct},
identifies the relative image of $\beta\eta$ with the action of
$\operatorname{ch}^{-}(\eta)$ on $t^{-1}b\in HC_2$, where $|t|=-2$.
Write a normalized representative of the former as
$c_0+tc_2+t^2c_4+\cdots$. Its rank term $c_0$ is zero. In the
shuffle-plus-cyclic-shuffle product
\cite[Section~2, Remark~2.2]{BokstedtOttosen}, every cyclic-shuffle
correction puts $b\in E$ in a normalized bar slot and hence is zero.
The remaining product in ordinary cyclic homology is exactly $bc_2$:
terms of positive $t$-degree vanish. Its two-form component is
\[
 -\lambda_E b\,c_1^{\mathrm{dR}}(I),\qquad\lambda_E\in E^\times,
\]
where the nonzero normalization follows from
\cite[Proposition~4.12]{AMMN}. In weight one, the map
$HC^-_0\to HP_0$ induces
$H^2(\Omega_B^{\geq1})\to H^2(\Omega_B^\bullet)$; both groups
are $\Omega_B^2/d\Omega_B^1$, so this is an isomorphism.
Consequently the periodic Chern component determines precisely the
two-form component used above. Thus no comparison with an \'etale
or syntomic regulator, and no additional exact normalization of
$K_0\to HC^{-}$, is needed.

The class $c_1^{\mathrm{dR}}(I)$ is represented, up to sign, by the
globally regular form
\[
 \omega_0=\frac{dX}{X}\wedge dZ
 =\operatorname{Tr}(P\,dP\wedge dP),\qquad
 P=\begin{pmatrix}-Z&X\\Y&Z+1\end{pmatrix}.
\]
To verify its algebraic nonvanishing, use the two root-deletion
bidiscs with $Z=Xt_0$ and $Z=-1+Xt_{-1}$. Their intersection is
$E\langle X,X^{-1},Z\rangle$. Local primitives are
\[
 q_0=-Z\,dX/X,\qquad q_{-1}=-(Z+1)dX/X,
 \qquad q_{-1}-q_0=-dX/X.
\]
If a global primitive existed, subtract it from these primitives,
restrict to $Z=0$, and take the Laurent residue in $X$. A closed
bidisc form $A(X,t)dX+B(X,t)dt$ has zero residue on $t=c/X$ for
every $|c|\le1$: its residue is
\[
 \sum_{m\ge0}(a_{m,m+1}-b_{m+1,m})c^{m+1}=0,
\]
by closedness and convergence of the Tate coefficients. Here the
two values of $c$ are $0$ and $1$, so the argument also applies at
$p=2$. The displayed difference has residue $-1$, a contradiction.
This also proves that $b\mapsto b[\omega_0]$ is injective.

The fiber square therefore detects $\beta\eta$ nontrivially in
$\pi_3K(B;\mathbf Q_p)$. Finally, localization preserves its rational
nonvanishing in $K_3(B[1/p])$: its possible kernel comes from
$K_3(C)$, a finite group. This is the image of $\beta(1-[J])$.
The argument is linear in the coefficient and proves the independence
assertion as well.
\end{proof}

\begin{remark}
The target in this proof is the ordinary de Rham quotient
$\Omega^2/d\Omega^1$ of the termwise completed complex, as in
\cite[Constructions~4.6 and~4.11]{AMMN}. It is not the quotient by
the closure of $d\Omega^1$. Indeed, with
$h_j(Z)=-(-Z)^{p^j}$, the regular primitives
$-(Z-h_j(Z))dX/X$ have derivatives converging to $\omega_0$, although
$\omega_0$ is not exact. This proof gives nonzero images of actual
$K$-classes in a rationalized completed spectrum; it does not assert
a prescribed order in every finite Artin approximation.
\end{remark}


\subsection{Infinite rank over a strict henselian base at every prime}
\label{d3:sec:directtcstrictallprimes}

We first extend the coefficient-cup calculation to the infinite
unramified base. Throughout this subsection, $K(-;\mathbf Q_p)$ means
the rationalization of the $p$-completed spectrum.

\begin{lemma}\label{d3:lem:directtcperfectcoeffcup}
Let $E$ be a complete discretely valued field of mixed characteristic
$(0,p)$ whose residue field $k$ is algebraic over $\mathbf F_p$, and
let $W=\mathcal O_E$. The AMMN point character
\[
 r_W:K_3(W)\longrightarrow\pi_3K(W;\mathbf Q_p)
       \xrightarrow{\sim}HC_2^{\mathrm{cts}}(W/W;\mathbf Q_p)=E
\]
is defined as before and is natural for continuous local maps of
these coefficient DVRs. Proposition~\ref{d3:prop:directtccoefficientcup}
holds for this $W$: any finite set of actual coefficients with
$\mathbf Q_p$-independent $r_W$-values gives rationally independent
products with the ruling class on the completed split node.
\end{lemma}
\begin{proof}
The standing hypotheses of \cite[Section~4.1]{AMMN} require a perfect
residue field, and Proposition~4.10 requires a smooth formal scheme;
neither requires a finite residue field. Example~4.3 applies to the
affine $p$-complete rings in question. Naturality when the coefficient
ring changes follows by keeping the canonical maps used in the proof
of Proposition~4.10: a map of pairs $(W,R)\to(W',R')$ gives a
commutative square
\[
\begin{matrix}
 HH(R;\mathbf Q_p)&\longrightarrow&HH(R/W;\mathbf Q_p)\\
 \downarrow&&\downarrow\\
 HH(R';\mathbf Q_p)&\longrightarrow&HH(R'/W';\mathbf Q_p).
\end{matrix}
\]
The absolute fiber square, cyclic constructions, and horizontal
fibers are natural as well. At the point, the cyclic scalar is sent
by $E\to E'$. Thus $r_{W'}(\beta|_{W'})=r_W(\beta)|_{E'}$.

Quillen's finite-field calculation and filtered-colimit continuity
show that $K_i(k)$ is uniquely $p$-divisible for $i>0$, and that
$K_3(k)$ is torsion. The coefficient exact sequence gives
$K(k)/p^a\simeq H\mathbf Z/p^a$ directly, hence
$K(k;\mathbf Z_p)\simeq H\mathbf Z_p$. In particular the two adjacent
positive homotopy groups needed to define the point fiber lift vanish.

For $B=W\langle X,Y,Z\rangle/(XY+Z(Z+1))$, its reduced special fiber
$C$ is an affine-line torsor over $\mathbf P^1_k$. Therefore
$K(C)\simeq K(k)\oplus K(k)$. The relative-to-absolute step in
Proposition~\ref{d3:prop:directtccoefficientcup} still has vanishing
$\pi_3$ and $\pi_4$ on this special fiber. Its ordinary de Rham
calculation and two-chart residue proof work over the complete field
$E$ verbatim. In particular $b\mapsto b[\omega_0]$ is injective,
where $\omega_0=(dX/X)\wedge dZ$ is globally regular. Finally
$K_3(C)$ is torsion, which is enough for localization to give
$K_3(B)_{\mathbf Q}\hookrightarrow K_3(B[1/p])_{\mathbf Q}$.
The module calculation and the nonzero scalar of
\cite[Proposition~4.12]{AMMN} therefore detect exactly the same
coefficient products. No Hausdorff quotient of de Rham cohomology
is taken.
\end{proof}

\begin{theorem}\label{d3:thm:directtcstrictinfiniteallp}
For every prime $p$, the complete strictly henselian regular local ring
\[
 A_{\infty,p}=W(\overline{\mathbf F}_p)[[x,y,z]]/(p+xy+z^2)
\]
has dimension three and infinite-dimensional rational generic
$K_3$-kernel. More precisely, for every $m\ge1$ there are $m$ actual
integral classes with zero generic image whose rational classes are
independent.
\end{theorem}
\begin{proof}
Put $V_m=W(\mathbf F_{p^m})$, $V_\infty=W(\overline{\mathbf F}_p)$,
$F_m=\operatorname{Frac}V_m$, and choose $u^2=-p$. Write
$W_m=V_m[u]$, $E_m=F_m(u)$, $W_\infty=V_\infty[u]$, and
$E_\infty=\operatorname{Frac}W_\infty$. Eisenstein's criterion makes
each $W_m$ and $W_\infty$ a complete DVR with uniformizer $u$.
The linear term $p$ of the equation in the regular ambient
power-series ring proves regularity and dimension three for
$A_{\infty,p}$. Its residue field is algebraically closed.

By Lemma~\ref{d3:lem:directtcactualcoefficients}, actual $r_{W_m}$-values
span $E_m$. Apply $1-\sigma$, where $\sigma(u)=-u$, and choose
$\beta_{m,1},\ldots,\beta_{m,m}\in K_3(W_m)$ with
$\sigma\beta_{m,j}=-\beta_{m,j}$ whose values are independent in
$E_m^-=uF_m$. This construction involves no division by two in
actual $K$-theory. Lemma~\ref{d3:lem:directtcperfectcoeffcup} carries
these values by the field inclusion into $E_\infty$, preserving
their $\mathbf Q_p$-independence.

Set $A_m=V_m[[x,y,z]]/(p+xy+z^2)$. Pull each coefficient to
$A_m/(x)=W_m[[y]]$ and push forward along $x=0$, obtaining
$c_{m,j}\in K_3(A_m)$. These vanish on $A_m[1/x]$ integrally.
The local map $A_m\to A_{\infty,p}$ is flat by miracle flatness:
both rings are regular of dimension three, and its closed fiber
is a field \cite[Tag~00R4]{StacksFlatness}. Thus their images
are the corresponding supported pushforwards on $A_{\infty,p}$.

After the finite free quadratic coefficient extension and inversion
of $p$, put
\[
 T=\bigl(W_\infty[[x,y,z]]/(xy+(z-u)(z+u))\bigr)[1/p],
 \qquad J_+=(x,z-u),\quad\eta_+=1-[J_+].
\]
The two components of $x=0$ are disjoint Cartier divisors. Their
structure-sheaf classes sum to zero because their union is the
principal divisor $x$. Flat base change and the projection formula
therefore give exactly
\[
 c_{m,j}|_T=\beta_{m,j}\eta_+
          +\sigma\beta_{m,j}\eta_-
          =2\beta_{m,j}\eta_+.
\]
Use $w=z-u$ and $d=2u$ in
Lemma~\ref{d3:lem:directtcperfectcoeffcup}. The substitution into the
good quadric is
$x=dX$, $y=dY$, $z=u+dZ$; all original formal variables have norm
at most $|u|<1$, including at $p=2$. The coefficient detector gives
$m$ independent images, so the original rational classes are
independent. Their generic images are zero by localization at $x$.
Since $m$ is arbitrary, the asserted kernel is infinite dimensional.
\end{proof}

\begin{corollary}\label{d3:cor:directtcstrictarithmeticallp}
For every prime $p$, the arithmetic strict henselization
\[
 S_p=\left(\left(\mathbf Z_{(p)}[x,y,z]/(p+xy+z^2)\right)
                         _{(x,y,z)}\right)^{sh}
\]
also has infinite-dimensional rational generic $K_3$-kernel,
represented by actual integral classes with integral generic vanishing.
\end{corollary}
\begin{proof}
Choose a finite \'etale local DVR $V_{m,0}/\mathbf Z_{(p)}$ with
number-field fraction field, residue $\mathbf F_{p^m}$, completion
$V_m$, and a compatible residue embedding. Set $W_{m,0}=V_{m,0}[u]$,
where $u^2=-p$. The fraction field of $W_{m,0}^h$ is the field of
constants in $E_m$, by Krasner's lemma.

The second assertion of \cite[Theorem~4.13]{LevineIndK3} and
$K_3^M(E_m)/p^a=0$ imply that
$K_3(\operatorname{Frac}W_{m,0}^h)$ surjects onto $K_3(E_m)/p^a$.
The Milnor vanishing follows from norm-residue
\cite[VI, Theorem~4.1]{Weibel} and local
$p$-cohomological dimension two, also at $p=2$
\cite[Chapter~I, discussion preceding Theorem~2.8]{MilneDuality}. Localization, using
$K_2(\mathbf F_{p^m})=0$, then gives
\[
 K_3(W_{m,0}^h)\twoheadrightarrow K_3(W_m)/p^a
 \qquad(a\ge1).
\]
Here the passage from $E_m$ back to $W_m$ does not change the finite
quotient because its kernel comes from prime-to-$p$ torsion
$K_3(\mathbf F_{p^m})$. The actual-completion argument of
Lemma~\ref{d3:lem:directtcactualcoefficients} now shows directly that
the henselian values of the AMMN point map span $E_m$. Applying
$1-\sigma$ supplies $m$ actual anti-invariant coefficients with
independent AMMN values. No comparison with an \'etale-selected
coefficient class is used.

Let $R_m=(A_{0,p}\otimes_{\mathbf Z_{(p)}}V_{m,0})^h$.
Its divisor quotient by $x$ receives $W_{m,0}^h$, its completion is
$A_m$, and it maps compatibly to $S_p$. Push the chosen coefficients
forward along $x=0$ and then map to $S_p$. Their images under
$S_p\to A_{\infty,p}$ are precisely the completed classes detected
in Theorem~\ref{d3:thm:directtcstrictinfiniteallp}; the map and completion
identification are supplied uniformly in $p$ by
Lemma~\ref{d3:lem:stricthenselcompletion}. Hence these $m$ rational
classes on $S_p$ are independent. They vanish after inverting $x$
integrally. Letting $m$ grow proves the corollary.
\end{proof}

\section{Higher point coordinates and unramified coefficients}
\label{d3:sec:higherpointcoordinates}


\begin{lemma}
\label{d3:lem:formalhigherTCcup}
Let $W$ be a complete mixed-characteristic discrete valuation ring
with uniformizer $\pi$, fraction field $E$, and residue field $k$
algebraic over $\mathbf F_p$. Let $B$ be a smooth affine formal
$W$-algebra, topologically of finite presentation and of relative
dimension two, and put $C=B/\pi B$.
Fix $n\geq2$ and suppose that
\[
 K_{2n}(C)=0.
\]
Let $\eta\in K_0(B)$ have rank zero. Write
\[
 \operatorname{ch}^-(\eta)=t\gamma,
 \qquad
 \gamma\in\widehat\Omega^2_{B/W}[1/p]
               /d\widehat\Omega^1_{B/W}[1/p],
 \qquad |t|=-2,
\]
for its negative-cyclic character in the finite mixed-HKR model,
and assume $\gamma\ne0$.

The AMMN point comparison gives a $\mathbf Q_p$-linear isomorphism
\[
 \pi_{2n-1}K(W;\mathbf Q_p)
       \xrightarrow{\ \sim\ }
 HC_{2n-2}^{\mathrm{cts}}(W/W;\mathbf Q_p)
       =t^{1-n}E.
\]
For an actual class $\beta\in K_{2n-1}(W)$, denote its value by
$t^{1-n}r_n(\beta)$. If $r_n(\beta)\ne0$, then the actual class
\[
                  \eta\,\beta_B\in K_{2n-1}(B)
\]
has infinite order. More generally, if a finite family of actual
classes $\beta_j$ has $\mathbf Q_p$-linearly independent values
$r_n(\beta_j)$, then the classes $\eta(\beta_j)_B$ are linearly
independent in $K_{2n-1}(B)\otimes_{\mathbf Z}\mathbf Q_p$.
\end{lemma}

\begin{proof}
All occurrences of $K(-;\mathbf Q_p)$ denote rationalized
spectral $p$-completion, and write
$\Omega_B^j=\widehat\Omega^j_{B/W}[1/p]$.
Since $k$ is algebraic over $\mathbf F_p$,
Quillen's finite-field computation and filtered colimits give
\[
 K_{2m}(k)=0\quad(m>0),
 \qquad K_{2m-1}(k)\text{ is prime-to-}p\text{ torsion}.
\]
The coefficient exact sequences also show that $K(k)^{\wedge}_p$
has no positive homotopy groups. The point isomorphism in the
statement therefore follows from the fiber square of
\cite[Proposition~4.10]{AMMN};
\cite[Example~4.3]{AMMN} identifies the continuous K-theory of
$\operatorname{Spf}W$ with the spectral completion used here.

First handle one actual class $\beta$. Its reduction in
$K_{2n-1}(k)$ need not be zero. Choose an integer $N$, prime to $p$,
which kills that reduction. The actual relative exact sequence,
together with $K_{2n}(k)=0$, gives a unique actual lift
\[
 \kappa\in K_{2n-1}(W,\pi W)
       \quad\text{of}\quad N\beta.
\]
By compatibility with completion, its relative AMMN character is
\[
          \chi_W(\kappa)=N t^{1-n}r_n(\beta).
\]
This constructs an actual relative class before using its
completed character; no completed relative lift is being treated
as an actual lift.

The AMMN fiber comparison is $K(B)$-linear, by the module
construction through
\cite[Proposition~2.8, Corollary~2.10, and Theorem~2.12]{AMMN},
the relative-to-$W$ projection in Proposition~4.10, and the
multiplicative cyclotomic trace
\cite[Theorem~7.4]{BGTTrace}.
More explicitly, with $H=THH(B)$, the objects $H\otimes D$ and
arrows $\mathrm{id}_H\otimes f$ in that construction are
$H$-modules and module maps in cyclotomic spectra. The lax
monoidal circle and TC functors preserve the induced $TC(B)$
actions, and the equivalences used there are inverted in the
module category. The relative Hochschild projection is also
$H$-linear. Thus this assertion uses the actual module maps in
the construction, rather than the naturality of an unspecified
additive equivalence.
Its target action is the usual action of negative cyclic homology
on the suspended ordinary cyclic fiber. The norm compatibility is
\cite[Proposition~2.3]{BokstedtOttosen}, with the spectral norm
identified with the Connes operator by
\cite[Theorem~2.1, Lemma~2.2, and Theorem~2.3]{HoyoisCircle}.

The finite mixed-HKR model in relative dimension two gives
\[
 HC^-_0(B/W;\mathbf Q_p)
    =H^0_{\mathrm{dR}}(B/W)\oplus tH^2_{\mathrm{dR}}(B/W),
\]
where $H^2_{\mathrm{dR}}=\Omega_B^2/d\Omega_B^1$ is the ordinary
quotient of continuous forms. Rank zero removes the first
component of $\operatorname{ch}^-(\eta)$, and higher form degrees
vanish. Thus this character is exactly $t\gamma$.

The pullback of $\chi_W(\kappa)$ is the scalar cyclic class
$Nt^{1-n}r_n(\beta)$ on $B$. Scalar multiplication has no cyclic
shuffle correction: in the normalized relative bar complex any
such correction contains a bar occupied by this scalar in $E$,
and hence is zero. Consequently the actual relative product
\[
       \xi=\eta\,\kappa_B\in K_{2n-1}(B,\pi B)
\]
has character
\begin{equation}\label{d3:eq:formalhigherTCcupvalue}
       \chi_B(\xi)=N t^{2-n}r_n(\beta)\gamma
       \quad\text{in }HC_{2n-2}^{\mathrm{cts}}(B/W;\mathbf Q_p).
\end{equation}
There is no constant component. For $n\geq2$, the displayed
two-form component is an ordinary de Rham cohomology summand;
since $r_n(\beta)\ne0$ and $E$ is a field, it is nonzero.
Thus $\xi$ has infinite order.

Now use the actual relative exact sequence
\[
 K_{2n}(C)=0\longrightarrow K_{2n-1}(B,\pi B)
                    \longrightarrow K_{2n-1}(B).
\]
It is injective at the middle term, and the absolute image of
$\xi$ is exactly $N\eta\beta_B$. Hence $\eta\beta_B$ has infinite
order. No inverse to multiplication by $N$ on actual K-groups is used.

For a finite family, choose one common prime-to-$p$ integer $N$
killing all reductions, and form the unique actual lifts
$\kappa_j$. Tensoring the last injection with the flat
$\mathbf Z$-module $\mathbf Q_p$ preserves injectivity.
Applying the $\mathbf Q_p$-linear extension of $\chi_B$ to a
relation between the relative products gives
\[
       N\left(\sum_j q_jr_n(\beta_j)\right)\gamma=0.
\]
The assumed independence forces every $q_j=0$.
\end{proof}

\begin{remark}
For the good all-prime quadric
\[
 B=W\langle X,Y,Z\rangle/(XY+Z(Z+1)),
 \qquad I=(X,Z),\qquad\eta=1-[I],
\]
the special fiber is an affine-line torsor over $\mathbf P^1_k$.
Homotopy invariance, descent, and the projective-bundle formula
give $K(C)\simeq K(k)\oplus K(k)$, so the required actual even
K-group vanishes. The idempotent
\[
 e=\begin{pmatrix}-Z&X\\Y&1+Z\end{pmatrix}
\]
has image $I$ and satisfies
\[
 \operatorname{Tr}(e\,de\wedge de)
                  =\frac{dX}{X}\wedge dZ.
\]
Thus $\gamma$ is, up to the consistent first-Chern sign, the
nonzero ruling form detected by the root-chart residue lemma.
The same formal coefficient calculation works in every fixed
degree $2n-1$: changing $n$ only changes the cyclic column.
For $n=1$, by contrast, the prospective term $t\gamma$ is killed
in ordinary cyclic homology, so this argument gives no
degree-one nonvanishing statement.
\end{remark}

\begin{remark}
Let $B_g$ be the completed Laurent test algebra obtained from a
smooth algebraic model of $B$ by adjoining
$T_1^{\pm1},\ldots,T_s^{\pm1}$ and inverting a polynomial $g$
in those variables with coefficients in $W$. Its relative
dimension is $D=s+2$. For the actual relative lift $\kappa$ above,
the same module calculation gives
\[
 \chi_{B_g}\!\left(\eta\,\kappa_{B_g}
                         [T_1]\cdots[T_s]\right)_{\mathrm{top}}
   =N t^{2-n}r_n(\beta)\,
           \gamma\wedge d\log T_1\wedge\cdots\wedge d\log T_s
\]
in the $t^{2-n}H^D_{\mathrm{dR}}$ component of
$HC_{2n+s-2}^{\mathrm{cts}}(B_g/W;\mathbf Q_p)$.

When $n\geq3$, the starting cyclic column is negative, so the
positive columns cannot simply be discarded on the basis of the
ordinary cyclic quotient. Instead use the following chain-degree
calculation. Write a negative-cyclic cycle for $\eta$ as
$\sum_{r\geq1}t^r e_{2r}$ and one for each unit as
$\sum_{a_i\geq0}t^{a_i}u_{2a_i+1}$.
The first sum starts at $r=1$: its degree-zero term represents the
rank in $HH_0(B_g/W)=B_g$, and $b:C_1\to C_0$ is zero for a
commutative algebra. Hence rank zero forces that term itself to
vanish.

The scalar point cup has zero cyclic-shuffle correction. After
the remaining $s$ unit multiplications, a term using indices
$r,a_1,\ldots,a_s$ and $c$ cyclic-shuffle corrections has
Hochschild degree
\[
                   2r+s+2\sum_i a_i+2c.
\]
The finite normalized HKR map in dimension $D=s+2$ kills this
term unless $r=1$, all $a_i=0$, and $c=0$. Thus only the first
Chern component, the lowest unit Dennis traces, and ordinary
shuffles survive. Their normalized image is exactly the
displayed wedge, with no additional degree-dependent factor.
This argument allows the unwanted terms to survive in the
ordinary cyclic complex; it kills them by their final form degree.

For a test polynomial $g$ chosen as the same-support unit lift of
its residual polynomial, the integral continuous Laurent residue
maps this wedge to $\gamma$. It is therefore nonzero whenever
$\gamma\ne0$. If the special fiber of $B_g$ is smooth of dimension
$D$ over $k$, Geisser--Levine's finite-coefficient vanishing
\cite[Theorem~8.4]{GeisserLevineCharP} implies
\[
 \pi_{2n+s}K(B_g/\pi;\mathbf Q_p)=0,
\]
since $2n+s>D$. The completed relative-to-absolute map is thus
injective in degree $2n+s-1$. In this Laurent extension one uses
this completed injection; no actual even K-group vanishing of
the higher-dimensional special fiber is asserted.
\end{remark}

\begin{remark}
The lemma does not assert surjectivity from actual K-theory to
spectral completion, and it does not identify $r_n$ with a
Bloch--Kato, syntomic, or polylogarithmic coordinate. Independently
known actual classes with independent \'etale Chern values may be
used once the separate factorization through spectral completion
has been established. The hypothesis on $k$ in this lemma is
algebraicity over $\mathbf F_p$; perfectness alone does not imply
the actual even K-group vanishing used in its proof.
\end{remark}


\subsection{Unramified actual coefficients at every prime}
\label{d3:sec:unramifiedhighercoefficients}

\begin{proposition}
\label{d3:prop:unramifiedhighercoefficients}
Let $p$ be any prime, $m\geq2$, and put
\[
 k_m=\F_{p^m},\qquad V_m=W(k_m),\qquad F_m=\Frac V_m.
\]
Take $u=i$ if $p=2$, and take $u$ to be a primitive $p$-th root
if $p>2$. Put $E_m=F_m(u)$, $O_m=V_m[u]$, and let $\sigma$
fix $F_m$ and invert $u$.
For each $n\geq2$ there exist Teichm\"uller roots
$\zeta_1,\ldots,\zeta_{m-1}\in\mu_{p^m-1}(V_m)\setminus\{1\}$
such that the actual rational classes
\[
 \beta_{n,j}=b_n(\zeta_ju)-b_n(\zeta_ju^{-1})
 \in K_{2n-1}(O_m)_{\Q},\qquad 1\leq j\leq m-1,
\]
are $\sigma$-anti-invariant and have $\Q_p$-linearly independent
images in $H^1(E_m,\Q_p(n))$ under $c_{n,1}$.
For each fixed $p,m,n$, a common integer multiple is represented
by strictly anti-invariant integral classes.
\end{proposition}
\begin{proof}
Put $\delta=u-u^{-1}\ne0$, $\epsilon=u^p=u^{-p}$, and
$F_{p,n}(z)=\operatorname{Li}_n(z)-p^{-n}\operatorname{Li}_n(z^p)$.
The equal $p$th powers make the depleted terms cancel exactly:
\[
 \operatorname{Li}_n(\zeta u)-\operatorname{Li}_n(\zeta u^{-1})
 =F_{p,n}(\zeta u)-F_{p,n}(\zeta u^{-1}).
\]
For $1\leq a\leq p-1$ the Laurent polynomial
\[
 Q_a(u)=\frac{u^a-u^{-a}}{u-u^{-1}}
       =\sum_{b=0}^{a-1}u^{a-1-2b}
\]
is integral and reduces to $a$.
Grouping denominators as $a+p\ell$ and using the binomial series
gives, with $S_j(t)=P_j(t)/(1-t)^{j+1}$ and $P_j\in\Z[t]$,
\begin{equation}\label{d3:eq:unramifiedhighercoordinate}
\begin{split}
 a_{n,p}(\zeta)
 &\coloneqq\frac{\operatorname{Li}_n(\zeta u)
                   -\operatorname{Li}_n(\zeta u^{-1})}{\delta}\\
 &=\sum_{j\geq0}(-1)^j\binom{n+j-1}{j}p^j
       \sum_{a=1}^{p-1}\frac{\zeta^aQ_a(u)}{a^{n+j}}
                              S_j(\epsilon\zeta^p)\in O_m.
\end{split}
\end{equation}
Here $S_0(t)=1/(1-t)$.
The series identity first holds for $|z|<1$ and extends uniformly
to $|z|\leq1$, $v_p(1-z^p)\leq2/3$.
The rigid identity principle and
\cite[Theorem~2.7(3)]{BesserDeJeuAlgorithm} identify it with Coleman
continuation. At the displayed points $1-\epsilon\zeta^p$ is a unit,
so every $j$th summand has valuation at least $j$.

If $t=\overline\zeta$, reduction of
\eqref{d3:eq:unramifiedhighercoordinate} is
\[
 h_n(t)=\frac{\sum_{a=1}^{p-1}a^{1-n}t^a}{1-t^p},
 \qquad t\in k_m^\times\setminus\{1\}.
\]
The numerator has $t$-coefficient one and degree $p-1$.
Thus $h_n$ is a nonconstant rational map of degree at most $p$.
If $p>2$, each finite fiber has at most $p$ elements, so
\[
 \#h_n(k_m^\times\setminus\{1\})
 \geq\left\lceil\frac{p^m-2}{p}\right\rceil=p^{m-1}.
\]
Its $\F_p$-span has dimension at least $m-1$.
If $p=2$, then $h_n(t)=t/(1+t)^2=w^2+w$, where
$w=(1+t)^{-1}$ ranges through $k_m\setminus\{0,1\}$.
Its image is exactly the nonzero part of
$\ker\operatorname{Tr}_{k_m/\F_2}$, whose span has dimension $m-1$.
Choose $m-1$ independent image values and Teichm\"uller lifts
$\zeta_j$ of corresponding inputs.
A nonzero $\Q_p$-linear relation among $a_{n,p}(\zeta_j)$ can
be normalized to coefficients in $\Z_p$, with one a unit.
Reduction would contradict the chosen $\F_p$-independence.
The coordinates are therefore $\Q_p$-independent; they need not
belong to $F_m$.

These are actual number-field coefficients. For $N=p^m-1$ put
$L_m=\Q(\mu_N,u)$.
The order of $p$ modulo $N$ is $m$, so the chosen completion
of $\Q(\mu_N)$ is $F_m$, and the completion of $L_m$ is $E_m$.
The shifted cyclotomic polynomial is Eisenstein in the ramified
factor, so its integer ring is $O_m$.
Theorem~1.12 of \cite{BesserDeJeu} gives $b_n(\zeta_ju)$ in
$K_{2n-1}(L_m)_{\Q}$.
Localization at the selected prime gives its rational lift to the
local integer ring, since the positive residue-field $K$-groups
are torsion. Pullback gives a class over $O_m$.
The global automorphism fixing $\mu_N$ and inverting $u$ preserves
the selected prime and induces $\sigma$.
Naturality proves exact anti-invariance of $\beta_{n,j}$.
Clearing denominators of the individual $b_n(\zeta_ju)$, then
subtracting their integral $\sigma$-images, gives the asserted
integral representatives.

Finally the common normalization
\eqref{d3:eq:highercyclotomicregulator} gives
\[
 c_{n,1}(\beta_{n,j})
 =\exp_{E_m,n}\bigl(\kappa_n\delta\,a_{n,p}(\zeta_j)\bigr).
\]
For $n\geq2$ this exponential is a $\Q_p$-linear isomorphism.
The independence of the coordinates and $\kappa_n\delta\ne0$
prove the assertion. No actual-to-completed $K$-theory
surjectivity is used.
\end{proof}

This proposition is a finite-level coefficient statement.
Its use over one fixed ring with infinite perfect residue field
requires the compatible ruling detector after that base change.


\subsection{One strict henselian ring in every odd degree}
\label{d3:sec:strictallprimeallodd}

Put $k=\overline{\F}_p$, $V_\infty=W(k)$, and define
\[
 f_p(z)=
 \begin{cases}
  \Phi_p(1+z),&p>2,\\
  \Phi_4(1+z)=z^2+2z+2,&p=2.
 \end{cases}
\]
The polynomial has Eisenstein degree $p-1$ for odd $p$ and degree
two for $p=2$.

\begin{theorem}\label{d3:thm:strictallprimeallodd}
For every prime $p$, the rings
\[
 R_{\infty,p}=V_\infty[[x,y,z]]/(xy+f_p(z)),\qquad
 S_p=\left(
  \bigl(\Z_{(p)}[x,y,z]/(xy+f_p(z))\bigr)_{(p,x,y,z)}
       \right)^{\mathrm{sh}}
\]
are strictly henselian regular local rings of dimension three,
and $\widehat S_p\cong R_{\infty,p}$ after fixing their residue
identifications with $k$.
For every $n\geq2$, each of these two rings $R$ satisfies
\[
 \dim_{\Q}\ker\left(
 K_{2n-1}(R)_{\Q}\longrightarrow
 K_{2n-1}(\Frac R)_{\Q}\right)=\infty.
\]
More precisely, for every $m\geq2$ the kernel contains $m-1$
rationally independent classes represented by actual integral
$K$-classes whose restrictions to $R[1/x]$ are zero integrally.
The ring for a fixed $p$ is independent of both $m$ and $n$.
\end{theorem}

\begin{proof}
Fix $n\geq2$ and $m\geq2$, and write $q=2n-1$,
$V_m=W(\F_{p^m})$, and $F_m=\Frac V_m$.
Take $\rho=i$ for $p=2$, and a primitive $p$th root for $p>2$.
Put $W_m=V_m[\rho]$, $E_m=F_m(\rho)$ and
$W_\infty=V_\infty[\rho]$, $E_\infty=\Frac W_\infty$.
These are complete coefficient DVRs and their fraction fields;
$\rho-1$ is a uniformizer in each ramified DVR.
Let $\sigma$ fix the unramified coefficient ring and invert $\rho$.

Proposition~\ref{d3:prop:unramifiedhighercoefficients} supplies
$m-1$ actual coefficients $\beta_j\in K_q(W_m)$, after one
common integer multiplication, with $\sigma\beta_j=-\beta_j$
and independent $\Q_p$-valued \'etale Chern classes.
Lemma~\ref{d3:lem:cherncompletionfactor} implies that their images
in $\pi_qK(W_m;\Q_p)$ are $\Q_p$-independent.
Under the AMMN point isomorphism used in
Lemma~\ref{d3:lem:formalhigherTCcup}, denote their coordinates by
\[
 a_j=r_{n,W_m}(\beta_j)\in E_m.
\]
These coordinates are independent and $\sigma a_j=-a_j$.
Naturality of that point calculation under $W_m\to W_\infty$
sends them by the field inclusion $E_m\hookrightarrow E_\infty$,
preserving independence. This argument only transports independence;
it does not identify the AMMN coordinates with a prescribed
polylogarithm normalization.

Set $R_m=V_m[[x,y,z]]/(xy+f_p(z))$.
Weierstrass division identifies its regular divisor $x=0$ with
$W_m[[y]]$, by $1+z\mapsto\rho$.
Pull back $\beta_j$ and push forward along this divisor to obtain
$g_j\in K_q(R_m)$. Localization makes $g_j|_{R_m[1/x]}=0$
integrally.
Both $R_m$ and $R_{\infty,p}$ are regular local of dimension three:
the equation has linear term $p$ in the regular ambient ring,
so $x,y,z$ are regular parameters.
The closed fiber of $R_m\to R_{\infty,p}$ is $k$; hence the map
is flat by \cite[Tag~00R4]{StacksFlatness}.
Its pullback therefore preserves the supported pushforward.

After the finite free extension $V_\infty\to W_\infty$ and
inversion of $p$, the polynomial splits with distinct roots
$a_g=g(\rho)-1$, indexed by
$G=\operatorname{Gal}(E_m/F_m)$.
Write $I_g=(x,z-a_g)$ and $\eta_g=1-[I_g]$ on this generic ring.
The root components of $x=0$ are disjoint Cartier divisors;
their structure-sheaf classes sum to zero because their union is
principal. Flat base change and the projection formula give
\[
 g_j=\sum_{g\in G}g(\beta_j)\eta_g.
\]
It suffices to retain the fixed inverse-root pair.
Put $d=\rho-\rho^{-1}\ne0$,
$H(z)=\prod_{g\ne1,\sigma}(z-a_g)$, and substitute
\[
 z=(\rho-1)+dZ,\qquad x=dX,\qquad
 y=dH((\rho-1)+dZ)Y
\]
into the smooth formal quadric
$B=W_\infty\langle X,Y,Z\rangle/(XY+Z(Z+1))$.
All original formal variables are topologically nilpotent, so the
substitution defines a ring map. It uses $z=t-1$, including at two.
The two selected root ideals pull to $J=(X,Z)$ and $J'=(X,Z+1)$,
up to the invertible generic scalar $d$.
Every omitted root ideal becomes trivial: if
$b=(a_g-a_1)/d\notin\{0,-1\}$, then
\[
 b(b+1)=-XY-(Z-b)(Z+b+1)\in(X,Z-b).
\]
No inversion of $H$ or flatness of this last substitution is used.
The pulled invertible module surjects onto the displayed invertible
ideal, hence is isomorphic to it.
The two remaining structure-sheaf classes sum to zero, so on
$B[1/p]$ the class becomes exactly
$2\beta_j(1-[J])$.
Lemma~\ref{d3:lem:formalhigherTCcup} detects these products by their
independent coordinates $2a_j$ in $K_q(B)$.
Here $K_q(B/\pi)$ is torsion, since the special fiber has the
$K$-theory of $\mathbf P^1_k$ and $k$ is algebraic over $\F_p$.
Localization therefore makes $K_q(B)_{\Q}\to K_q(B[1/p])_{\Q}$
injective, so these products remain independent on the generic
test ring.
Thus the $g_j$ are rationally independent on $R_{\infty,p}$.

For the arithmetic assertion choose $N=p^m-1$ and a prime $v$
above $p$ in $\Q(\mu_N)$ compatible with $k$.
The essentially \'etale local DVR
$V_{m,0}=(\Z_{(p)}[\zeta_N])_v$ has completion $V_m$, since
$p$ has multiplicative order $m$ modulo $N$; put $W_{m,0}=V_{m,0}[\rho]$.
The coefficients in Proposition~\ref{d3:prop:unramifiedhighercoefficients}
come from the actual rational $K_q$ of
$\Frac W_{m,0}=\Q(\mu_N,\rho)$.
The localization map from $K_q(W_{m,0})$ is surjective
\cite[Sections~5--7]{Quillen}, since
$K_{q-1}(\F_{p^m})=0$ by
\cite{QuillenFinite}\cite[IV, Corollary~1.13]{Weibel}.
Clear denominators before lifting and applying $1-\sigma$.
This produces the same independent finite family from actual
classes on $W_{m,0}$.
The chosen pointed \'etale neighborhood gives $V_{m,0}\to S_p$,
and $\rho\mapsto1+z$ gives
$W_{m,0}\to S_p/(x)$.
Pull the coefficients to this regular divisor and push to $S_p$.
They vanish after inverting $x$ integrally.

The universal property gives a compatible local map
$S_p\to R_{\infty,p}$.
Both rings have residue $k$ and regular parameters $x,y,z$;
the induced associated-graded map is the identity on $k[x,y,z]$.
It follows that $\widehat S_p\cong R_{\infty,p}$.
Completion is flat, so these arithmetic supported classes map to
the already detected completed classes and are independent.
Letting $m$ grow proves both assertions.
\end{proof}


\section{One fixed ring in every degree at least three}
\label{d3:sec:singlefixedalldegree}

We use the ordinary continuous de Rham quotient throughout.
First we record why introducing one Laurent variable preserves both
the coefficient--ruling products and their products with that unit.

\begin{lemma}\label{d3:lem:onelaurentcoefficientdetector}
Let $W$ be the integers of a finite extension $E/\Q_p$, with
residue field $k$ and uniformizer $\pi$. Put
\[
 B=W\langle X,Y,Z\rangle/(XY+Z(Z+1)),\qquad
 J=(X,Z),\qquad \eta=1-[J].
\]
Let $0\ne g\in W[T]$ have unit nonzero coefficients, and put
$B_g=(B[T^{\pm1},1/g])^\wedge_\pi$.
For $n\geq2$, suppose actual classes
$\beta_1,\ldots,\beta_r\in K_{2n-1}(W)$ have
$\Q_p$-linearly independent AMMN point coordinates
$a_1,\ldots,a_r\in E$.
Then each of the two families
\[
 \eta\beta_j\in K_{2n-1}(B_g[1/p])_{\Q},\qquad
 \eta\beta_j[T]\in K_{2n}(B_g[1/p])_{\Q}
\]
is rationally independent.
\end{lemma}

\begin{proof}
Choose a common integer $N$, prime to $p$, which kills the
reductions of the $\beta_j$ in $K_{2n-1}(k)$.
Their unique actual relative lifts have point characters
$Nt^{1-n}a_j$, where $|t|=-2$.
Write $\gamma=\operatorname{ch}^-_1(\eta)$ for the nonzero ruling
two-form class of Lemma~\ref{d3:lem:formalhigherTCcup}.
The module construction of that lemma, in the finite mixed-HKR
model of relative dimension three, gives the relative characters
\begin{equation}\label{d3:eq:onelaurentcoefficientcharacters}
 Nt^{2-n}a_j\gamma,\qquad
 Nt^{2-n}a_j\gamma\wedge d\log T.
\end{equation}
Here is the degree check for the second formula, including $n>2$.
The rank-zero negative Chern cycle has terms $t^r e_{2r}$ with
$r\geq1$; a unit cycle has terms $t^a u_{2a+1}$ with $a\geq0$.
A cyclic-shuffle correction increases Hochschild degree by two.
Dimension three therefore retains only $r=1$, $a=0$, and the
ordinary shuffle. Normalized HKR gives exactly the second term
in \eqref{d3:eq:onelaurentcoefficientcharacters}.
This argument does not discard positive cyclic columns
merely because the initial column might be negative.

The degree-zero projection $P_0$ of
Lemma~\ref{d3:lem:continuouslaurentresidue} sends the first form
coefficient in \eqref{d3:eq:onelaurentcoefficientcharacters}
back to $Na_j\gamma$; its residue sends the second coefficient there.
The independence of the $a_j$ thus detects every nonzero rational
linear combination of either relative family.
The special fiber $C_g=B_g/\pi$ is smooth of dimension three.
Lemma~\ref{d3:lem:laurentspecialvanishing} gives
$\pi_{q+1}K(C_g;\Q_p)=0$ for every $q\geq3$.
Consequently the relative-to-absolute map is injective in all the
degrees used here after spectral completion and rationalization.
This proves independence in the actual groups of $B_g$.

The final generic localization requires a separate actual
$K$-theory observation. Let
$U_g=\Spec k[T^{\pm1},1/\bar g]$.
The ruling and projective-bundle formula give
$K(C_g)\simeq K(U_g)\oplus K(U_g)$.
For $q\geq2$, localization on $\A^1_k$
\cite[Sections~5--7]{Quillen} puts $K_q(U_g)$ between
a quotient of $K_q(k)$ and a subgroup of the finite direct sum
$\bigoplus_vK_{q-1}(k(v))$, where the omitted points $v$ have
finite residue fields. These groups are finite by
\cite{QuillenFinite}\cite[IV, Corollary~1.13]{Weibel}.
Thus $K_q(C_g)$ is torsion, and localization makes
\[
 K_q(B_g)_{\Q}\longrightarrow K_q(B_g[1/p])_{\Q}
\]
injective. This proves the stated generic assertions.
\end{proof}

\begin{theorem}\label{d3:thm:singlefixedalldegree}
Fix a prime $p$, let $S=\Z_{(p)}^{\mathrm{sh}}$ with residue
$k=\overline\F_p$, and set
\[
 V_0=(S[T])_{(p)},\qquad
 f_p(z)=
 \begin{cases}
  \Phi_p(1+z),&p>2,\\
  \Phi_4(1+z)=z^2+2z+2,&p=2,
 \end{cases}
\]
\[
 R_p=\bigl(V_0[x,y,z]/(xy+f_p(z))\bigr)_{(x,y,z)}.
\]
This is one Noetherian regular local ring of dimension three,
with residue field $k(T)$, and for every $j\geq3$,
\[
 \dim_{\Q}\ker\left(K_j(R_p)_{\Q}
            \longrightarrow K_j(\Frac R_p)_{\Q}\right)=\infty.
\]
Each of the finite independent families used to establish this
assertion consists of actual integral classes which vanish after
inverting $x$ integrally.
\end{theorem}

\begin{proof}
The strict henselization $S$ is a Noetherian DVR with uniformizer
$p$; see \cite[Tags~06LJ and~0AP3]{StacksHenselization}.
The localization $V_0$ is a DVR with the same uniformizer and
residue field $k(T)$.
The polynomial $f_p$ is Eisenstein of degree at least two.
In the regular ambient local ring of dimension four, the equation
has nonzero linear term $p$, proving the assertions about $R_p$.
Equivalently,
\begin{equation}\label{d3:eq:singlefixedringlocalization}
 R_p=\bigl(S[T,x,y,z]/(xy+f_p(z))\bigr)_{(p,x,y,z)}.
\end{equation}

Fix $n\geq2$ and $m\geq2$, and take $\rho=i$ at two and a
primitive $p$th root otherwise.
Choose a compatible prime $v$ of $\Q(\mu_{p^m-1})$ and put
\[
 V_{m,0}=(\Z_{(p)}[\zeta_{p^m-1}])_v\subset S,
 \qquad W_{m,0}=V_{m,0}[\rho].
\]
The actual number-field classes of
Proposition~\ref{d3:prop:unramifiedhighercoefficients} lift to
$K_{2n-1}(W_{m,0})$ after a common integer multiplication:
the obstruction is $K_{2n-2}(\F_{p^m})=0$.
Choose the integral lifts before applying $1-\sigma$, where
$\sigma$ fixes $V_{m,0}$ and inverts $\rho$.
We obtain $m-1$ strictly anti-invariant classes $\beta_j$.
Their finite-local-field \'etale Chern classes are independent;
Lemma~\ref{d3:lem:cherncompletionfactor} and the AMMN point isomorphism
therefore give independent point coordinates in
$E_m=\Frac W(\F_{p^m})(\rho)$.

The regular divisor $x=0$ in $R_p$ receives $W_{m,0}$ by
$\rho\mapsto1+z$.
Pull back and push forward to obtain $a_j\in K_{2n-1}(R_p)$,
and put $b_j=a_j[T]\in K_{2n}(R_p)$.
Here $T$ is a unit. Both families vanish on $R_p[1/x]$
integrally, by localization and the projection formula.

We check independence after any putative finite relation.
Using \eqref{d3:eq:singlefixedringlocalization} and
filtered-ring continuity of $K$-theory
\cite[IV, Section~6.4]{Weibel}, the coefficients, the classes,
and an actual integral relation all descend to
\[
 \bigl(V_{\lambda,0}[T,x,y,z]/(xy+f_p(z))\bigr)[1/F],
 \qquad \bar F(T)\ne0,
\]
where $V_{\lambda,0}\subset S$ is an essentially \'etale local
DVR containing $V_{m,0}$ and all constants in the finite data.
We include $T$ in the inverted product $F$ when necessary.
Its completion is an unramified coefficient DVR $V_\ell$ with
finite residue field; put $W_\ell=V_\ell[\rho]$ and
$E_\ell=\Frac W_\ell$.
Point-character naturality sends the independent coordinates by
the field inclusion $E_m\hookrightarrow E_\ell$.

Over $W_\ell$, retain the two roots
$\rho-1,\rho^{-1}-1$ of $f_p$.
With $d=\rho-\rho^{-1}$ and
$H(z)=\prod_{a\ne\rho-1,\rho^{-1}-1}(z-a)$, use
\[
 z=(\rho-1)+dZ,\qquad x=dX,\qquad
 y=dH((\rho-1)+dZ)Y.
\]
Choose $g(T)$ by lifting precisely the nonzero coefficients of
$\bar F(T)$ to units in $W_\ell$ and form $B_g$ as in
Lemma~\ref{d3:lem:onelaurentcoefficientdetector}.
All three displayed images are topologically nilpotent.
The image of $F$ is $g+\pi h$, hence a unit of $B_g$.
Thus this finite relation has an actual image in the test ring;
no map from all of $R_p$ to a fixed test algebra is asserted.

Flat coefficient base change and the projection formula, followed
by inversion of $p$, split the supported class as
$\sum_\tau\tau(\beta_j)(1-[I_\tau])$.
The pair substitution makes every omitted root ideal trivial.
The two remaining disjoint components have structure-sheaf
classes summing to zero, and anti-invariance therefore gives
\[
 a_j\longmapsto2\eta\beta_j,\qquad
 b_j\longmapsto2\eta\beta_j[T]
 \quad\text{in }K_*(B_g[1/p]).
\]
Lemma~\ref{d3:lem:onelaurentcoefficientdetector} excludes the asserted
relation in either degree.
Hence both families have rational rank $m-1$ in $R_p$.
As $m$ is arbitrary and the degrees $2n-1,2n$ with $n\geq2$
cover every integer at least three, the theorem follows.
\end{proof}


\subsection{A single complete ring in every degree at least three}
\label{d3:sec:completedfixedalldegree}

The finite-witness proof of Theorem~\ref{d3:thm:singlefixedalldegree}
does not itself pass to arbitrary completed coefficients.  We now
give a complete-ring construction, using
Lemma~\ref{d3:lem:essentiallysmoothAMMN} and a constant-term map
defined on all primitive denominators.

\begin{lemma}
\label{d3:lem:completedcoefficientresidue}
Let $W$ be a complete DVR with uniformizer $\pi$, and set
\[
 V_W=\widehat{W[T]_{(\pi)}}^{\,\pi},\qquad
 B'=W\langle X,Y,Z\rangle/(XY+Z(Z+1)),
\]
\[
 \mathcal B=V_W\langle X,Y,Z\rangle/(XY+Z(Z+1)).
\]
There are continuous $W$-linear maps on the relative form complexes,
a degree-zero chain projection $P_0$ and a degree-minus-one residue,
which fix base forms and satisfy
\[
 P_0(\omega)=\omega,\qquad
 \operatorname{Res}_T(\omega\wedge d\log T)=\omega
 \quad(\omega\in\widehat\Omega^\bullet_{B'/W}).
\]
They preserve actual de Rham boundaries. The complexes here retain
the differential $dT$: their base ring is $W$, not $V_W$.
\end{lemma}

\begin{proof}
For $a\geq1$, put $L_a=(W/\pi^a)((T))$.
Every $g\in W[T]\setminus(\pi)$ is a unit in $L_a$.
Indeed its nonzero reduction is invertible in the field
$k((T))$, and inverses lift across the nilpotent ideal $(\pi)$.
Uniqueness makes these inverses compatible in $a$. Hence
\[
 W[T]_{(\pi)}\longrightarrow\varprojlim_a L_a
\]
is an integral, $\pi$-adically continuous ring homomorphism.
Constant-term extraction is compatible with these reduction maps,
so extends to
\[
 \operatorname{CT}:V_W\longrightarrow W,\qquad
 \operatorname{CT}(1)=1,\qquad
 \operatorname{CT}(T\partial_T f)=0.
\]
The identity with the logarithmic derivative is coefficientwise
at every level and then passes to the inverse limit.

Tensor with the quadric algebra before completion and use
$\varprojlim_a(B'/\pi^a)((T))$ as target.  The completed relative
forms decompose into a component with no $d\log T$ and one
with $d\log T$ at the right. Apply constant term to the first
to define $P_0$, and to the second to define
$\operatorname{Res}_T$.
The derivative identity gives compatibility with the de Rham
differentials (with the conventional shift sign for a residue).
In particular an actual primitive maps to an actual primitive.
All operations are integral and bounded before inverting $p$;
there is no completion of the image of the differential.
\end{proof}

\begin{theorem}
\label{d3:thm:completedfixedalldegree}
For a prime $p$, put $k=\overline{\F}_p$, $W_0=W(k)$, and
\[
 V=\widehat{\,W_0[T]_{(p)}\,}^{\,p},\qquad
 \mathcal R_p=V[[x,y,z]]/(xy+f_p(z)),
\]
where $f_p(z)=\Phi_p(1+z)$ for $p>2$, and
$f_2(z)=\Phi_4(1+z)=z^2+2z+2$.
Then $\mathcal R_p$ is a complete Noetherian regular local
domain of dimension three, with residue field $k(T)$.
For every $j\geq3$,
\[
 \dim_{\Q}\ker\left(K_j(\mathcal R_p)_{\Q}
          \longrightarrow K_j(\Frac\mathcal R_p)_{\Q}\right)
 =\infty.
\]
The finite independent families proving this assertion consist
of actual integral classes with zero image in
$K_j(\mathcal R_p[1/x])$.
\end{theorem}

\begin{proof}
The ring $W_0[T]_{(p)}$ is a Noetherian DVR with uniformizer
$p$ and residue $k(T)$. Its completion $V$ has the same
properties. Since $f_p$ is Eisenstein of degree at least two,
the quotient in the statement has nonzero linear term $p$
in a regular ambient local ring of dimension four.
It is therefore regular of dimension three, with maximal
ideal $(x,y,z)$ and the asserted residue field.

Let $\rho=i$ for $p=2$, and otherwise let $\rho$ be a primitive
$p$th root of unity. Put
\[
 W=W_0[\rho],\qquad \pi=\rho-1,\qquad E=\Frac W.
\]
There are compatible identifications
\begin{equation}\label{d3:eq:completedfixedcoefficients}
 W_0[T]_{(p)}[\rho]=W[T]_{(\pi)},\qquad
 V[\rho]=\widehat{W[T]_{(\pi)}}^{\,\pi}.
\end{equation}
To verify the first, the norm of a polynomial
$g\in W[T]\setminus(\pi)$ reduces to $\bar g^{[W:W_0]}$
and is primitive over $W_0$. The norm and adjugate identity
therefore invert $g$ in the left side. The reverse inclusion
is immediate. Finite freeness over $W_0$ gives the second
identity after completion; the $p$-adic and $\pi$-adic
topologies are cofinal.

Fix $n\geq2$ and an arbitrary $m\geq2$.
Proposition~\ref{d3:prop:unramifiedhighercoefficients} and
Lemma~\ref{d3:lem:cherncompletionfactor} supply $m-1$ actual
classes $\beta_1,\ldots,\beta_{m-1}$ over
$W_m=W(\F_{p^m})[\rho]$, exactly anti-invariant under
$\sigma(\rho)=\rho^{-1}$, with independent AMMN point
coordinates. Their images in $K_{2n-1}(W)$ have coordinates
$b_1,\ldots,b_{m-1}\in E$ still independent over $\Q_p$:
the point character is natural under the coefficient inclusion
$\Frac W_m\hookrightarrow E$.
No surjectivity from actual $K$-groups to all point coordinates
is asserted.

The regular divisor $x=0$ of $\mathcal R_p$ has coordinate ring
\[
 \mathcal R_p/(x)\simeq V[\rho][[y]],
 \qquad \rho\longmapsto1+z.
\]
Pull back the $\beta_\nu$ and push forward to obtain
$a_\nu\in K_{2n-1}(\mathcal R_p)$.
Also put $a_\nu[T]\in K_{2n}(\mathcal R_p)$.
These actual classes vanish after inverting $x$.

Consider the single test algebra
\[
 \mathcal B=V[\rho]\langle X,Y,Z\rangle/(XY+Z(Z+1)),
 \qquad J=(X,Z),\qquad\eta=1-[J].
\]
By \eqref{d3:eq:completedfixedcoefficients} it is the completion
of an essentially smooth $W$-algebra of relative dimension
three. Lemma~\ref{d3:lem:essentiallysmoothAMMN} applies with
this perfect-constant coefficient ring $W$.
Its continuous de Rham model includes the variable $T$.

Retain the inverse pair of roots
$a=\rho-1$, $a'=\rho^{-1}-1$ of $f_p$.
Set $d=\rho-\rho^{-1}$ and
$H(z)=\prod_{\alpha\ne a,a'}(z-\alpha)$.
The substitution
\begin{equation}\label{d3:eq:completedfixedpairmap}
 z=a+dZ,\qquad x=dX,\qquad y=dH(a+dZ)Y
\end{equation}
defines an actual map $\mathcal R_p\to\mathcal B$.
The equation becomes $d^2H(a+dZ)(XY+Z(Z+1))=0$.
All three images lie in $\pi\mathcal B$, so arbitrary formal
power series converge. The coefficient map is the actual
inclusion $V\to V[\rho]$; no finite-stage descent of completed
coefficients is required.

Flat coefficient base change first splits the divisor after
inverting $p$. Pulling its root-ideal line bundles along
\eqref{d3:eq:completedfixedpairmap} leaves only $J$ and
$J'=(X,Z+1)$ on $\mathcal B[1/p]$.
For an omitted root with normalized value
$b=(\alpha-a)/d\notin\{0,-1\}$, the ideal
$(X,Z-b)$ is the unit ideal, since
\[
 b(b+1)=-XY-(Z-b)(Z+b+1).
\]
The two surviving disjoint components have classes summing
to zero, because their union is the principal divisor $X=0$.
The projection formula and exact anti-invariance therefore give
\begin{equation}\label{d3:eq:completedfixedimages}
 a_\nu\longmapsto2\beta_\nu\eta,\qquad
 a_\nu[T]\longmapsto2\beta_\nu\eta[T]
 \quad\text{on }\mathcal B[1/p].
\end{equation}

It remains to check these classes in actual groups.
The reduction of each $\beta_\nu$ lies in
$K_{2n-1}(k)$, a torsion group of order prime to $p$
element by element. Choose one common prime-to-$p$ integer
$N$ killing the finite list of reductions.
Since $K_{2n}(k)=0$, each $N\beta_\nu$ has a unique actual
relative lift at the coefficient point.
Its character is $Nt^{1-n}b_\nu$.
The module comparison and the finite dimension-three HKR
calculation of Lemma~\ref{d3:lem:formalhigherTCcup} give
\[
 Nt^{2-n}b_\nu\gamma,\qquad
 Nt^{2-n}b_\nu\gamma\wedge d\log T,
\]
where $\gamma$ is the nonzero ordinary ruling two-form class
on $B'=W\langle X,Y,Z\rangle/(XY+Z(Z+1))$.
Lemma~\ref{d3:lem:completedcoefficientresidue} sends these
respectively to $Nb_\nu\gamma$.
The independence of the $b_\nu$ detects every nonzero
rational combination of either actual relative family.

Finally,
\[
 C=\mathcal B/\pi
   =k(T)[X,Y,Z]/(XY+Z(Z+1)),\qquad
 K(C)\simeq K(k(T))\oplus K(k(T)).
\]
For $q\geq2$, the actual group $K_q(k(T))$ is torsion.
Every element descends first to a finite field of constants
and then to a finite open of its affine line.
Quillen localization expresses the $K_q$-group of that open
as an extension of a quotient of the finite group $K_q(k_m)$
by a subgroup of a finite sum of $K_{q-1}(k(v))$.
These groups are finite, including at $q=2$.
Continuity proves the claim for $k(T)$.

The algebra $\mathcal B$ is Noetherian and flat over the complete
DVR $V[\rho]$, and its special fiber is smooth.
Consequently $\pi$ is a nonzerodivisor and $\mathcal B$ is regular:
at every maximal ideal the regular quotient by $\pi$ gives the
regularity criterion, and $\pi$ belongs to the Jacobson radical.
Thus $K_{q+1}(C)_{\Q}=K_q(C)_{\Q}=0$ for every $q\geq3$.
The actual relative sequence and the regular-divisor
localization sequence respectively give injections
\[
 K_q(\mathcal B,\pi)_{\Q}\hookrightarrow K_q(\mathcal B)_{\Q},
 \qquad
 K_q(\mathcal B)_{\Q}\hookrightarrow
                     K_q(\mathcal B[1/p])_{\Q}.
\]
The detected families therefore remain independent in the
generic test ring. Equation~\eqref{d3:eq:completedfixedimages}
proves their independence on $\mathcal R_p$.
Their ranks are at least $m-1$ for arbitrary $m$, and the
degrees $2n-1$ and $2n$ cover every integer at least three.
\end{proof}


\section{Milnor, motivic, and higher-Chow consequences}
\label{d3:sec:milnormotivicconsequences}

Put
\[
 A=\bigl(\mathbf Z_{(3)}[x,y,z]/(3+xy+z^2)\bigr)_{(x,y,z)},
 \qquad R\in\{A,A^h,\widehat A,A^{\mathrm{sh}}\},
 \qquad F_R=\operatorname{Frac}(R).
\]
In each ring use the units
\[
 a=\frac{1+z}{2},\qquad \lambda=1+x,\qquad
 \mu=1+\frac{(1+x)y}{4},\qquad
 s_0=a(1-a)=1+\frac{xy}{4}.
\]
Thus $1+x\mu=\lambda s_0$ and
$c_R=2[a][\lambda][\mu]\in K_3(R)$ is the explicit class already
constructed. The superscript $M,\mathrm{naive}$ below means the
tensor algebra on $R^\times$ modulo the two-sided Steinberg ideal;
graded commutativity is not added to its definition.

\begin{corollary}\label{d3:cor:explicitmilnormotivic}
For every such $R$, each of the four generic restriction maps
\[
\begin{gathered}
 K_3^{M,\mathrm{naive}}(R)_{\mathbf Q}
       \longrightarrow K_3^M(F_R)_{\mathbf Q},\qquad
 \widehat K_3^M(R)_{\mathbf Q}
       \longrightarrow K_3^M(F_R)_{\mathbf Q},\\
 H_B^3(R,\mathbf Q(3))\longrightarrow H_B^3(F_R,\mathbf Q(3)),
 \qquad
 CH^3(R,3)_{\mathbf Q}\longrightarrow CH^3(F_R,3)_{\mathbf Q}
\end{gathered}
\]
has a nonzero kernel. Explicit witnesses are, respectively,
$2\{a,\lambda,\mu\}$, its image in improved Milnor theory,
$2[a]_B\cup[\lambda]_B\cup[\mu]_B$, and twice the cubical graph
cycle specified below. The Milnor symbols and the graph cycle vanish
over $F_R$ integrally.
\end{corollary}

\begin{proof}
We first record the stronger Quillen input
\begin{equation}\label{d3:eq:explicitnonzeropopen}
 c_R|_{R[1/3]}\ne0\quad\text{in }K_3(R[1/3])_{\mathbf Q}.
\end{equation}
Proposition~\ref{d3:prop:directtcoriginal} detects the actual image of
$c_R$ on
$B=W\langle X,Y,Z\rangle/(XY+Z^2-1)$, where
$W=\mathbf Z_3[u]/(u^2+3)$, by $(x,y,z)=(uX,uY,uZ)$.
For $A^{\mathrm{sh}}$, first map to its completion and use the same
calculation over
$W_\infty=W(\overline{\mathbf F}_3)[u]/(u^2+3)$:
the nonzero endpoint difference remains nonzero in the larger
coefficient field, and Lemma~\ref{d3:lem:tcreviewresidue} applies there.
The perfect-residue-field scope of
\cite[Proposition~4.10]{AMMN} permits this base change.
In either case $C=B/u$ is an affine-line torsor over
$\mathbf P^1_k$, with $k=\mathbf F_3$ or
$\overline{\mathbf F}_3$. Hence $K_3(C)$ is torsion by the projective
bundle formula, Quillen's finite-field computation, and filtered
continuity. Localization and d\'evissage therefore make
$K_3(B)_{\mathbf Q}\to K_3(B[1/3])_{\mathbf Q}$ injective.
The nonzero detector image on $B[1/3]$ factors through $R[1/3]$,
proving~\eqref{d3:eq:explicitnonzeropopen}.

\emph{Milnor theory.}
There are natural symbol maps
\[
 K_3^{M,\mathrm{naive}}(R)
 \xrightarrow{\eta_R}\widehat K_3^M(R)
 \xrightarrow{\iota_R}K_3(R),
\]
whose composite sends a symbol to the product of its unit classes.
Kerz's universal construction, Proposition~4.2.8(4)--(6), and
Theorem~4.3.2 give these maps, surjectivity of $\eta_R$, and a map
$r_R:K_3(R)\to\widehat K_3^M(R)$ with
$r_R\iota_R=2$; these assertions hold for arbitrary local rings
\cite[Theorem~4.2.5, Proposition~4.2.8, Theorem~4.3.2]{KerzThesis}.
Thus $\iota_R$ is split injective rationally. Both displayed Milnor
classes have infinite order because their Quillen image is $c_R$.
For the first three rings the residue field is $\mathbf F_3$ and no
equality of naive and improved theory is asserted. For
$A^{\mathrm{sh}}$ the residue field is infinite, so $\eta_R$ is an
isomorphism integrally.

In the field $F_R$, Steinberg relations give
\[
 \{\lambda s_0,\mu\}
 =\{-x,\lambda s_0\}=\{-x,s_0\},
 \qquad 2\{a,s_0\}=0.
\]
The first equality uses $\{-x\mu,1+x\mu\}=0$, and the second
cancellation uses $\{-x,1+x\}=0$. Consequently
\[
 2\{a,\lambda,\mu\}
 =2\{a,\lambda s_0,\mu\}
 =-2\{a,s_0,-x\}=0.
\]
These calculations use $-x$ only in the field, where it is a unit.

\emph{Beilinson cohomology.}
All four rings are regular, noetherian, and finite dimensional.
Cisinski--D\'eglise identify
\[
 H_B^q(X,\mathbf Q(p))
 \simeq\operatorname{Gr}_\gamma^pK_{2p-q}(X)_{\mathbf Q}
\]
for every regular $X$, compatibly with the multiplicative Adams
decomposition \cite[Corollaries~14.2.14 and~14.2.19]{CisinskiDegliseMotives}.
A unit has Adams weight one, so $c_R$ has pure weight three and
corresponds to the asserted nonzero cup product. Its generic
restriction is zero by the preceding symbol calculation.

\emph{Ordinary cubical higher Chow groups.}
Let $\square=\mathbf P^1\setminus\{1\}$, with faces $0,\infty$.
Intersect the graph of
$(a,\lambda,\mu):\operatorname{Spec}R\to(\mathbf P^1)^3$ with
$\operatorname{Spec}R\times\square^3$, and call the resulting
closed cycle $\Gamma_R$. It has codimension three and
\[
 \Gamma_R\simeq\operatorname{Spec}R[1/(xy)],
 \qquad \gamma_R=2[\Gamma_R]\in CH^3(R,3).
\]
Indeed $a-1$ is a unit, $\lambda-1=x$, and
$\mu-1=\lambda y/4$; the three graph equations form a regular
sequence. This is not a morphism from all of $\operatorname{Spec}R$
to $\square^3$. Every face intersection is empty because the three
coordinates are units, so $\Gamma_R$ is admissible and has zero
boundary. Over $F_R$, Totaro's field isomorphism sends the Milnor
symbol to this graph; the integral vanishing already proved therefore
gives $\gamma_R|_{F_R}=0$
\cite[Theorem~1 and Sections~1--2]{TotaroMilnor}.

Set $U=\operatorname{Spec}R[1/3]$. This is a regular equal-characteristic
zero scheme, and the compatible comparisons give
\[
 CH^3(U,3)_{\mathbf Q}\simeq H_B^3(U,\mathbf Q(3))
       \simeq K_3(U)_{\mathbf Q}^{(3)}.
\]
Here $U$ is affine, excellent, regular, noetherian, finite dimensional,
and of equal characteristic zero. The comparison with the actual
Bloch cycle complex in precisely this affine scope is recorded in
\cite[Remark~1.3.2]{BondarkoIvanovChowWeights}; see also
\cite[Introduction~B, p.~xxii; Theorem~11.1.24 and
Examples~11.1.25 and~11.2.3]{CisinskiDegliseMotives}.
The cubical and simplicial cycle complexes give the same higher Chow
groups \cite[Theorem~4.7]{LevineHigherChowRevisited}.
The rational monoidal comparison with Beilinson motives is
\cite[Theorem~16.1.4 and Corollary~16.1.7]{CisinskiDegliseMotives}.
The graph-symbol normalization and its compatibility with the product
of unit classes can be checked on a smooth finite-type
characteristic-zero model and then pulled back: each map $A\to R$
is flat, and flat pullback takes the graph cycle to the displayed
graph cycle. Thus the restricted class $\gamma_R|_U$ corresponds to
$c_R|_U$.
Thus~\eqref{d3:eq:explicitnonzeropopen} implies
$\gamma_R|_U\ne0$, and hence $\gamma_R\ne0$.
Only open restriction on $R$ is used in this last step; no comparison
between the full mixed-characteristic cycle complex and Beilinson
cohomology on $R$ is assumed.
\end{proof}


% End input: sections/04-three-dimensional-families.tex


% Begin input: sections/05-two-dimensional-families.tex
\section{Higher odd degrees for the arithmetic quartic}
\label{d2ext:higher-odd-quartic}

The arithmetic root lattice also detects coefficients in every higher
odd degree. In this section write $A_2=A$ for the quartic ring of
Theorem~\ref{d2:thm:geometric}, let
$W_{\mathrm{ar}}=\mathbf Z_{(5)}[\rho]$ and
$E_{\mathrm{ar}}=\mathbf Q(\rho)$, and let
$i:V(x)\hookrightarrow\operatorname{Spec}A_2$ be its regular divisor.
We retain the root classes $\eta_g$ and their real Chern classes $e_g$
on the split affine surface $\mathcal Q$. Throughout this section
$g\in\mathbf F_5^\times$, and $\sigma_g$ sends $\rho$ to $\rho^g$.

\begin{theorem}
\label{d2ext:twodimalloddgeometric}
For the quartic arithmetic local ring $A_2$ and every $n\ge2$, the kernel of
$K_{2n-1}(A_2)\to K_{2n-1}(\Frac A_2)$ contains an actual infinite-order
class. The rational kernel has dimension at least two for even $n$
and at least one for odd $n$.
\end{theorem}
\begin{proof}
Borel's theorem gives $\dim_{\Q}K_{2n-1}(E_{\mathrm{ar}})_{\Q}=2$ and
$\dim_{\Q}K_{2n-1}(\Q)_{\Q}=0$ for even $n$, respectively $1$
for odd $n$ \cite[Proposition~12.2]{Borel}. Denote restriction and transfer by $\operatorname{res}$
and $\operatorname{Tr}$, respectively; they satisfy
$\operatorname{res}\operatorname{Tr}=\sum_g\sigma_g$ and
$\operatorname{Tr}\operatorname{res}=4$. Thus the invariant subspace
is exactly the restriction from $\Q$, and its trace-zero complement
$V_n$ has the asserted dimension. Quillen localization \cite{Quillen,Weibel} gives
$K_{2n-1}(W_{\mathrm{ar}})_{\Q}\simeq K_{2n-1}(E_{\mathrm{ar}})_{\Q}$.
Choose an actual $\gamma\in K_{2n-1}(W_{\mathrm{ar}})$ with noninvariant Borel
image and put $\beta_n=4\gamma-\sum_g\sigma_g\gamma$.
Its regulator is nonzero and its Galois sum is zero exactly.

Set $c_n=i_*\beta_n$. Flat base change again identifies its image
on $\mathcal T$ with $\sum_g(\sigma_g\beta_n|_{E_{\mathrm{ar}}})\eta_g$.
At a complex embedding, represent the point regulators
$b_g\in\C/\R(n)$ in the complementary real line $i\R(n)$
\cite[Section~6]{BunkeTamme}.
They form a nonzero vector with sum zero, hence a nonconstant vector.
For each finite open $U=\mathcal Q_s$, the product regulator is
represented by $2\pi i\sum_gb_ge_g$ in
\[
 \HD^3(U,\R(n+1))=H^2(U,\C)/H^2(U,\R(n+1)),
 \qquad n+1\ge3.
\]
Multiplication by $2\pi i$ induces an isomorphism
$\C/\R(n)\simeq\C/\R(n+1)$. The root lattice and
Proposition~\ref{d2:prop:geometry} therefore make this class
nonzero. If an integer killed $c_n$, continuity for the filtered
arithmetic localization would make that equality hold on some
$\mathcal Q_s$, contradicting the regulator. Localization at the
source divisor gives integral vanishing on $A_2[1/x]$.
Applying the argument to every nonzero vector of $V_n$ proves
injectivity on $V_n$ and gives the asserted rank bounds.
\end{proof}

The restriction $n\ge2$ is essential here: at $n=1$, the target
twist is two and $F^2$ on a complex affine surface is not zero.
The displayed quotient formula cannot be used in that degree.

\section{An explicit product of three units}
\label{d2ext:sec-explicit-triple}

Retain $A_2=A$ from Setup~\ref{d2:setup:ring} and its
completion $\widehat A_2=\widehat A$ from Theorem~\ref{d2:thm:completed}.
We identify an explicit product of three units with the weighted
root class after pullback to the split surface and the completed test algebra.

Let $N$ be a complex manifold of dimension two. Write $A^\bullet(N)$
for complex smooth forms and
\[
 I(p)^k=\{\xi\in A^k([0,1]\times N):
 \xi|_0\in (2\pi i)^pA^k_{\mathbf R}(N),\quad
 \xi|_1\in F^pA^k(N)\}.
\]
Thus $H^k(I(p))=H^k_{\mathcal D,\mathrm{an}}(N,\mathbf R(p))$,
with multiplication induced by wedge product.

\begin{lemma}\label{d2ext:interval-product}
Let $g$ be a nowhere-vanishing holomorphic function with a holomorphic
logarithm $L$, and let $z\in H^2(I(2))$. If $\xi$ is a closed
representative of $z$, put $\omega=\xi|_1$. Then $\omega$ is a
closed holomorphic two-form, is independent of the representative,
and
\[
 r_1(g)z=[-L\omega]\quad\hbox{in}\quad
 H^3_{\mathcal D,\mathrm{an}}(N,\mathbf R(3))
 =H^2(N,\mathbf C)/H^2(N,\mathbf R(3)).
\]
Here $r_1(g)$ is the analytic Deligne first regulator of $g$, with
the interval and cone conventions of
\cite[Definition~2.7 and Proposition~2.8]{BunkeTamme}.
\end{lemma}

\begin{proof}
Write $L=a+i\phi$, with $a,\phi$ real, and let $\tau$ denote the
interval coordinate. The usual representative of $r_1(g)$ is
\[
 i\,d\phi+d(\tau a).
\]
Its difference from $d(\tau L)$ is
$d(i(1-\tau)\phi)$. The primitive $i(1-\tau)\phi$ belongs to
$I(1)^0$: its value at $0$ is purely imaginary, hence in
$\mathbf R(1)$, and its value at $1$ is zero. Consequently
$d(\tau L)$ represents $r_1(g)$ in $I(1)$.

The endpoint $\omega=\xi|_1$ is closed and lies in $F^2A^2(N)$,
so it is a closed holomorphic two-form. Changing $\xi$ by an
$I(2)$-boundary cannot change this endpoint, because
$F^2A^1(N)=0$.

The product has representative
\[
 \upsilon=d(\tau L)\wedge\xi=d(\tau L\xi).
\]
Its endpoint at $0$ vanishes. Its endpoint at $1$ is
$dL\wedge\omega$, which vanishes by complex dimension two.
Thus the cone map sends its class to
$[-\int_0^1\upsilon]$. Fiberwise Stokes gives
\[
 -\int_0^1\upsilon
 =-L\omega+d_N\!\left(\int_0^1\tau L\xi\right),
\]
and hence gives the asserted representative. It is closed, since
$d(L\omega)=dL\wedge\omega=0$. The primitive $\tau L\xi$
has endpoint $L\omega$ at $1$; as $F^3A^2(N)=0$, it generally
does not belong to $I(3)^2$.

Finally $F^3A^\bullet(N)=0$, so the cone long exact sequence gives
the stated quotient: the next real-to-complex cohomology map is
injective. This proves the formula.
\end{proof}

\begin{proposition}
\label{d2ext:dimtwoExplicitTriple}
The triple
\begin{equation}\label{d2ext:eq-explicit-triple}
 c_{\rm exp}=\left\{2+y,\ 1-x,\
 1+(1-x)(x^3+Ty^3)y\bigl((1+y)^5+1\bigr)\right\}\in K_3(A_2)
\end{equation}
has infinite order in $A_2$ and in $\widehat A_2$, and its
restriction to the localization at $x$ is zero integrally in both cases.
\end{proposition}

\begin{proof}
\emph{Integral vanishing.}
Put $t=-1-y$, $v=x^3+Ty^3$, and
\[
 H(t)=\frac{t^{10}-1}{\Phi_{10}(t)}=(t^5-1)(t+1)
      =y\bigl((1+y)^5+1\bigr),\quad b=vH(t).
\]
Then $xv=-\Phi_{10}(t)$, so $1-xb=t^{10}$.
Set $g=1-t$, $\lambda=1-x$, and $\mu=1+\lambda b$.
These are units: their residues are $2,1,1$ respectively, and
$t$ is also a unit. In the convention
$\langle r,s\rangle=\{r,1-rs\}$ when $r$ is a unit, the
Dennis--Stein relations give
\[
 \delta=\langle x,b\rangle
 =\langle x,b+1-xb\rangle=\langle x,\mu\rangle
 =\{\lambda t^{10},\mu\}.
\]
Here $\langle x,1\rangle=0$ and $1-x\mu=\lambda t^{10}$;
see \cite[Chapter~III, Definition~5.11 and relations (D1)--(D3)]{Weibel}.
Since $\{g,t^{10}\}=10\{1-t,t\}=0$, multiplication gives
\begin{equation}\label{d2ext:dimtwoDSidentity}
 [g]\delta=\{g,\lambda,\mu\}=c_{\rm exp},\qquad
 [g]\delta|_{D(x)}=\{g,x,t^{10}\}=0.
\end{equation}
All identities are integral and persist under every ring map in use.

\emph{The universal supported coefficients.}
Write $E_{\mathrm{ar}}=\mathbf Q(\rho)$, where $\rho$ is primitive of order five,
and $t_g=-\rho^g$ for $g\in\mathbf F_5^\times$. On the smooth surface
\[
 Q^0=\operatorname{Spec}E_{\mathrm{ar}}[X,V,t]/(XV+\Phi_{10}(t)),\qquad
 U=Q^0[1/(t(1-t))],
\]
the class $\alpha=[1-t]\langle X,VH(t)\rangle$ is defined and
vanishes on $U[1/X]$. The entire divisor $D=V(X)$ is contained
in $U$. Localization and d\'evissage therefore give an actual
supported lift $\kappa\in K_3(Q^0)$ of $\alpha$.
Put $\eta_g=1-[(X,t-t_g)]$. Since the four components of $D$
are affine lines, $\kappa=\sum_g\gamma_g\eta_g$ with
$\gamma_g\in K_3(E_{\mathrm{ar}})$, and $\sum_g\eta_g=0$.

\emph{Independence of the supported lift.}
This lift is unique rationally as a class on $Q^0$.
Here is a direct $K$-theoretic check, independent of its periods,
using localization and the Laurent fundamental theorem \cite{Weibel}.
Put $Z=\mathbf A^1_{E_{\mathrm{ar}}}\setminus\{0,1\}$.  Then
$U[1/X]\simeq\mathbf G_m\times Z$, and localization on the
affine line, split by the units $t$ and $1-t$, gives
\[
 K_r(Z)\simeq K_r(E_{\mathrm{ar}})\oplus
       K_{r-1}(E_{\mathrm{ar}})[t]\oplus K_{r-1}(E_{\mathrm{ar}})[1-t].
\]
Together with the Laurent fundamental theorem this gives
\[
\begin{split}
 K_4(U[1/X])\simeq{}&
 K_4(E_{\mathrm{ar}})\oplus K_3(E_{\mathrm{ar}})[t]\oplus K_3(E_{\mathrm{ar}})[1-t]\\
 &{}\oplus K_3(E_{\mathrm{ar}})[X]
 \oplus K_2(E_{\mathrm{ar}})[X][t]\oplus K_2(E_{\mathrm{ar}})[X][1-t].
\end{split}
\]
The first three summands extend to $U$, so their boundary in
$K_3(D)=K_3(E_{\mathrm{ar}})^4$ is zero.  For a number field,
$K_2(E_{\mathrm{ar}})_{\mathbf Q}=K_4(E_{\mathrm{ar}})_{\mathbf Q}=0$: localization
reduces this to the rational even $K$-group vanishing for its
number ring \cite[Proposition~12.2]{Borel} and the torsion odd
$K$-groups of finite fields
\cite{QuillenFinite}\cite[IV, Corollary~1.13]{Weibel}.
The only remaining boundary is that of
$\beta[X]$, namely the diagonal tuple
$\pm(\beta,\beta,\beta,\beta)$ with a common convention-dependent
sign, since $X$ has order one on each root component. Localization therefore
identifies the kernel as
\[
 \ker\bigl(K_3(D)_{\mathbf Q}\longrightarrow K_3(U)_{\mathbf Q}\bigr)
       =\Delta K_3(E_{\mathrm{ar}})_{\mathbf Q}.
\]
The diagonal pushes forward to zero on $Q^0$, by the
length-one free resolution of the principal divisor $X=0$.
Consequently any two actual supported lifts of $\alpha$ give
the same class in $K_3(Q^0)_{\mathbf Q}$.
Moreover, since
$Q^0[1/X]\simeq\mathbf G_m\times\mathbf A^1_{E_{\mathrm{ar}}}$, the same
Laurent calculation gives
$\ker(K_3(D)\to K_3(Q^0))=\Delta K_3(E_{\mathrm{ar}})$ already integrally.
Thus its supported coefficients are unique modulo the diagonal.

Write $\operatorname{bl}$ for the Bloch-symbol map to algebraic
$K_3$ in \cite[Theorem~4.9]{BunkeTamme}. For a root $a$ of order ten define
$\beta_{\rm BT}(a)=\frac1{10}\operatorname{bl}(10[a])$
using \cite[Theorem~4.9 and (4.8)]{BunkeTamme}. Both $a$ and
$1-a$ are global units, and $10[a]$ has zero wedge boundary.
The normalization at every complex embedding $\sigma$ is
$r_2(\beta_{\rm BT}(a))_\sigma=-iD_{\rm BW}(\sigma a)$.
We claim the exact rational identity
\begin{equation}\label{d2ext:dimtwoUniversalCoefficients}
 \kappa=-10\sum_g\beta_{\rm BT}(t_g)\eta_g
          \quad\hbox{in }K_3(Q^0)_{\mathbf Q}.
\end{equation}

\emph{Determination of the supported coefficients.}
Fix a complex embedding $\sigma:E_{\mathrm{ar}}\hookrightarrow\mathbf C$ and
write $a_g=\sigma t_g$. The four roots are
$e^{\pm i\pi/5}$ and $e^{\pm3i\pi/5}$. For any distinct pair
$a_i,a_j$, their chord avoids $0$, $1$, and the other roots. Choose
one neighborhood of the chord on which the logarithm
$L=\operatorname{Log}(1-t)$ and the classical dilogarithm
$F=\operatorname{Li}_2(t)$ are single-valued. We use the branches
continued from their power series on the open unit disc; they satisfy
$dF=-L\,dt/t$.

Over this chord the matching sphere is parametrized by
\[
 X=\sqrt{|\Phi_{10}(t)|}\,e^{i\vartheta},\qquad
 V=-\frac{\Phi_{10}(t)}{|\Phi_{10}(t)|}\,\overline X,
\]
with $X=V=0$ at the two endpoints. This parametrization satisfies
$XV=-\Phi_{10}(t)$ and is smooth at each pole: $|X|^2$ is a
smooth positive multiple of the distance to that endpoint along
the chord. Orient the sphere by $ds\wedge d\vartheta$, where $s$
increases from $a_i$ to $a_j$.

On $X\ne0$, the weight-two holomorphic endpoint of the regulator
of $\delta=\langle X,VH(t)\rangle$ is
\[
 \omega=10\frac{dX}{X}\wedge\frac{dt}{t}
        =-\frac{dX\wedge d(VH(t))}{t^{10}}.
\]
The last expression extends holomorphically across both poles.
It therefore identifies the endpoint on the entire neighborhood.
Lemma~\ref{d2ext:interval-product} now gives
$r_3(\alpha)=[-L\omega]$. In particular this calculation retains
the full interval regulator, including its real endpoint.
On the punctured sphere,
$L\omega=10d(F\,dX/X)$. The boundary circles at the two poles
contribute $2\pi iF(a_j)$ and $-2\pi iF(a_i)$, and the integrals
over shrinking smooth caps tend to zero. The resulting period is
\begin{equation}\label{d2ext:matching-period}
 \int_{\Sigma_{ij}}r_3(\alpha)
   =-20\pi i\bigl(F(a_j)-F(a_i)\bigr)
       \pmod{\mathbf R(3)}.
\end{equation}
Since $|a_i|=|a_j|=1$, the real part is
$20\pi(D_{\rm BW}(a_j)-D_{\rm BW}(a_i))$.
The root divisor has local normal coordinate $X$. Its intersection
numbers with the oriented sphere are $+1$ at the initial pole
and $-1$ at the terminal pole. Thus the Chern periods of
$\eta_i$ and $\eta_j$ are $2\pi i$ and $-2\pi i$, respectively.
Multiplicativity gives, for every $\sigma$,
\[
 r_2(\gamma_i-\gamma_j)_\sigma
   =-10r_2\bigl(\beta_{\rm BT}(t_i)-\beta_{\rm BT}(t_j)\bigr)_\sigma.
\]
Borel injectivity \cite[Section~6]{BunkeTamme}, applied to all
embeddings, proves equality in $K_3(E_{\mathrm{ar}})_{\mathbf Q}$. The remaining coefficient is common to
all four components and contributes zero because $\sum_g\eta_g=0$.
This proves \eqref{d2ext:dimtwoUniversalCoefficients}.

\emph{Comparison with the actual cyclotomic coefficient.}
Write $\beta_{\rm BDJ}(a)$ for the field class supplied by the
symbol map of \cite[Theorem~1.10(1)--(2) and Remark~1.7]{BesserDeJeu}.
The complex formula of \cite[Proposition~4.1 and Remark~5.2]{DeJeuWedge}
compares its regulator with that of $\beta_{\rm BT}(a)$ by one
nonzero real scalar, independent of the field, root, and embedding.
This scalar is rational. Indeed, take $a=\zeta_6$ in
$L=\mathbf Q(\sqrt{-3})$ and put
$\beta_{\rm BT}(a)=\frac16\operatorname{bl}(6[a])$.
Both classes belong to the one-dimensional rational vector space
$K_3(L)_{\mathbf Q}$ \cite[Proposition~12.2]{Borel}
and are nonzero, since $D_{\rm BW}(\zeta_6)>0$. Consequently
$\beta_{\rm BDJ}(\zeta_6)=s\beta_{\rm BT}(\zeta_6)$ for a
unique $s\in\mathbf Q^\times$. Taking the complex regulator
identifies $s$ with the universal real scalar. Borel injectivity
at all embeddings \cite[Section~6]{BunkeTamme} then yields
\[
 \beta_{\rm BDJ}(a)=s\beta_{\rm BT}(a)
\]
for every root in use. The same field symbol map occurs in the
complex and $p$-adic formulas of the cited theorem.
Finite-residue localization identifies
$K_3(W_{\rm ar})_{\mathbf Q}$ with $K_3(E_{\mathrm{ar}})_{\mathbf Q}$,
so this comparison also holds before local completion.
Consequently the coefficient and character argument of
Proposition~\ref{d2:prop:coefficient} applies to
\[
 \Gamma=\operatorname{bl}(10[-\rho])\in K_3(\mathcal O_{E_{\mathrm{ar}}}),
 \qquad \beta=\Gamma-\sigma_{-1}\Gamma\in K_3(W_{\rm ar}).
\]
This is an actual integral antisymmetrization. Its residue class
vanishes because $\sigma_{-1}$ acts trivially on the residue field.
The identities $D_{\rm BW}(a^{-1})=-D_{\rm BW}(a)$ at every
embedding and Borel injectivity \cite[Section~6]{BunkeTamme} give
\[
 \beta_{\mathbf Q}=20\beta_{\rm BT}(-\rho)
                  =\frac{20}{s}\beta_{\rm BDJ}(-\rho).
\]
The syntomic symbol formula therefore makes the point regulator
of $\beta$ a common nonzero rational multiple of
$\operatorname{Li}_2(-\rho)$. The two congruences in
the proof of Proposition~\ref{d2:prop:coefficient} imply that its $e_1$ and
$e_3$ components are nonzero. The finite Chern factorization and
point comparison in Proposition~\ref{d2:prop:coefficient}
transfer this statement to its AMMN point coordinates.
Equation~\eqref{d2ext:dimtwoUniversalCoefficients} becomes
\begin{equation}\label{d2ext:dimtwoExactHalf}
 \kappa=-\frac12\sum_g\sigma_g(\beta)\eta_g.
\end{equation}

\emph{Detection on the arithmetic ring.}
On the chord $e^{-i\pi/5}\to e^{i\pi/5}$, the real period in \eqref{d2ext:matching-period} is
\[
 40\pi D_{\rm BW}(e^{i\pi/5})>0.
\]
The inequality follows from
$D_{\rm BW}(e^{i\theta})=-\int_0^\theta
\log(2\sin(v/2))\,dv$ for $0<\theta<\pi/3$.
This sphere lies in $Q^0[1/(t+1)]$, identified with $\mathcal Q$
by $X=x$, $V=x^3+Ty^3$, and $t=-1-y$.
The supported expression for $\kappa$ makes its regulator a real
root class. Proposition~\ref{d2:prop:geometry} keeps this class
nonzero on every permitted finite localization. If a nonzero
integer killed $c_{\rm exp}$, continuity would realize that
vanishing at one such finite localization, contradicting its
regulator. This proves infinite order in $A_2$ directly.

\emph{Detection after completion.}
The same substitutions define a ring map
$\mathcal O(Q^0)\to B[1/\pi]$,
with $X\mapsto\pi X$, $t\mapsto-1-\pi Y$.
Both $t$ and $1-t$ are units already in $B$, and the root line
$(X,t-t_g)$ pulls back to $J_g[1/\pi]$.
The denominator and completion maps of
Section~\ref{d2:sec:assembly} therefore give
\[
 c_{\rm exp}|_{B[1/\pi]}
 =\widehat c_{\rm exp}|_{B[1/\pi]}
 =-\frac12\sum_g\sigma_g(\beta)(1-[J_g])\big|_{B[1/\pi]}.
\]
Proposition~\ref{d2:prop:weighted} and the actual assembly
detect the right side as nonzero in rationalized completed
$K$-theory. Thus both actual source classes have infinite order.
The equality used here is on the split surface and its test pullback.
\end{proof}

\begin{corollary}
\label{d2ext:twodimexplicitMilnor}
For either $A=A_2$ or $A=\widehat A_2$, the ordinary Milnor symbol
\[
 m=\left\{2+y,\ 1-x,\
 1+(1-x)(x^3+Ty^3)y\bigl((1+y)^5+1\bigr)\right\}
       \in K_3^M(A)
\]
has infinite order and has zero image in
$K_3^M(\Frac A)$, integrally. In particular the ordinary
Milnor Gersten augmentation also fails after rationalization.
\end{corollary}

\begin{proof}
Write
\begin{align*}
 t&=-1-y,\qquad g=1-t,\qquad \lambda=1-x,\qquad q=t^{10},\\
 b&=(x^3+Ty^3)y\bigl((1+y)^5+1\bigr),\qquad \mu=1+\lambda b.
\end{align*}
The defining equation gives $1-xb=q$ and hence
$1-x\mu=\lambda q$. The elements $g,\lambda,\mu,q$
are units in $A$: their residues are respectively $2,1,1,1$.
In its fraction field, $x$ and $x\mu$ are also nonzero.
The two Steinberg relations and skew-symmetry in field
Milnor $K$-theory give
\begin{align*}
 0&=\{x\mu,\lambda q\},\qquad \{x,\lambda\}=0,\\
 \{x,q\}&=-\{\mu,\lambda\}-\{\mu,q\}
            =\{\lambda q,\mu\}.
\end{align*}
Furthermore $\{g,q\}=10\{1-t,t\}=0$. Therefore, without
rationalizing,
\[
 m=\{g,\lambda,\mu\}
   =\{g,\lambda q,\mu\}
   =\{g,x,q\}
   =-\{g,q,x\}=0
       \quad\text{in }K_3^M(\Frac A).
\]
The natural homomorphism $K_3^M(A)\to K_3(A)$ induced
by products of units sends $m$ to the already detected
infinite-order class $c_{\rm exp}$. Thus no nonzero integer
can kill $m$. 
\end{proof}

\section{A tame family at every residue prime}
\label{d2ext:sec-tame-family}

\begin{theorem}
\label{d2ext:twodimtamefamily}
Let $p$ be any prime and let $d\ge3$ satisfy $p\nmid d$. Put
\[
 A_{p,d}=\left(\Z_{(p)}[T,x,y]/
       (p+x^d+y^d+Txy^{d-1})\right)_{(p,x,y)}.
\]
This is a regular local ring of dimension two, with residue field
$\F_p(T)$, and
\[
 \dim_{\Q}\ker\bigl(K_3(A_{p,d})_{\Q}
       \longrightarrow K_3(\Frac A_{p,d})_{\Q}\bigr)
       \ \ge\ \lfloor d/2\rfloor.
\]
In particular these dimensions are unbounded with the residue
prime and local dimension fixed. The theorem concerns the
uncompleted rings displayed above.
\end{theorem}

\begin{proof}
\emph{The actual source and the flat splitting ring.}
Write $A=A_{p,d}$. Its regular ambient local ring has parameters
$p,x,y$, and its defining equation has linear term $p$.
Thus $A$ is regular of dimension two with maximal ideal $(x,y)$.
Solving the equation for $T$ gives $\Frac A=\Q(x,y)$.

Choose $\theta^d=-p$ and set
\[
 K=\Q(\theta),\qquad W_0=\Z_{(p)}[\theta],\qquad
 A/(x)=W_0[T]_{(\theta)}.
\]
Eisenstein's criterion gives $[K:\Q]=d$; $W_0$ is a DVR with
uniformizer $\theta$ and residue field $\F_p$.
Quillen localization gives
$K_3(W_0)_{\Q}\simeq K_3(K)_{\Q}$, and every actual
$K_3(K)$ class lifts to $K_3(W_0)$, since $K_2(\F_p)=0$.
Pull back such a coefficient $\beta$ to $A/(x)$ and push it
forward to obtain $c_\beta\in K_3(A)$. Its restriction to
$A[1/x]$ is zero. We will prove that the resulting rational map
$K_3(K)_{\Q}\to K_3(A)_{\Q}$ is injective.

Let $E=K(\mu_d)$ and choose a prime $\mathfrak P$ above $p$.
Put $W=(\mathcal O_E)_{\mathfrak P}$ and let $k$ be its finite
residue field. The completion of $E$ at this prime is
$\Q_p(\mu_d,\theta)$. Since $p\nmid d$, adjoining $\mu_d$
is unramified, and $Z^d+p$ remains Eisenstein over that
unramified extension. Consequently $\theta$ is a uniformizer
of $W$, and the $d$ roots of unity $(\zeta_i)_{i=1}^d$ in $W$ have distinct
reductions in $k$.
The ring $W$ is flat over $\Z_{(p)}$, but need not be finite.
Use the explicit local ring
\[
 B=(A\otimes_{\Z_{(p)}}W)_{(\theta,x,y)}
  =\left(V[x,y]/
      (x^d+y^d+Txy^{d-1}-\theta^d)\right)_{(\theta,x,y)},
 \qquad V=W[T]_{(\theta)}.
\]
The ideal indicated is maximal with residue field $k(T)$, and
$A\to B$ is flat by tensor product followed by localization.

\emph{The exceptional curve and its root points.}
The projectivized tangent cone of $B$ is
\[
 C:\ X^d+Y^d+TXY^{d-1}=U^d
       \subset\mathbf P^2_{k(T)}.
\]
It is geometrically smooth. A singular point must have $U=0$;
then $Y\ne0$ and $q(s)=s^d+Ts+1$ would have a multiple root.
The identity $dq-sq'=(d-1)Ts+d$ shows this can occur only at
finitely many algebraic values of $T$, or never if $p\mid d-1$.
Thus it cannot occur over $k(T)$.

The smooth plane curve $C$ is geometrically integral, so the
associated graded ring of $B$ is a domain. Krull intersection
then shows that $B$ is a domain. The blow-up
$\mathcal Y=\operatorname{Bl}_{(\theta,x,y)}\Spec B$
is regular and has the single reduced exceptional curve $C$.
Here are the local details needed for this assertion. Its
$\theta$-chart has equation
$X^d+Y^d+TXY^{d-1}=1$, with $x=\theta X$, $y=\theta Y$;
the other two strict-transform charts reduce to the corresponding
charts of the smooth $C$. A nonzero differential tangent to
the ambient exceptional plane therefore proves regularity at
every point of $C$. The nonregular locus is closed by excellence.
Its proper image in the local scheme $\Spec B$, if nonempty,
would contain the closed point, a contradiction. This proves
regularity everywhere.
Off the closed point the blow-up is an isomorphism; the
hypersurface $S_2$ property then also shows $B$ is normal.

The tautological ideal gives
$\mathcal O_{\mathcal Y}(C)|_C=\mathcal O_C(-1)$.
The root primes $(x,y-\theta\zeta_i)$ have strict transforms
meeting $C$ at $P_i=[1:0:\bar\zeta_i]$, each with multiplicity
one. Indeed, on the $\theta$-chart the $Y$-derivative is a unit
at $X=0,Y=\zeta_i$, and the two regular parameters there are
$\theta,X$.

\emph{The exceptional Picard lattice.}
On $C^*=C\cap\{UY\ne0\}$ the root-point classes have only
the diagonal relation. To prove this, extend $k$ to $\bar k$ and
consider
\[
 M=\Spec\bar k[X,Z,Y^{\pm1}]/(XZ-(1-Y^d)),\qquad
 T=(Z-X^{d-1})/Y^{d-1}.
\]
On inverting $X$ its ring is the Laurent UFD
$\bar k[X^{\pm1},Y^{\pm1}]$. Divisor localization and its units
$cX^aY^b$ give
$\Pic(M)=\Z^d/\Z(1,\ldots,1)$, generated by
$(X,Y-\bar\zeta_i)$.

Every finite fiber $T=t$ is geometrically integral. Under
$s=X/Y$, $w=1/Y$ it is
$w^d=s^d+ts+1$, $w\ne0$.
The polynomial $q_t=s^d+ts+1$ has $\deg\gcd(q_t,q_t')\le1$
by $dq_t-sq_t'=(d-1)ts+d$ and $d\ne0$ in $\bar k$.
If it were an $\ell$th power for a prime $\ell\mid d$,
that degree would be at least $d-d/\ell\ge2$, a contradiction.
A rational-function power would already be a polynomial power.
Since $\bar k$ contains $\mu_d$ and $p\nmid d$, this also
proves irreducibility of $w^d-q_t$: the root extension is
cyclic, and a proper degree would express $q_t$ as such a power.
Each removed $T$-fiber is therefore one principal prime divisor,
of multiplicity one. Divisor localization gives
$\Pic(M)\simeq\Pic(M_{\bar k(T)})$.
The generic fiber of $M$ is $C^*_{\bar k(T)}$, proving the root-point
assertion.

\emph{Arithmetic Picard and Betti persistence.}
The generic affine surface and test localization are
\begin{align*}
 Q&=\Spec E[T,x,y,y^{-1}]/(p+x^d+y^d+Txy^{d-1})\\
  &\simeq\Spec E[x,v,y^{\pm1}]/(xv+y^d+p),\qquad
 \mathcal T=\Spec B[1/\theta,1/y],
\end{align*}
where $v=x^{d-1}+Ty^{d-1}$. Inverting $x$ again gives a
Laurent UFD. Thus $\Pic(Q)=\Z^d/\Z(1,\ldots,1)$, generated
by $D_i=(x,y-\theta\zeta_i)$; this remains true after every
coefficient field extension.
The map $\Pic(Q)\to\Pic(\mathcal T)$ is injective.
Indeed, a root combination made principal on $\mathcal T$
can acquire on $\mathcal Y$ only strict divisors over $\theta=0$
or $y=0$, and a multiple of $C$. The former miss $C^*$ and
$\mathcal O_C(-1)$ is trivial on $U\ne0$.
Restriction to $C^*$ gives the same combination of the $P_i$,
so all coefficients are equal, which was already the zero
relation on $Q$.

Write $\mathcal O(\mathcal T)=\varinjlim_s\mathcal O(Q_s)$,
where $s$ runs over the permitted arithmetic denominators.
Each $E$-prime divisor removed by such an $s$ has zero Picard
class by the injection just proved. Its geometric components
are one Galois orbit. The absolute Galois group over $E$ acts trivially on the torsion-free
root Picard lattice, so each geometric component also has zero
class. The degree-two support sequence consequently gives
\[
 H^2(Q(\C),\R)\lhook\joinrel\longrightarrow
 H^2(Q_s(\C),\R)\qquad\text{for every permitted }s.
\]
The root Chern classes $e_i$ span an injected copy of
$\R^d/\R(1,\ldots,1)$: in the Gysin sequence of the $D_i$,
the angular class of $x$ has residue the diagonal and that of
$y$ has residue zero. This proves precisely the needed real
root-class persistence.

\emph{Detection of every nonzero coefficient.}
Let $\iota_i:K\hookrightarrow E$ send $\theta$ to
$\theta\zeta_i$, and put $\eta_i=1-[\mathcal O_Q(-D_i)]$.
Flat base change for the actual Gysin class and the projection
formula, followed by restriction to the regular $\mathcal T$,
give
\[
 c_\beta|_{\mathcal T}
    =\left.\sum_i\iota_i(\beta)\eta_i\right|_{\mathcal T}.
\]
At one complex embedding of $E$, the point regulators
$b_i\in\C/\R(2)$ run through all embeddings of $K$.
Choose their representatives in $i\R$. If $\beta\ne0$ in
$K_3(K)_{\Q}$, Borel injectivity \cite[Section~6]{BunkeTamme}
makes this vector nonzero.
Conjugate embeddings have opposite coefficients and real
embeddings have coefficient zero, so it is not diagonal.

The absolute multiplicative regulator of Bunke--Tamme
\cite[Theorem~2.31]{BunkeTamme} of the supported sum
is the real root class $2\pi i\sum_i b_i e_i$, modulo
$H^2(Q_s(\C),\R(3))$. It is nonzero on every permitted finite
open by the preceding injection. Here the analytic Deligne
target is $H^2(Q_s,\C)/H^2(Q_s,\R(3))$ because its complex
dimension is two and $F^3=0$; a nonzero real class survives this
quotient. Each individual regulator calculation spreads to a
regular finite type arithmetic model after inverting one
integer. The standard continuity
$K_3(\mathcal T)=\varinjlim_s K_3(Q_s)$ now rules out any
nonzero rational coefficient in the kernel of the Gysin map.

Finally $Z^d+p$ has one real root when $d$ is odd and none
when $d$ is even. Borel's rank calculation
\cite[Proposition~12.2]{Borel} therefore gives
$\dim_{\Q}K_3(K)_{\Q}=r_2(K)=\lfloor d/2\rfloor$.
Choosing actual lifts of a rational basis yields the asserted
number of independent classes of infinite order in the generic kernel.
All uses of localization, base change and regulators are the
standard inputs cited above \cite{Quillen,QuillenFinite,Borel,BunkeTamme,Weibel}.
\end{proof}

For example, $(p,d)=(2,3)$ gives
$\bigl(\Z_{(2)}[T,x,y]/(2+x^3+y^3+Txy^2)\bigr)_{(2,x,y)}$.
For any fixed $p$, arbitrarily large odd $d$ prime to $p$ give
the unbounded rank assertion without enlarging the residue
field $\F_p(T)$. 

\begin{remark}
The condition $d\ge3$ is used in the vertical-fiber calculation.
When $d=2$ and $p\ne2$, the fiber at $t=2$ has
$s^2+2s+1=(s+1)^2$, so the Kummer irreducibility argument is
unavailable and the vertical fiber splits. Localization can then
kill a root divisor difference. The proof therefore supplies no
degree-two analogue of the displayed tame family.
\end{remark}

\begin{corollary}
\label{d2ext:twodimtameallodd}
For $A_{p,d}$ as in Theorem~\ref{d2ext:twodimtamefamily} and every
$n\ge2$, the rational kernel of
$K_{2n-1}(A_{p,d})\to K_{2n-1}(\Frac A_{p,d})$ has dimension at least
\[
 \begin{cases}
 \lfloor d/2\rfloor,&n\text{ even},\\
 \lceil d/2\rceil-1,&n\text{ odd}.
 \end{cases}
\]
These bounds are realized by actual classes of infinite order.
For odd $d$ the bound is $(d-1)/2$ in every odd $K$-degree at least three.
\end{corollary}

\begin{proof}
Retain $K=\Q(\theta)$, $W_0$, the splitting field $E$, and the
$d$ embeddings $\iota_i:K\hookrightarrow E$ from the theorem.
Use the specified rational coefficient space
\[
 V_n=\ker\left(\operatorname{Tr}_{K/\Q}:
       K_{2n-1}(K)_{\Q}\longrightarrow K_{2n-1}(\Q)_{\Q}\right).
\]
For an actual $\gamma\in K_{2n-1}(K)$ set
$\beta=d\gamma-\operatorname{res}_{K/\Q}
                 \operatorname{Tr}_{K/\Q}\gamma$.
The identity $\operatorname{Tr}\operatorname{res}=d$ shows
that $\operatorname{Tr}\beta=0$ integrally. These actual
classes span $V_n$ rationally, because the displayed operator
acts as multiplication by $d$ on $V_n$.
Each $\beta$ lifts to $K_{2n-1}(W_0)$, since
$K_{2n-2}(\F_p)=0$. Localization, flat base change and the
projection formula give an actual class vanishing generically whose
test image is $\sum_i\iota_i(\beta)\eta_i$.

At a fixed complex embedding of $E$ let $b_i$ be the point
regulators, represented in the complementary real line
$i\R(n)$ to $\R(n)$ in $\C$. Base change for field transfer
gives $\sum_i\iota_i(\beta)=
\operatorname{res}_{E/\Q}\operatorname{Tr}_{K/\Q}\beta=0$;
therefore $\sum_i b_i=0$. If $\beta$ is nonzero rationally,
Borel injectivity \cite[Section~6]{BunkeTamme} makes the vector
$(b_i)$ nonzero and hence nondiagonal. Its line-bundle product is
$2\pi i\sum_i b_i e_i$. The preceding Betti persistence
keeps this root class nonzero on every permitted finite open.
The relevant analytic Deligne group is
\[
 H^3_{\mathcal D,\mathrm{an}}(Q_s,\R(n+1))
   =H^2(Q_s,\C)/H^2(Q_s,\R(n+1)),
\]
because $n+1\ge3$ and $Q_s(\C)$ is a surface.
Multiplication by $2\pi i$ identifies
$\C/\R(n)$ with $\C/\R(n+1)$, so the product regulator
is nonzero. The same continuity argument proves injectivity
of the supported map on $V_n$; applying it to every rational
relation gives independent actual representatives.

Rational transfer is surjective, since its composition with
restriction is multiplication by $d$. Borel's ranks
\cite[Proposition~12.2]{Borel} give
\[
 \dim V_n=
 \begin{cases}
 r_2(K),&n\text{ even},\\
 r_1(K)+r_2(K)-1,&n\text{ odd}.
 \end{cases}
\]
Here $r_2(K)=\lfloor d/2\rfloor$, while $r_1(K)$ is one for odd $d$ and zero for even $d$; the bounds follow.
\end{proof}

\section{A fixed arithmetic ring with infinite rank}
\label{d2ext:sec-strict-coefficients}

\begin{theorem}
\label{d2ext:twodimstrictcoefficientinfinite}
Fix a prime $p$ and an integer $d\ge3$ prime to $p$. Let
$V=\Z_{(p)}^{\mathrm{sh}}$, with residue field $\overline{\F}_p$,
and put
\[
 R=\left(V[T,x,y]/(p+x^d+y^d+Txy^{d-1})\right)_{(p,x,y)}.
\]
Then $R$ is a Noetherian regular local ring of dimension two,
with residue field $\overline{\F}_p(T)$. For every $n\ge2$,
\[
 \dim_{\Q}\ker\!\left(K_{2n-1}(R)_{\Q}
       \longrightarrow K_{2n-1}(\Frac R)_{\Q}\right)=\infty.
\]
The coefficient DVR $V$ is strictly henselian; no henselianity
of the displayed $R$ is asserted. This theorem concerns this
single uncompleted ring.
\end{theorem}

\begin{proof}
\emph{The coefficient and splitting rings.}
The strict henselization $V$ is a DVR with uniformizer $p$,
and $F=\Frac V$ is algebraic over $\Q$; see the Stacks Project,
\href{https://stacks.math.columbia.edu/tag/0AP3}{Tag 0AP3}
and the proof of
\href{https://stacks.math.columbia.edu/tag/09ET}{Tag 09ET}.
It is excellent because it is a DVR of fraction characteristic
zero
(\href{https://stacks.math.columbia.edu/tag/07QW}{Tag 07QW(4)}).
The simple roots of $Z^N-1$ in $\overline{\F}_p$ lift by
Hensel's lemma whenever $p\nmid N$, so $\mu_N\subset V$.
The ambient local ring defining $R$ has parameters $p,x,y$,
and the defining equation has initial term $p$. Thus $R$ has
the stated regularity, dimension and residue field.

Choose $\theta^d=-p$ and set
\[
 W=V[\theta],\qquad E_\infty=F(\theta),\qquad
 G=\operatorname{Gal}(E_\infty/F)=\{\sigma_\zeta:\zeta\in\mu_d\},
 \quad \sigma_\zeta(\theta)=\zeta\theta .
\]
Eisenstein shows that $W$ is a finite free DVR extension of
$V$, with uniformizer $\theta$, residue field $\overline{\F}_p$,
and degree $d$. Its displayed Galois action preserves $W$.
Moreover
\[
 D=\Spec(R/(x))=\Spec W[T]_{(\theta)}
\]
is regular. Pullback to $D$ followed by its Gysin map
$i_*:K_{2n-1}(D)\to K_{2n-1}(R)$ therefore constructs
actual classes vanishing on $R[1/x]$.

\emph{The geometric detector over $E_\infty$.}
Use the flat splitting map to
\[
 B=(R\otimes_VW)_{(\theta,x,y)},\qquad
 \mathcal T=\Spec B[1/\theta,1/y],
\]
and the smooth surface
\[
 Q=\Spec E_\infty[x,v,y^{\pm1}]/(xv+y^d+p),
 \qquad v=x^{d-1}+Ty^{d-1}.
\]
The geometric proof of Theorem~\ref{d2ext:twodimtamefamily}
applies over these coefficient rings. The required geometric calculations are as follows.
The blow-up still has the single smooth exceptional curve
\[
 C:\ X^d+Y^d+TXY^{d-1}=U^d
       \subset\mathbf P^2_{\overline{\F}_p(T)}.
\]
The same three charts prove its regularity along $C$, and
excellence and properness prove regularity everywhere. The
root strict transforms meet it transversely at
$[1:0:\bar\zeta]$. On $C\cap\{UY\ne0\}$ their Picard
classes have the lattice $\Z^d/\Z(1,\ldots,1)$:
the auxiliary surface
$XZ=1-Y^d$, $Y\ne0$, has that lattice, and every finite
$T=(Z-X^{d-1})/Y^{d-1}$ fiber is integral by the
$dq-sq'$ argument. These calculations already take place
over $\overline{\F}_p$ and do not require finite residue
constants.

Consequently, restriction of divisors to this exceptional open
gives
\[
 \Pic(Q)=\Z^d/\Z(1,\ldots,1)
       \lhook\joinrel\longrightarrow\Pic(\mathcal T),
\]
with generators $D_\zeta=(x,y-\zeta\theta)$.
The Picard calculation is unchanged by coefficient field
extension. Each $E_\infty$-prime divisor removed by a permitted
finite localization $Q_s$ has zero Picard class. Its geometric
components form one Galois orbit and all have the same class,
since the absolute Galois group over $E_\infty$ fixes its root
generators and the lattice is
torsion-free. Each component class is therefore zero.
The support sequence gives, for every embedding
$\jmath:E_\infty\hookrightarrow\C$,
\begin{equation}
 H^2(Q_\jmath(\C),\R)
      \lhook\joinrel\longrightarrow H^2((Q_s)_\jmath(\C),\R).
 \label{d2ext:strictcoefficientbetti}
\end{equation}
The Gysin sequence with complement $(\mathbf G_m)^2$ identifies
the root subspace with $\R^d/\R(1,\ldots,1)$, since the
angular classes of $x$ and $y$ have residues respectively the
diagonal vector and zero. This argument uses no regulator on
an infinite arithmetic base.

\emph{Finite actual coefficient spaces of unbounded rank.}
For $N>2$ divisible by $d$ and prime to $p$, choose compatible
roots of unity in $V$ and set
\[
 L_N=\Q(\mu_N)\subset F,\qquad K_N=L_N(\theta)\subset E_\infty .
\]
The chosen $p$-adic place of $L_N$ is unramified over $\Q_p$.
Eisenstein there gives $[K_N:L_N]=d$, and
$\operatorname{Gal}(K_N/L_N)$ is the restriction of $G$. Both fields
are totally imaginary. For fixed $n\ge2$, Borel's ranks
\cite[Proposition~12.2]{Borel} and restriction--transfer identities give
\begin{equation}
 V_{N,n}=
 \ker\!\left(\sum_{\sigma\in G}\sigma:
       K_{2n-1}(K_N)_{\Q}\to K_{2n-1}(K_N)_{\Q}\right),
 \qquad
 \dim_{\Q}V_{N,n}=\frac{(d-1)\varphi(N)}2.
 \label{d2ext:strictcoefficientrank}
\end{equation}
Indeed, the invariant subspace is exactly the restriction of
$K_{2n-1}(L_N)_{\Q}$, because
$\operatorname{res}\operatorname{Tr}=\sum_\sigma\sigma$ and
$\operatorname{Tr}\operatorname{res}=d$.
The two Borel ranks are $d\varphi(N)/2$ and $\varphi(N)/2$
in every degree under consideration.

Let $W_N=(\mathcal O_{K_N})_{\mathfrak P_N}$ be the DVR
at the prime induced by $W$. It maps to $W$, and $G$ preserves
it: above the chosen $L_N$ place the extension is totally
ramified and has a unique prime. Finite-field localization
gives
\[
 K_{2n-1}(W_N)_{\Q}\simeq K_{2n-1}(K_N)_{\Q},
\]
as well as actual surjectivity to the field, since
$K_{2n-2}$ of the finite residue field is zero.
Thus $V_{N,n}$ defines a rational linear map into
$K_{2n-1}(R)_{\Q}$ by coefficient pullback and $i_*$.
Every element of $V_{N,n}$ has an actual representative after clearing
denominators. If exact integral trace zero is wanted, replace
an actual lift $\gamma$ by
$d\gamma-\sum_{\sigma\in G}\sigma\gamma$; this operation
acts as multiplication by $d$ on $V_{N,n}$.

\emph{Detection, with both continuity passages explicit.}
Suppose $0\ne\beta\in V_{N,n}$ mapped to zero in
$K_{2n-1}(R)_{\Q}$. Borel injectivity \cite[Section~6]{BunkeTamme} supplies an embedding
$j_0:K_N\hookrightarrow\C$ at which its point regulator
is nonzero. Extend $j_0$ to an embedding
$\jmath:E_\infty\hookrightarrow\C$, using that
$E_\infty/K_N$ is algebraic.
For $\sigma\in G$, let $b_\sigma$ denote the point regulator
at $j_0\sigma$, represented in the complementary real line
$i\R(n)$ to $\R(n)$ in $\C$. These coefficients satisfy
\[
 b_1\ne0,\qquad \sum_{\sigma\in G}b_\sigma=0,
\]
so this vector is not diagonal. The detecting embedding is
chosen for this coefficient; no single embedding is required
to detect every element of $V_{N,n}$.

Put $\eta_\zeta=1-[\mathcal O_Q(-D_\zeta)]$. Flat Gysin
base change and the projection formula identify the image of
the supported class on $\mathcal T$ with the restriction of
\[
 \alpha=\sum_{\zeta\in\mu_d}(\sigma_\zeta\beta)\eta_\zeta
       \in K_{2n-1}(Q)_{\Q}.
\]
The assumed vanishing therefore makes $\alpha$ zero on
$\mathcal T$. First, $K$-theory continuity for the filtered
localization
$\mathcal O(\mathcal T)=\varinjlim_s\mathcal O(Q_s)$
gives a finite permitted denominator set $s$ for which
$\alpha|_{Q_s}=0$ rationally.
Second, $E_\infty$ is a filtered union of number fields.
All coefficients of these finitely many denominators, the
root ideals, and the chosen $K$-classes descend to a number
field $E_{\mathrm{ar}}\subset E_\infty$ containing $K_N$.
Continuity for this filtered coefficient extension implies
that the asserted vanishing already holds on the corresponding
finite open over some larger number field
$E_1\subset E_\infty$.

Apply the published absolute multiplicative regulator to that
finite stage, after spreading it to a regular finite-type
arithmetic model by inverting one integer. Evaluate at
$\jmath|_{E_1}$. The complex open of this finite model is exactly
$(Q_s)_\jmath$, and naturality gives the product class
\[
 2\pi i\sum_{\zeta\in\mu_d}b_{\sigma_\zeta}e_\zeta
 \quad\text{in}\quad
 H^3_{\mathcal D,\mathrm{an}}((Q_s)_\jmath,\R(n+1))
 =\frac{H^2((Q_s)_\jmath,\C)}
        {H^2((Q_s)_\jmath,\R(n+1))}.
\]
The equality holds because $n+1\ge3$ and the complex
dimension is two. Multiplication by $2\pi i$ identifies
$\C/\R(n)$ with $\C/\R(n+1)$. The non-diagonal root
vector and \eqref{d2ext:strictcoefficientbetti} make this
regulator nonzero, contradicting the finite-stage vanishing.
This proves injectivity on every $V_{N,n}$.
No commutation of analytic cohomology with an infinite
coefficient colimit has been used.

Finally, take for example $N=d^a$ as $a\to\infty$.
The dimensions in \eqref{d2ext:strictcoefficientrank} are
unbounded, while $R$ remains fixed. This proves infinite
rational rank in every degree $2n-1\ge3$.
All coefficient lifts, transfers and products used above are
actual Quillen $K$-theory operations
\cite{Quillen,QuillenFinite,Weibel,Borel,BunkeTamme}.
\end{proof}

For each $p$ one may fix $d=3$ if $p\ne3$, and $d=4$
if $p=3$. This gives one ring at each residue prime
with the stated infinite-rank conclusion in all higher odd
degrees. The construction here uses the displayed arithmetic local ring;
the complete rings considered below use a separate continuous detector.

\section{Visible coefficient rank and all higher odd degrees}
\label{d2ext:dimtwoVisibleHigherRank}

The elliptic contraction detects the $e_3$-component of the point
coefficient. We construct arbitrarily large families whose projections
to this component are $\mathbf Q_5$-linearly independent, and then use
them on a single complete local ring.
We write $\mathcal K(-)=K(-)^\wedge_5[1/5]$ for spectral
$5$-completion followed by inversion of $5$.

For the continuous detector, write $W'=W(\overline{\mathbf F}_5)[\rho]$
for the coefficient ring denoted by $R$ in
Proposition~\ref{d2:prop:weighted}, and put $E_\infty=W'[1/5]$.
The automorphisms $\sigma_g$ fix $W(\overline{\mathbf F}_5)$.
We use the normalization of \eqref{d2:eq:linear-weight}: the bounded
functional has $\ell(\omega)=1$, and the polarization fixes the common
sign. Trace, contraction, and formal replacement are $E_\infty$-linear
on the differential quotients. The same four-term calculation therefore
gives, for every $b\in E_\infty$ with $\sigma_{-1}b=-b$,
\[
 L\operatorname{Tr}_h\left(\sum_{g\in\mathbf F_5^\times}
                  \sigma_g(b)\gamma_g\right)
       =4e_3(b)[\widetilde Q^*\eta].
\]
Here the operators on the scaled model are transported by formal
replacement, as in the proof of Proposition~\ref{d2:prop:weighted}.
The vector $[\widetilde Q^*\eta]$ is nonzero in the ordinary Cohen
quotient by Proposition~\ref{d2:prop:graph}. This formula applies to
coefficients in the finite local fields below and requires only their
$e_3$-projections. It uses the linear operators over the fixed scalar
field, with no passage to a colimit of completed $K$-theory.

For each complete coefficient discrete valuation ring $O$ below and
$n\geq2$, let $r_{n,O}(\beta)$ denote the scalar whose AMMN point
character is $t^{1-n}r_{n,O}(\beta)$, with $|t|=-2$, using the point
comparison in Lemma~\ref{d3:lem:formalhigherTCcup}.

\begin{lemma}
\label{d2ext:dimtwoVisibleHigherCoefficientRank}
Fix $n\geq2$ and $m\geq2$. Put
\[
 V_m=W(\mathbf F_{5^m}),\quad F_m=V_m[1/5],\quad
 O_m=V_m[\rho],\quad E_m=F_m(\rho),\quad \pi=\rho-1,
\]
where $\rho$ is a primitive fifth root of unity.
Let $\sigma_g$ fix $F_m$ and send $\rho$ to $\rho^g$.
Let $\omega_T:\mathbf F_5^\times\to\mathbf Z_5^\times$ be the
Teichm\"uller character, and put
\[
 e_3=\frac14\sum_{g=1}^4\omega_T(g)^{-3}\sigma_g,\qquad
 \tau_3=\sum_{g=1}^4\omega_T(g)^{-3}\rho^g.
\]
There are actual classes
$\beta_j\in K_{2n-1}(O_m)$, indexed by $1\leq j\leq m-1$, each of the form
$(1-\sigma_{-1})\beta_{0,j}$, with the following properties:
their reductions in $K_{2n-1}(\mathbf F_{5^m})$ vanish integrally,
and their AMMN point coordinates $b_j=r_{n,O_m}(\beta_j)\in E_m$
have $\mathbf Q_5$-linearly independent projections $e_3b_j$.
\end{lemma}

\begin{proof}
Define
$S_j(z)=(z\,d/dz)^j(1-z)^{-1}\in\mathbf Z[z,(1-z)^{-1}]$.
The expansion and continuation argument in
the proof of Proposition~\ref{d2:prop:coefficient}, now with
$\binom{n+j-1}{j}$ in place of $j+1$, gives
\[
 F_n(z)\coloneqq \operatorname{Li}_n(z)-5^{-n}\operatorname{Li}_n(z^5)
 =\sum_{j\geq0}(-1)^j\binom{n+j-1}{j}5^j S_j(z^5)
             \sum_{a=1}^4\frac{z^a}{a^{n+j}}.
\]
The series is uniformly convergent on $|z|\leq1$,
$|1-z^5|=1$, and agrees there with Coleman continuation by
\cite[Theorem~2.7(3)]{BesserDeJeuAlgorithm}.
The fifth power acts on the polylogarithm variable.

For a Teichm\"uller root
$\zeta\in\mu_{5^m-1}(V_m)\setminus\{1\}$, the value
$(\zeta\rho)^5=\zeta^5$ is fixed by every $\sigma_g$.
The exact Gauss identity
$e_3(\rho^a)=\frac14\omega_T(a)^3\tau_3$ therefore gives
\begin{equation}
\label{d2ext:dimtwoVisibleHigherDepletion}
\begin{split}
 U_n(\zeta)
 &\coloneqq \frac{4e_3\operatorname{Li}_n(\zeta\rho)}{\tau_3}\\
 &=\sum_{j\geq0}(-1)^j\binom{n+j-1}{j}5^j S_j(\zeta^5)
        \sum_{a=1}^4\frac{\zeta^a\omega_T(a)^3}{a^{n+j}}
       \in V_m.
\end{split}
\end{equation}
The terms with $j\geq1$ belong to $5V_m$.  Thus, writing
$t=\overline\zeta$, its reduction is
\[
 h_{3,n}(t)
   =\frac{P_n(t)}{1-t^5},\qquad
 P_n(t)=\sum_{a=1}^4a^{3-n}t^a\in\mathbf F_5[t],
 \qquad t\in\mathbf F_{5^m}^{\times}\setminus\{1\}.
\]
The polynomial $P_n$ is nonzero, has coefficient one at $t$,
and has degree four.  Hence $h_{3,n}$ is nonconstant and has
rational-map degree at most five.  Every finite fiber has at most
five elements, so
\[
 \#h_{3,n}\bigl(\mathbf F_{5^m}^{\times}\setminus\{1\}\bigr)
 \geq\left\lceil\frac{5^m-2}{5}\right\rceil=5^{m-1}.
\]
Its $\mathbf F_5$-span has dimension at least $m-1$.
Choose $m-1$ independent image values and corresponding
Teichm\"uller roots $\zeta_j$.  Reduction of a normalized
$\mathbf Q_5$-linear relation shows that the $U_n(\zeta_j)$
are $\mathbf Q_5$-independent.  Since $\tau_3\ne0$, the values
$e_3\operatorname{Li}_n(\zeta_j\rho)$ are independent as well.

These are regulators of actual rational number-field classes.
Indeed, all $\zeta_j\rho$ lie in
$L=\mathbf Q(\mu_{5^m-1},\rho)$ and are roots of unity.
At the chosen prime above $5$ they and their complements are
units, since $\overline\zeta_j\ne1$.  The rational symbols
$[\zeta_j\rho]_n$ are cycles, and
\cite[Theorem~1.10(2)]{BesserDeJeu} maps them to a common
nonzero rational multiple of
$L_{\mathrm{mod},n}(\zeta_j\rho)
 =\operatorname{Li}_n(\zeta_j\rho)$.
The point normalization in \eqref{d3:eq:pointnormalization}
identifies the normalized Besser--de Jeu coordinate with the
ordinary syntomic coordinate, so the scalar is common to the
chosen roots. The number-field case of the cited theorem applies
unconditionally. The local completion at the chosen prime is $E_m$.

Localization at that coefficient prime lifts these rational
classes to its integer ring, since
$K_{2n-2}(\mathbf F_{5^m})=0$.
Clear one common integer denominator and then form the actual
differences
$\beta_j=(1-\sigma_{-1})\beta_{0,j}$, mapping them to $K_{2n-1}(O_m)$.
Their residue reductions are zero because $\sigma_{-1}$ fixes
the residue field.  Applying $e_3$ to their syntomic regulators
multiplies each of the preceding values by the same nonzero
rational scalar, including the factor two from the difference.

The ramified point comparison
\cite[Example~2.2.4 and Propositions~2.2.7, 2.2.9, 2.3.4]{HuberKings}
is Galois-equivariant and identifies syntomic and \'etale point
cohomology in every twist $n\geq2$.  Thus the $e_3$-projections
of the \'etale Chern values are $\mathbf Q_5$-independent.
Lemma~\ref{d3:lem:cherncompletionfactor} gives a natural
Galois-equivariant $\mathbf Q_5$-linear map
\[
 \pi_{2n-1}\bigl(K(O_m)^\wedge_5[1/5]\bigr)
       \longrightarrow H^1(E_m,\mathbf Q_5(n))
\]
through which the actual Chern class factors. The finite Chern
maps are applied over $E_m$, where $5$ is invertible, after the
canonical projections to $K_{2n-1}(O_m;\mathbf Z/5^\nu)$ and
restriction to the field. Compatibility with reduction gives the
map to the inverse limit of $H^1$. The finite $H^0$ groups are
Mittag--Leffler, identifying this target with continuous cohomology.
This construction uses the canonical source projections and
allows a nonzero source $\lim^1$.

The point case of the AMMN comparison identifies the source with
$t^{1-n}E_m$, naturally under $\sigma_g$.  A relation among the
$e_3b_j$ would therefore induce the same relation among the
independent projected Chern values.  This proves the claim
without identifying the AMMN and polylogarithm coordinates.
\end{proof}

\begin{remark}
The rational-map degree bound cannot uniformly be replaced by four.
When $n\equiv3\pmod4$ one has
\[
 h_{3,n}(t)=(t-1)^{-5}-(t-1)^{-1}.
\]
Its image lies in the trace-zero hyperplane of
$\mathbf F_{5^m}/\mathbf F_5$.  The lower bound $m-1$ accommodates
this actual loss of one dimension.
\end{remark}

Put
\[
 k=\overline{\mathbf F}_5,\qquad W=W(k),\qquad
 V=\widehat{W[T]_{(5)}}^{\,(5)},\qquad
 R=V[[x,y]]/(\Phi_5(1+y)+x^4+Txy^3).
\]

\begin{theorem}
\label{d2ext:dimtwoCompleteAllOddInfiniteRank}
The ring $R$ is a complete regular local domain of dimension two,
with maximal ideal $(x,y)$ and residue field $k(T)$.
For every $n\geq2$,
\[
 \dim_{\mathbf Q}\ker\left(
 K_{2n-1}(R)_{\mathbf Q}\longrightarrow
 K_{2n-1}(\operatorname{Frac}R)_{\mathbf Q}\right)=\infty.
\]
For each fixed $n,m\geq2$, the kernel contains $m-1$ independent
classes represented by actual integral classes that vanish on
$R[1/x]$ integrally.  The ring $R$ is independent of $n,m$.
\end{theorem}

\begin{proof}
The regularity argument is the same linear-term calculation as in
Section~\ref{d2:sec:preliminaries}.
Set $W'=W[\rho]$ and $E_\infty=W'[1/5]$.
For fixed $n,m$, take the actual coefficients of
Lemma~\ref{d2ext:dimtwoVisibleHigherCoefficientRank}.
Naturality of the AMMN point calculation sends their coordinates
under $E_m\hookrightarrow E_\infty$; hence their visible
projections remain $\mathbf Q_5$-independent.

Use the same finite parameter cover and scaled formal test ring
$B$ as in \eqref{d2:eq:test-algebra}, now over $V$.
It is the completion of an essentially smooth $W'$-algebra
of relative dimension two.  The spectrum comparison underlying
Lemma~\ref{d3:lem:essentiallysmoothAMMN} gives, for each fixed $n$,
the relative character to $HC_{2n-2}^{\mathrm{cts}}(B/W';\mathbf Q_5)$.
In its finite mixed-HKR model the top-form summand is
\begin{equation}
\label{d2ext:dimtwoHigherTopColumn}
 t^{2-n}\left(
 \Omega^2_{B/W'}[1/5]/d\Omega^1_{B/W'}[1/5]\right).
\end{equation}
Indeed its cycles in this summand are all two-forms, and its
boundaries are precisely $t^{2-n}d\Omega^1$; form degrees above
two vanish.

The zero residue reductions of the coefficients provide actual
relative lifts at the point. Pull them to $B$ and multiply by
$\eta_g=1-[J_g]$.
With the common scalar $\lambda_{W'}\in E_\infty^\times$ of
Proposition~\ref{d2:prop:cyclic}, the negative-cyclic module product is
\[
 (t^{1-n}\sigma_g(b_j))\,(\lambda_{W'}t\gamma_g)
       =\lambda_{W'}t^{2-n}\sigma_g(b_j)\gamma_g.
\]
Rank zero removes the scalar term, higher Chern forms vanish in
dimension two, and the cyclic-shuffle terms vanish because the
point coefficient is a scalar relative to $W'$.
The scalar-bar calculation applies in every cyclic column, including
the negative columns that occur for $n>2$.
The scalar occurs to the first power because it belongs to the
first Chern component.

Trace and contraction in Proposition~\ref{d2:prop:weighted} are
$E_\infty$-linear. They send the coefficient of
\eqref{d2ext:dimtwoHigherTopColumn} to
\[
 4\lambda_{W'}e_3(b_j)[\widetilde Q^*\eta].
\]
The fixed vector $[\widetilde Q^*\eta]$ is nonzero over $E_\infty$
after the generic Cohen completion, by
Lemma~\ref{d2:lem:scalar} and Proposition~\ref{d2:prop:graph}.
Multiplication by the common nonzero $\lambda_{W'}$ is an
injective $\mathbf Q_5$-linear map on this coefficient field.
Thus these character images are $\mathbf Q_5$-independent,
without any rationality or Galois-invariance assumption on
$\lambda_{W'}$.

Write $C=B/\pi$.  As before, $C$ is a filtered colimit of smooth
affine surfaces over $k$.  By
\cite[Theorem~8.4]{GeisserLevineCharP},
$\pi_q\mathcal K(C)=0$ for every $q\geq3$.
The relative and regular-divisor localization sequences therefore
give, in the present degree,
\[
 \pi_{2n-1}\mathcal K(B,(\pi))
 \hookrightarrow\pi_{2n-1}\mathcal K(B)
 \hookrightarrow\pi_{2n-1}\mathcal K(B[1/\pi]).
\]
These use the vanishings in degrees $2n$ and $2n-1$, respectively.
These are vanishings of the rationalized completed spectrum.

The source divisor $x=0$ has ring $V[\rho]$.
Pull each $\beta_j$ to that divisor and push forward to define
$c_j\in K_{2n-1}(R)$.  Localization gives $c_j|_{R[1/x]}=0$
integrally.  Finite flat coefficient extension $V\to V[\rho]$,
followed by inversion of $5$, splits the divisor into the four
root components with coefficients $\sigma_g(\beta_j)$.
Pullback of their line-bundle classes under the scaled map gives
\[
 c_j|_{B[1/\pi]}
   =\sum_g\sigma_g(\beta_j)(1-[J_g])\big|_{B[1/\pi]}.
\]
This is exactly the actual product detected above; the scaled map
itself is not required to be flat.
Consequently the $c_j$ are rationally independent.
Since $m$ is arbitrary while $R$ is fixed, the kernel has infinite
dimension for this $n$.  Applying the argument separately to each
$n\geq2$ proves the theorem.
\end{proof}

\section{One complete two-dimensional ring in every degree at least three}
\label{d2ext:dimtwoAllDegreesFixedRing}

The additional parameter below is a unit of the source ring.
Its logarithmic differential raises the degree of the detected
class by one.  Detecting the even products also proves independence
of their odd factors, including degree three.

\begin{lemma}
\label{d2ext:dimtwoFullCohenLaurentResidue}
Let $W'$ be the coefficient discrete valuation ring, with
uniformizer $\pi$, used in the quartic construction.  Let $D_0$
be a smooth affine $W'$-curve whose special fiber is integral,
and put
\[
 A_T=\widehat{(D_0)_{(\pi)}}^{\,(\pi)},\qquad
 A_{TS}=\widehat{(D_0[S])_{(\pi)}}^{\,(\pi)}.
\]
Here localization at $(\pi)$ inverts every element with nonzero
reduction.  Let $B_T$ be the $\pi$-adic completion of a smooth
affine relative curve over $A_T$, defined by a model over $D_0$,
and let $B_{TS}$ be the corresponding completed model over
$A_{TS}$.  Suppose their continuous differentials over $W'$
are finite projective of ranks two and three, respectively.
There is a $W'$-linear operator
\[
 \operatorname{Res}_S:
 \Omega^{r,\mathrm{cts}}_{B_{TS}/W'}
       \longrightarrow \Omega^{r-1,\mathrm{cts}}_{B_T/W'}
\]
which is integral and $\pi$-adically continuous and satisfies
\[
 \operatorname{Res}_S(d\alpha)=-d\operatorname{Res}_S(\alpha),
 \qquad
 \operatorname{Res}_S(\theta\wedge d\log S)=\theta
       \quad\text{for every }\theta\in\Omega^2_{B_T/W'}.
\]
In particular it sends actual exact three-forms to actual exact
two-forms. The induced map
\[
 \frac{\Omega^2_{B_T/W'}[1/5]}
      {d\Omega^1_{B_T/W'}[1/5]}
 \longrightarrow
 \frac{\Omega^3_{B_{TS}/W'}[1/5]}
      {d\Omega^2_{B_{TS}/W'}[1/5]},
 \qquad [\theta]\longmapsto[\theta\wedge d\log S],
\]
is injective; these are ordinary continuous quotients.
\end{lemma}

\begin{proof}
For every $a\geq1$ there is a compatible ring map
\[
 B_{TS}/\pi^a
       \longrightarrow (B_T/\pi^a)((S)).
\]
Indeed, if $g\in D_0[S]$ has nonzero reduction, its image modulo
$\pi$ is a nonzero polynomial over the field
$\operatorname{Frac}(D_0/\pi)$.  It is therefore a unit in that
field's Laurent series ring.  The ideal generated by $\pi$ in
$(B_T/\pi^a)((S))$ is nilpotent, so this inverse lifts.  This
argument applies to every primitive denominator, including
denominators whose literal lowest coefficient is not a unit.
The lifts of the inverses are unique and commute with reduction
in $a$.  The maps extend to the completed algebra because the
inverse-limit target is complete.  Thus the relevant target is
\[
 \mathcal L=\varprojlim_a(B_T/\pi^a)((S)).
\]
The negative Laurent bound may depend on $a$, which is why this
inverse-limit target is used.

Use Laurent series of continuous base forms, with differential
defined coefficientwise and with the usual $S$-derivative.
These operations are compatible modulo $\pi^a$ and preserve
the $\pi$-adic filtration.  The universal continuous derivation
of $B_{TS}$ therefore gives a map into this Laurent differential
complex.  Finite projectivity of the source form modules
justifies extension to their completions and exterior powers.

Write a form of degree $r$ in the target as
$\alpha(S)+\beta(S)\wedge dS$, with base-form degrees $r$
and $r-1$.  Define
\[
 \operatorname{Res}_{S,r}
       \bigl(\alpha(S)+\beta(S)\wedge dS\bigr)
      =(-1)^{r-1}[S^{-1}]\beta(S).
\]
Coefficient extraction is compatible with reduction and is
integral.  The coefficient of $S^{-1}$ in the derivative of a
Laurent series is zero.  It follows directly that
$\operatorname{Res}_{S,r+1}d=-d\operatorname{Res}_{S,r}$.
For $r=3$ its sign is positive, which gives the second displayed
identity.  Inverting $5$ preserves these identities.  All
quotients in the conclusion use the ordinary image of the
differential.
\end{proof}

Put
\[
 k=\overline{\mathbf F}_5,\qquad W_0=W(k),\qquad
 V_{TS}=\widehat{W_0[T,S]_{(5)}}^{\,(5)},
\]
and
\[
 R_{TS}=V_{TS}[[x,y]]/
       \bigl(\Phi_5(1+y)+x^4+Txy^3\bigr).
\]

\begin{theorem}
\label{d2ext:dimtwoCompleteAllDegreesInfiniteRank}
The ring $R_{TS}$ is a complete regular local domain of dimension
two, with maximal ideal $(x,y)$ and residue field $k(T,S)$.
For every integer $j\geq3$,
\[
 \dim_{\mathbf Q}\ker\!\left(
 K_j(R_{TS})_{\mathbf Q}\longrightarrow
 K_j(\operatorname{Frac}R_{TS})_{\mathbf Q}\right)=\infty.
\]
For every $n,m\geq2$, the odd degree $2n-1$ and the even degree
$2n$ kernels contain $m-1$ independent classes represented by
actual integral classes which vanish on $R_{TS}[1/x]$ integrally.
The ring $R_{TS}$ is independent of $n$ and $m$.
\end{theorem}

\begin{proof}
The ring $W_0[T,S]_{(5)}$ is a discrete valuation ring with
uniformizer $5$ and residue field $k(T,S)$.  Its completion
$V_{TS}$ has the same properties.  The defining equation of
$R_{TS}$ is congruent to $5$ modulo $(5,x,y)^2$ in the regular
local ambient ring of dimension three.  Its quotient is
therefore regular of dimension two, and is complete with
maximal ideal $(x,y)$.  In particular it is a domain.
The element $S$ is a unit of $R_{TS}$.

Set $W'=W_0[\rho]$, $\pi=\rho-1$, and
$E_\infty=W'[1/5]$.  Fix $n,m\geq2$, and take the $m-1$
actual coefficients $\beta_j$ from
Lemma~\ref{d2ext:dimtwoVisibleHigherCoefficientRank}.
Map them to $K_{2n-1}(W')$ and write $b_j\in E_\infty$
for their natural AMMN point coordinates.  Their visible
projections $e_3b_j$ remain $\mathbf Q_5$-independent.
Their reductions vanish integrally, since each coefficient
is an actual difference $(1-\sigma_{-1})\beta_{0,j}$.

Choose the same finite parameter cover, independent of $S$,
as in the one-parameter quartic detector.  In the notation of
Lemma~\ref{d2ext:dimtwoFullCohenLaurentResidue}, it gives the old
coefficient ring $A_T$ and the new one $A_{TS}$.  The latter
receives $V_{TS}[\rho]$: every primitive polynomial in $T,S$
has nonzero image in $k(U)(S)$, and hence is a unit.  Completion
extends this map.  Put
\[
 B_{TS}=A_{TS}\langle X,Y\rangle/
 \left(X^4+\frac{\Phi_5(1+\pi Y)}{\pi^4}+TXY^3\right),
 \qquad C_{TS}=B_{TS}/\pi.
\]
The old algebra $B_T$ is given by the same formula over $A_T$.
They are the completions of essentially smooth $W'$-algebras
of relative dimensions three and two, respectively.
Smoothness holds near the entire special fiber; localizing a
model at an element congruent to one modulo $\pi$ leaves the
completion unchanged.  In particular $B_{TS}$ and $C_{TS}$
are regular and $\pi$ is a non-zero-divisor.

For $g\in\mathbf F_5^\times$, let
\[
 a_g=\frac{\rho^g-1}{\pi},\qquad
 J_g=(X,Y-a_g),\qquad
 \eta_g=1-[J_g],\qquad \gamma_g=-c_1(J_g).
\]
These invertible ideals are pulled from $B_T$.  We can therefore
represent their Chern classes by two-forms pulled from the old
continuous complex, independently of $S$.

The actual residue sequence at the point supplies actual
relative lifts of all $\beta_j$; here
$K_{2n}(k)=0$ by Quillen's finite-field computation and filtered
colimits.  Pull these lifts to $B_{TS}$ and form their relative
products
\[
 \xi_j=\sum_g\sigma_g(\beta_j)\eta_g[S]
       \quad\text{in }K_{2n}(B_{TS},(\pi)),
\]
where the coefficient symbols in this formula denote the
specified relative lifts.

Apply the essentially smooth relative fiber comparison of
Lemma~\ref{d3:lem:essentiallysmoothAMMN} in each fixed degree.
Its $K(B_{TS})$-module compatibility gives the negative-cyclic
action on the ordinary cyclic target.  In relative dimension
three the normalized finite mixed-HKR map uses only the
denominators $2$ and $6$, both units at $5$.
Let $\lambda_{W'}\in E_\infty^\times$ be the common scalar of
Proposition~\ref{d2:prop:cyclic}.
The line bundles are pulled from the rank-two algebra $B_T$,
so naturality pulls back their negative-cyclic characters
$\lambda_{W'}t\gamma_g$.
The point coefficient is $t^{1-n}\sigma_g(b_j)$, and the
unit character starts with $d\log S$.
The top character component is
\[
 \chi_{B_{TS}}(\xi_j)_{\mathrm{top}}
    =\lambda_{W'}t^{2-n}
       \sum_g\sigma_g(b_j)\gamma_g\wedge d\log S
       \quad\text{in }HC^{\mathrm{cts}}_{2n-1}(B_{TS}/W';
                                                    \mathbf Q_5).
\]
Here $|t|=-2$.  This formula also holds when the initial cyclic
column is negative.  To check it at chain level, the rank-zero
line has underlying Hochschild degrees at least two; the
lowest unit term has degree one.  Every higher line or unit
term, or cyclic-shuffle correction in their product, has final
Hochschild degree at least five, and is killed by the
dimension-three HKR map.  The scalar point cup has no
cyclic-shuffle correction because its scalar lies in the
relative base.  Thus precisely the displayed ordinary
shuffle survives.

The top-form summand is
\[
 t^{2-n}\bigl(
 \Omega^3_{B_{TS}/W'}[1/5]/
                  d\Omega^2_{B_{TS}/W'}[1/5]\bigr).
\]
All its three-forms are cycles, and its boundaries are precisely
the indicated ordinary exact forms.  Apply
Lemma~\ref{d2ext:dimtwoFullCohenLaurentResidue} to their coefficients.
It recovers
$\lambda_{W'}\sum_g\sigma_g(b_j)\gamma_g$ on $B_T$.
The weighted trace and elliptic contraction of
Proposition~\ref{d2:prop:weighted} are $E_\infty$-linear and send
these classes to
\[
 4\lambda_{W'}e_3(b_j)[\widetilde Q^*\eta].
\]
The fixed base vector $[\widetilde Q^*\eta]$ survives the old generic
Cohen completion by Lemma~\ref{d2:lem:scalar}, and remains
nonzero after the finite constant extension adjoining $\rho$ by
Proposition~\ref{d2:prop:graph}.
Multiplication by the common nonzero $\lambda_{W'}$ preserves
$\mathbf Q_5$-linear independence.
Consequently the characters of the $\xi_j$ are
$\mathbf Q_5$-independent.

The special fiber $C_{TS}$ is a filtered colimit of smooth
affine threefolds over the perfect field $k$.  Geisser--Levine
\cite[Theorem~8.4]{GeisserLevineCharP} gives
\[
 K_q(C_{TS};\mathbf Z/5^\nu)=0
      \quad(q>3,\ \nu\geq1).
\]
The Milnor sequence for spectral completion, using degrees
$q$ and $q+1$, therefore gives
$\pi_q\mathcal K(C_{TS})=0$ for $q\geq4$.
For the even degree under consideration the relative and
regular-divisor localization sequences give
\[
 \pi_{2n}\mathcal K(B_{TS},(\pi))
   \hookrightarrow\pi_{2n}\mathcal K(B_{TS})
   \hookrightarrow\pi_{2n}\mathcal K(B_{TS}[1/\pi]).
\]
The first injection uses the vanishing in degree $2n+1$,
and the second that in degree $2n$; both are valid for $n\geq2$.
Here $\mathcal K=K^\wedge_5[1/5]$ denotes the rationalized
spectrally completed $K$-theory.

It remains to match these detected even classes with actual
classes of the fixed source.  Its regular Cartier divisor
$i:x=0\hookrightarrow\operatorname{Spec}R_{TS}$ has ring
$V_{TS}[\rho]$.  Pull the coefficients to this divisor and set
\[
 c_{n,j}=i_*(\beta_j)\in K_{2n-1}(R_{TS}),\qquad
 a_{n,j}=c_{n,j}[S]\in K_{2n}(R_{TS}).
\]
Both classes vanish on $R_{TS}[1/x]$ integrally.
Make the finite flat coefficient extension
$V_{TS}\to V_{TS}[\rho]$, and then invert $5$.
The divisor $x=0$ splits into its four root components with
coefficients $\sigma_g(\beta_j)$.
Their length-one Cartier resolutions and the projection
formula express the resulting odd class as
$\sum_g\sigma_g(\beta_j)(1-[I_g])$, where
$I_g=(x,y-(\rho^g-1))$.
Apply the continuous scaling
\[
 x\longmapsto\pi X,\qquad y\longmapsto\pi Y,\qquad
 R_{TS}\longrightarrow B_{TS}.
\]
It exists because the defining equation maps to $\pi^4$ times
the test equation and the maximal ideal maps into $(\pi)$.
After inverting $\pi$, pullback of the invertible ideals $I_g$
gives $J_g$. The resulting actual class is
\[
 a_{n,j}|_{B_{TS}[1/\pi]}
   =\sum_g\sigma_g(\beta_j)\eta_g[S]\big|_{B_{TS}[1/\pi]}
       \quad\text{in actual }K_{2n}.
\]
Only the initial coefficient extension is asserted to be flat;
the subsequent substitution pulls back line bundles.
This equality and the two completed injections prove that
the actual classes $a_{n,j}$ are rationally independent.

Finally, a rational relation among the actual odd classes
$c_{n,j}$ would, after multiplication by the same actual unit
class $[S]$ in the same ring, give a relation among the
$a_{n,j}$.  Thus the odd classes are independent as well.
This argument includes $n=2$ and requires no vanishing of
$\pi_3\mathcal K(C_{TS})$.
Letting $m$ grow proves infinite rank separately in degrees
$2n-1$ and $2n$.  These pairs cover every $j\geq3$.
\end{proof}

\begin{remark}
No general injectivity of multiplication by $[S]$ is used:
it is proved only on the displayed finite spans by detecting
their products.  In particular the odd degree-three conclusion
does not use a degree-three generic-localization injection for
the new special fiber.  The residue field in this fixed-ring
statement is $\overline{\mathbf F}_5(T,S)$.
\end{remark}


% End input: sections/05-two-dimensional-families.tex


\appendix

% Begin input: sections/06-comparisons.tex
\section{Low degrees and comparison with other Gersten statements}\label{d3:sec:comparisons}
We use the arithmetic-quadric notation of Section~\ref{d3:sec:geometry} in this appendix.

\subsection{A low-degree stress test}
The analogous $K_1$ construction does not supply a counterexample.
Its product regulator would have degree three and twist two.
Here the filtration is the analytic truncation of the holomorphic
de Rham complex, which need not vanish in filtration degree two
on a complex surface.

For a smooth affine complex surface $U$, write $X=U^{\mathrm{an}}$.
An affine embedding realizes $X$ as a closed complex submanifold
of $\C^N$. Coordinate functions separate points and give local
coordinates. They bound the holomorphic hull of each compact set
inside a compact coordinate polydisc; the hull is closed there
and hence compact. Together with second countability, these
properties make $X$ Stein.
Cartan's theorem~B
\cite[\S5, proof of Theorem~B, pp.~19-14--19-15]{CartanStein1952}
gives $H^q(X,\Omega_X^a)=0$ for every $q>0$.
The holomorphic Poincar\'e lemma follows on coordinate balls by
radial integration and identifies $\C_X$ with the holomorphic
de Rham complex. Its hypercohomology is therefore computed by
global holomorphic forms, in degrees zero through two. Thus
\[
 H^2(X,\C)=\Gamma(X,\Omega_X^2)/d\Gamma(X,\Omega_X^1),
 \qquad H^3(X,\C)=0.
\]
For
$F^2_{\mathrm{an}}\Omega_X^\bullet=\Omega_X^{\ge2}
=\Omega_X^2[-2]$, the map from degree-two hypercohomology to
$H^2(X,\C)$ is the displayed quotient map and is surjective.
Moreover,
$\mathbb H^3(X,F^2_{\mathrm{an}}\Omega_X^\bullet)
=H^1(X,\Omega_X^2)=0$.
The constant real sheaf $\R(2)$ is a direct summand of the
underlying real sheaf $\C$, so $H^3(X,\R(2))=0$ as well.
The analytic Deligne cone sequence contains
\[
\begin{split}
 H^2(X,\R(2))\oplus
 \mathbb H^2(X,F^2_{\mathrm{an}}\Omega_X^\bullet)
 &\longrightarrow H^2(X,\C)\longrightarrow\HD^3(X,\R(2))\\
 &\longrightarrow H^3(X,\R(2))\oplus
 \mathbb H^3(X,F^2_{\mathrm{an}}\Omega_X^\bullet).
\end{split}
\]
Its first arrow is surjective and its final group is zero.
Hence $\HD^3(U,\R(2))=0$.
This calculation uses the analytic truncation just defined;
it makes no corresponding assertion about the logarithmic
mixed-Hodge filtration.
The elementary quotient in Lemma~\ref{d3:lem:Deligne} is
unavailable in twist two. This agrees with the fact that $K_1$
of a commutative local ring is its unit group and injects into
the fraction field.

There is also an explicit algebraic version of this test. On the split quadric in coordinates
\[
uv=t(1-t),
\]
the matrix
\[
e=\begin{pmatrix}t&u\\v&1-t\end{pmatrix}
\]
is a rank-one idempotent. Up to choosing the ruling, its image is the line bundle $P=\OO(D_+)$. For $\alpha\in E^\times$, $\alpha\neq1$, the determinant-one matrix
\[
M=\begin{pmatrix}\alpha^{-1}&0\\0&1\end{pmatrix}
\bigl(I+(\alpha-1)e\bigr)
\]
represents $\alpha([P]-1)$ in $K_1$. Its upper-left entry is $s/\alpha$ with $s=1+(\alpha-1)t$. On $D(s)$ Gaussian elimination makes $M$ elementary. The removed curve has
\[
t=-\frac1{\alpha-1},\qquad uv=-\frac{\alpha}{(\alpha-1)^2}\neq0,
\]
so it is a geometrically irreducible principal curve isomorphic to $\mathbb G_m$. Hence ordinary Betti injectivity alone cannot protect a $K_1$ product. The condition $3>\dim_{\C}U=2$ in the main proof is essential.

\subsection{Equal characteristic and other localizations}
Replacing the arithmetic DVR by $k[t]_{(t)}$ and adjoining $u$
with $u^2=-t$ does not reproduce the same coefficient class.
For every field $k$, restriction gives an isomorphism
\[
 K_3(k)_{\Q}^{(2)}
       \xrightarrow{\sim}K_3(k(u))_{\Q}^{(2)}.
\]
To prove this, put $F=k(u)$ and compare the Milnor and Quillen
rational-function residue sequences
\cite[III, Theorem~7.4; V, Corollary~6.7.1]{Weibel}.
They have respective kernels $K_3^M(k)$ and $K_3(k)$ and
surjective residue maps to
$\bigoplus_P K_2^M(k[u]/P)$ and
$\bigoplus_P K_2(k[u]/P)$, where $P$ runs over monic
irreducible polynomials.
The right-hand terms agree by Matsumoto's theorem
\cite[III, Theorem~6.1]{Weibel}.

Their residue maps are compatible with the symbol maps.
At each polynomial valuation, a product of unit classes extends
over the valuation ring and has zero boundary, while for a
uniformizer $\pi$ and units $a,b$ the boundary-product formula gives
\[
 \partial\{\pi,a,b\}=\{\bar a,\bar b\}.
\]
Multilinearity and $\{\pi,\pi\}=\{-1,\pi\}$ reduce every Milnor
symbol to unit symbols and symbols with exactly one uniformizer.
These formulas therefore give its tame residue
\cite[III, Theorem~7.3; V, equation~(6.1.1),
Example~6.1.2, and Lemma~6.7.3]{Weibel}.

Given $z\in K_3(F)$, surjectivity of the Milnor residue map
supplies an image $m$ of $K_3^M(F)$ with the same residues as
$z$. Exactness of the Quillen sequence gives
$z=\operatorname{res}(\beta)+m$ for some $\beta\in K_3(k)$.
The rational Milnor image has pure Adams weight three
\cite[IV, Example~5.9.1]{Weibel}.
Naturality of the rational Adams decomposition
\cite[IV, Theorem~5.11]{Weibel} therefore makes projection
to weight two surjective from $K_3(k)_{\Q}^{(2)}$.
Injectivity follows from the injective constant-field map in
the Quillen sequence, proving the displayed isomorphism.
The involution $u\mapsto-u$ fixes $k$ and hence acts trivially
on this space; in characteristic two it is already the identity.
Thus the needed rational weight-two anti-invariant regulator
source is absent. The arithmetic number-field class is an
essential ingredient.

The scheme $T$ is regular and contains a field, but it is not semilocal: its nonzero Picard group already rules that out. Local Gersten theorems over fields therefore do not give an injection $K_3(T)\to K_3(\Frac T)$ for this global scheme.

Sections~\ref{d3:sec:node}--\ref{d3:sec:Delignedetection} prove persistence for the precise noncomplete localization; they do not by themselves justify completion. Section~\ref{d3:sec:completed} supplies a separate arithmetic $p$-adic proof for completed examples. Propositions~\ref{d3:prop:explicitcomplete}, \ref{d3:prop:explicittwocomplete}, and \ref{d3:prop:explicitoddcomplete} verify that the particular explicit classes at every residue prime survive both completion and henselization. These arguments use arithmetic detectors after completion.

\subsection{The distinction between the generic fiber and the fraction field}
The rings considered here are regular and flat over $\Z$, and are consequently $F$-smooth in the sense relevant to \cite{BouisMotivic}. Corollaries~5.12--5.13 of that reference apply to them: in degree $i$, the motivic comparison with $\Z/p^\nu(i)$-coefficients injects into the cohomology of the entire generic fiber $A[1/p]$. This is compatible with the construction here, whose nonzero classes are detected on that generic fiber and become zero only after a further restriction to the fraction field. The generic fiber has positive dimension and cannot be replaced by its generic point in that statement.

The special fiber $A/(p)$ of the basic examples is the singular domain $k[x,y,z]/(xy+z^2)$, or its corresponding completion. Thus these maps are not smooth over their coefficient DVR. They cannot become smooth through a different local DVR coefficient presentation either: primality of $(p)$ forces the base uniformizer to be associated to $p$, and the same singular special fiber remains. Absolute purity also does not assert that a principal-divisor Gysin map in higher $K$-theory is zero; triviality of the normal bundle gives the weaker self-intersection identity $i^*i_*=0$.

\subsection{Rational and finite coefficients}
The regulator proves that the kernel contains an element of infinite order. It does not by itself prove that any fixed finite-coefficient Gersten map fails. In the explicit $p=3$ case, Section~\ref{d3:sec:modthree} supplies a separate algebraic proof for $\Z/3$-coefficients. That coefficient characteristic equals the residue characteristic; no prime-to-residue-characteristic conclusion is asserted here. A rationally nonzero class can otherwise become divisible in a sufficiently large group. In particular, one may not interchange a filtered colimit over shrinking opens with an inverse limit over coefficient exponents without a separate argument.

Feld's Theorem~1.1 \cite{FeldFiniteGersten2026} constructs a nonzero
generic-kernel class in $K_2(A_{\mathrm F};\mathbb{Z}/3)$, where
$A_{\mathrm F}=V[[x,y]]/(3+x^2-y^3)$ and $V$ is a complete mixed-characteristic
$(0,3)$ discrete valuation ring with uniformizer $3$.
In that example, an integral lift acquires divisibility by three over
the fraction field while the integral $K_2$ restriction is injective
\cite[Corollary~5.1]{FeldFiniteGersten2026}.

\subsection{An explicit number-field source when the prime is three}
Let $u^2=-3$ and $\zeta=(1+u)/2$. Then $\zeta=e^{i\pi/3}$ at the chosen embedding and
\[
1-\zeta=\zeta^{-1}.
\]
Thus the pre-Bloch symbol $[\zeta]$ has zero wedge boundary. Its Bloch--Wigner dilogarithm is positive, since
\[
D(e^{i\pi/3})=-\int_0^{\pi/3}\log\bigl(2\sin(t/2)\bigr)\,dt>0.
\]
The rational Bloch-group description of $K_3^{\mathrm{ind}}(E)$ \cite[VI, Theorem~5.2]{Weibel} supplies an explicit possible rational source class. The main proof does not require fixing this comparison or its integral normalization: Borel's theorem and Lemma~\ref{d3:lem:beta} already produce an integral anti-invariant source with nonzero regulator.

\section{A mixed-square obstruction to a simplicial-functor construction}
\label{d3:sec:mochizukisquare}

Theorem~5 of \cite{MochizukiParameters} asserts that, for a noetherian
local ring $B$ and a non-zero-divisor $g$, the functor
$\mathcal P_{B/gB}\to\mathcal M_B(1)$ induces the zero map on
$K$-theory. Its application to $B=A$ and $g=x$ would kill the supported
class used above. We therefore examine a specific step in its proof.
The following calculation concerns Lemma~13(III) as written; it does
not, by itself, establish that Theorem~5 is false.

Let $B$ be a local ring, let $g$ be a non-zero-divisor in its maximal
ideal, and put $t=(1+g)^{-1}$. Write
\[
G=\operatorname{Typ}(g)=[B\xrightarrow{g}B],\qquad
I=\operatorname{Typ}(1)=[B\xrightarrow{1}B],\qquad M=G\oplus I,
\]
with complexes in homological degrees $1,0$. These are respectively
the objects $(1,0)_B$, $(0,1)_B$, and $(1,1)_B$ of the category $\mathcal C$
in \cite[Definition~9]{MochizukiParameters}. For $b\in B$ define
chain maps $\alpha_b:I\to M$ and $\beta_b:M\to G$ by
\begin{align*}
(\alpha_b)_1&=\binom{b}{1},&
(\alpha_b)_0&=\binom{gb}{1},\\
(\beta_b)_1&=\begin{pmatrix}1&-b\end{pmatrix},&
(\beta_b)_0&=\begin{pmatrix}1&-gb\end{pmatrix}.
\end{align*}
They give a split short exact sequence of complexes
\[
\mathcal S_b:\qquad 0\longrightarrow I\xrightarrow{\alpha_b}
M\xrightarrow{\beta_b}G\longrightarrow0.
\]
Indeed $s:G\to M$, given by $\binom10$ in both degrees, is a
chain section of $\beta_b$, while $r:M\to I$, given by
$\begin{pmatrix}0&1\end{pmatrix}$ in both degrees, satisfies
$r\alpha_b=1$ and $\alpha_b r+s\beta_b=1$. In the block notation
of \cite[Conventions~11]{MochizukiParameters}, the arrows are
$\binom b1$ and $\begin{pmatrix}1&-b\end{pmatrix}$, respectively.
Both are upper triangular. Thus $\mathcal S_b$ is an object of
$S_2^{\triangle}\mathcal C$ in that paper's notation.

There is an isomorphism of short exact sequences
$\mathcal S_1\to\mathcal S_t$ whose components on $I,M,G$ are,
respectively,
\[
\ell=(1+g)\operatorname{id}_I,\qquad
v_1=\begin{pmatrix}1&0\\g&1\end{pmatrix},\quad
v_0=\begin{pmatrix}1&0\\1&1\end{pmatrix},\qquad
q=t\operatorname{id}_G.
\]
The relation $\operatorname{diag}(g,1)v_1
=v_0\operatorname{diag}(g,1)$ verifies that $v$ is a chain map.
All three components are isomorphisms. The four naturality equations
can be checked without the block convention:
\begin{align*}
v_1\binom11&=\binom{1}{1+g}
             =\binom t1(1+g),&
v_0\binom g1&=\binom{g}{1+g}
             =\binom{gt}{1}(1+g),\\
\begin{pmatrix}1&-t\end{pmatrix}v_1
 &=\begin{pmatrix}t&-t\end{pmatrix},&
\begin{pmatrix}1&-gt\end{pmatrix}v_0
 &=\begin{pmatrix}t&-gt\end{pmatrix}.
\end{align*}
In the paper's block notation $v$ is
$\begin{pmatrix}1&0\\1&1\end{pmatrix}$, so every component is
lower triangular. This is therefore a morphism in
$i_{\mathrm{lower}}S_2^{\triangle}\mathcal C$.

The assignments $\mu'_1,\mu'_2$ preceding Definition~12 take an
object $(n,m)_B$ to $\operatorname{Typ}(g)^{\oplus n}$ or
$\operatorname{Typ}(1)^{\oplus n}$, and a morphism to its upper-left
$n'\times n$ block in both degrees. Each $\beta_b$ consequently
maps to the identity of $\operatorname{Typ}(g)$ or
$\operatorname{Typ}(1)$. The component $v$ maps to the identity,
whereas $q$ maps to multiplication by $t$. Naturality of the
resulting right-hand square would require $1=t$. This is false,
since $1-t=g/(1+g)\ne0$. The assignments thus fail to send this
morphism of $S_2$-objects to a morphism, contrary to
\cite[Lemma~13(III)]{MochizukiParameters}.

The failure already occurs with $B=\Z_{(3)}$, $g=3$, and $t=1/4$.
The calculation uses a regular DVR with regular quotient. It identifies
a failure of the stated simplicial-functor construction used in the
subsequent additivity argument.



% End input: sections/06-comparisons.tex


% Begin input: verification-scope.tex
\section*{Authoring and verification}
The OpenAI byline was requested by the user to attribute the preparation
of this consolidated article to ChatGPT/Codex. It does not identify an
official OpenAI publication or institutional endorsement.

This article consolidates the source manuscripts and the subsequent proof
audit with assistance from ChatGPT and Codex. The accompanying source
repository records the statements, proofs, and outcomes of independent
model-based reviews of the central constructions and their dependencies.
A separate Lean~4.34.1 development verifies ten polynomial identities
and two inequalities in natural numbers used in the calculations.
The Lean development does not formalize the complete algebraic
$K$-theory, regulator, or comparison arguments. The mathematical proofs
are given in the article; the repository records the precise scope of
the auxiliary checks and preserves the original sources.

% End input: verification-scope.tex


% Begin input: references.tex
\begin{thebibliography}{99}

\bibitem{Achinger}
P.~Achinger, \emph{Wild ramification and $K(\pi,1)$ spaces},
Invent. Math. \textbf{210} (2017), no.~2, 453--499.
Section~6.2, Theorem~6.2.1, records the Gabber--Fujiwara comparison
for every torsion abelian sheaf.
\href{https://arxiv.org/abs/1701.03197v2}{arXiv:1701.03197v2}
(4 May 2017);
\href{https://doi.org/10.1007/s00222-017-0733-5}{doi:10.1007/s00222-017-0733-5}.

\bibitem{AMMN}
B.~Antieau, A.~Mathew, M.~Morrow, and T.~Nikolaus,
\emph{On the Beilinson fiber square},
Duke Math. J. \textbf{171} (2022), no.~18, 3707--3806.
Propositions~2.8, 4.10, 4.12, Corollary~2.10, Theorem~2.12,
and Constructions~4.6, 4.11.
Numbered locators follow
\href{https://arxiv.org/abs/2003.12541v2}{arXiv:2003.12541v2}
(29 September 2021).

\bibitem{BerkovichEtale}
V.~G.~Berkovich, \emph{\'Etale cohomology for non-Archimedean analytic spaces},
Publications Math\'ematiques de l'IH\'ES \textbf{78} (1993), 5--161.
Theorem~6.1.1.
\href{https://www.numdam.org/article/PMIHES_1993__78__5_0.pdf}{Published article}.

\bibitem{Berthelot1997Duality}
Pierre Berthelot,
\emph{Dualit\'e de Poincar\'e et formule de K\"unneth en cohomologie rigide},
Comptes Rendus de l'Acad\'emie des Sciences de Paris,
S\'erie I, Math\'ematique \textbf{325} (1997), no.~5, 493--498.
\href{https://wstein.org/people/berthelo/publis/Poincare_Kunneth.pdf}{Author's text}.

\bibitem{BesserRegulators}
A.~Besser, \emph{Syntomic regulators and $p$-adic integration I:
rigid syntomic regulators}, Israel Journal of Mathematics
\textbf{120} (2000), part B, 291--334.
Proposition~8.6(2)--(3), its proof and equation~(8.1);
Proposition~9.9, Corollary~9.10, and Proposition~9.11.
Numbered locators refer to the author's version.
\href{https://www.math.bgu.ac.il/~bessera/reg/reg.html}{Author's original version};
\href{https://doi.org/10.1007/BF02834843}{doi:10.1007/BF02834843}.

\bibitem{BesserDeJeuAlgorithm}
A.~Besser and R.~de~Jeu, \emph{Li$^{(p)}$-service? An algorithm for computing $p$-adic polylogarithms}, Math. Comp. \textbf{77} (2008), no.~262, 1105--1134.
\href{https://doi.org/10.1090/S0025-5718-07-02027-3}{doi:10.1090/S0025-5718-07-02027-3};
\href{https://research.vu.nl/ws/portalfiles/portal/2459860/220267.pdf}{Published article}.

\bibitem{BesserDeJeuCurves}
A.~Besser and R.~de Jeu, \emph{The syntomic regulator for $K_4$ of curves}, Pacific J. Math. \textbf{260} (2012), no.~2, 305--380. Equation~(1.8), p.~307, and p.~377 record the Coleman dilogarithm normalization and functional equations.
\href{https://msp.org/pjm/2012/260-2/pjm-v260-n2-p03-s.pdf}{Published article}.

\bibitem{BesserDeJeu}
A.~Besser and R.~de Jeu, \emph{The syntomic regulator for the $K$-theory of fields}, Ann. Sci. \'Ecole Norm. Sup. (4) \textbf{36} (2003), no.~6, 867--924. Equation~(1.2) and Theorems~1.6, 1.10, 1.12.
\href{https://www.numdam.org/article/ASENS_2003_4_36_6_867_0.pdf}{Published article};
\href{https://doi.org/10.1016/j.ansens.2003.01.003}{doi:10.1016/j.ansens.2003.01.003}.

\bibitem{BGTTrace}
A.~J.~Blumberg, D.~Gepner, and G.~Tabuada,
\emph{Uniqueness of the multiplicative cyclotomic trace},
Advances in Mathematics \textbf{260} (2014), 191--232, Theorem~7.4.
\href{https://arxiv.org/abs/1103.3923v3}{arXiv:1103.3923v3}.

\bibitem{BokstedtOttosen}
M.~B\"okstedt and I.~Ottosen,
\emph{A Hochschild--Kostant--Rosenberg theorem for cyclic homology},
J. Pure Appl. Algebra \textbf{221} (2017), no.~6, 1458--1493.
Section~2, especially Definition~2.1, Remark~2.2, and Proposition~2.3.
\href{https://arxiv.org/abs/1601.07412v2}{arXiv:1601.07412v2}.

\bibitem{BondarkoIvanovChowWeights}
M.~V.~Bondarko and M.~A.~Ivanov,
\emph{On Chow weight structures for $cdh$-motives with integral coefficients},
arXiv preprint, version~2, 18 August 2015. Remark~1.3.2.
\href{https://arxiv.org/abs/1506.00631v2}{arXiv:1506.00631v2}.

\bibitem{Borel}
A.~Borel, \emph{Stable real cohomology of arithmetic groups}, Ann. Sci. \'Ecole Norm. Sup. (4) \textbf{7} (1974), no.~2, 235--272. In particular, Proposition~12.2 gives the relevant rank calculation.
\href{https://doi.org/10.24033/asens.1269}{doi:10.24033/asens.1269}.

\bibitem{BouisMotivic}
T.~Bouis, \emph{Motivic cohomology of mixed characteristic schemes},
arXiv preprint, version~3, 7 September 2026, Corollaries~5.12--5.13.
\href{https://arxiv.org/abs/2412.06635v3}{arXiv:2412.06635v3}.

\bibitem{BunkeTamme}
U.~Bunke and G.~Tamme, \emph{Multiplicative differential algebraic $K$-theory and applications}, Annals of $K$-Theory \textbf{1} (2016), no.~3, 227--258. See Remark~2.6, Theorem~2.31, and Section~6.
Numbered locators refer to the published version.
\href{https://doi.org/10.2140/akt.2016.1.227}{doi:10.2140/akt.2016.1.227};
\href{https://msp.org/akt/2016/1-3/akt-v1-n3-p01-s.pdf}{Published article}.

\bibitem{CaroDAddezio2025}
D.~Caro and M.~D'Addezio,
\emph{Injectivity failure in crystalline comparisons},
arXiv preprint, version~1, 14 November 2025.
\href{https://arxiv.org/abs/2511.11444v1}{arXiv:2511.11444v1}.

\bibitem{CartanStein1952}
H.~Cartan,
\emph{Faisceaux analytiques sur les vari\'et\'es de Stein :
d\'emonstration des th\'eor\`emes fondamentaux},
S\'eminaire Henri Cartan \textbf{4} (1951--1952),
expos\'e~19, 1--15 (1952).
\href{https://numdam.org/item/SHC_1951-1952__4__A19_0/}{Published article}.

\bibitem{CisinskiDegliseMotives}
D.-C.~Cisinski and F.~D\'eglise,
\emph{Triangulated categories of mixed motives},
Springer Monographs in Mathematics, Springer, 2019.
Introduction~B, p.~xxii; Theorem~11.1.24;
Examples~11.1.25 and~11.2.3; Corollaries~14.2.14 and~14.2.19;
Theorem~16.1.4 and Corollary~16.1.7.
\href{https://arxiv.org/abs/0912.2110v4}{arXiv:0912.2110v4}.

\bibitem{ColmezNiziol}
P.~Colmez and W.~Nizio\l, \emph{On $p$-adic comparison theorems for rigid analytic varieties, I},
M\"unster J. Math. \textbf{13} (2020), 445--507.
Theorem~1.1(4), Theorem~1.3(2), Theorem~7.9, and
Sections~3.1.1, 3.2.3, 5.1, and~7.2.4.
Numbered locators follow
\href{https://arxiv.org/abs/1905.04721v2}{arXiv:1905.04721v2}
(5 October 2019).

\bibitem{DLZ2011}
Christopher Davis, Andreas Langer, and Thomas Zink,
\emph{Overconvergent de Rham--Witt cohomology},
Annales Scientifiques de l'\'Ecole Normale Sup\'erieure (4)
\textbf{44} (2011), no.~2, 197--262.
\href{https://www.numdam.org/item/10.24033/asens.2143.pdf}{Published article}.

\bibitem{DeJeuCurves}
R.~de Jeu, \emph{Towards Regulator Formulae for the $K$-Theory of Curves over Number Fields}, Compositio Math. \textbf{124} (2000), 137--194. Theorem~2.3 restates the all-embedding complex regulator formula for the symbol map.
\href{https://doi.org/10.1023/A:1026440915009}{Published article};
\href{https://arxiv.org/abs/math/9811192}{Earlier preprint}.

\bibitem{DeJeuWedge}
R.~de Jeu, \emph{Zagier's conjecture and wedge complexes in algebraic $K$-theory}, Compositio Math. \textbf{96} (1995), no.~2, 197--247. Proposition~4.1 and Remark~5.2.
\href{https://www.numdam.org/article/CM_1995__96_2_197_0.pdf}{Published article}.

\bibitem{DennisStein}
R.~K.~Dennis and M.~R.~Stein, \emph{$K_2$ of discrete valuation rings}, Advances in Mathematics \textbf{18} (1975), no.~2, 182--238. The modern sign convention used here is specified explicitly in Section~\ref{d3:sec:explicit}.
\href{https://doi.org/10.1016/0001-8708(75)90157-7}{doi:10.1016/0001-8708(75)90157-7}.

\bibitem{EsnaultViehweg}
H.~Esnault and E.~Viehweg, \emph{Deligne--Beilinson cohomology}, in \emph{Beilinson's Conjectures on Special Values of $L$-Functions}, Perspectives in Mathematics \textbf{4}, Academic Press, 1988, 43--91. See Remark~2.13, Section~3 for products, and Section~8 for Chern classes.
\href{https://page.mi.fu-berlin.de/esnault/preprints/helene/13-deligne_beilinson.pdf}{Author's institutional copy}.

\bibitem{FeldFiniteGersten2026}
N.~Feld, \emph{Finite-coefficient Gersten injectivity fails in ramified mixed
characteristic}, preprint, 2026.
\href{https://arxiv.org/abs/2608.05005v2}{arXiv:2608.05005v2},
version~2, 24 September 2026; first submitted 5 August 2026.
Theorem~1.1 and Section~5.

\bibitem{GeisserLevineCharP}
T.~Geisser and M.~Levine, \emph{The $K$-theory of fields in
characteristic $p$}, Inventiones Mathematicae \textbf{139}
(2000), no.~3, 459--493. Theorem~8.4.
\href{https://www2.rikkyo.ac.jp/web/geisser/p-part.pdf}{Published article};
\href{https://doi.org/10.1007/s002220050014}{doi:10.1007/s002220050014}.

\bibitem{Ginot}
G.~Ginot, \emph{Formules explicites pour le caract\`ere de Chern en $K$-th\'eorie alg\'ebrique}, Annales de l'Institut Fourier \textbf{54} (2004), no.~7, 2327--2355.
\href{https://www.math.univ-paris13.fr/~ginot/papers/flecarch.pdf}{Author's preprint};
\href{https://doi.org/10.5802/aif.2081}{doi:10.5802/aif.2081}.

\bibitem{GrosseKlonne}
E.~Grosse-Kl\"onne, \emph{Rigid analytic spaces with overconvergent structure sheaf}, Journal f\"ur die reine und angewandte Mathematik \textbf{519} (2000), 73--95. Proposition~3.1 gives coherent acyclicity for dagger affinoids.
\href{https://arxiv.org/abs/1408.3329v1}{arXiv:1408.3329v1}.

\bibitem{HoyoisCircle}
M.~Hoyois, \emph{The homotopy fixed points of the circle action on
Hochschild homology}, arXiv preprint, version~2, 21 April 2018.
Theorem~2.1, Lemma~2.2, and Theorem~2.3.
\href{https://arxiv.org/abs/1506.07123v2}{arXiv:1506.07123v2}.

\bibitem{HuberKings}
A.~Huber and G.~Kings, \emph{A $p$-adic analogue of the Borel regulator and the Bloch--Kato exponential map},
J. Inst. Math. Jussieu \textbf{10} (2011), no.~1, 149--190.
Section~2 gives the Chern-compatible comparison with Besser's
syntomic cohomology and the exponential, without an unramifiedness restriction.
Numbered citations refer to
\href{https://arxiv.org/abs/math/0612611v1}{arXiv:math/0612611v1}.

\bibitem{Jannsen}
U.~Jannsen, \emph{Continuous \'etale cohomology}, Mathematische Annalen \textbf{280} (1988), 207--245. Equations~(0.2), (0.4), and Section~6.
\href{https://epub.uni-regensburg.de/26684/1/jannsen12.pdf}{Published article}.

\bibitem{KerzThesis}
M.~Kerz, \emph{Milnor $K$-theory of local rings}, doctoral thesis, Universit\"at Regensburg, 2008. Theorem~4.2.5, Proposition~4.2.8(4)--(6), and Theorem~4.3.2 give the improved theory, the ordinary-to-improved surjection, and the comparison with Quillen $K$-theory whose composite is $(n-1)!$ for general local rings.
\href{https://d-nb.info/98935329x/34}{Full thesis}.

\bibitem{LevineHigherChowRevisited}
M.~Levine,
\emph{Bloch's higher Chow groups revisited},
Ast\'erisque \textbf{226} (1994), 235--320. Theorem~4.7.
\href{https://www.numdam.org/item/AST_1994__226__235_0/}{Published article}.

\bibitem{LevineIndK3}
M.~Levine, \emph{The indecomposable $K_3$ of fields}, Ann. Sci. \'Ecole Norm. Sup. (4) \textbf{22} (1989), no.~2, 255--344. Theorem~4.12 is the finite-coefficient Chern comparison used here; see also Theorem~4.13 and Section~5.1.
\href{https://doi.org/10.24033/asens.1585}{doi:10.24033/asens.1585}.

\bibitem{MerkurjevTorsion}
A.~S.~Merkurjev, \emph{On the torsion in $K_2$ of local fields}, Annals of Mathematics \textbf{118} (1983), no.~2, 375--381.
\href{https://annals.math.princeton.edu/1983/118-2/p06}{Publisher's page};
\href{https://doi.org/10.2307/2007033}{doi:10.2307/2007033}.

\bibitem{MilneDuality}
J.~S.~Milne, \emph{Arithmetic Duality Theorems}, second edition, BookSurge, 2006. Chapter~I, Corollary~2.3 and the discussion preceding Theorem~2.8 give local Galois finiteness.
\href{https://www.jmilne.org/math/Books/ADTnot.pdf}{Author's edition}.

\bibitem{MilneEtale}
J.~S.~Milne, \emph{Lectures on \'Etale Cohomology}, author's lecture notes, version~2.21, 22 March 2013. Theorem~19.1(b) gives geometric finite-coefficient finiteness.
\href{https://www.jmilne.org/math/CourseNotes/LEC.pdf}{Author's notes}.

\bibitem{MochizukiParameters}
S.~Mochizuki, \emph{Local Gersten's conjecture for regular system of parameters}, preprint, version~8, 6 August 2020. Section~\ref{d3:sec:mochizukisquare} examines Lemma~13(III) in this version.
\href{https://arxiv.org/abs/1503.07966v8}{arXiv:1503.07966v8}.

\bibitem{NikolausScholze}
T.~Nikolaus and P.~Scholze, \emph{On topological cyclic homology},
Acta Mathematica \textbf{221} (2018), no.~2, 203--409, Corollary~I.4.3.
\href{https://arxiv.org/abs/1707.01799v2}{arXiv:1707.01799v2}.

\bibitem{Petrequin2003}
D.~Petrequin, \emph{Classes de Chern et classes de cycles en
cohomologie rigide}, Bulletin de la Soci\'et\'e Math\'ematique
de France \textbf{131} (2003), no.~1, 59--121.
\href{https://www.numdam.org/article/BSMF_2003__131_1_59_0.pdf}
{Published article};
\href{https://doi.org/10.24033/bsmf.2437}{doi:10.24033/bsmf.2437}.

\bibitem{Quillen}
D.~Quillen, \emph{Higher algebraic $K$-theory I}, in \emph{Algebraic $K$-Theory I: Higher $K$-Theories}, Lecture Notes in Mathematics \textbf{341}, Springer, 1973, 85--147.
\href{https://doi.org/10.1007/BFb0067053}{doi:10.1007/BFb0067053}.

\bibitem{QuillenFinite}
D.~Quillen, \emph{On the cohomology and \(K\)-theory of the general
linear groups over a finite field}, Ann. of Math. (2)
\textbf{96} (1972), no.~3, 552--586.
\href{https://doi.org/10.2307/1970825}{doi:10.2307/1970825}.

\bibitem{SaitoCompletion}
T.~Saito, \emph{Graded quotients of ramification groups of local fields with imperfect residue fields}, Amer. J. Math. \textbf{145} (2023), no.~5, 1389--1464, Lemma~4.2.1.
\href{https://preprint.press.jhu.edu/ajm/sites/default/files/AJM-saito.pdf}{Author's version};
\href{https://arxiv.org/abs/2004.03768v2}{arXiv:2004.03768v2};
\href{https://doi.org/10.1353/ajm.2023.a907702}{doi:10.1353/ajm.2023.a907702}.

\bibitem{SilvermanAEC2009}
Joseph H. Silverman,
\emph{The Arithmetic of Elliptic Curves},
second edition, Graduate Texts in Mathematics, vol.~106,
Springer, New York, 2009.
\href{https://doi.org/10.1007/978-0-387-09494-6}{doi:10.1007/978-0-387-09494-6}.

\bibitem{StacksBlowup}
The Stacks Project Authors, \emph{The Stacks Project}, Section~31.33, ``Blowing up,'' especially the standard chart construction and tautological ideal identity in Lemmas~31.33.2 and~31.33.4.
\href{https://stacks.math.columbia.edu/tag/01OF}{Tag 01OF};
\href{https://stacks.math.columbia.edu/tag/0804}{Tag 0804};
\href{https://stacks.math.columbia.edu/tag/02OS}{Tag 02OS}.

\bibitem{StacksFlatness}
The Stacks Project Authors, \emph{The Stacks Project}, Lemma~10.128.1 (Miracle flatness), Tag~00R4.
\href{https://stacks.math.columbia.edu/tag/00R4}{Tag 00R4}.

\bibitem{StacksHenselization}
The Stacks Project Authors, \emph{The Stacks Project}, Section~10.156, Henselization and quasi-finite ring maps (Tag~0GIN), and Section~15.46, Permanence of properties under henselization (Tag~07QL).
\href{https://stacks.math.columbia.edu/tag/0GIN}{Tag 0GIN}; \href{https://stacks.math.columbia.edu/tag/07QL}{Tag 07QL}. Also Section~10.153, \href{https://stacks.math.columbia.edu/tag/04GK}{Tag 04GK}, and Section~10.155, \href{https://stacks.math.columbia.edu/tag/0BSK}{Tag 0BSK}.
For the cited permanence and quotient results, see
\href{https://stacks.math.columbia.edu/tag/07QM}{Tag 07QM}
(Lemma~15.46.1),
\href{https://stacks.math.columbia.edu/tag/06LJ}{Tag 06LJ}
(Lemma~15.46.3),
\href{https://stacks.math.columbia.edu/tag/06LN}{Tag 06LN},
\href{https://stacks.math.columbia.edu/tag/0AP3}{Tag 0AP3}
(Lemma~15.46.11),
\href{https://stacks.math.columbia.edu/tag/05WQ}{Tag 05WQ}
(Lemma~10.156.2), and
\href{https://stacks.math.columbia.edu/tag/05WS}{Tag 05WS}
(Lemma~10.156.4).

\bibitem{Tamme}
G.~Tamme, \emph{Karoubi's relative Chern character, the rigid syntomic regulator, and the Bloch--Kato exponential map},
Forum Math. Sigma \textbf{2} (2014), e20.
Corollary~5.19 and its proof give the syntomic-to-\'etale compatibility used for $\Spec\OO_E$.
\href{https://arxiv.org/abs/1111.4109v4}{arXiv:1111.4109v4}
(21 July 2014).

\bibitem{TotaroMilnor}
B.~Totaro, \emph{Milnor $K$-theory is the simplest part of algebraic
$K$-theory}, $K$-Theory \textbf{6} (1992), 177--189.
Theorem~1 and Sections~1--2.
\href{https://www.math.ucla.edu/~totaro/papers/public_html/milnor.pdf}
{Author's copy}.

\bibitem{WeibelChern}
C.~Weibel, \emph{\'Etale Chern classes at the prime 2},
in P.~G.~Goerss and J.~F.~Jardine (eds.),
\emph{Algebraic $K$-Theory and Algebraic Topology},
NATO ASI Series C \textbf{407}, Kluwer, 1993, 249--286.
Theorem~3.11(i) gives the line-bundle product formula in $K$-degree three;
Section~5 recalls the indecomposable $K_3$ comparison.
Page locators 10 and 12 refer to the author's version.
\href{https://sites.math.rutgers.edu/~weibel/archive/papers-dir/chernclass.pdf}{Author's paper};
\href{https://doi.org/10.1007/978-94-017-0695-7_14}{doi:10.1007/978-94-017-0695-7\_14}.

\bibitem{Weibel}
C.~A.~Weibel, \emph{The $K$-book: An Introduction to Algebraic $K$-theory}, Graduate Studies in Mathematics \textbf{145}, American Mathematical Society, 2013.
Chapter~III, Definition~5.11 and relations (D1)--(D3), pp.~43--44,
fix the Dennis--Stein convention used here; Chapter~III,
Theorem~6.2.4, p.~50, gives the local-field $K_2$ decomposition.
See also Chapter~IV, Theorem~1.12 and Corollary~1.13, p.~10,
for finite fields, Section~2 for coefficients, and Section~6.4
for filtered colimits; Chapter~V, 3.3, 3.3.2, 3.7.2, 3.12,
4.1, and 6.1 for the cited transfers, base change, projection,
d\'evissage, and localization results; and Chapter~VI,
Theorems~4.1 and~5.2 for norm-residue and the rational
Bloch-group identification. Page locators refer to the author PDFs.
\href{https://sites.math.rutgers.edu/~weibel/Kbook/Kbook.III.pdf}{Author's Chapter~III};
\href{https://sites.math.rutgers.edu/~weibel/Kbook/Kbook.IV.pdf}{Author's Chapter~IV};
\href{https://sites.math.rutgers.edu/~weibel/Kbook/Kbook.V.pdf}{Author's Chapter~V};
\href{https://sites.math.rutgers.edu/~weibel/Kbook/Kbook.VI.pdf}{Author's Chapter~VI}.
\end{thebibliography}

% End input: references.tex

\end{document}
